\\documentclass[12pt]{article}
\\usepackage[T1]{fontenc}
\\usepackage[utf8]{inputenc}
\\usepackage[brazil]{babel}
\\usepackage[a4paper,margin=2.5cm]{geometry}
\\usepackage{setspace}
\\setstretch{1.15}
\\title{INDÚSTRIA DO PETRÓLEO PARA NÃO ENGENHEIROS}
\\date{}
\\begin{document}
\\maketitle

\\section*{Pagina 1}

INDOSTRIA\\
DO PETROLEO\\
\\
AO-ENGENHEIROS\\
\\
eRe\\
\\
Digitalizada com CamScanner\\

\\section*{Pagina 2}

INDUSTRIA DO PETROLEO\\
\\
PARA NAO-ENGENHEIROS\\
\\
Digitalizada com CamScanner\\

\\section*{Pagina 3}

INDUSTRIA DO PETROLEO\\
\\
PARA NAO-ENGENHEIROS\\
\\
Geraldo André Raposo Ramos\\
\\
i=\\
\\
EDITORA\\
\\
DOISO3\\
\\
Digitalizada com CamScanner\\

\\section*{Pagina 4}

Aos meus pais (in memoriam),\\
esposa e filhos.\\
\\
———————————\\
\\
Digitalizada com CamScanner\\

\\section*{Pagina 5}

Sumario\\
\\
1 Conceitos gerais\\
\\
1\\
1.1 O que é Engenharia de Petréleos? 1\\
1.2 O que é 0 petrdleo? 3\\
1.3 O primeiro pogo de petrdleo 10\\
1.4 Historia do Petrdleo no mundo 41\\
1.5 — Histérico do petréleo em Angola 13\\
1.6  Cadeia produtiva do Petrdleo e Gas Natural 16\\
1.7. Contratos de exploragao e producao petrolifera 18\\
2  Cadeia de valor da industria petrolifera 25\\
2.1. Segmento da cadeia produtiva de Upstream 26\\
2.2. Segmentos de Midstream e Downstream 39\\
3  Exploragao de hidrocarbonetos 41\\
3.1 Teorias das placas tectonicas 41\\
3.2 Conceito das rochas 47\\
3.3 Geologia de petrdleos 53\\
34  Exploragao Petrolifera 61\\
4 Perfuracdo e completacao de pocos 67\\
4.1 Perfuracdo de pogos de petréleo 68\\
4.2  Tipos e selecdo de sondas ou plataformas de perfuragao 73\\
4.3 Completacao de pogos 106\\
4.4 Operacao Start-up \\& First Oil 109\\
4.5. Abandono de pogos 110\\
46 Risco na atividade de perfuracao e completagao 110\\
47 Custos e contratos de perfuragao e completacgao de pocos 411\\
Avaliacao de Formacoes 115\\
5.1 Avaliagdo de macroestruturas 115\\
vii\\
\\
Digitalizada com CamScanner\\

\\section*{Pagina 6}

vill SUMARIO\\
\\
2 6 ‘S 1\\
5.2 Técnicas de avaliagao de fone 3\\
5.3 Estudos de reservatorio\\
éri 133\\
jos\\
§\\_Reserveses cas de um engenheiro de reservatérios 133\\
eee ed rochas e fluidos de reservatérios 134\\
6.2 Propriedades basicas das 1a\\
6.3  Tiposde reservat6rios ;\\
Métodos de recuperagao de hidrocarbonetos 157\\
‘s i s 161\\
65 Estimativas de reserva\\
zs é6 ubmarinas 169\\
7 Producao e operacoes S\\
1\\
7.1 Producao em Angola ue\\
72 Sistema de Produgao F\\
73 Vendas e armazenamento do petrdleo 182\\
7.4 Descarte de agua salgada 183\\
75  Tratamento de gas 184\\
7.6 — Plataformas de produgao ; 185\\
7.7 Paises da Organizacao dos Paises Exportadores de Petrdleo\\
(OPEP) 194\\
8 Midstream e Downstream 197\\
8.1 Caracterizagao de petrdleo bruto 198\\
8.2 Processos de refinacao de petrdleo bruto 203\\
8.3 Transferéncia e Armazenamento do Petréleo Bruto 223\\
84  Armazenamento e distribuigao dos derivados do petréleo 229\\
9 Seguranca, meio ambiente e satde na industria 233\\
9.1 Principais impactos ambientais 234\\
9.2 Relacéo humana com a natureza 238\\
9.3 Aspectos ambientais de perfuracao terrestre de pocos 239\\
9.4 Aspectos de seguranca e satide ocupacional 247\\
9.5 Principais acidentes registrados 249\\
9.6 — Fontes alternativas do petréleo 250\\
Acrénim\\
ii 253\\
R\\
eferénclas Bibliograficas 275\\
\\
Digitalizada com CamScanner\\

\\section*{Pagina 7}

LISTA DE FIGURAS\\
\\
1.1 Configuracao tipica de uma jazida de petrdleo 3\\
1.2 Esquema de um poco de extracao 3\\
13 Exemplo de butano 5\\
14 Exemplo de eteno 5\\
1.5 Exemplo de etino 6\\
1.6 Estrutura de Butadieno 6\\
17 Estrutura de Ciclohexano 6\\
1.8 Estrutura de Ciclobuteno 6\\
19 Estrutura de Anel Benzeno 7\\
1.10 Resumo de tipos dos hidrocarbonetos 7\\
4411 Reservas mundias de petréleo 11\\
1.12 Mapa da antiga Babilénia 12\\
1.13 Coronel Drake a esquerda e a direita 0 primeiro pogo 12\\
1.14 Engenheiro mecanico, Anthony Francis Lucas 13\\
1.15 Refinaria de Luanda 14\\
\\
1.16 Sanha Floating production storage and offloading /Unidade flutuante\\
de armazenamento e transferéncia (FPSO) 15\\
\\
1.17 Cadeia produtiva da industria petrolifera adotada pela\\
\\
Sonangol, E.P. 16\\
\\
1.18 Estrutura da divisao do segmento de upstream 16\\
1.19 Ilustragdo grafica da partilha de producao baseada nos termos\\
\\
fiscais 22\\
\\
21 Cadeia produtiva da industria petrolifera - A 25\\
\\
ix\\
\\
Digitalizada com CamScanner\\

\\section*{Pagina 8}

x LISTA DE FIGURAS\\
\\
22 Cadeia produtiva da industria petrolifera - B 26\\
industria petrolifera - C 26\\
2.3 Cadeia produtiva da industria p\\
24 principais etapas do segmento de Upstream 27\\
25 Segmento de Upstream .\\
2.6 Plateau de produgao a\\
27 Investimento tipico para campos or apr es *\\
28 Custos de sondas e equipamentos de perfuracao 39\\
2.9 Custos de perfuragao de pogos 39\\
3.1 Tipos de Limites de Placas 42\\
3.2 Deriva continental 43\\
3.3 Placas Tecténicas 45\\
34 Sondas eletroacusticas para investigagao submarina 4\\%\\
3.5 Algumas tecnologias utilizadas para exploracao do fundo\\
oceanico 46\\
3.6 Relevo submarino dos oceanos 47\\
3.7 Ciclo de transformacao de rochas 48\\
3.8 Rochas igneas ou magmaticas 49\\
3.9 Exemplo de rochas metamorficas 52\\
3.10 Rochas sedimentares clasticas 52\\
3.11 Rochas sedimentares quimicas 53\\
3.12 Rochas sedimentares organicas 53\\
3.13 Técnicas para a busca de petréleo 54\\
3.14 Posicionamento dos fluidos no subsolo 54\\
3.15 Maturacdo da matéria Organica 55\\
3.16 Esquematica da estrutura com a profundidade de pré-sal 56\\
3.17 \\_Jazida de Petréleo 57\\
3.18 Acumutacao do Petrdéleo 39\\
319 Trapa Estrutural 60\\
\\
4\\
Digitalizada com CamScanner\\

\\section*{Pagina 9}

USTADEFIGURAS =X\\
\\
3.20 Trapa Estratigrafica 60\\
3.21 Trapa mista 61\\
3.22 Deformagées da Crosta Terrestre 62\\
3.23 Mapa de Bouguer 64\\
3.24 Mapa aeromagnético 64\\
3.25 Métodos Sismicos 65\\
41 Cacamba de perfuracao 68\\
42 Exemplo de alguns pocos quanto 4 finalidade e trajetéria 70\\
43 Sonda terrestre de alta perfuracdo 75\\
44 Barcaca de pantano de perfuracao 76\\
4.5 Plataformas de perfuragao maritimas: fixa, autoelevaveis e\\
\\
semissubmersivel 78\\
4.6 Plataformas maritimas de perfuragao: FPSO Monocoluna,\\
\\
Navio-Sonda e SPAR 79\\
47 Ceriménia tradicional de batismo do Libongos na Coreia do Sul 80\\
4.8 Principais equipamentos que comp6em uma sonda de\\
\\
perfuracao de petrdleo 82\\
49 Esquema de uma sonda rotativa com uma torre e 0 Mastro de\\
\\
perfuragao 83\\
4.10 Sistema de movimentacao de cargas 84\\
4.11 Mesa rotativa, Kelly, bucha de Kelly, bucha da mesa rotativa 86\\
412 — Cabeca de injecao (swivel) 87\\
4.13 Top drive do sistema de rotagdo da sonda de perfuracao 88\\
4.14 Esquema de uma sonda mecdnica com cinco motores diesel 88\\
4.15  Esquema de uma sonda AC/DC 89\\
4.16 Sistema de circulacao de lama 90\\
417 Conjunto tipico de um Blowout Preventer/ Preventor de Erup¢ao,\\
\\
Obturador de Seguranca (BOP) 92\\
4.18 Painel do sistema de monitoramento de uma sonda de\\
\\
perfuracao 4\\
\\
Digitalizada com CamScanner\\

\\section*{Pagina 10}

LISTA DE FIGURAS\\
\\
xii\\
419 Comando de perfuragao (Drill mum) 58\\
4.20 Tubos de perfuragao pesados "\\
4.21 Tubos de perfuracao i\\
4.22 Brocas de perfuracao de pocos 95\\
4.23 Substitutos \\%6\\
4.24 Estabilizadores 7\\
4.25 Escareadores 7\\
4.26 Alargadores 97\\
4.27 Chaves flutuantes e iron roughneck 98\\
4.28 Cunhas e colar de seguranga 98\\
4.29 Classificagao dos fluidos de perfuragao 99\\
4.30 Revestimento de poco 103\\
4.31 Cimentacao de revestimento 103\\
4.32 Categorias basicas de completacgao 107\\
4.33 Arvore de natal convencional 107\\
4.34 Condicionamento do poco 108\\
4.35 Canhoneio 109\\
436 Tipos de canhoneio utilizado 109\\
4.37 Resumo de custos de atividades de perfuracao e completacao\\
de pocos 111\\
5A Comparacao entre resolucao sismica e perfis 123\\
52 Grafico de gradientes de pressdo 125\\
6.1 Composicao da rocha-reservatério 135\\
6.2 Classificacao de Porosidade 135\\
6.3 Porosidade 136\\
mA cee hg origem: A, D porosidade primaria; B, 137\\
6.5 Saturacao de fluidos 139\\
\\
Digitalizada com CamScanner\\

\\section*{Pagina 11}

LISTA DE FIGURAS xiii\\
6.6 Permeabilidade de uma rocha 140\\
6.7 Esquema do experimento de Henry Darcy 141\\
68 Itustragao da definicao do valor de 1 Darcy 141\\
69 Velocidade de Darcy 142\\
6.10 Fluxo linear 143\\
6.11 Fluxo linear com inclinagao com fluxo para cima 144\\
6.12 Fluxo linear com inclinagao com fluxo para baixo 144\\
6.13 Intervalo de aplicabilidade da Lei de Darcy 145\\
6.14 Fluxo radial permanente 145\\
6.15 Diagrama esquematico da medi¢ao de permeabilidade do tipo\\
\\
Hassler 147\\
6.16 Efeito Klinkenberg No ensaio experimental com hidrogénio,\\
\\
nitrogénio e gas carbénico 148\\
6.17 Grafico de kg em funcao de 1/P,, 149\\
6.18 Efeito da sobrecarga sobre a permeabilidade 149\\
6.19 Curva de permeabilidade relativa em funcao da saturacao de\\
\\
agua 150\\
6.20 Influéncia da saturagao sobre as curvas de permeabilidade\\
\\
relativa 151\\
6.21 Molhabilidade de fluidos 151\\
6.22 Pressao capilar entre agua e ar 153\\
6.23 Diagrama de fases de uma mistura de hidrocarbonetos 154\\
6.24 Diagrama de fases misturas liquidas 155\\
6.25 Grafico de fase de envelopes ilustrando os reservatérios de\\
\\
; 6leo e gas 156\\
\\
6.26 Métodos de recuperagao de hidrocarbonetos 158\\
6.27. \\_ Exemplo de injecao em malhas e periférica 160\\
6.28. - Tipos de reservas de petrdleo segundo a SEC 163\\
6.29 Volume original reservatério de dleo 165\\
7A Sistema de producio completo 170\\
\\
Digitalizada com CamScanner\\

\\section*{Pagina 12}

xiv\\
7.2\\
ra)\\
74\\
\\
75\\
76\\
\\
77\\
\\
78\\
79\\
7.10\\
711\\
712\\
713\\
714\\
715\\
7.16\\
TAT\\
"7:18\\
719\\
7.20\\
721\\
7.22\\
723\\
744\\
7.25\\
81\\
82\\
\\
LISTA DE FIGURAS\\
\\
Ilustragao das tras etapas do sist\\
\\
i de\\
Fluidos produzidos com um sistema\\
\\
Pocos equipados com o sistema de gas lift\\
\\
Elevacao arti\\
\\
ema de producao\\
\\
elevacao natural\\
\\
172\\
173\\
174\\
\\
ificial através do bombeio mecanico com hastes 475\\
\\
‘al com o sistema de Bombeamento Centrifugo\\
\\
Baie EE Submersible Pumping (BCS) 176\\
Sistema de elevacao artificial com Bombeamento de Cavidades\\
Progressivas/ Progressive Cavity Pumping (BCP) M7\\
Esquema do processamento primario de petrdleo 178\\
Separador bifasico horizontal 180\\
Separador trifasico vertical 181\\
Separador de teste 181\\
Free water knockout horizontal 182\\
Free water knockout vertical 182\\
Aquecedor-tratador horizontal 182\\
Processo de adogamento de gas 184\\
Configuracao da planta de processamento 185\\
Plataformas de producao fixas e flutuantes 187\\
Plataformas de producao flutuantes 188\\
Sistema submarino de producao 190\\
Cabegca de poco submarino 190\\
Veiculo de Pperacao Remota (ROV) 191\\
Arvore de Natal Molhada (ANM) 191\\
Umbilical 192\\
Linhas flexiveis e risers flexiveis 193\\
Manifolde 193\\
Exemplo de petréteo bruto 198\\
Relacao entre o grau Ap| do dte0 e os derivados extraides 200\\
\\
Digitalizada com CamScanner\\

\\section*{Pagina 13}

LSTA DE FIGURAS xv\\
\\
8.3 co da relagao entre a densidade e a densidade API\\
201\\
8.4 Relacao entre a viscosidade e a densidade do éleo em grau API 201\\
8.5 Esquema geral de uma refinaria 204\\
8.6 Processo de decantacao do petréleo 205\\
8.7 Esquematica do processo de pré-aquecimento do dleo bruto 205\\
8.8 \\textbackslash{}lustra 0 vaso separador para o processo de dessalinacao 206\\
8.9 Esquematica do processo eléctrico de dessalgagao 207\\
8.10 Esquematica do processo quimico de dessalgacao 207\\
8.11 Fornalha e torre de destilacao 208\\
8.12 Coluna de destilacao atmosférica 209\\
8.13 Processo de destilagao atmosférica 209\\
8.14 Asfalto 210\\
8.15 Coluna de destilacao a vacuo e seus produtos 210\\
8.16 Fluxograma de reforma catalitica 212\\
8.17 Esquema do percurso dos derivados do petréleo 214\\
8.18 Coluna de destilagdo com diferentes derivados do petroleo 214\\
8.19 Derivados do petréleo bruto 215\\
8.20 Refinaria de Luanda 217\\
8.21 Importancia e utilidade do petroleo bruto 218\\
8.22 Chicletes e refrigerantes 219\\
8.23 Planta de Angola LNG no Soyo, Angola 220\\
8.24 Plataforma de Sanha-FPSO Bloco 0 220\\
8.25 Opes A e B do desenho de uma UPGN 221\\
8.26  Opcdes C e D do desenho de uma UPGN 221\\
8.27 Produtos finais de uma refinaria ou UPGN 222\\
‘etroleo do poco ao produto final 222\\
\\
8.28 Esquema da industria do p\\
\\
8.29 Oleodutos para transferéncia de petrdleo bruto e seus derivados 224\\
\\
Digitalizada com CamScanner\\

\\section*{Pagina 14}

XVi-—\\_LISTA DE FIGURAS\\
\\
8.30 Navio petroleiro Jahre Vicking e Irati 224\\
8.31 Navio para o transporte de Gas Natural Liquefeito (GNL) 205\\
8.32 Transporte rodoviario dos derivados do petréleo 225\\
\\
8.33 Vagio ferroviario para 0 transporte de derivados de petréleo 26\\
\\
8.34 Tanque atmosférico 226\\
8.35 Tanques cilindricos e esféricos 227\\
8.36 Tanque cénico e curvo 228\\
8.37 Tanque flutuante 229\\
8.38 Carrosséis de enchimento de gas butano - Instalagdes Carlos\\
\\
Pinto Nogueira (ICPN) 230\\
8.39 Linha de enchimento das garrafas levita 230\\
8.40 Placas de abastecimento de camides em IBV-1 231\\
8.41 Lubrificante NGOL - IMUL 232\\
9.1 Escopo de atuacao de metodologia da producao mais limpa 239\\
9.2\\
\\
Esquema de processo de separacao térmica de fases 245\\
\\
Digitalizada com CamScanner\\

\\section*{Pagina 15}

LISTA DE TABELAS\\
a ee\\
\\
11\\
1.2\\
13\\
1.4\\
15\\
1.6\\
17\\
1.8\\
19\\
\\
2.1\\
2.2\\
23\\
24\\
2.5\\
3.1\\
3.2\\
3.3\\
41\\
\\
42\\
\\
43\\
\\
44\\
\\
Fragoes tipicas do petréleo\\
\\
4\\
Contaminantes ou nao-hidrocarbonetos 5\\
Exemplo de alguns hidrocarbonetos do tipo Alcanos e Alcenos 7\\
Classificacao do petréleo em Grau API 8\\
Taxa Interna de Rentabilidade (TIR) 21\\
Partilha de Producao do Bloco Y, baseada na TIR 21\\
\\
Partilha de Produgao do Bloco X, baseada na Producao passada 21\\
\\
Tipos de contratos 23\\
Comparacao entre o Contrato de Concessao (CC) e Contrato de\\
Partilha de Producgao (CPP) 23\\
Exemplo de CPP - Bloco AA 29\\
Trés fases que caracterizam a producdo de um campo 34\\
Ramas e qualidade do petréleo produzido em Angola 35\\
Ramas e qualidade do petréleo produzido no mundo 36\\
Principais diferencas entre gas convencional e gas de xisto 40\\
Placas principais 45\\
Placas menores 45\\
Classificagao das rochas 48\\
Tipos de sondas ou plataformas de perfuragao 75\\
Comparacao entre sondas terrestres e barcaca de pantanos 79\\
Comparacao das sondas ou plataformas de perfuracao de\\
petrdleo e gas 80\\
Principais custos na implementacao de um projeto de campo “\\
offshore\\
\\
xvii\\
\\
\\_\\
\\
Digitalizada com CamScanner\\

\\section*{Pagina 16}

xviii\\
\\
5.1\\
5.2\\
6.1\\
62\\
63\\
7A\\
7.2\\
73\\
7.4\\
75\\
7.6\\
77\\
78\\
8.1\\
8.2\\
8.3\\
8.4\\
8.5\\
8.6\\
8.7\\
8.8\\
91\\
9,2\\
\\
93\\
\\
9.4\\
\\
LSTA DE TABELAS\\
\\
incipais perfis\\
Resumo da aplicagao dos principals pel\\
\\
Exemplos de gradientes \\&\\
\\
Classificagdo de reservatorio CO!\\
\\
Métodos de medigao de permeabilidade\\
\\
Resumo de diferentes mecanismos de produgao\\
\\
Anomenclatura dos campos petroliferos\\
\\
producao de Angola de 2017 a 2020\\
\\
Composigao tipica de petrdleo bruto\\
\\
Composicao tipica do gas natural\\
\\
Petroleo de referéncia utilizado nas transagoes\\
Comparacao das atividades realizadas pelas plataformas\\
Numero de estruturas de produgao existentes em Angola\\
Reservas provadas da OPEP\\
\\
Rama de petrdéleo bruto produzido em Angola\\
\\
Ramas produzidas internacionalmente\\
\\
Tipos de petrdleo quanto a predominancia\\
\\
Exemplo de uma unidade de processo\\
\\
Resumo dos processos da refinacao de petréleo\\
\\
Indice de octano de alguns componentes\\
\\
Projetos para construcao de refinarias em Angola\\
Especificagdes de Gas de Petrdleo Liquefeito (GPL)\\
\\
Componentes ambientais dos diferentes meios\\
\\
Os principais aspectos ambientais potenciais causadores de\\
\\
impactos negativo\\
\\
Principais im A\\
pactos no meio i\\
ocorréncia antrépico e suas formas de\\
\\
Classificacado de res\\{duos sélidos - NBR 10.004\\
\\
123\\
\\
massas especificas tipicas de fluidos 124\\
\\
m base na permeabilidade 143,\\
\\
146\\
159\\
171\\
171\\
178\\
179\\
183\\
188\\
189\\
195\\
201\\
202\\
202\\
203\\
213\\
213\\
218\\
222\\
235\\
\\
235\\
\\
237\\
240\\
\\
7 |\\
\\
Digitalizada com CamScanner\\

\\section*{Pagina 17}

PREFACIO\\
nna\\
\\
O petréleo é 0 principal produto de exportacao de Angola. Portanto, adquirir conheci-\\
mentos sobre Petréleo e Gas Natural é crucial Para todos os angolanos, principalmente\\
aqueles que trabalham direta ou indiretamente nas atividades relacionadas 4 industria pe-\\
trolifera. Nesse contexto, este livro é introdutério e pretende capacitar os leitores com\\
conhecimentos basicos da cadeia de valor da industria do petréleo,\\
de upstream, downstream e midstream.\\
\\
E importante que os leitores tenham em mente que a cadeia produtiva do petréleo é\\
bastante complexa e envolve etapas com a aplicacao de técnicas e conhecimentos diferen-\\
Giados, exigindo uma formacio especializada. Este livro nao tem a intengao de abordar\\
em profundidade os processos complexos da industria do petréleo, mas sim fornecer uma\\
visio geral da cadeia de valor da industria do petréleo. Para obter conhecimentos espe-\\
cializados, os leitores podem consultar livros dedicados a cada capitulo abordado neste\\
livro. O gas natural nao é amplamente explorado neste livro, por isso, incentivamos os\\
leitores a consultar o Volume 1 e 2 do livro “Industria do Gas Natural”, do mesmo autor.\\
\\
Este livro é composto por nove capitulos, além da introdugio. O Capitulo 1 aborda\\
\\
conceitos gerais da industria do petréleo, enquanto o Capitulo 2 descreve a cadeia pro-\\
dutiva da industria petrolifera, bem como alguns contratos na exploracao e produgio.\\
Os contetidos relacionados ao segmento de upstream sao abordados do Capitulo 3 ao 7,\\
incluindo exploragao de hidrocarbonetos, perfuragao e completacao de pocos, avaliacio\\
de formacées, reservat6rios, producao e operac6es submarinas. O Capitulo 8 aborda as-\\
suntos relacionados ao segmento de midstream e downstream. O ultimo capitulo do livro\\
apresenta contetidos relacionados 4 seguranca, meio ambiente e satide na industria do\\
petréleo, incluindo tépicos como principais impactos ambientais, acidentes registrados e\\
fontes alternativas de energia.\\
\\
Este livro pode ser utilizado em cursos de introducio a industria do petréleo, assim\\
como em cursos modulares para capacitar os participantes interessados em adquirir co-\\
nhecimentos sobre a industria petrolifera. Os materiais deste livro so baseados em varias\\
fontes referenciadas ao longo do texto. Fizemos o possivel para dar os devidos créditos\\
4s fontes utilizadas, aos quais somos imensamente gratos, e pedimos desculpas por quais\\
quer omiss6es. Como se trata de um livro introdutério, evitamos fazer referencias cruza\\
das, e também gostariamos de apresentar nossas sinceras desculpas por isso.\\
\\
Qu seja, os segmentos\\
\\
xix\\
\\
—s\\
Digitalizada com CamScanner\\

\\section*{Pagina 18}

xx PREFACIO\\
" laboracao dos lei\\
ib este livroe contamos coma = '€20 dos leito,\\
apace nee aren utor espera que este livro seja uma fonte util para a\\
por meio de criticas\\
\\
eitores em seus estudos.\\
\\
Geraldo Ramos, Ph.D,\\
\\
Digitalizada com CamScanner\\

\\section*{Pagina 19}

AGRADECIMENTOS\\
\\
a\\
\\
Agradeco especialmente ao meu Deus, que me concedeu sabedoria, inteligéncia, vontade\\
e satide para tornar esta obra uma realidade,\\
\\
Agradeco 4 minha querida esposa e aos meus filhos, por seu encorajamento constante\\
para que eu seguisse em frente com a escrita, dedicando seu tempo de lazer para apoiar-\\
me totalmente neste trabalho. Também agradeco aos parentes e amigos que oraram por\\
mim e me deram suporte durante a elaboracao deste livro.\\
\\
Expresso meu agradecimento especial ao Departamento de Responsabilidade Social\\
da Sonangol, cujo patrocinio viabilizou esta obra e proporcionou maior espaco para dar\\
voz aos académicos dispostos a compartilhar conhecimentos, tao necessarios em Angola.\\
\\
Gostaria também de agradecer a todos os autores referenciados neste livro, cujos tra-\\
\\
balhos foram fonte de pesquisa para a composi¢\\&io desta obra. Desde ja, peco desculpas\\
sinceras por eventuais omissGes.\\
\\
xxi\\
\\
Digitalizada com CamScanner\\

\\section*{Pagina 20}

INTRODUCAO\\
\\
A procura, descoberta e extrac¢io do petrdleo e¢ gds natural, assim como a sua transferén-\\
cia para as refinarias e a transformacao dos derivados até ao consumidor final, exigem\\
a intervengio de muitos especialistas, respeitando uma ordem cronolégica. Essa ordem\\
cronolégica ¢ sistematica de desenvolvimento de atividades, objetivando suprir um mer-\\
cado consumidor final com os produtos do sistema, chamamos de cadeia de valor ou\\
produtiva da industria do petrédleo,\\
\\
Embora o petréleo seja um recurso mineral nao-renovavel, é um produto estratégico\\
ea principal fonte de energia ¢ aquisigao de receitas de Angola, em particular, e do mundo\\
em geral, Nesse contexto, os paises nfo industrializados e dependentes do petréleo so-\\
frem graves consequéncias econémicas e politicas com a reducio abrupta do prego do\\
barril de petréleo no mercado.\\
\\
Tendo em conta a sua importancia para a economia de Angola, é crucial a dissemina-\\
gio do conhecimento do petréleo e gés natural para os angolanos, de forma a capacita-los\\
nos tépicos basicos dos diferentes segmentos da cadeia de valor da industria petrolifera.\\
\\
A cadeia de valor da industria petrolifera esta dividida em trés grandes segmentos prin-\\
cipais conhecidos como upstream, midstream e downstream. A atividade relacionada com\\
a explora¢io, desenvolvimento e producio de petréleo e gas natural € chamada de ups-\\
tream. Ba fase de maiores riscos de investimentos, que é compensada pela possibilidade\\
de obter elevados lucros.\\
\\
Geralmente, as companhias petroliferas criam parcerias para compartilhar os riscos\\
€, consequentemente, os lucros. Os grandes players deste segmento em Angola sao a\\
Sonangol, E.P., empresa nacional de Petréleo, conhecida como NOC (do inglés, National\\
Oil Company). As empresas internacionais de petréleo, conhecidas como IOC (do inglés,\\
International Oil Company), sio a TOTAL EP, Chevron, ESSO Angola, BP Angola e ENI.\\
\\
As operacées so realizadas com o suporte de empresas de prestacao de servicos, onde\\
destacamos Halliburton e Schlumberger. Além dessas duas empresas, existem outras\\
como Baker, Cameron, GE, Seadrill, Transocean, Sonasurf, Sonamet, Aker, FMC, Ma-\\
ersk, Saipem, Hercules, Weatherford, Operatec, Argonauta e Fugro, entre outras.\\
\\
A primeira licenga de exploragao, ou seja, 0 inicio das atividades de exploracéo em\\
Angola, teve lugar em 1910 na Bacia de Cuanza e na regiao maritima do Congo, pela\\
Companhia Canha \\& Formigal. O primeiro pogo comercial foi declarado em 1955 (poco\\
Benfica 2), na mesma bacia, pela Petrofina.\\
\\
O midstream compreende as atividades de refinago, transporte, importagao e expor-\\
tagdo de petréleo e gis natural, Essa fase da cadeia de valor destaca-se por ser uma das\\
atividades menos lucrativas e com a possibilidade de obtencao de grandes prejuizos.\\
\\
xxiii\\
\\
Digitalizada com CamScanner\\

\\section*{Pagina 21}

xxiv INTRODUCAO\\
\\
cionado com as atividades de logistica, distribuicag «\\
\\
comerdializagao é conhecido como downstream. Eo nia ati MENOS TISCOS € com\\
maior margem de ganhos, assim como um vasto mercado consumi loraser explorado, As\\
atividades de refinagdo em Angola tiveram inicio em 1958, com a inaugura¢ao da Tefinaria\\
de Luanda, cuja construcao teve inicio em 1956. .\\
\\
Nesse contexto, sendo o petréleo 0 principal recurso de exportacao de Angola, é im.\\
portante apresentarmos este livro com a finalidade de familiarizar 0 leitor, mesmo nig\\
especializado, com a linguagem e técnicas da industria do petrdleo. Esta obra é transyer.\\
sal para todos os estratos sociais.\\
\\
Os leitores terao a possibilidade de adquirir conhecimentos basicos da industria do pe-\\
tréleo e poderao consultar a bibliografia referenciada para aprofundar os conhecimentos\\
adquiridos, A esses, recomendamos a aquisi¢ao de livros com estudos mais aprofundados\\
de cada capitulo, que de forma geral é uma disciplina ou livro de especializacao,\\
\\
\\& Este , éo resultado de materiais de um curso modular da introdugao da industria\\
\\
patron experiéncia na industria do petréleo e material disponivel no dominio Publico\\
\\
hi ae neste ba Esperamos que seja uma fonte de aprendizado para todos os\\
que possam adquirir esta obra.\\
\\
O segmento da cadeia de valor rela\\
\\
Digitalizada com CamScanner\\

\\section*{Pagina 22}

CAPITULO 1\\
\\
CONCEITOS GERAIS\\
\\
See\\
\\
Neste capitulo, apresentamos de forma resumida os conceitos basicos da Engenharia de\\
Petréleos e as fungGes do engenheiro nesse sector. O principal recurso, que é 0 foco desse\\
profissional, também foi abordado de maneira concisa. O objetivo é familiariza-lo com\\
esses conceitos transversais a todos os outros capitulos e fornecer ferramentas para uma\\
melhor compreensao dos tépicos abordados nos capitulos subsequentes.\\
\\
1.1 O que é Engenharia de Petréleos?\\
\\
A Engenharia de Petrdleos é uma disciplina de engenharia relacionada As atividades de\\
produgio de hidrocarbonetos, como petréleo bruto e gis natural. Ela é considerada um\\
setor de upstream (montante) da industria do petréleo, envolvendo as atividades de explo-\\
racio e producio de petréleo e gas natural.\\
\\
1.1.1 O que é um engenheiro de petrdleos?\\
\\
Um engenheiro de petréleos est4 envolvido em quase todas as etapas da avaliacio, de-\\
senvolvimento e producao de campos de petréleo e gas. O objetivo desse profissional é\\
maximizar a recuperacao de hidrocarbonetos com 0 minimo custo, com uma forte énfase\\
na redugio dos impactos ambientais associados. Embora a Engenharia de Petrdleos seja\\
\\
marcada por sua caracteristica multidisciplinar, ela pode ser dividida em trés areas basicas\\
de atuagao:\\
\\
Engenheiros de reservatério: os profissionais dessa area trabalham para otimizar a pro-\\
ducio de petréleo e gas natural por meio do posicionamento adequado de pogos,\\
niveis de producio e técnicas avangadas de recuperacio de petréleo.\\
\\
Engenheiria de perfuragio e completagio: os profissionais dessa area gerenciam os as\\
Pectos técnicos da perfuracao e conclusio de pogos exploratorios, de produgio e\\
de injegdo. Isso também inclui engenheiros de lama, que gerenciam a qualidade do\\
fluido de perfuracao.\\
\\
Indistrnia do\\
\\
1\\
Petréleo para ndo-engenheiros, 1* elisa.\\
By Genias nt Le Ee\\
\\
André Raposo Ramos Copyright © 2025 Bditora Dois03\\
\\
—\\
Digitalizada com CamScanner\\

\\section*{Pagina 23}

2 —\\_ CONCEITOS GERAIS\\
iros de subsuperficie: 08 profissionais dessa\\
\\
o reservatorio € 0 Poses incluindo perfuragées, con,\\
lo do pogo \\& equipamentos de monitora.\\
étodos de elevagio artificial ¢ selecionam\\
\\
© equipamento de superficie que Sep: fluidos produzidos (leo, gas ¢ agua),\\
+ ou subdividir a engenharia de petréleos em nove\\
\\
engenharia de perfuracao, engenharia de com.\\
\\
do petrdleo, tecnologia offshore, transporte de\\
as e distribuigao de combustiveis e outros\\
\\
Jugio e engenhe!\\
a interface entre\\
\\
controle de fluxo no fund\\
avaliam m\\
\\
ara OS\\
\\
Engenheiros de prod\\
area gerenciam\\
trole de areia, ct\\
mento do fundo do pogo. Bles\\
\\
Algumas literaturas tentam amplia\\
areas basicas: engenharia de reservatorios,\\
pletacio, processo de produgao, economia\\
petrdleo, processamento de petroleo em refinari\\
\\
derivados.\\
\\
1.1.2 Disciplinas auxiliares\\
roducao de petrdleo € gas natural requer 0 auxilio de varias\\
paleontologia, geologia, geofisica, engenharia\\
oceanografia, contabilidade e meio am-\\
descrever algumas dessas disciplinas\\
\\
Aatividade de exploracao € P!\\
freas da ciéncia, incluindo geoquimica,\\
quimica, fisica, engenharia naval, meteorologia,\\
biente, devido a complexidade da atividade. Vamos\\
\\
auxiliares:\\
* Gedlogos de petréleo encontra:\\
neas com métodos geoldgicos e geofisicos.\\
\\
m hidrocarbonetos analisando estruturas subterra-\\
\\
. Geofisicos de petréleo aplicam um conjunto de ferramentas investigativas indispen-\\
siveis para a prospeccao de dleo e gas, seja em qualquer nivel de detalhe ou econo-\\
micamente viavel. Esse conjunto de ferramentas é popul i\\
Seas: popularmente conhecido como\\
\\
. parsers 0 estudo das propriedades fisicas das rochas relacionadas 4 distribui¢ao\\
rm luidos ‘ea Bey espagos porosos. Atualmente, o conjunto de disciplinas que trata\\
propriedades fisicas das rochas e minerais é chamado de Petrofisica.\\
\\
* Petroquimica é a ciéncia e a técni y\\
sie cnica que utilizam o\\
\\
matéria- 3 petréleo e o gas natu:\\
ia-prima para o desenvolvimento de diversos produtos qui - al como\\
\\
1.1.3 Campos de trabalho\\
\\
O engenheiro de\\
\\
critérios, ee cad em diversos campos de atividade, incluindo es\\
(onshore), industrias cstiatednids ae (offshore), plataformas terrestres\\
res desse profissional podem ser 0 eens \\_ Pesquisa € universidades, Os empregado-\\
empresas de apoio, instituicdes académicas, mon siprusas de nerriges\\
\\
1.1.4 C\\
‘ompeténcias de um engenheiro de petrdle\\
os\\
\\
Projeto: i .\\
jetos eficientes que visem A exploragio ¢ 4\\
\\
produgao do reservatéri\\
rid sem riscos e :\\
Prejuizo ao meio ambiente. Além disso, ele cuida\\
\\
4\\
Digitalizada com CamScanner\\

\\section*{Pagina 24}

© QUE E 0 PETROLEG? 3\\
\\
do transporte do petréleo seus derivados,\\
\\
desde o |\\
Com habilidades nas areas de local\\
\\
engenharia, geofisica,\\
i berta de Teservatérios\\
tes de realizar perfuragSes nesses reservatérios.\\
\\
de exploracgao até a refinaria.\\
mineracao e geologia, o engenheiro\\
de petréleo e busca maneiras eficien-\\
\\
1.2 O que 6 0 petréleo?\\
\\
O petréleo é um dleo formado pela decom\\
por bactérias. Ele € encontrado no fundo\\
coberto por sedimentos. Ao longo de um\\
carbonetos, compostos de hidrogénio e c:\\
\\
posicao de matéria Organica e mineral atacada\\
dos mares e lagoas em terrenos sedimentares,\\
periodo, essa matéria se transforma em hidro-\\
\\
‘ c at ‘arbono. Quando as condigGes geolégicas sao\\
favordveis, 0 dleo fica aprisionado entre camadas de rochas impermedveis, formando re-\\
\\
servatérios de dleo e gas natural. Nas jazidas, a camada de éleo é coberta por uma camada\\
de gas ¢ flutua sobre uma camada de Agua salgada (Figura 1.1).\\
\\
Figura 1.1: Configuracao tipica de uma jazida de petréleo.\\
\\
A extracao de petréleo € realizada por meio da perfuracio de pocos, nos quais o petré-\\
leo é liberado devido  pressio dos fluidos presentes no reservatério. Em alguns casos,\\
quando a pressao natural nao é suficiente, o petréleo precisa ser bombeado para a super-\\
\\
ficie (Figura 1.2).\\
\\
Figura 1.2: Esquema de um poco de extracao.\\
\\
Durante o processo de extracao do petréleo bruto, é realizado 0 processamento ou trata\\
mento primario para a separacao gas-leo-Agua livre (1° etapa do processamento prims-\\
tio) e desidratagao do éleo (2 etapa do processamento primirio), para eliminar a aguae\\
os sais do petréleo processado em baterias, permitindo alcangar os niveis de Basic Sed iment\\
and Water/ Sedimento basico e gua (BS\\&W) que aumentam o valor de comercializagio\\
do fluido mitigam problemas relacionados a seguranga (corrosao ou explosio). lngeste\\
realcar que a BS\\&W \\& medida pela quantidade de agua ¢ sedimentos na amostra do dleo.\\
\\
Digitalizada com CamScanner\\

\\section*{Pagina 25}

4 CONCEITOS GERAIS\\
aria, onde passa por transfoy\\
\\
Uma vez extraido, © petrdleo é transportado e ¢ terrestre, S40 utilizade,\\
Para © ‘\\
magodes que resultam em seus subprodutos. werge™ paraum ponto especialmente\\
©\\
\\
inhi 5 c ; “\\
oleodutos (pipelines), camntnsbes* EE tee troleo, conhecido como terminal petrol;\\
preparado para 0 transporte ¢ comércio do pe bém chamado de petroleiro oy,\\
\\
i isterma,\\
fero, onde ele é embarcado em um navio cis'\\
\\
navio-tanque.\\
ode\\
O proceso de refino p ay UPGN tem CO\\
\\
Gas Natural (UPGN)e refinari . Al * aplicaco\\
dutos especificados para atender diversa P. ito de petro)\\
\\
«20, Gas Liquefeit\\
gasolina, gas combustivel, reinjegao, Gas Liquele\\
petréleo bruto é tratado nas\\
\\
corrente, como gas liquefeito,\\
produzidos tecidos na industria petro!\\
\\
para@ refin\\
anspor\\
\\
tam\\
\\
iras; Unidade de Processamento do\\
rrer de duas 1 50 basica fracionar o gis em pro.\\
das pelo mercado, incluindo\\
Jeo (GLP), entre outros,\\
diversos produtos de uso\\
asfalto, e também sig\\
\\
gasolina, 6!\\
quimica.\\
\\
1.2.1 Constituintes do petroleo bruto\\
\\
lo no subsolo é chamad\\
uimicos chama\\
quenas plantas € anit\\
\\
lo de petrdleo bruto. B composto por uma\\
dos hidrocarbonetos. Esses hidrocarbone-\\
mais se decompuseram sob camadas\\
\\
de areia ¢ lama. O petréleo bruto nem sempre tem a mesma aparéncia, isso depende de\\
sua origem. As vezes, ele é quase incolor, enquanto em outras situagOes pode ser espesso\\
€ preto. No entanto, geralmente 0 petréleo se assemelha a um melago fino e marrom.\\
Quando é extraido de um pogo, especialmente um pogo submarino, 0 petréleo bruto esta\\
\\
uentemente misturado com gases, Agua ¢ areia.\\
Opetrdleo é normalmente separado em frases conforme a faixa de ebuli¢ao dos com-\\
\\
postos. A Tabela 1.1 ilustra as fracdes tipicas obtidas do petrdleo.\\
\\
O petréleo encontrad\\
mistura de diferentes produtos q)\\
tos foram produzidos quando pe:\\
\\
Tabela 1.1: FragGes tipicas do petrdleo.\\
\\
‘Temperatura de ebu- Composi¢io aproxi- Usos\\
\\
Fracio\\
ligdo (°C) mada\\
Gis By\\
ot ps Cy -Cz Gis combustivel\\
40 Z\\
C3-Cq Gas combustivel engarrafado, uso domés-\\
tico ¢ industrial\\
Gasolina 40-175\\
ie ra Cio Combustivel de automével, solvente\\
ae are bere By “C12 Iuminagiio, combustivel de avides ¢ jato\\
vanes Saban zn 3°Ci7 Diesel, fornos\\
* -¢;\\
Gastleo pesado 18 C25 Combustivel, matéria-prima para lubrifi-\\
cantes\\
Lubnficantes 400-510 C26 - Cy\\
ibis ie e 8 Oleos lubrificantes\\
— Asfalto, picl\\
\\
O petréleo bruto\\
nao é composto a\\
a nas ;\\
tes conhecidos como pit iin arias eta mas também contém con\\
netos, Esses contaminantes podem incluir\\
\\
Digitalizada com CamScanner\\

\\section*{Pagina 26}

O QUE £ O PETROLEO? 5\\
\\
Tabela 1.2: Contaminantes ou nio-hidrocarbonetos.\\
— SSS\\
\\
Contaminante Elemento quimico presente Prejuizo\\
\\
Compostos onginicos sulfurados——Enxofie (\\$) Corrosio, toxicidade e poluigio.\\
\\
Compostos orginicos nitrogenados \\_ Nitrogénio (N) Retengio de Agua emulsionada, contami\\
nagio de catalisadores, alteragio de colo\\
racio de produtos finais.\\
\\
Compostos orginicos oxigenados —\\_Oxigénio (0) Acidez, corrosividade, formagio de go-\\
mas, odor,\\
Compostos orginicos metilicos. —\\_\\_Metais (principalmente Ni e V) Agressio a materiais, contaminagio de ca:\\
\\
talisadores.\\
\\
1.2.2 Hidrocarbonetos\\
\\
O petréleo bruto é uma mistura de hidrocarbonetos. Sao frequentemente cadeias de\\
atomos de carbono com hidrogénios ligados. As cadeias mais longas tém pontos de ebu-\\
licdo mais altos, o que permite sua separago por destilagao. Os grupos mais simples so\\
os alcanos e alcenos. Eles terminam com “ano” e “eno”, respectivamente. O primeiro\\
trecho do nome depende do ntimero de atomos de carbono. Os hidrocarbonetos estao\\
resumidos na Figura 1.10.\\
\\
Alcanos: sao hidrocarbonetos alifaticos saturados, isto é, apresentam cadeia aberta com\\
simples ligag6es apenas. Férmula geral: C,,H2,2. Exemplo: Butano: C4Hjo (Fi-\\
gura 1.3). Na Tabela 1.3 ilustramos alguns exemplos de Alcanos.\\
\\
H oH H oH\\
ee d\\
H—C.—¢—c—c¢—H\\
Ee ae I\\
HH H H\\
\\
Figura 1.3: Exemplo de butano.\\
\\
Alcenos ou olefinas: sao hidrocarbonetos alifaticos nao saturados que apresentam dupla\\
ligagao. Formula geral: C,,H,,. Exemplo: Bteno: C)H4 (Figura 1.4). Na Tabela 1.3,\\
ilustramos alguns exemplos de hidrocarbonetos do tipo Alcenos.\\
\\
Hy at\\
i” Ni\\
\\
Figura 1.4: Exemplo de eteno.\\
Alcinos: sio hidrocarbonetos alifaticos ndo saturados por uma tripla ligaco. Formula\\
\\
geral: C,H2,\\_9. Exemplo: Etino: C,H também conhecido como Acetileno (Fi-\\
Sura 1.5),\\
\\
—\\
Digitalizada com CamScanner\\

\\section*{Pagina 27}

6 —\\_CONCEITOS GERAIS\\
H-C= C—H\\
\\
Figura 1.5: Exemplo de etino.\\
\\
os nao saturados por duas ligacGes duplas, p bn\\
\\
cadienos: sio hidrocarbonetos alifatic ,\\
Alcadienos: sao hidroca' ho eM eee.\\
\\
mula geral: C,, H2n-2- Exemplo: Bu\\
\\
HoH\\
H\\
APS\\
a? SC Ser\\
1\\
HoH\\
\\
Figura 1.6: Estrutura de Butadieno.\\
\\
Cicloalcanos: apresentam cadeia fechada com apenas simples ligacdes. Formula geral:\\
C,H. Exemplo: Ciclohexano: CoHj2 (Figura 1.7).\\
\\
CH\\
\\
yr,\\
\\
H.C CH,\\
i |\\
\\
H.C CH,\\
KF\\
\\
CH,\\
\\
Figura 1.7: Estrutura de Ciclohexano.\\
\\
Grnskean Sto hidrocarbonetos ciclicos nao saturados por uma dupla ligagao. For\\
mula geral: C,,H2,\\_2. Exemplo: Ciclobuteno: C4Hg (Figura 1.8).\\
\\
CH,\\
\\
fo\\
H.C CH\\
\\
Figura 1.8: Estrutura de Ciclobuteno.\\
\\
y\\
\\
Digitalizada com CamScanner\\

\\section*{Pagina 28}

OQUEEOPETROLEO? = 7\\
Aromaticos ou Arenos: Sio hidrocarbonetos em cuja estrutura existe pelo menos um\\
\\
anel benzénico (aromatico), Formula geral: C,H). Exemplo: Anel Benzeno: CH,\\
(Figura 1.9), 6fl6\\
\\
Ss\\
A\\
\\
Figura 1,9: Estrutura de Anel Benzeno.\\
\\
Os hidrocarbonetos podem apresentar-se no estado Mquido (éleo bruto ou condensado),\\
gasoso ou fragGes leves (gas natural), sdlidos ou fragSes pesadas (asfalto e betume). O\\
gas pode ser classificado em gis associado (ao petréleo) ¢ gas nfio associado (com ou sem\\
fragGes liquidas). Ambos podem ser usados nas operagées como combustivel, reinjetado\\
para recuperacao secundaria, terciéria ou armazenagem, queimado (falta de Opsao) ou\\
\\
para comercializagio. Na Figura 1.10 e na Tabela 1.3, resumimos alguns tipos de hidro-\\
carbonetos abordados deste capitulo.\\
\\
Figura 1.10: Resumo de tipos dos hidrocarbonetos.\\
\\
‘Tabela 1.3: Exemplo de alguns hidrocarbonetos do tipo Alcanos e Alcenos.\\
\\
Hidrocarboneto Férmula Simbolo\\
\\
Alcanos Metano CHy Ci\\
Etano CoH, C2\\
Propano C3Hg C3\\
Eteno CoHy Co\\
Alcenos i\\
Buteno CyHy Cy\\
\\
Digitalizada com CamScanner\\

\\section*{Pagina 29}

B \\_CONCEITOS GERAIS\\
\\
le Petrdleo\\
, ‘ 50 classificados pelo American Petroleum\\
idade (gravity), OS éleos sa\\
Dependendo de sua densi raus er densidade na escala de API), sendo que 0s com maior\\
Institute (API) em varios opto, th see a a\\
\\
. A 3 melhores. Com =\\
\\
graduacao Pe Ate gma ee equagio 1.1 foi criada pelo API e serve como pari.\\
\\
pesado e um de am de petrbleas derivados. A mesma equagao pode ser representada\\
at\\
\\
metro de classifica¢:\\
na equacao 1.2.\\
\\
1.2.3 Classificagao d\\
\\
141,5 ;\\
Yo= 731,54 API a)\\
131, 5+ API\\
ou\\
apie SABE = 18i,5 (1.2)\\
\\
0\\
Onde yp é a densidade relativa medida em condi¢oes padrao. Agua (agua doce), por\\
definicao, tem densidade de 1,0. Isso implica que usando a equacao 1.2, a API desta agua\\
€ 10. As condicdes padrao de temperatura ¢ pressio referenciadas sao 60°F (15,56°C) e\\
14,696 psi (101,325 kPa). Alguns fatores podem afetar o grau API dos éleos, tais como:\\
\\
A idade geoldgica: as rochas antigas tendem a ter maior graduagao, mas as rochas ter-\\
Giarias podem ter cerca de 40n° API, como as do Mar do Norte;\\
\\
Profundidade do reservatério: quanto maior a profundidade, maior a graduacao;\\
\\
Tectonismo: altas graduacGes sio mais comuns em regides com muitas tenses nas ca-\\
madas geolégicas;\\
\\
Salinidade: os reservatorios de origem marinha tendem a ter maiores graduacées do que\\
os de origem de ambientes com Agua salobra ou fresca;\\
\\
Teor de enxofre: este teor é alto em éleos de baixa graduacio.\\
\\
\\& 5 i i sons ae rar\\
Tibelaasa: do petréleo em Grau API pelas diferentes instituicdes esta ilustrada na\\
\\
Tabela 1.4: Classificagio do petréleo em Grau API.\\
eee\\
Orgio API (Grau API)\\
\\
Oleo Leve Oleo Médio Oleo Pesado Oleo Ultrapesado\\
Governo de Alberta\\
rs / Canada 234 25-34 10-25 S10\\
Departamento de Energia dos EUA 2 35,1 25.35.) <\\
OPEP :\\
\\
10-25 <10\\
32 26-32 10,5 - 26 \\$105\\
Petrobras offshore hs 19-32 14-19 si\\
— amid ne 18-32 13-18 su\\
\\
Digitalizada com CamScanner\\

\\section*{Pagina 30}

© QUE EO PETROLEO? 9\\
\\
A classificagao do petréleo, de acordo com seus constituintes, é de interesse tanto para ge-\\
oquimicos quanto para refinadores. Os primeiros visam caracterizar 0 6leo para relaciona-\\
lo a rocha mae e medir seu grau de degradacao. Os refinadores querem saber a quanti-\\
\\
dade das diversas fragdes que podem ser obtidas, assim como sua composi¢ao e proprie-\\
dades fisicas.\\
\\
Assim, os 6leos parafinicos s4o excelentes para a producao de Querosene de Aviacao\\
(QAV), diesel, lubrificantes e parafinas. Os dleos nafténicos produzem frac6es significati-\\
vas de gasolina, nafta petroquimica, QAV e lubrificantes, enquanto os éleos aromaticos\\
so mais indicados para a produgao de gasolina, solventes e asfaltos.\\
\\
Classe parafinica (75\\% ou mais de parafinas): nesta classe esto os éleos leves, fluidos\\
ou de alto ponto de fluidez, com densidade inferior a 0,85, teor de resinas e asfaltenos\\
menor que 10\\% e viscosidade baixa, exceto nos casos de elevado teor de n-parafinas\\
com alto peso molecular (alto ponto de fluidez).\\
\\
Os aromiaticos presentes so dos anéis simples ou duplos e 0 teor de enxofre é baixo.\\
Este tipo de petréleo produz subprodutos com as seguintes propriedades: (1) Ga-\\
solina de baixo indice de octanagem; (2) Querosene de alta qualidade; (3) Oleo die-\\
sel com boas caracteristicas de combustio; (4) Oleos lubrificantes de alto indice de\\
viscosidade, elevada estabilidade quimica e alto ponto de fluidez; (5) Residuos de\\
refinacdo com elevada porcentagem de parafina.\\
\\
Classe parafinico-nafténica (50-70\\% parafinas, >20\\% de nafténicos): os dleos desta clas-\\
se apresentam um teor de resinas e asfaltenos entre 5 e 15\\%, baixo teor de enxo-\\
fre (menos de 1\\%), teor de nafténicos entre 25 e 40\\%. A densidade e viscosidade\\
apresentam valores maiores do que os parafinicos, mas ainda sao moderados.\\
\\
Classe nafténica (>70\\% de nafténicos): nesta classe enquadra-se um niimero muito pe-\\
queno de dleos. Apresentam baixo teor de enxofre e se originam da alteraco bi-\\
oquimica de dleos parafinicos e parafinico-nafténicos. Alguns éleos da América do\\
Sul, da Russia e do Mar do Norte pertencem a esta classe.\\
\\
O petréleo do tipo nafténico produz subprodutos com as seguintes propriedades\\
principais: (1) Gasolina de alto indice de octanagem; (2) Oleos lubrificantes de baixo\\
residuo de carbono; (3) Residuos asfalticos na refinagao.\\
\\
Classe aromatica intermediaria (>50\\% de hidrocarbonetos aromiaticos): sao dleos fre-\\
quentemente pesados, contendo de 10 a 30\\% de asfaltenos e resinas e teor de enxofre\\
acima de 1\\%. O teor de monociclicos é baixo e, em contrapartida, o teor de tiofe-\\
nos e dibenzotiofenos é elevado. A densidade usualmente é maior que 0,85. Alguns\\
6leos do Oriente Médio, Africa Ocidental, Venezuela, California e Mediterraneo (Si-\\
cilia, Espanha e Grécia) so desta classe.\\
\\
Classe aromatico-nafténica (>35\\% de nafténicos): os dleos deste grupo sofreram pro-\\
cesso inicial de biodegradacao, no qual foram removidas as parafinas. Eles sao de-\\
tivados dos éleos parafinicos e parafinico-nafténicos, podendo conter mais de 25\\%\\
de resinas e asfaltenos, ¢ teor de enxofre entre 0,4 e 1\\%. Alguns dleos da Africa\\
Ocidental sio desse tipo.\\
\\
Digitalizada com CamScanner\\

\\section*{Pagina 31}

10 — CONCEITOS GERAIS:\\
\\
\\$ eos \\$40 Oriu\\
\\
Classe aromatico-asfaltica (>35\\% de asfaltenos ¢ eons pnt — rn omni\\
um processo de biodegradagio avangada, sa qua’ ites GAG vardedettoonane\\
alcenos ¢ oxidagio. Podem também incluir rig oy cider, Ainweneren, ant\\
aromaticos nio degradados da Venezuela e da Afri recone, tomtattton Uc wees\\
classe compreende principalmente éleos pesados € ents uAieecne\\
clo dos éleos aromaticos intermedirios. Desta Le vii vir Gh 180\\% ee =\\textasciitilde{}\\
é elevado, havendo equilibrio entre ambos. Oteor Seo uid Otidenta Vonacasy\\
sos extremos. Nessa classe, encontram-se 0S dleos do\\
\\
¢ Sul da Franga. .\\
contetidos sobre a classificacdo do petréleo que nao\\
\\
No capitulo 8, vamos abordar alguns\\
foram mencionados neste capitulo.\\
\\
1.3 O primeiro poco de petrdleo\\
\\
A moderna indastria petrolifera remonta cerca de 150 anos. O primeiro poco it dies\\
do mundo foi perfurado em Titusville, Pensilvania, em 1859. Ele atingiu petréleo a 21\\
metros abaixo do solo e produziu 3.000 litros de dleo por dia. Conhecido como 0 poco\\
Drake, em homenagem ao “Coronel” Edwin Drake, 0 homem responsavel pelo poco,\\
iniciou uma busca internacional por petréleo e, de varias maneiras, acabou mudando a\\
\\
maneira como vivemos.\\
\\
1.3.1 Quanto tempo durara o petréleo do mundo?\\
\\
O petréleo levou milhdes de anos para se formar e o suprimento de petréleo no solo nao\\
durard para sempre. Os campos de petrdleo j4 descobertos possuem mais de 1 bilhdo de\\
barris de petrdleo (1.000.000.000.000). Embora estejamos usando petréleo rapidamente,\\
as reservas aumentam a cada ano. Isso ocorre porque mais petrdéleo é descoberto e sio\\
encontradas novas maneiras de extrair 0 dleo que antes nao podia ser extraido. Mesmo\\
assim, nosso petréleo nao vai durar para sempre.\\
\\
Atualmente, o mundo consome cerca de 26 bilhdes de barris por ano. A esse ritmo,\\
deve haver petréleo suficiente por pelo menos mais 40 anos. provavel que mais petréleo\\
seja descoberto nesse periodo. As empresas de petréleo estao sempre procurando novos\\
campos de petrdleo ¢ ainda ha muito mais dreas do fundo do mar para explorar.\\
\\
1.3.1.1 Reservas mundiais de Petréleo\\
\\
fontes de petréleo do mundo sio:\\
\\
Venezuela: 0 pais latino-americano é 0 campeio em\\
, Teservas de petrdleo no mundo.\\
Com cerea de 303,2 bilhGes de barris, Tepresenta 17,9\\% da Participagdo mundial de\\
\\
petréleo,\\
\\
Digitalizada com CamScanner\\

\\section*{Pagina 32}

HISTORIA DO PETROLEO NO MUNDO 11\\
\\
Arabia Saudita: o pais do Oriente Médio esta e:\\
barris em suas reservas e Tepresenta 17,2\\%\\
maiores exportadores de petréleo, integra\\
tégica no setor de dleo e gis,\\
\\
™m segundo lugar, com 266,2 bilhdes de\\
das reservas mundiais totais. \\& um dos\\
a OPEP e tem grande importancia estra-\\
\\
Canada: as reservas canadenses POssuem cerca de\\
10\\% do total mundial. A maioria do petréleo\\
EUA e Europa.\\
\\
168,9 bilhdes de barris, representando\\
Produzido no pais é exportada para os\\
\\
Ira: 0 pais, localizado na Asia Ocidental, també:\\
lugar na lista dos maiores produtores, Es\\
bilhGes de barris em suas reservas, equival\\
\\
'm integra a OPEP e aparece em quarto\\
tima-se que o Ira Possua cerca de 157,2\\
lente a 9,3\\% das reservas mundiais.\\
Iraque: outro integrante da OPEP, o Traque apresenta 148,8 bilhdes de barris de petréleo\\
€ representa 8,8\\% da participacao mundial.\\
Ilustramos as reservas mundiais de petréleo por ti\\
gura 1.11.\\
\\
po de petréleo e por regides na Fi-\\
\\
"Ars Chensas eteme) » Gice Convencional + Gieo Pesado = Gieotitra Pexado\\
\\
(a) Reservas mundiais por tipo de petréleo. (b) Reservas mundiais por re ides.\\
\\
Figura 1.11: Reservas mundias de petréleo.\\
\\
1.4 Histéria do Petrédleo no mundo\\
\\
Existem relatos da existéncia e utilizagao do petréleo que remontam a antiguidade. Mui-\\
0s povos utilizavam-se dos vazamentos naturais e os registros dao conta da utilizagao na\\
Torre de Babel e na Arca de Noé, no embalsamento de mortos ilustres pelos egipcios, na\\
Pavimentacio de estradas pelos Incas, como aglutinante de tijolos pelos Sumérios, para\\
fins bélicos Por gregos e romanos, entre outros. a\\
O registro da Participacdo do petréleo na vida do homem remonta a tempos eaten\\
antiga Babil6nia (Figura 1.12), os tijolos eram assentados com 0 asfalto, e o betume\\
a largamente utilizado pelos fenicios na calafetacao de embarca¢ées. iaieeame\\
Desde o século XVI, 0 principal motivo das expansdes maritimas e das —\\_ le : be °\\
némicas Curopeias foi a busca do ouro. Reis, navegantes, soldados e wanna le on\\
gal, Espanha, Holanda e Inglaterra se lancaram na busca desse precioso mineral ao Fe\\
\\
Na\\
a\\
\\
——— mSca\\
\\
Digitalizada com CamScanner\\

\\section*{Pagina 33}

12 CONCEITOS GERAIS\\
\\
Figura 1.12: Mapa da antiga Babilonia.\\
\\
do mundo, No entanto, a partir do século XIX, outro tipo de “ouro” comegou a despertar\\
\\
a cobica humana.\\
\\
Por volta de 1847, 0 petréleo comegou a ser utilizado comercialmente quando um co-\\
merciante de Pittsburgh, na Pensilvania, EUA, comegou a engarrafar ¢ vender petréleo\\
proveniente de vazamentos naturais (6leo que emergia naturalmente do solo) como lu-\\
brificante. Cinco anos depois, em 1852, um quimico canadense descobriu que aquecendo\\
¢ destilando o petréleo era possivel produzir querosene, um Iiquido que poderia ser usado\\
\\
em limpadas.\\
Em 1859, o Coronel Drake descobriu petréleo em Titusville, Pensilvania, utilizando o\\
\\
método de perfuragao com percussao. Esse poco produzia cerca de 2 metros ciibicos de\\
6leo por dia (Figura 1.13).\\
\\
Figura 1.13: Coronel Drake a esquerda ¢ a direita 0 primeiro pogo.\\
\\
Em 1900, o americano a3,\\
auma profundidade rh ka utilizou o processo rotativo para encontrar leo\\
deodreddianienge tibet Seok, ca Nos anos seguintes, a perfuragao rotativa\\
Figura 1.14, le perfuracao Por percussdo, conforme se vé na\\
A década de 60 foi marcada pel: .\\
\\
consumo desenfreado. "ma abe opin de petréleo no mundo, o que estimulou 0\\
Médio ¢ gis na Unido Sovittica. Nessa déeada. Arig se Ae Pettdleo no Oriente\\
‘a década, Ardbia Saudita Kuwait, Ira, Iraque ¢\\
\\
economicamente vidveis exploragées Precos do petré\\
\\
no M petroleo, o que tornoU\\
\\
outras grandes descobertas em territé lar do Norte © no México, Houve também\\
Nas décadas de 80 ¢ 90, os linen do terceiro mundo, .\\
\\
criando um novo ciclo econémico para a indiisti, reduziram os custos de produgio:\\
\\
Ustria petrolifera. Em 1996, as reservas P\\&\\
\\
4\\
Digitalizada com CamScanner\\

\\section*{Pagina 34}

HISTORICO DO PETROLEO EM ANGOLA 13\\
\\
Figura 1.14: Engenheiro mecanico, Anthony Francis Lucas.\\
\\
troliferas mundiais provadas eram 60\\% maiores do que em 1980, e os custos médios de\\
prospeccao e producao cairam cerca de 60\\% nesse mesmo periodo.\\
\\
Ao longo do tempo, o petréleo se tornou uma fonte de energia predominante. Como\\
advento da petroquimica, centenas de novos compostos passaram a ser produzidos, como\\
plisticos, borrachas sintéticas, tintas, corantes, adesivos, solventes, detergentes, explosi-\\
vos, produtos farmacéuticos, cosméticos, entre outros. O petréleo, além de ser utilizado\\
como combustivel, tornou-se imprescindivel na vida moderna.\\
\\
1.5 Histérico do petréleo em Angola\\
\\
A atividade de prospecco e pesquisa de hidrocarbonetos em Angola teve inicio em 1910.\\
Nesse ano, foi concedida 4 Companhia Canha \\& Formigal uma area de 11.400 km? no\\
offshore do Congo e na Barra do Kwanza, sendo o primeiro poco perfurado em 1915.\\
Pesquisas Mineiras de Angola (PEMA) e a Sinclair dos EUA também estiveram envolvidas,\\
desde cedo, na atividade de prospeccao e pesquisa em Angola.\\
\\
Apés breve paragem, em 1952, reiniciou-se a atividade com a concessao 4 Purfina da\\
mesma drea, adicionada a sua extens4o na Plataforma Continental em 1955. Em 1962, foi\\
efetuado o primeiro levantamento sismico do offshore de Cabinda pela Cabinda Gulf Oil\\
Company e, em setembro desse ano, surgiu a primeira descoberta.\\
\\
Em 1767, o entio governador Francisco Sousa Coutinho retirou algumas amostras\\
de petréleo na regiio dos Libongos (Barra do Dande) e levou para Lisboa. Em 1910,\\
as autoridades portuguesas concederam 4 PEMA a primeira licenga para exploragéo em\\
terra nas bacias do Congo e Cuanza.\\
\\
Em 1915, foi perfurado o primeiro poco na Bacia do Cuanza, Em 1919, iniciou-se 0\\
envolvimento das companhias estrangeiras em Angola, como a Sinclair Oil dos BUA, mas\\
\\
apesar das boas prospeccdes, nao houve interesse comercial.\\
\\
Em 1954, a Gulf Oil formou a Cabinda Gulf Oil Company Limited (CABGOC), em Ca-\\
binda e, nos anos seguintes, chegou a Companhia Belga Petrofina, que formou uma subsi\\
didria chamada Purfina, a qual posteriormente se expandiu para Petrangol. Em 1955, foi\\
feita a primeira descoberta de petréleo no campo Benfica, seguida por Luanda, Cacuaco,\\
Galinda e Tobias pela Petrofina,\\
\\
|\\
\\
Digitalizada com CamScanner\\

\\section*{Pagina 35}

14 CONCEITOS GERAIS\\
em Luanda (Figura 1.15). Em 1966, foi con,\\
Ba ustiveis SARL (ANGOL), que se AS5OCio4\\
\\
Em 1957, a Petrangol construiu a refina\\
odutos. Em 1962, iniciou-se a exploraci,,\\
\\
dade de Lubrificantes € Comb\\
\\
tituida a Socie\\
ir os seus pre\\
\\
4 Petrangol para refinar e distribu\\
offshore\\
\\
Figura 1.15: Refinaria de Luanda.\\
\\
1968: Inicia-se a primeira producao em offshore, (4guas rasas — Bloco 0) operado pela\\
\\
Chevron.\\
\\
1970; Outras companhias americanas, como Esso, Sun Oil, Amoco € Amerada Hess, as-\\
sinam um acordo com Angola para o terminal de Malongo, Cabinda, do Liquefied\\
Petroleum Gas/ Gas Natural Liquefeito (LPG). A Total também chega a Angola em\\
1975, na data da independéncia.\\
\\
1976: Criacio do Ministério dos Petréleos (MINPET), que atualmente é 0 Ministério dos\\
Recursos Minerais e Petréleos (MIRENPET). Nesse mesmo ano, a SONANGOL é\\
criada a partir da nacionalizacao de ANGOL Sociedade, por meio do Decreto-Lein”\\
52/76, de 9 de junho. ASONANGOL passa a coordenar, regular e controlar todas\\
as atividades relacionadas a indistria petrolifera em Angola. Além disso, torna-se a\\
concessionéria exclusiva para as operacées de exploracao e produgio de hidrocarbo-\\
\\
netos no pais. Sao estabelecidas joint ventures e contratos, como concessées, partilhas\\
¢ contratos de risco. ,\\
\\
OP eats producao de LPG no terminal de Malongo, Cabinda (Figura 1.16). 1991:\\
io da exploracao em Aguas profundas nos Blocos 14, 15, 16, 17, 18 e 20.\\
\\
1996: Descoberta do Gi itai\\
J cara “mon ' irassol pela Elf Aquitaine Petroleum nas aguas profundas\\
1999: Instalacdo da primeira\\
oe x primeira unidade de FPSO (Kuito) no offshore angolano, no Bloco 14.\\
? Inicio da produgio em Aguas Profundas no Bloco 14 pela Chevron.\\
\\
2001: Inicio da producdo no cam G\\
ir P\\
utilizada foi a maior do Reade ‘assol, no Bloco 17, Nesse ano, a plataforma FPSO\\
\\
2005: A producio de bleo de\\
a la ultrapassa a marca de 1.000.000 de barris pof\\
\\
]\\
\\
Digitalizada com CamScanner\\

\\section*{Pagina 36}

HISTORICO DO PETROLEO EM ANGOLA 15\\
\\
Figura 1.16: Sanha FPSO,\\
\\
2007: Angola torna-se membro da OPEP,\\
2010: Comissionamento do primeiro FPSO, Xicomba.\\
2011: Realizagao de um leilao de 11 blocos de pré-sal em Angola.\\
\\
2012: Inicio da producao em Aguas profundas no Bloco 31 pela BP.\\
\\
2012: Descoberta do primeiro poco de pré.\\
\\
-sal nas Aguas profundas, Pogo Azul 1, no Bloco\\
23 pela Maersk.\\
\\
2013: Primeira producdo de GNL na planta de ALNG, com a primeira venda para o Bra-\\
sil.\\
2013: A empresa Cobalt anuncia descobertas si;\\
\\
ignificativas de petréleo e gas no pré-sal\\
nos blocos 20 € 21.\\
\\
2016: Inicia-se o processo de reestrutura¢ao da Sonangol, com a nomeagao de um novo\\
Conselho de Administracao.\\
\\
2017: O Presidente da Republica cria um grupo de trabalho para a reestruturagio do\\
setor petrolifero, por meio do Decreto n° 290/17.\\
\\
2018: Lancamento do Programa de Regeneracao da Sonangol E.P.\\
\\
2019: Criacao da Agéncia Nacional de Petréleo, Gas e Biocombustiveis (ANPG) e aprova-\\
Gao do seu Estatuto Organico pelo Decreto Presidencial n° 49/19, de 6 de fevereiro.\\
A ANPG tem como objetivo regular, fiscalizar e promover a execucio das atividades\\
petroliferas, incluindo operacées e contratagées no setor de petroleo, gas e biocom-\\
\\
bustiveis.\\
\\
2020: Criagdo da Macroestrutura da Sonangol no ambito do Programa de Regeneracio\\
em curso na Sonangol, E.P. A nova estrutura tem como objetivo tornar a empresa\\
mais 4gil ¢ eficiente. Foram estabelecidas sete areas de apoio, quatro comités, cinco\\
unidades de negécios e a Sonangol Holdings, que engloba os negécios nio nucleares\\
da empresa. A privatizacio desses ativos esta prevista de acordo com o Decreto\\
Presidencial n° 250/19, de 5 de agosto.\\
\\
2022: Constituig\\&o de Azule Energy, através da fusio dos negécios angolanos das em\\
Presas BP e ENI. O acordo segue o memorando de entendimento entre as empresas,\\
firmado em maio de 2021.\\
\\
Digitalizada com CamScanner\\

\\section*{Pagina 37}

16 CONCEITOS GERAIS\\
\\
as Natural\\
1.6 Cadeia produtiva do petroleo \\& G\\
\\
deia de valor foi desenvolvido como um instrumen,,\\
a ou cadeia vnissa de que @ produgao de bens Pode\\
Parte da ee atores estéo interconectados por flu,\\
gpa a de atender a um mercado Consumido,\\
c\\
\\
Oconceito de cadeia produtiv\\
para ter uma visio sistémica.\\
representada como um sistema,\\
de materiais, capital e informagao,\\
final com os produtos do sistema. ar aiadliiteas as petreo ‘ tans .\\
aie ee oe a ito da cadeia produtiva. Existem varias classificacges\\
co consumidor final, aplica-se 0 conce! 5 interesses das empresas em maximizar g,\\
\\
\\_ endem do: ; 4\\
disponiveis na literatura, que dep financeiro de cada uma. Os trés gran.\\
\\
ray condmico \\&\\
€cnicos e o poder e' : piesa 7 -\\_\\
recursos humanos, enna \\& cadeia produtiva da industria petrolifera sao: Upstream,\\
des segmentos que ¢\\
\\
downstream, conforme a estrutura adotada pela Sonangol, E.P, ilustrado ng\\
ic low a pie, rear\\} ‘\\
\\
—. e ‘Algumas empresas, pot questdes econdmicas, dividem a cadeia apenas em\\
Figura 1.17. ”\\
\\
downstream, onde as atividades de midstream esto incluidas no segmento de\\
upstream e ?\\
\\
downstream.\\
\\
Figura 1.17: Cadeia produtiva da industria petrolifera adotada pela Sonangol, E.P.\\
\\
O acesso ou etapa de ganho de entrada é o primeiro passo para uma empresa de petrdleo\\
iniciar no setor de exploracao e producio de hidrocarbonetos é determinar as regides do\\
mundo que devem ser consideradas interessantes em termos técnicos,\\
micos, sociais e ambientais.\\
O objetivo desta etapa é obter direitos de concessiao,\\
plorar e perfurar em areas com licengas para atividades\\
sao realizadas atividades de defini\\
\\
politicos, econd-\\
\\
ou seja, obter permissao para ex-\\
petroliferas. Durante esta etapa,\\
\\
Digitalizada com CamScanner\\

\\section*{Pagina 38}

CADEIA PRODUTIVA DO PETROLEO \\& GAS. NATURAL 17\\
\\
© segmento de upstream engloba as ativi\\
fontes de dleo, bem como o transpo:\\
ende as atividades de exploracao,\\
para muitas empresas e literatura\\
¢ produgao:\\
\\
idades de busca, identificagio e localizagio das\\
rte desse dleo extraido até as refinarias. Compre\\
\\
Perfuracio e Produgio. O conceito de upstream,\\
\\
€ subdividido em exploragio, desenvolvimento\\
\\
. — omen buscar, identificar e quantificar Novas reservas\\
» Por meio das seguintes atividades Principais: garantir\\
acesso a reservas por meio de Negociacdes com entidades publicas ou privadas;\\
analisar a geologia dos subsolos; identificar potenciais reservatérios de PRG:\\
confirmar a existéncia do reservatério, :\\
. Desenvolvimento: tem como finalidade planejar a abordagem e definir os recur-\\
sos necessarios para a producao, de modo a maximizar a rentabilidade de uma\\
reserva. Inclui toda a preparacio para a etapa de produgdo. As atividades prin-\\
cipais dessa etapa sido: avaliar, por meio de PoO¢os, a extensio, o potencial de\\
producio e a viabilidade econémica da reserva; investigar as caracteristicas do\\
subsolo que possam afetar a producio; avaliar possiveis cendrios de producio,\\
desde a localizagao das perfuracées até as especificagSes da infraestrutura a ser\\
utilizada; implementar a infraestrutura de producio.\\
\\
* Produgdo: extrair 0 petréleo e gas de uma reserva com 0 intuito de maximizar sua\\
vida util. Suas atividades mais importantes sio: extrair petrdleo e gis utilizando\\
diversas técnicas de recuperacao (primaria, secundaria e avangada); manter os\\
niveis de producao da reserva otimizados (workover); encerrar as atividades de\\
\\
producio, por exemplo, desativagao de infraestrutura e descarte de residuos t6-\\
xicos.\\
\\
O segmento de midstream € a fase em que as matérias-primas sao transformadas em pro-\\
\\
dutos prontos para uso especifico, como gasolina, diesel, querosene, GLP, nafta,\\
6leo lubrificante etc. Sao as atividades de refinacio:\\
\\
" Transporte: devido ao fato de os campos petrol\\{feros nem sempre estarem lo-\\
calizados préximos aos terminais e refinarias de petréleo e gas, é necessirio o\\
\\
transporte da producao por meio de embarcagées, caminhGes, vagdes ou tubu-\\
lagées (oleodutos e gasodutos).\\
\\
* Processamento e Refino: embora a separacao da Agua, dleo, gis ¢ sdlidos produ-\\
zidos ocorra em estagdes ou na prépria unidade de produgio, é necessirio o\\
processamento e refino da mistura de hidrocarbonetos proveniente da rocha re-\\
servatério para a obtencio dos componentes que sero utilizados nas mais di-\\
versas aplicacdes, como combustiveis, lubrificantes, plisticos, fertilizantes, me-\\
\\
dicamentos, tintas, tecidos etc.\\
\\
O segmento de downstream \\&a parte logistica, ou seja, o transporte dos produtos da reli\\
Naria até os locais de consumo. Envolve o transporte, distribuigio e comercializagio\\
dos derivados do petréleo:\\
\\
Digitalizada com CamScanner\\

\\section*{Pagina 39}

18 —\\_ CONCEITOS GERAIS\\
s estagdes ¢ refinarias (gis nay\\
« provenientes das es Ural\\
. — apenas“ querosene, er err rene Fesiduos pesade .\\
“ ror r 8) io comercializados pelas distribuidoras, wes Se encarregan,\\
—\\_= \\_ a - ein forma original ou aditivada, ao consumidor final\\
de oferecé-los, 0\\
\\
ffer:\\
1.7 Contratos de exploragao \\& produgao petrolifera\\
Jadas por meio de licengas ou contratos concedidos\\
a a abidodls econémica das operagoes de exploracio ¢\\
\\
produco nao depende apenas de fatores geoldgicos ou tecnolégicos, mas também do\\
impacto dos sistemas fiscais € contratuais dos respectivos paises ou entidades locais. No\\
capitulo 4, sobre perfuracao ¢ completagao, abordaremos outros contratos que nao serio\\
\\
i itulo.\\
\\
ag ty nie aye de exploracao e producio petrolifera é regulada por legislacio\\
propria, adotada em funcao da realidade e dos interesses do pais, levando em consideracio\\
as diversas legislacées existentes no mundo. Angola aprovou leis ¢ decretos que servem\\
como instrumentos legais para a regulacao do setor petrolifero nacional.\\
\\
O principal instrumento de regulagao do setor petrolifero é a Lei 10/94 das ativida-\\
des petroliferas e os Decretos de Concessio Petrolffera, aos quais toda atividade petrol\\
fera deve estar em conformidade. Além desse decreto, existem outros decretos auxiliares,\\
\\
destacando-se os seguintes:\\
\\
Decreto Presidencial n° 52/19: Estratégia Geral de Atribuicao de Concessées Petrolife-\\
tas. Este decreto visa assegurar a substituicao de reservas, promovendo a atividade\\
de exploragao de forma racional e adequada, desencadear medidas adequadas 4 con-\\
firmagio do potencial petrolifero do pais e fornecer petréleo bruto suficiente para\\
satisfazer a ny = de refinacdo, mediante ponderacao econdmica da ex-\\
\\
As operacdes petroliferas s0\\
pelo Governo. Nesse contexto,\\
\\
Decreto Presidencial n°\\
n° 86/18, de 2 de Abril: atualiza as Tegras e procedimentas dos con-\\
\\
cursos piblicos no setor petro! z\\
48/06. Petrolifero, revogando o anterior Decreto Presidencial n°\\
\\
Decreto Presidencial n* 9\\}/\\} .\\
vidades de abandono de pores ogo APT: estabelece as rgrase pocekimenics dasa\\
\\
Pogos e desm; ri\\
antelamento de instalagdes de petrdleo e gas\\
\\
no territério angolano,\\
Decreto Legislativo Presidencial n°\\
5/1 5\\
bre as atividades de pesquisa dane de 18 de Maio; estabelece o regime jurihoo 5\\
\\
petroliferas, revogando o anterior Decresa reas de desenvolvimento de concess6es\\
ia ds areas de desenvolvimento alvo, Presidencial n° 211/15, que fazia refere™\\
\\
Decreto Legislativo Presidencial n°\\
mento 6/18, de 18\\
peewee Para a adequacdo dos termos de Maio: define os incentives eo procedt\\
\\
contratuais € fise ;\\
revogando 0 anterior Deer © Aseais apliciveis As zonas mary!\\
13 de Junho. Legislative Presidencial n® 2/16,\\
\\
Digitalizada com CamScanner\\

\\section*{Pagina 40}

CONTRATOS DE EXPLORACAO E PRODUGAO PETROLIFERA 19\\
\\
Decreto Legislativo Presidencial n° 7/18, de 18 d\\
fiscal aplicavel as atividades de Prospeccio,\\
producao ¢ venda de gas natural em Angola.\\
\\
le Maio: estabelece o regime juridico e\\
Pesquisa, avaliacao, desenvolvimento,\\
\\
1.7.1 Contratos de exploracdo e Producdo em Angola\\
\\
Os principais contratos vigentes na indtistria petrolifera angolana sio, nomeadamente, o\\
Contrato de Partilha de Produgao (CPP), o Contrato de Concessao (CC), o Contrato de\\
\\
Associacao (CA) e 0 Contrato de Servico com Risco (CSR). Em Angola, esses contratos\\
esto estabelecidos da seguinte forma:\\
\\
® O contrato comum estabelecido em Angola € 0 CPP, que é aplicado em todos os\\
blocos de offshore (Aguas rasas, profundas e ultra-profundas).\\
\\
= OCC é estabelecido para o Bloco 0.\\
* OCA é estabelecido para os Blocos FS e FST.\\
* OCSR €é estabelecido para os blocos de pré-sal.\\
1.7.1.1 Contrato de Partilha de Producaéo (CPP)\\
\\
Este tipo de contrato € o mais popular na industria petrolifera angolana e tem como pre-\\
ceito fundamental a partilha da Producao (Petréleo Lucro) entre o Grupos Empreiteiros\\
(GE) ea Concessionaria Nacional (CN). No passado, a Sonangol Concessionaria eraa CN,\\
mas atualmente é a ANPG, que atua como representante legal e contratual do governo\\
da Repiblica de Angola no negocio petrolifero do pais, detendo a propriedade exclusiva\\
€ 0 direito de exploracao dos hidrocarbonetos em Angola.\\
\\
O CPP sungiu na Indonésia na década de 1960 e hoje é utilizado como referéncia por\\
Varios paises produtores de petréleo e gas, incluindo Angola, China, Egito, Libia, Filipi-\\
nas, Malasia, Guatemala, Trinidad e Tobago, Quénia, Costa do Marfim, Guiné Equatorial,\\
fntre outros. A ideia por tras do surgimento do CPP esté relacionada A necessidade de\\
teduzir os desequilibrios entre os paises do Oriente Médio e as empresas petroliferas in-\\
ternacionais em relac4o aos contratos de concessao.\\
\\
Os contratos de partilha de produco inverteram a légica da propriedade dos hidrocar-\\
bonetos explorados, Antériormente, as empresas petroliferas internacionais eram propri-\\
¢tarias dos hidrocarbonetos explorados, enquanto agora a propriedade é transferida para\\
\\& estados hospedeiros, O Estado deixa de ser apenas remunerado por meio de royalties\\
© Outros tributos pelo direito outorgado as empresas petroliferas para explorar exclusi-\\
vamente esse recurso, Em vez disso, os hidrocarbonetos passam a ser propriedade do\\
Pe uma parte deles ¢ entregue as empresas internacionais em troca de suas ativida-\\
\\
‘ee envolvido na exploragao ¢ producio. ae\\
\\
mente, os contratos de partilha de producio tém sido amplamente utilizados na\\
\\
Mic ena Asia como modelos de regulacio da atividade de exploracio e producao nesses\\
Continentes,\\
\\
Digitalizada com CamScanner\\

\\section*{Pagina 41}

20 CONCEITOS GERAIS\\
i do (CPP,\\
Regime fiscal dos Contrato de Partilha de Produgdo (CPP)\\
\\
ime fiscal aplicdvel as atividades petroliferas na area de concessio é definido ely\\
Beanto de Gerens De acordo com esse decreto, as associadas da concessions, \\%\\
ae eons. de um unico imposto de rendimento do pirsleg A taxa \\&\\
sth, ales do iboats de rendimento do petrdleo, os i ne amen adem\\
dos a pagar 4 concessionaria o bonus de assinatura dos con oe “Pr oducin,\\
o bénus de produgao e as contribuigées sociais. EB sipatiac “ fanieet ot ee\\
pagamentos feitos 4 concessionaria sao integralmente Pn as Finan\\
¢as, conforme estabelecido no n° 2 do artigo 55° da Lei 13/04.\\
\\
Mecanismo do Contrato de Partilha de Produ¢do (CPP)\\
\\
No CPP, o Grupos Empreiteiros (GE) assume integralmente 0 risco de Pesquisa © arca\\
com todos os custos de desenvolvimento e operagdes. O CPP determina que todo 9\\
Petréleo produzido e obtido na area de concessao seja dividido em duas Partes: uma\\
denominada petréleo para recuperacdo de custos (Cost Oil) e outra denominada petrdleo\\
lucro (Profit Oil). O petrdleo de custo é a parte da producio destinada a Tecuperacio\\
dos custos elegiveis, e pode ser exportado livremente até o limite maximo de 50\\% dos\\
custos elegiveis. O petréleo lucro é a parte da producao que é compartilhada entre a\\
Concessionéria e o Grupos Empreiteiros (GE), com base na TIR da area de concessio,\\
Geralmente, o petrdleo lucro (Profit Oil) e a recuperacao de custos (Cost Oil) sao divi-\\
didos igualmente (50\\% / 50\\%) no inicio da producao. A medida que os custos de i\\
\\
investi-\\
mento forem amortizados, a Porcentagem do petréleo lucro aumentar4,\\
\\
Recuperacao de custos\\
\\
* Para efeitos de Tecuperacao, os cu.\\
\\
; Stos de A\\&S sig incorporados anualmente nas\\
outras categorias de custos com ba:\\
\\
Se na proporcaio correspondente;\\
\\
« Existe um limite de recuperacdo de custos de 50\\%;\\
\\
* Caso os custos nao sejam totalmente\\
\\
aumenta para 65\\%; epee eit aes\\
\\
o limite de recuperagi?\\
* Hé um incentivo ao investimento de 33\\% (Uplift,\\
\\
* O periodo de amortizacio é de 4 anos, confo i cretorLe!\\
a rm elo Decreto\\
13/04 de 24 de Dezembro de 2004, mimibalaatts pola: Dee\\
\\
Digitalizada com CamScanner\\

\\section*{Pagina 42}

CONTRATOS DE EXPLORAGAO E PRODUGAO PETROLIFERA 21\\
\\
partilha do Petrdleo Lucro\\
\\
+ O Petrdleo Lucro remanescente,\\
\\
; a apds a dedugio do Cost Oil, é partilhado entre o\\
Governo (concessionaria) e o GE;\\
\\
» A partilha do Petréleo Lucro é baseada em uma TIR fixa, conforme indicado na\\
Tabela 1.5.\\
\\
Tabela 1.5: TIR.\\
\\
Estado (Concessionario) Grupos Empreiteiros (GE) (As-\\
\\
sociados)\\
\\
70\\% A70\\%\\
\\
llustragao da aplica¢ao do Contrato de Partilha de Producao (CPP)\\
\\
Nesta secao, iremos ilustrar a aplicacdo do CPP com base na producio passada e na TIR\\
utilizando as Tabelas 1.6 e 1.7.\\
\\
Tabela 1.6: Partilha de Producao do Bloco Y, baseada na TIR.\\
\\
Estado (Concessionario) Grupos Empreiteiros (GE) (As-\\
sociados)\\
\\
Petrdleo para recuperacio de custos (Cost Oil) 50\\%\\
\\
Petréleo Lucro (produgao passada)\\
\\
<15\\% 30\\% 70\\%\\
\\
De 15.25\\% 40\\% 60\\%\\
\\
De 25430\\% 60\\% 40\\%\\
\\
De 30840 \\% 80 \\% 20\\%\\
\\
> 40\\% 90\\% 10\\%\\
\\
Tabela 1.7: Partilha de Produco do Bloco X, baseada na produgao passada.\\
\\
—\\
\\
Estado (Concessionério) Grupos Empreiteiros (GE) (As-\\
sociados)\\
Petréleo para recuperacio de custos (Cost Oil) 50\\%\\
Petréleo Lucro (Producio passada)\\
<75MMbbls 60\\% 40\\%\\
De 75 4100 MMbbls 80\\% 20\\%\\
>100 MMbbls 90\\% 10\\%\\
\\
aN ent e § RUNG ein ii\\
\\
Digitalizada com CamScanner\\

\\section*{Pagina 43}

22 —\\_ CONCEITOS GERAIS\\
Representacdo gréfica do Contrato de Partilha de Produ¢ao (CPP)\\
produgio bruta ser partilhada entre o GE € 0 Gove\\
dugaio liquida do GE versus a produgao do Govern\\
fiscais do CPP, que incluem: 0\\
\\
A Figura 1.19 ilustra como a\\
(concessionaria) para determinar a pro\\
(concessionaria), com base nos termos\\
\\
* Recuperacao dos custos:\\
\\
— Limite inicial de 4 anos;\\
maxima de 65\\%.\\
\\
— Extensfo no quinto ano para uma taxa\\
\\
= Partilha do Petréleo Lucro com base na TIR.\\
\\
‘=Petréleo para recuperagao de custos |\\
(ost 018\\
\\
sInelui 33\\% Uplift sobre o Capex\\
*Capex seré depreciado em 4 anos\\
\\
0 principio a cedéncia, pelos governos dos paises hos\\
0 dos hidrocarbonetos aos GE. Nesse tipo de contrato,\\
operacGes petroliferas, enquanto os governos s40\\
Ostos de transacio petrolifera e outras taxas. EM\\
\\
do como Risk Service Agreement (RSA)\\
deiro contrata uma empre\\
o. Nesse tipo de contrat\\
\\
ll\\
\\
Digitalizada com CamScanner\\

\\section*{Pagina 44}

CONTRATOS DE EXPLORACAO E PRODUGAO PETROLIFERA 23\\
\\
a empresa contratada recebe uma compensag:\\
\\
‘io financeira do pais hos; i ic.\\
eceb pedeiro pelos ser-\\
rasil fe S paisi ili SSE ti ‘ontrato en\\
vigos prestados. O Brasil foi um dos paises que utilizou esse tipo de con 1979,\\
\\
Contrato de joint ventures\\
\\
No contrato de joint ventures, 0 governo hospedeiro forma uma companhia de operagées\\
em parceria com © operador. Nesse tipo de contrato, cabe ao operador assumir todos\\
os riscos, desde a exploracio até a descoberta. O reembolso geralmente ocorre a partir\\
\\
da produgao do pais hospedeiro. \\& importante ressaltar que esse tipo de contrato nao é\\
atualmente utilizado em Angola.\\
\\
1.7.2, Comparacao dos principais contratos\\
\\
Nesta seco, faremos uma comparacio resumida dos principais contratos para elucidar\\
\\
os leitores sobre as vantagens de cada um deles. Os resultados dessa comparagio estéo\\
resumidos nas Tabelas 1.8 € 1.9.\\
\\
Tabela 1.8: Tipos de contratos.\\
\\
——————SSS—SS—SsS—\\
\\
Tipologia de Contratos Grupos Empreiteiros (GE) Governo\\
\\
cc Todo risco/Recompensa Recompensa em fungio da producio\\
edo prego\\
\\
cpp Risco de Exploragio/Partilha de re- Partilha de recompensa\\
\\
compensa\\
\\
Joint Venture (Parcerias) Partilha de riscoe de recompensa \\_Partilha de risco e de recompensa\\
\\
Contrato de Prestagao de Servicos Sem risco Todo risco\\
\\
Contrato de Risco Todo risco/Recompensa Recompensa\\
\\
——\\
\\
Tabela 1.9: Comparacao entre o CC e CPP.\\
\\
0 —————— eee\\
\\
Atividade Concessio Partilha de Produgio\\
\\
Pesquisa, desenvolvi- Financiada pela Associacio \\_Financiado pelo Grupos Empreiteiros (GE)\\
mento e produgio\\
\\
Producio Toda produgio é comerciali- As despesas do Grupos Empreiteiros (GE) sio reembolsadas (Re-\\
zada pela Associagio cuperacio de custos) ¢ o grupo é remunerado com uma parte\\
do petrdleo bruto restante (Petrdleo Lucro)\\
\\
Reyaliies cimpostos Paga royalties e impostos O Grupos Empreiteiros (GE) paga imposto sobre o rendimento\\
(Petroleum Income Tax (PIT),\\
Petroleum transaction Tax (PTT))\\
\\
ertes So propriedades da associagio io propricdade do Estado representado pela ANPG (no pas\\
\\
sado era Sonangol Concessioniria)\\
\\
Mento do rendi- PTT Produgio acumulada para alguns casos ou escalio da TIR para\\
\\
— ‘outros\\
\\
Digitalizada com CamScanner\\

\\section*{Pagina 45}

24 CONCEITOS GERAIS\\
\\
Conclusao\\
\\
resentamos os conceitos basicos e abrangentes da Indtistria de Petra,\\
leo, proporcionando aos leitores conhecimentos sobre os sae CONCE:TOS teenicn,\\
relacionados ao petréleo, a cadeia de valor da industria do petrdleo (upstream, midstreany\\
¢ downstream), assim como os principais contratos vigentes em Angola. Também foram\\
apresentados os principais marcos histéricos da industria petrolifera mundial ¢ em 4,\\
\\
gola.\\
\\
Neste capitulo, ap\\
\\
1\\
\\
Digitalizada com CamScanner\\

\\section*{Pagina 46}

CAPITULO 2\\
\\
CADEIA DE VALOR DA INDUSTRIA\\
PETROLIFERA\\
\\
As operagées da industria do petréleo sio divididas em trés principais segmentos: Ups-\\
tream (montante), Midstream e Downstream (jusante). Upstream engloba as atividades de\\
busca, identificagao e localizagao das fontes de dleo e o transporte deste dleo extraido até\\
as refinarias. Midstream € 0 segmento em que as matérias-primas sao transformadas em\\
produtos prontos para uso especifico. Downstream é a parte logistica, ou seja, o transporte\\
dos produtos da refinaria até os locais de consumo. Este conceito depende da literatura\\
utilizada, assim como as necessidades da propria empresa operadora, conforme ilustrado\\
nas Figuras 2.1, 2.2 2.3. O conceito de Upstream para muitas empresas ou literaturas é\\
subdividido em exploragao, desenvolvimento e produgao.\\
\\
Downstream\\
\\
Figura 2.1: Cadeia produtiva da industria petrolifera - A.\\
\\
Exploracio: tem como finalidade buscar, identificar e quantificar novas reservas de pe-\\
tréleo e gis, através das seguintes atividades principais: garantir acesso a reservas\\
Por meio de negociagées com entes publicos ou privados; analisar a geologia dos\\
\\
subsolos; identificar potenciais reservatérios de petréleo e gis; confirmar a existén-\\
cia do reservatério.\\
\\
Desenvolvimento: tem como finalidade planejar a abordagem e definir os recursos ne-\\
cessarios para a produgio, que maximizem a rentabilidade de uma reserva. Inclui\\
toda a preparacao para a etapa de produgio. As atividades principais dessa etapa sio:\\
avaliar, com auxilio de pocos, a extensio, o potencial de produgio e a viabilidade\\
€conémica da reserva; investigar as caracteristicas do subsolo que possam afetar a\\
\\
Indisina do\\
\\
rt petrileo para ndo-engenheiros, 1 edigdo. 25\\
\\
André Ramos Copyright © 2023 Editora Dois03\\
PY\\
\\
Digitalizada com CamScanner\\

\\section*{Pagina 47}

26 CADENA DE VALOR DA INDUSTRIA PETROUFERA\\
\\
; 50: planejar a melhor forma de ex\\
on ivei nirios de produgao; P 5 i =\\
\\
Tar, do.\\
\\
utilizada; implementar 4 infraestrutura de producto\\
\\
oo aS\\
\\
jastria petrolifera - B. Cortesia: Sonangol, E.P.\\
\\
Figura 2.2: Cadeia produtiva da in\\
\\
de uma reserva com intuito de maximizar sua vida\\
til, Suas atividades mais importantes so; extrair petréleo e gas com as mais \\&\\
versas técnicas de recuperacao (primaria, secundaria e enhanced); manter niveis de\\
produco da reserva otimizados (workover); encerrar as atividades de producio (ex:\\
desativacao de infraestrutura € descarte de residuos t6xicos).\\
\\
Producio: extrair 0 petréleo e gas\\
\\
Transporte: pelo fato dos campos petroliferos nao serem localizados préximos dos ter-\\
minais e refinarias de éleo e gas, € necessario 0 transporte da producao através de\\
embarcacoes, caminhGes, vagoes ou tubulaces (oleodutos e gasodutos).\\
\\
Processamento e Refino: apesar da separacdo da agua, dleo, gas e sdlidos produzidos,\\
ocorrer em estacSes ou na propria unidade de produclo, é necessario o processe-\\
mento ¢ refino da mistura de hidrocarbonetos proveniente da rocha reservat6rio,\\
para a obtencao dos componentes que serao utilizados nas mais diversas aplicacées\\
il iveis, lubrificantes, plasticos, fertilizantes, medicamentos, tintas, tecidos\\
etc.).\\
\\
Distribuigao: os produtos finais das estacdes e refinarias (gis natural, gis residual, GLP\\
gasolina, nafta, querosene, lubrificantes, residuos pesados e outros destilados) so\\
\\
comercializados com as distribuidoras 4\\
= * , que se encarregarao de oferecé-los, na Su\\
forma original ou aditivada, ao consumidor final.\\
\\
i\\
So Ge\\
\\
Figura 2.3: i i i\\
Cadeia produtiva da indtistria petrolifera - C.\\
\\
2.1\\
Segmento da cadeia Produtiva de Upstri\\
eam\\
\\
de forma A s\\
diversas etapas da cadeia de valor da ind senunaida as diferentes atividades realizadas ¥\\
Petrolifera no segmento de Upstrea" -\\
\\
——— y |\\
\\
Digitalizada com CamScanner\\

\\section*{Pagina 48}

SEGMENTO DA CADEIA PRODUTIVA DE UPSTREAM 27\\
\\
forme ilustrado na Figura 2.4, Este segmento de Upstream esta dividido em\\
pais etapas: acesso, exploragio, desenvolvimento, Produgiio e abandono\\
\\
1-9 anos ———*-— 1-2 anos + 2.\\textasciitilde{} 10-20 anos —+ [em]\\
\\
5-10 anos ——+[Oxew]\\
\\
Figura 2.4: principais etapas do segmento de Upstream.\\
\\
cinco princi-\\
\\
———\\_ 0- Llanos ———+— 0 anos ———————— 01 anos\\
\\
| 2.1.1 Etapa de acesso\\
\\
Esta etapa também é conhecida como etapa de ganho de entrada. Ba etapa para adquirir\\
os direitos de concessao. O objetivo principal desta etapa de Upstream é obter os direitos\\
de explorar e perfurar. Uma das principais atividades é a defini¢ao de bacias para a pros-\\
peccao, negociar com o governo e garantir a licenga de exploragdo, desenvolvimento e\\
producao do campo. Nesta etapa, determina-se as regiées com maior potencial quanto\\
aos aspectos técnicos, politicos, econdmicos, sociais e de meio ambiente:\\
\\
" Aspectos técnicos: avaliar a quantidade do potencial de hidrocarbonetos que podem\\
ser descobertos e produzidos, incluindo os desafios técnicos. Esta avaliacao é feita\\
através de estudos de reconhecimento que serdo feitos usando os estudos publica-\\
mente disponiveis.\\
\\
* Aspectos politicos e econémicos: incluem regime politico e estabilidade do go-\\
verno, potencial nacionalizacao da industria do petréleo, embargos atuais, estabi-\\
lidade fiscal e niveis de taxaciio, restrigdes 4 remessa de lucros, seguranca pessoal,\\
custos locais, previsdes de inflacio e taxas de cambio.\\
\\
* Aspectos sociais: disponibilidade de forca de trabalho especializada e treinamento\\
local requerido, o grau de esforco requerido para estabelecer a presenca de um con-\\
tetido local capacitado para diversas areas da industria petrolifera.\\
\\
* Aspectos ambientais: estudo sobre o impacto ambiental e as precaugGes necessarias\\
Para proteger o meio ambiente e qualquer legislacao local especifica contra danos\\
ambientais durante as opera¢6es petroliferas.\\
\\
2.1.2 Licitagées de blocos\\
\\
Oconvite para a participacao é anunciado publicamente, sob a forma de rodada de licita-\\
S60. O pais tem a 4rea a licitar com hidrocarbonetos, mas nenhum recurso financeiro,\\
\\
lumano ¢/ou tecnolégico, enquanto as empresas investidoras possuem © ee\\
Cursos necessdrios a atividade petrolifera. O governo, por meio da concessionaria ure\\
4gencia reguladora, oferecer4 certa area “bloco” para empresas durante o licencia\\
\\
Digitalizada com CamScanner\\

\\section*{Pagina 49}

28 CADEIA DE VALOR DA INDUSTRIA PETROLIFERA\\
\\
para atrair investimentos para as ativi.\\
sas de petroleo competirao entre si n,\\
bloco. O vencedor receberd o direite\\
riodo e obter uma compensacio com,\\
timentos oferecendo “Pacotes\\
\\
O governo “oferecera” dados brutos ¢ contratos\\
\\&P). As empre\\
\\
dades de Exploragao e produgio (E\\
fase de licitagio para obter o direito de E\\&P de um\\
\\
de gerir as atividades de E\\&P para determinado pe\\
\\
‘, ia invest\\
ontrato. Nesta fase, o governo precisa atrair in\\
\\
base noc ; n\\
: m a protecao do interesse nacional.\\
\\
interessantes” enquanto manté\\
2.1.3 Contratos internacionais na industria petrolifera\\
oderao ser adotados pelo governo e as empresas investido.\\
\\
Existem varios contratos que P' is\\
usados sao:\\
\\
ras. Ao nivel internacional, os contratos mais\\
= Contrato de Partilha de Producao (Production Sharing Agreement): contrato assi-\\
nado entre o estado e uma empresa contratada que vai pesquisar uma area a sua\\
conta ¢ risco. Em caso de descoberta, 0 estado assume a receita do petréleo produ-\\
\\
zido e faz os pagamentos devidos ao governo € A empresa contratada;\\
\\
* Acordo de Licitagdo Conjunta (Joint Operating Agreement): organizacao de empre-\\
sas a fim de participarem de uma licitagdo sobre uma area especifica para explorare\\
produzir, formalizando regras entre si dessa participa¢ao;\\
\\
= Acordo de Operacées Conjuntas (Joint Operating Agreement): acordo formulado\\
pelas empresas exploradoras e produtoras de petréleo para que se desenvolva a sua\\
operacao, dividindo fungées;\\
\\
. Acordo de Farmout: contratos de cessao de direitos de exploracao e producio dos\\
concessionarios com terceiros;\\
\\
. nn by yale (Lifting) ou Acordo de participacdes (Offtake Agreements) € 05\\
\\
oes wii he de Gas (Gas Balance Agreements): celebrados por empre-\\
res que re; i 4 6\\
\\
ceauael cme gem a retirada ordenada da producio de petréleo\\
\\
* Acordo internaci "\\
Agreements): shang. 9 vendas de petréleo bruto (International Crude Oil Sales\\
venda de petréleo bruto e de gés natural entre export\\
\\
dores e importadores,\\
» Produtores e oui\\
ici itros agentes da . Aunt entre\\
Participantes de uma mesma joint venture. ee! cadeia econémica ou\\
\\
2.1.3.1\\
Contratos de partitha de producao (CPP)\\
\\
lucro di »\\
antes que a propredade de todo g cg Pet le®" © periodo tipico varia de 5 308"\\
bide direitos minerais, Uma vez Spssendieat seja transferida de volta ao proP™\\
de producio dividem a es ad 98 custos de investimento, 0s a¢ os\\
"sa0 entre o §0verno ¢ a companhia de petrole?\\
\\
4\\
Digitalizada com CamScanner\\

\\section*{Pagina 50}

SEGMENTO DA CADEIA PRODUTIVA DE UPSTREAM 29\\
\\
recuperagao de custos\\
\\
oGrupo Empreiteiro deve recuperar todas as despesas de exploragao, desenvolvime a;\\
operagao \\& administracao e servicos (A\\&S) sequencialmente: mg\\
\\
= Custos de operacio, desenvolvimento (valor depreciado) e exploracio;\\
\\
» Para efeitos de recuperacao, A\\&S sao incorporados anualment\\
\\
= € nas demais catego-\\
rias de custos com base na proporc4o correspondente;\\
\\
* 50\\% Limite de recuperacao de custos;\\
\\
= 65\\% caso os custos nao forem recuperados em 4 anos;\\
\\
= 33\\% de incentivo ao investimento (Uplift);\\
\\
* Periodo de amortizacao de 4 anos: Decreto Lei 13/04 de 24 de Dezembro de 2004.\\
\\
Partilha de lucro\\
\\
Cost oil (petréleo para recuperacao de custos) é o petréleo em que o Grupo Empreiteiro\\
(GE) deve recuperar as despesas de pesquisa, de desenvolvimento, de producio e de ad-\\
ministra¢ao e servicos decorrentes da aplicacao do contrato no termo da Lei sobre a Tri-\\
butagao das Atividades Petroliferas, através do levantamento e livre disposi¢ao, até ao\\
maximo por ano de 50\\%. O Profit Oil (petréleo lucro) remanescente apés deducio do\\
Cost oil (petrdleo para recuperacio de custos), sera partilhado entre o Governo (Conces-\\
sionaria) e o Grupo Empreiteiro. A Partilha do Profit Oil é baseada na TIR fixa. Exemplo\\
de CPP - BLOCO AA baseada na TIR est descrito na Tabela 2.1.\\
\\
Tabela 2.1: Exemplo de CPP - Bloco AA.\\
\\
Estado (Concessiondrio) Grupo Empreiteiro (GE)\\
\\
Petréleo para recuperacao de custos (cost oil) 50\\%\\
\\
Petréleo lucro (Profit Oil)\\
\\
Menos 10\\% 35\\% 65\\%\\
De 10\\% 4 menos de 25\\% 45\\% 55\\%\\
De 25\\% \\& menos de 30\\% 70\\% 30\\%\\
De 30\\% a menos de 35\\% 80\\% 20\\%\\
De 35\\% ou mais 90\\% 10\\%\\
\\
2. x\\
14 Etapa de exploracdo\\
s . servas de P\\&G.\\
Exploragio tem como finalidade buscar, identificar € quantificar novas reser ss le PS\\
ravi das Seguintes atividades principais: garantir acesso a reservas Por me rid\\
pear com entes publicos ou privados; analisar a geologia dos subsolos; i\\
“enciais Teservatérios de P\\&G; confirmar a exis'\\
\\
io de ne\\
\\
ncia do reservatorio.\\
\\
Digitalizada com CamScanner\\

\\section*{Pagina 51}

30 ——\\_CADEIA DE VALOR DA INDUSTRIA PETROLIFERA\\
\\
Figura 2.5: Segmento de Upstream.\\
\\
Depois que um contrato de exploracio e producio é estabelecido, a equipe de exploracao\\
comeca a reunir os dados de subsuperficie que serdo usados para selecionar locais paraa\\
Perfuracao de um ou mais pocos de exploracao, ou wildcat,\\
\\
As principais técnicas ou atividades utilizadas na exploracao sao:\\
\\
* Estudos geol6gicos: colher e analisar as amostras de forma a identificar as possiveis\\
areas de existéncia dos hidrocarbonetos;\\
\\
* Processamento sismico: criar imagens na subsuperficie de forma a identificar as\\
areas potenciais de hidrocarbonetos antes de perfuracio;\\
\\
. \\$40: perfurar pocos para termos certeza da existéncia de hidrocarbonetos;\\
\\
Teservatorios e rochas selantes, ,\\
Ps Ppt pip ape 2 Potencial de dleo e gas antes de perfurar um pos0, «\\
direcionada para a terra, A ee terrane, O método sismico usa energia sonora hw ;\\
volta pelas camadas subterriness. pi ae das Tochas subterraneas e ¢ ere\\
ou, mais Tecentemente, uma visio Ad lanad Sismicas oferecem uma visio 2D 0\\
Uma vez que Pesquisas geoldg, ou 'magens sismicas com lapso de tempo. —\\
Pogos de exploracio sig Perfura ied € sismicas determinam uma zona produit\\
bonetos a serem desenvolvidos em gc. bem-sucedidos, seré a descoberta de hidro\\
Esta fase de <™ quantidades comerciais, ;\\
€ a fase onde certificamos a vena mente conhecida como fase de descoberta, pom”\\
\\
ta\\
fase ¢ somente conhecida de hi tbonetos, Em muitas literatura \\&\\
£Omo exploracig ¢ Seguir-se-4 a fase de avaliagio. Em outras\\
\\
Digitalizada com CamScanner\\

\\section*{Pagina 52}

SE\\
‘GMENTO DA CADEIA PRoDUTIVA DE UPSTREAM =. 3.\\
\\
yreraturas, como a que adotamos neste livro, a exploracio é com;\\
\\
osta por descoberta\\
pi scobi e\\
avaliag20-\\
\\
2.1.4.1 Etapa de Avaliacgao\\
\\
0 objetivo desta etapa é avaliar o potencial do reservatério. As atividades principais sao:\\
cio de avaliacdo e simulacao intensa de reservatérios. Bsta fase \\& posterior a ise\\
\\
da descoberta. Uma vez identificados os hidrocarbonetos apés a perfuragao do poco de\\
exploraco, ainda sera necessario um esforco consideravel para avaliar com precisio o\\
potencial da descoberta, e o papel da avaliacao é fornecer informagdes econdmicas.\\
\\
Durante a avaliacdo, normalmente sao perfurados mais pocos para coletar informa-\\
goes e amostras do reservatério, e outra pesquisa sismica também pode ser adquirida, a\\
fim de melhor delinear o reservatério. Essa fase do processo visa reduzir a faixa de in-\\
certeza nos volumes de hidrocarbonetos existentes, definir a extensio e a configuracio\\
do reservatério e coletar dados para a previsao do desempenho do reservatério durante\\
a vida util prevista da producio.\\
\\
Nos casos em que a quantidade de dados obtidos nao fornecem um quadro preciso do\\
tamanho, forma produtiva das reservas, as op¢Ges seguintes devero ser consideradas:\\
\\
* Continuar desenvolvendo os campos para gerar renda em um periodo de tempo\\
curto. No entanto, é preciso ter em conta os possiveis riscos, como subestimar ou\\
superestimar as dimensdes do campo, o que pode afetar a lucratividade do projeto;\\
\\
* Planejamento de avaliaco com objetivo de otimizar o desenvolvimento técnico do\\
projeto. Os riscos dessa opcao sao retardar a data do “primeiro dleo” para muitos\\
anos e que pode exceder 0 investimento inicial do projeto. No entanto, a rentabili-\\
dade do projeto pode melhorar através da otimizacao técnica;\\
\\
* Avaliar o custo da descoberta e comercializ4-la. Existem empresas qualificadas para\\
exploracio e sem intencdo de investir na etapa de desenvolvimento. Os lucros gera-\\
dos na venda das 4reas descobertas so usados para explorar outras oportunidades;\\
\\
\\&\\
t\\
f\\
5\\
3\\
f\\
\\
* Nada fazer. Esta foi sempre uma opcAo principalmente quando os campos ou proje-\\
tos sio marginais, ou seja, o volume das descobertas nao € economicamente vidvel.\\
Contudo, pode ser uma op¢do inadequada para governos que querem gerar lucros,\\
\\
ig " "uma vez que campos desenvolvidos com sucesso podem gerar receitas significativas.\\
\\
i juan-\\
Em suma, como j4 mencionamos acima, nesta fase, quantificam-se os volumes ou q\\
tidades de hidrocarbonetos existentes no reservatorio e serve para:\\
\\
* Avaliar com precisio o potencial da descoberta;\\
reservatorio, €\\
\\
. = a rat s do\\
Mais 5 rmagées e amostra'\\
Pocos sao perfurados para obter info: i delinear o reser\\
\\
Outra pesquisa sfsmica também pode ser adquirida, a fim de melhor\\
vatério;\\
\\
A xii es;\\
* Reduzir a taxa de incertezas nos volumes de hidrocarbonetos existent\\
\\
Digitalizada com CamScanner\\

\\section*{Pagina 53}

32 CADEIA DE VALOR DA INDUSTRIA PETROLIFERA\\
\\
19 e obter dados para a preyigs\\
ie\\
\\
reservatori\\
prevista da produ¢io\\
\\
* Definir o tamanho € a configuracao do reserve\\
do desempenho do reservatorio durante 4 vida util\\
\\
2.1.5 Etapa de Desenvolvimento\\
toe construir instalagées. As principaiy\\
\\
¢ sancionar 0 proje\\
\\
nar os conceitos, definir o plano de desenvolvimento de\\
rfuracao de desenvolvimento.\\
é realizado € considerando que 0 projeto é econo.\\
micamente viavel, um plano de desenvolvimento é elaborado ¢ executado. O Plano de\\
\\
Desenvolvimento de Campo é0 documento reitor usado para aplicagdo € monitorar as\\
esenvolvimento de novo campo ou expandir projetos de\\
\\
atividades necessarias para 0 d\\
desenvolvimento ja existentes. O plano de desenvolvimento deve incluir 0 seguinte:\\
\\
O objetivo desta etapa é definir\\
atividades desta etapa sio selecio\\
campo, construir instalagdes € pe\\
\\
Uma vez 0 estudo de viabilidade\\
\\
* Objetivos do desenvolvimento;\\
\\
* Dados de engenharia de petréleo;\\
\\
* Principios operacionais ¢ de manutengao;\\
* Descricio das instalacGes de engenharia;\\
\\
\\# Estimativas de custos ¢ de mio de obra;\\
\\
* Plano de projeto;\\
\\
* Resumo econdmico-financeiro do projeto;\\
\\
* Proposta do orgamento;\\
\\
* Caracteristicas do reservatério joi\\
téno, contaminantes etc.); (volume original de éleo e gas existente no reserv\\
\\
* Volume\\
de reserva recuperivel em P90-P50-P1 0;\\
\\
* Taxa de Produgio ¢ Instalagbe:\\
de are\\
lages de superficie); \\$s de Producao (numero e tipo de pogos, tipo de inst\\
\\
* Estudo de impacto ambiental;\\
* Avaliacio de riscos;\\
* Health, Safety and Envi\\
ronment (HSE,\\
) ou Seguranga, Satide Ocupacional e Met A™\\
\\
Se '\\
T aprovado pelas autoridades: Um ver\\
\\
Pamentos de\\
rides,» compen usta, perfuracdo © comple ‘ + instalagdes\\
4 decisio de —. Se forem cumpridos os cntene® ue\\
dinheiro no projeto, Decisie Final \\&\\
\\
Digitalizada com CamScanner\\

\\section*{Pagina 54}

SEGMENTO DA CADEIA PRODUTIVA DE UPSTREAM 33\\
\\
Investimento (FID). Uma vez aprovado o Plano de Desenvol\\
\\
b3 ; vimento, existem vari a\\
yidades que deverao ser realizadas antes do inicio da produca stasratt:\\
\\
0:\\
» Plano de desenvolvimento de campo;\\
\\
» Projeto detalhado das instalaces;\\
\\
+ Aquisicao de materiais para construcao;\\
\\
« Fabricagao das instalagdes;\\
\\
» Efetivacao das instalacdes;\\
\\
* Ativacdo de toda unidade petrolifera e equipamentos.\\
\\
Equipes multidisciplinares intervém nesta fase do projeto. Depois que uma descoberta\\
promissora é feita, essas equipes de especialistas preparam um plano de desenvolvimento\\
que integra o cronograma de producao em campo as necessidades do mercado. Abaixo\\
vamos mencionar alguns intervenientes nesta etapa de projeto:\\
\\
« A equipe técnica analisa as condi¢Ges de reservatério em detalhes e prepara varias\\
opg6es de desenvolvimento, juntamente com seu impacto ambiental, selecionando\\
o nimero e a localizacao dos pocos de desenvolvimento e especificando as instala-\\
ges de superficie necessdrias para processar a producdo e exportacao de petréleo,\\
o plano de desenvolvimento do projeto, avaliacao de riscos e orcamento geral;\\
\\
* Aequipe comercial (marketing) elabora uma carta de intengdes e/oua venda a longo\\
prazo dos hidrocarbonetos produzidos, especialmente para o gas natural, que ainda\\
nao é uma mercadoria internacional comercializada de maneira semelhante ao pe-\\
tréleo;\\
\\
* Os analistas financeiros incorporam o trabalho das equipes técnicas e de marketing\\
nos modelos financeiros e preparam um relatério detalhado da economia do projeto\\
para cada op¢4o de desenvolvimento proposta;\\
\\
* Os termos fiscais do Contrato de Petréleo de Upstream sao incluidos, de modo a\\
alocar o fluxo de caixa futuro ao proprietario do mineral e ao contratado ou 4 joint\\
venture;\\
\\
* Por meio de um processo iterativo, os especialistas tecnicos, de marketing e finan-\\
ceiros decidem sobre um plano de desenvolvimento ideal ¢ o recomendam para a\\
decisio final de Investimento e para financiamento.\\
\\
2.1.5.1 Instalacgées Petroliferas\\
\\
Na etapa de desenvolvimento, inclui também a construgao de equipamentos que serio\\
\\
stlizados Para producao de hidrocarbonetos, ou seja, instalaces petroliferas é um a\\
\\
Hnto de equipamentos para extrair (elevar) petréleo e gas de subsuperficie a aypeciice,\\
\\
rome exportar para o cliente, conforme o requisito especificado no contrato. Con!\\
a localizagio, consiste em subsuperficie (pogos) € instalagdes de su\\
\\
por Ou dimensdes ideais das instalagdes dependera do tamanho do res\\
\\
° do contrato, que definira a taxa étima de producao.\\
\\
perficie. Ot\\
ervatorio \\& da\\
\\
Digitalizada com CamScanner\\

\\section*{Pagina 55}

34 CADEIA DE VALOR DA INDUSTRIA PETROLIFERA\\
\\
2.1.6 Etapa de Producgao “\\
ei rimeiro éleo se torna comercial mente vii\\
Lane rea recuperar © dinheiro investido em outras\\
\\
as procuram reduzir o peri\\
\\
da cadeia produtiva da industria petrolifera. As empres: a ag ea Periods\\
\\
ses da ei y exploragio ¢ 0 primeiro éleo. Esta é uma Pp es ques,\\
entre 0 inicio da\\
\\
i vem buscar.\\
\\
Mm) iteiros deve! SCs\\
empresas operadoras ou emprell 3 ie er -\\
Os planos de desenvolvimento ¢ de produgao baseiam-se nO perfil de p ‘oducig de\\
\\
i 6 essarias, a quantidade € o ajustamento de\\
hidrocarbonetos. Os tipo! de instalagdes neces .\\
\\
r tos. Os tipos s : d\\
fa es dos pocos a serem perfurados so determinados pelo perfil de P roducao. A producig\\
OS P\\
s i seakiealé atinge um plateau sustentado, que pode continuar por muitos anos,\\
de atingir o nivel em que nio é mais econdmico, como ilustrado na Figura 2.6 ena\\
ant al\\
\\
Tabela 2.2.\\
\\
Esta etapa comeca quando a\\
E nesta fase que se comega a geral\\
\\
Figura 2.6: Plateau de producao: < 5 anos (dleo), 5 — 10 anos (gas natural).\\
\\
O pessoal de engenharia, operacGes e manutencAo, projetam, constroem, gerem e man-\\
tém os pocos de producio e as instalacées de superficie, aplicando as melhores praticas\\
para maximizar a recuperacao de hidrocarbonetos, cumprindo os padrdes de seguran?\\
€ as restriges ambientais, para atingir os objetivos de producao esperados.\\
\\
Tabela 2.2: Trés fases que caracterizam a producao de um campo.\\
a a eee\\
1. Periodo de\\
construgdo —Pocos produtores recém-perfurados so colocados progressiva\\
mente em funcionamento,\\
Os pogos antigos comecam a declinar enquanto os pocos novos\\
Bs, ete ser colocados em produgao com as instalagdes de\\
pu ze funcionarem em plena capacidade. A produgio se\\
my emer até um certo periodo. Para um campo de pro-\\
aah rf slag Periodo perdura até cinco anos e de gis de\\
A\\
ier ind i hidrocarbonetos neste periodo entra em declinio.\\
mente © periodo mais longo da fase da produgio de um\\
\\
Taree\\
O pessoal da drea comercial\\
Pessoal fina: Bere a venda\\
i resultados ae hidrocarbonetos produzidos, ao pass° e\\
\\
Tn especie Se pen consoan; ‘ ibui pagamentos ao governo anfitn\\
"© 08 acordos contratuais em vigor.\\
\\
2, Perfodo de plateau\\
\\
3. Periodo de declis\\
\\
Digitalizada com CamScanner\\

\\section*{Pagina 56}

SEGMENTO DA CADEIA PRODUTIVA DE UPSTREAM = -355\\
Nem todos 0 hidrocarbonetos descobertos no reservatdy\\
omicamente com a tecnologia atual. Fatores de recupe!\\
de dleo podem variar de 15 a 50\\%, ao passo que para ca:\\
altos, da ordem de 60 a 85\\%.\\
\\
Mais de 90\\% da producio de hidrocarbonetos em An\\
maritimos (offshore), com os Blocos 0 e 14 operados pel\\
Cabinda Gulf Oil Company, Blocos 17 e 32 operados pel:\\
rado pela ESSO Angola, Blocos 3/05 e 4 operados pela Sonangol Pesquisa e Producio,\\
Bloco 15/05 ENI Angola e Bloco 32 operado pela BP. Nos campos de producio em terra\\
(Onshore), temos os Blocos FS e FST operados pela Somoil e Cabinda Sul operado pela\\
PlusPetrol.\\
\\
A qualidade e as ramas do petréleo produzido nestes campos esto ilustrados na Ta-\\
bela 2.3. A qualidade do petréleo produzido em Angola est entre as melhores qualidades\\
de dleos produzidos no mundo. Esta comparacio pode ser feita através da Tabela 2.4,\\
onde estao ilustradas as diferentes ramas do petréleo produzido no mundo.\\
\\
tio podem ser Tecuperados eco-\\
Tacao tipicos para Teservatérios\\
impos de gas sao geralmente mais\\
\\
gola é proveniente dos campos\\
la Chevron através da sua filial\\
la TOTAL Angola, Bloco 15 ope-\\
\\
Tabela 2.3: Ramas e qualidade do petrdleo produzido em Angola.\\
|)\\
Rama \\_Densidade API Acidez (TAN, mg KOH/g) Enxofre (\\% Peso)\\
\\
Palanca 37,0 0,03 0,18\\
Cabinda 32,5 0,06 0,12\\
Nemba 39,7 Ol 02\\
Dalia 23,6 1,39 0,51\\
Girassol 29,9 0,36 0,42\\
Gimboa 24,9 0,49 0,63\\
Hungo 28,3 0,72 0,65\\
Kissange 31,8 0,3 0,39\\
Mondo 30,3 0,77 0,39\\
Paz-Flor 25,3 1,8 0,39\\
Pluténio 33,6 0,37 0,15\\
Satumo 24,4 0,77 0,85\\
Saxi/Bamque 32,7 0,65 0,25\\
\\
Amelhor rama da producdo petrolifera de Angola é Nemba, com a densidade em grau\\
API de 39,7, que fica muito préxima do petréleo produzido nos Estados Unidos, West\\
Tat iate (WTI), assim como o petréleo de referéncia do Reino Unido, Brent,\\
coma densidade em grau de API de 40,0.\\
\\
2.7 Abandono ou desmantelamento do campo\\
\\
Otempo de vida econdmica de um projeto é definido como 0 ponto em que a sy yes\\
uo cobre mais os Custos operacionais (e os royallties). Tecnicamente, ainda € Pc =\\
ir, Mas haverd perda financeira. Ou seja, as despesas serdo maiores em comp®!\\
\\
. utras estratégias\\
otoues td 0 de campos pode ser retardado adotando o'\\
\\
\\textasciitilde{}\\textasciitilde{} —\\
\\
Digitalizada com CamScanner\\

\\section*{Pagina 57}

36 CADEIA DE VALOR DA. INDUSTRIA PETROLIFERA\\
\\
leo produzido no mundo,\\
\\
do petrél\\
\\
Tabela 2.4: Ramas € qualidade\\
\\
pI Enxofre (\\% Peso)\\
\\
Rama Origem Grau Al\\
xtra Lt Arabia Saudita 38.) 11\\
Arabia Lt Aribia Saudita 34,0 1,9\\
Brent Reino Unido 40,0 0,5\\
Cano Limon Colémbia 25,2 09\\
Daging China 33,0 0,1\\
Forcados Nigéria 29,5 0,2\\
Kwait Blend Kuwait 30,9 2,5\\
Marlim Brasil 20,1 07\\
Maya México 21,3 34\\
1 BUA 39,8 0,3\\
\\
Wo i\\
\\
* Reduzir os custos operacionais;\\
= Aperfeigoar o processamento de hidrocarbonetos;\\
\\
= Investir ou aplicar outras técnicas de recuperacao de hidrocarbonetos que sejam ef-\\
Gientes e baratas;\\
\\
* Aplicar incentivos fiscais.\\
\\
Nesta etapa do ci a\\
tam as cig id cs um campo, Os Custos operacionais € de manutengio represen-\\
al peeeeeniis scam co Esses custos estao relacionados com o numero de\\
\\
F er funcionar Os equipamentos, a quantidade de equipamentos\\
\\
€ Os custos diarios das platafor\\
‘mas de 4 Pi 2\\
alugados, producao em casos em que os equipamentos sio\\
\\
Além da condic’o econémica mencionad:\\
\\
a, exi é\\
um poco pode ser abandonado; » €xistem outras raz6es para declarar que\\
\\
* Cessacao da Producio:\\
lucros com sua producio\\
lo em operacio;\\
\\
© poco nao is *\\
aks pa é wuals €conomicamente rentavel, ou seja, \\%\\
m mais os gastos, e portanto nao é vidvel mante-\\
\\
0\\
essa técnica consi\\
mpletado e reperfur4-| Onsiste em abandonar parte do pogo ja perfurado¢\\
\\
0 Seguindo +\\
um bri\\
@ trajetéria lateral para atingir outras areas\\
\\
. Nessa \\textasciitilde{}\\
peracao, retira-se a parte superior da coluné\\
\\
* Abandono de Pe,\\
Tfuracde,\\
rad ¢5es Pilot,\\
Pados e abandonados imedlatameene Exploratérios: Nesse caso, pogos sio a"\\
\\
completacio d, apé d\\
“Stes Pocos, POS serem perfurados e/ou testados. Ni? hi\\
\\
Digitalizada com CamScanner\\

\\section*{Pagina 58}

SEGMENTO DA CADEIA PRODUTIVA DE UPSTREAM\\
\\
37\\
2.1.7.1 Opsdes de abandono ou desmantelamento de Pocos\\
\\
Os custos de abandono sao muito elevados e nao ha Teceitas, j4 que a producio est\\
ralisada. Recomenda-se que os custos de desmantelamento Sejam pré considerados\\
\\
rojeto do pogo, antes de sua perfuracdo. Um projeto comprometido com © fstar\\
mantelamento garante que 0 abandono seja feito mais rapidamente e de forma s\\
e ainda possibilita reducao de custos. Os seguintes fatores como Prioridades do opera-\\
dor durante 0 desmantelamento: seguranc¢a (completar o projeto sem prejuizo a pessoas\\
eao meio ambiente), custos (operagGes economicamente viaveis sempre minimizando\\
custos), reputagao (operador com responsabilidades perante todas as partes envolvidas),\\
responsabilidade futura (minimizar responsabilidade pos desmantelamento), minimizar\\
horas de trabalho humana no projeto. Entre estes fatores, custo e seguranca serao direta-\\
mente afetados por decisdes tomadas durante o projeto do campo.\\
\\
4 pa-\\
jano\\
‘0 des-\\
egura\\
\\
2.1.7.2 Abandono de pocos e instalagées em Angola\\
\\
O abandono de pogos e instalagdes em Angola segue as regras internacionais, no qual\\
as companhias operadoras passam a alimentar um fundo que ir4 garantir despesas de\\
desmantelamento e abandono de pogos e instalagGes petroliferas. Toda a informacao\\
sobre abandono e desmantelamento de pogos e instalacGes petroliferas esta escrita de\\
forma detalhada no Diario da Republica de 10 de Abril de 2018, I Série - N° 46.\\
\\
Neste Diario, est4 detalhado o plano de desenvolvimento, principios gerais, plano de\\
abandono previsional, plano de abandono definitivo ou permanente, abandono de pocos\\
e abandono de instalagées petroliferas.\\
\\
O despacho presidencial publicado dia 10 de Abril no Diario da Republica estabelece\\
que os fundos de abandono e desmantelamento sejam constituidos pelas empresas ope-\\
radoras, depositados em contas abertas em nome da concessionaria, que apenas podera\\
\\
utilizar esse dinheiro no final da concess4o, para pagar as operagées de desmantelamento\\
dos blocos.\\
\\
2.1.8 Custos de exploracdo, desenvolvimento e producao de\\
campos\\
\\
40 anos.\\
\\
O ciclo de investimento de um campo convencional normalmente varia de 15 a\\
plorar ¢\\
\\
Isto é desde o momento em que o operador come¢a a obter os direitos de ex\\}\\
erfurar, até o desmantelamento do ativo como ilustrado na Figura 2.4.\\
\\
Como sabemos, o desenvolvimento de um campo é um empreendimento de\\
custo, Principalmente em Areas offshore, e as incertezas sio particularmente altas durante\\
4 aploracio ¢ avaliac¢io. Um investimento de capital tipico para um campo paraliecn\\
offhore de tamanho médio (ou seja, reservas de 100 Mbbl) seria da ordem de USS | As\\
a Na Figura 2.7 esta ilustrado o investimento tipico para campos convencionais. 4\\
\\
Borias de custos envolvidos sio:\\
\\
alto\\
\\
cdo e\\& avaliagao, sao cus\\
\\
"Os A s de explor\\
Custos da descoberta: inclui os investimentos de explora 0 de poses\\
\\
‘0S para a Tealizaco de estudos geologicos ¢ sismicos € para 4 pe rfuras\\
\\
Digitalizada com CamScanner\\

\\section*{Pagina 59}

IA PETROLIFERA\\
\\
38 CADEIA DE VALOR DA INDUSTR\\
\\
Figura 2.7: Investimento tipico para campos convencionais.\\
\\
de exploracao. Eles também sao obrigados para delinear e testar o reservatério, an-\\
tes de tomar a decisdo de desenvolvimento;\\
\\
* Os custos de desenvolvimento: também denominados Capital Expenditure/Despesr\\
de Capital (CAPEX), so os custos de desenvolvimento do campo, incluindo a per\\
furacao dos pocos de exploracao e desenvolvimento, e a constru¢ao das instalagées\\
de superficie e pipeline necessarias para colocar 0 campo em opera¢ao;\\
\\
* Os custos de produgao: sao os custos fixos e variaveis anuais necessarios para getit\\
as operacGes de producao;\\
\\
Embora a maioria dos custos possa ser considerada Operational Expenditure/ Despes\\
\\
Operacionais (OPEX), os campos de producdo também esto sujeitos a despesas de brow\\
\\
field, que sao consideradas CAPEX.\\
\\
oat See e ae de perfuracio de pocos sfio custos elevados oe\\
\\
Dentantaean Artie pa Entre os custos de exploracio estao o apoio logistic\\
\\
de perfuragio, cabeca de so € Montagem (DMM), revestimento, cimentacio, fut\\
, cabeca de poco, brocas e outros custos, como ilustrado na Figura ase\\
\\
2.1.9 Dif\\
pA idleliad entre desenvolvimento de campos cone\\
nals € néo-convencionais\\
\\
Na Figura 2.1 estdo ilus\\
no-convencionais, Em Hisham ‘apas de desenvolvimento de campos conv’\\
podem ser amplamente UMbesiicada s Estados Unidos da América, os ativos 4\\
\\
enciom®*\\
e Ups\\
sede\\
15s.\\
em trés tipos de desenvolvimento de dleo\\
\\
Digitalizada com CamScanner\\

\\section*{Pagina 60}

SEGMENTOS DE MIDSTREAM — DOWNSTREAM 39\\
\\
DMM — -Revest. Fluido—-Brocas\\_ Cabecade Apolo Sondac\\
Poco Log. Equipe\\
\\
Figura 2.8: Custos de sondas e equipamentos de perfuracio.\\
\\
6-7 US\\$M/poco 5= 10 US\\$M/pogo 30 - 50 US\\$M/poso lguas rasas)\\
100 - 150 US\\$M/poso [équas.\\
profundas) !\\
\\
Figura 2.9: Custos de perfuragao de pocos.\\
\\
* Campos onshore nao convencionais, incluindo leo de xisto e gas de xisto, com uma\\
participacao de 50\\% na produgao nos EUA ¢ uma participacao de 18\\% no mundo;\\
\\
* Desenvolvimentos no exterior, representando 29\\% da producao.\\
\\
Na Tabela 2.5, ilustramos a comparagao entre desenvolvimentos convencionais ¢ de xisto.\\
Ao contrario dos recursos convencionais, que residem em reservat6rios altamente poro-\\
\\$08 e permeéveis, nos quais 0 petréleo pode fluir e podem ser facilmente aproveitados por\\
Po¢os verticais padrao, os recursos de xisto permanecem presos em sua rocha original, °\\
xisto rico em organicos que se formou a partir da deposicao sedimentar de lama, silte,\\
argila © matéria organica no fundo dos mares rasos.\\
\\
2.2 Segmentos de Midstream e Downstream\\
ates Segmentos, como mencionamos na se¢io introdutéria deste médulo, se os -\\
sancarTegam do transporte ¢ refinagio do petréleo bruto (Midstream), distri Oe de-\\
‘omercializacio de derivados do petrdleo (Downstream). Estes dois segmentos ser\\
Se] a\\
\\
"wolvidos de forma detalhada no capitulo 8.\\
\\
ns J\\
\\
Digitalizada com CamScanner\\

\\section*{Pagina 61}

40 CADEIA DE VALOR DA INDUSTRIA PETROLIFERA\\
\\
gis convencional e gas de xisto.\\
\\
‘Tabela 2.5: Principais diferengas entre\\
Xisto\\
\\
a Convencional\\
Paw Descrigio\\
atores IC: Requer migra¢do através da Nao requer migracio, pos\\
e-\\
\\
Migragio de hidrocarbon' rocha geradora para 0 reser- os hidrocarbonetos estio 4,\\
tol \\& mieeadiatio 2 SPFS™ yatorio onde sio armazena- mazenados na rocha ger.\\
Geolégicos namento A dora\\
Ivo para perfuragao Rocha reservat6rio Rocha geradora\\
Zona alvo\\
do Norte: Porosidade baixa: 4 a 15\\%\\
\\% Campos do mar\\
Porosidade (\\%) 20a 25\\%\\
ili- Permeabilidade ultra 6:\\
\\
‘li 5 Expresso em Mili-Darcy axa\\
\\
Permeabilidade (darcy-D) ao3) expressa\\_ em Nano-Darcy\\
\\
o-\\%)\\
\\
Estimulagio (fraturamento Somente em alguns casos, Sim, cerca de 100\\% de pocos\\
hidrdulico) para ser comer- menos de 20\\% de pocos con- de xistos devem ser fratura-\\
\\
J Ot al vencionais sio fraturados dos\\
comerciais\\
\\
Perfil tipico do poco produ- Plateau: A produgao se es- Declinio imediato: Pico de\\
tabiliza no plateau por al- producao nos primeiros trés\\
guns anos, depois um decli- meses e depois um acentu-\\
nio constante ado declinio\\
\\
tor\\
\\
Numero de novos pogos Baixo. Vida util de pocos-al- Muito elevado - conceito\\
Para manter o nivel de gumas décadas. Reservas u- de “fabrica de perfuracio”.\\
\\
producio Picas de 50 - 300 Bcf/poco e Vida util de poco - alguns\\
0,1 - 2 Mbb/poco anos. Producdo acumu-\\
lada de 0,5 Bcf/pogo e 150\\
kbbl/poco\\
Habilidades de planejar o Sim N§o - Aprende e ajusta 20\\
desenvolvimento de campo perfurar\\
com antecedéncia\\
\\
Fonte: BCG experience,\\
\\
Conclusaéo\\
Neste capitulo, abordam\\
Posicio dos ‘08 a cadeia de valor do setor petrolifero, Abordamos que 4 ¢e"\\
\\
Downst talhes - perreall\\
. Team serio desenvolvides .- médulos subsequentes. Os segmentos de Midstre\\#\\
no Wie\\
Capitulo 8; Midstream e Downstream.\\
\\
Digitalizada com CamScanner\\

\\section*{Pagina 62}

CAPITULO 3\\
\\
EXPLORAGAO DE HIDROCARBONETOos\\
Se a eae,\\
\\
Neste capitulo, abordaremos inicialmente as condi¢ées necessérias Para a existéncia de\\
acumula¢ao de hidrocarbonetos. Antes disso, faremos uma revisio dos conceitos da for-\\
macio dos continentes ¢ das placas tecténicas, bem como dos tipos de rochas existentes\\
que sao de interesse para a engenharia de petréleos. As técnicas utilizadas para a locali-\\
zacio de hidrocarbonetos serao apresentadas de forma resumida. £ importante ressaltar\\
que as atividades de exploracao tém como objetivo descobrir novos volumes de hidro-\\
carbonetos, substituindo aqueles que estao sendo Produzidos. O sucesso dos esforcos\\
de exploracao de uma empresa determinara suas perspectivas de permanéncia no nego-\\
cio a longo prazo. Por esse motivo, muitas empresas investem bastante na formacao de\\
profissionais para as atividades de exploracao.\\
\\
O objetivo geral deste capitulo é assegurar que o leitor seja capaz de identificar inicial-\\
mente as condi¢des necessérias para a existéncia de acumulacio de hidrocarbonetos. Os\\
objetivos especificos sio os seguintes:\\
\\
* Adquirir conhecimentos basicos de geociéncia relacionados com a crosta terrestre,\\
teoria das placas tecténicas e morfologia dos fundos oceanicos;\\
\\
* Identificar os diferentes tipos de rochas e quais sao aquelas que armazenam os hidro-\\
\\
carbonetos;\\
\\
* Conhecer os fatores que contribuem para a formacio de uma acumulacio de hidro-\\
\\
carbonetos;\\
\\
* Conhecer as técnicas geoldgicas e geofisicas utilizadas para localizar depdsitos de\\
Petréleo gis natural.\\
\\
31 Teorias das placas tecténicas\\
\\
‘ oe , Origindria da palavra grega TEKTOVIKOG, “relativo 4 -\\_ spay\\
nance geologia que descreve os movimentos de grande escala que 0¢¢\\
tosferg terrestre\\
os | 41\\
na\\
1) Cera Para ndo-engenheivos, 1* edigio\\
Aniet Repos Ramos Copyright © 2028 Bdnora Do's)\\
\\
Digitalizada com CamScanner\\

\\section*{Pagina 63}

az expLonacAo DE HIDROCARBONETOS\\
\\
a ‘om 1S\\
as tectonicas, a parte mais externa \\#8 bipsichd “ re per das an,\\
° . , lidificada na pal erna do mango\\
sali ue inclui a crosta ¢ a , : ‘ on\\
poles ende a parte mais interna viscosa do manto. Em uma a\\
re ce se comportar como um liquido Superagiy\\
Ie.\\
\\
a astenosfera, que COMP \\textasciitilde{}—\\
ilhd nto\\
\\
temporal de milhoes de anos, o manto P: est eonuporti-ee Co orn et\\
\\
lo\\
\\
cido, mas em resposta 4 forgas subitas,\\
\\
Na teoria das plac\\
\\
lacas tectonicas, que se deslocam sobre a ate\\
\\
ear" varias pl\\
tosfera esta fragmentada em Ao de dois fendmenos geoldgicos distintos\\
\\
Ali\\
nosfera. Essa teoria surgiu ap «tig do século XX por Alfred Wegener, e a im\\
\\
i i 1, jdentificada ;\\
‘te tut soe ocelnic, detectada pela primeira na are he ” A teotia en\\
si foi desenvolvida no final dos anos 60 POF Rober? Pal Fe ‘in i onzie ¢ desde\\
entao tem sido universalmente aceita pelos cientistas, tendo revolucionado as Ciéncias d,\\
ants das camadas exteriores da Terra em litosfera astenosfera baseia-se em suas\\
diferencasmecanicas. A litosfera é mais fria € rigida, enquanto a astenosfera é mais quente\\
¢ mecanicamente mais fraca. Essa diviséo nao deve ser confundida com a subdivisio\\
quimica da Terra, do interior para a superficie, em nucleo, manto e crosta.\\
\\
O principio chave da tectonica de placas é a existéncia de uma litosfera constituida\\
\\
placas tectonicas separadas € distintas, que flutuam sobre a astenosfera. A relativa\\
fluidez da astenosfera permite que as placas tectonicas se movimentem em diferentes di-\\
regdes. As placas entram em contato umas com as outras a0 longo dos limites de placa\\
(Figura 3.1), frequentemente associados a eventos geolégicos como terremotos e \\& forma\\
lo de caracteristicas topograficas como cadeias montanhosas, vulcées € fossas oceanicas.\\
‘A maioria dos vulc6es ativos do mundo encontra-se a0 longo dos limites de placas, sendo\\
is conhecida e ativa. Esses limites\\
\\
a regiio do Circulo de Fogo do Oceano Pacifico a mais\\
serao detalhados mais adiante.\\
\\
Digitalizada com CamScanner\\

\\section*{Pagina 64}

TEORIAS DAS PLACAS TECTONICAS 43\\
icas comum a todas elas é a camada s6lida superior do manto, situada sob\\
oe ocntal e oceanica, constituindo, juntamente com a Crosta, a litosfera, Sob as crostas\\
\\textasciitilde{} distingdo entre crosta continental € crosta oceanica baseia-se na diferenga de densida-\\
des dos materiais que constituem cada uma delas; a crosta oceanica é mais densa devido\\
js diferentes proporgoes dos elementos Constituintes, em particular do silicio, A crosta\\
oceanica mais pobre em silica e mais rica ¢m minerais maficos, que geralmente so mais\\
densos, enquanto a crosta continental apresenta uma maior percentagem de minerais fél-\\
sicos, que geralmente sao menos densos.\\
Como consequéncia, a crosta oceanica geralmente esta abaixo do nivel do mar,\\
maior parte da placa do Pacifico, enquanto a crosta continental situa-:\\
conforme principio da isostasia elaborado por John Pratt,\\
\\
comoa\\
se acima desse nivel,\\
\\
3.1.1 Acrosta terrestre\\
\\
Acrosta terrestre (Figura 3.2) foi concebida pelo meteorologista alemao Alfred Wegener,\\
que afirmou que ha aproximadamente 200 milhdes de anos nao havia separacSo entre\\
os continentes. Ou seja, existia uma nica massa continental chamada Pangeia e um\\
tinico oceano, o Pantalassa. Ao longo de milhGes de anos, ocorreu a fragmentacio desses\\
continentes, resultando em dois megacontinentes chamados Laurisia e Gondwana. A\\
partir desse ponto, os continentes comecaram a se mover e se adaptar as configuracdes\\
atuais.\\
\\
O ponto crucial para o desenvolvimento da crosta terrestre, que em esséncia envolve\\
0 movimento dos continentes, ou seja, placas que se movem, é a constatacdo de que a\\
Terra nao é estatica. Wegener observou que a costa da Africa tinha um contorno que\\
se encaixava na costa da América do Sul. Outra evidéncia que corroborou com a teoria\\
foi a descoberta de fosscis de animais da mesma espécie em continentes diferentes. Se-\\
ria impossivel que esses animais tivessem atravessado o Oceano Atlantico, sendo a unica\\
explicacao plausivel de que os dois continentes ja estiveram juntos no passado.\\
\\
Pangeia Laurdsia © Godwana Mundo Modemo\\
\\
Figura 3.2: Deriva continental.\\
\\
3,\\
12 Tipos de limites de placas\\
\\
— ttés tipos de limites de placas, caracterizados pelo modo como as placas we i”\\
perf ©m relacdo as outras. Cada tipo esta associado a diferentes fenome|\\
\\
 —\\textasciitilde{}— \\_\\
\\
Digitalizada com CamScanner\\

\\section*{Pagina 65}

ON\\
\\
rem quando as placas q\\
\\
44 EXPLORACAO DE HIDROCARBONETOS,\\
rvativos Ocor! cali\\
* Limites transformantes ou ag eiongotelfillias transformantes, © ou,\\
i i a na ou Sad - Ong\\
cisamente, rogam um dirs Ot Pata a eaten\\
a id lativo das duas placas pode ser para a ervador aicinn md “a\\
dendo . oman 3 direita ou a esquerda de um obs 40 longo \\&\\
lendo se o\\
falha.\\
\\
as placas se af;\\
Limites divergentes ou construtivos surgem quando duas p astam uma dy\\
= Limites\\
outra.\\
\\
ce rgentes ou destrutivos ocorrem quando duas placas se Movem em,\\
“Bee ch 8 ais formando uma zona de subduc¢ao (quando uma das placas\\
Rs Satta,\\
fests es outra) ou uma cadeia montanhosa (quando as placas Colidem e ®\\
mergulha sob a\\
comprimem uma contra a outra).\\
\\
Existem limites de placas cuja situacao é mais complexa, especialmente NOS casos em que\\
és ou mais placas se encontram. Nesses casos, ocorre uma combinacio dos trés tipos\\
\\
trés\\
\\
de limites mencionados anteriormente.\\
\\
3.1.3 Teorias das placas tecténicas\\
\\
entanto, a prioridade deve ser atribuida a Hes,\\
nao publicado de seu artigo de 1962 em 1960\\
Moviam através da crosta oce:\\
\\
‘Anica (como P\\
bacias ocednicas e Os continentes adjacentes\\
crustal ou placa.\\
\\
s, ja que ele compartilhou um manuscrito\\
). Hess sugeriu que os continentes nio s¢\\
toposto pela deriva continental), mas que as\\
Moviam-se juntos como uma tinica unidade\\
\\
€poca (e até ridicularizada), essa Publicag: i\\
\\
damental e Perspicaz, Em 1967, Jason Morgan Propés que a superficie da Terra consis\\
em 12 Placas rigidas que se moviam em elaco umas as outras, Dois meses depois\\
1968, Xavier Le Pichon apresentou um modelo completo baseado em 6 placas princips\\
\\
com seus movim, .\\
eatos relativos, figura 3.3 ilustra as placas tectonicas.\\
\\
3.1.3.1 Classificacao das Placas\\
\\
Existem 55 Placas tecténj\\
\\
ras i ‘ a “f menor\\
Essas placas estig Tesumidas nese Quais 15 sao principais ¢ 40 sao\\
le\\
\\
3.2,\\
\\
Digitalizada com CamScanner\\

\\section*{Pagina 66}

TEORIAS DAS PLACAS\\
\\
Te CTONICAS, 45\\
\\
Figura 3.3: Placas Tecténicas.\\
\\
Tabela 3.1: Placas principais.\\
\\
1. Placa Africana 6. Placa de Cocos 11. Placa de Nazea\\
2. Placa da Antirtida 7. Placa Buroasistica 12, Placa Notte Americana\\
3. Placa Arabica 8. Placa das Filipinas 13, Placa do Pacifico\\
\\
4, Placa Australiana 9. Placa Indiana 14, Placa de Scotia\\
\\
5. Placa das Carafbas 10. Placa Juan de Fuca\\
\\
5, Placa Sul-americana\\
\\
Tabela 3.2: Placas menores.\\
——E———— ee\\
\\
1, Placa da Anatélia 11, Placadas Galapagos 21. Placa do Mar Egeu 31. Placa do Recife de\\
Conway\\
2. Placa do Altiplano 12. Placa de Gorda 22. Placa do Mar das Mo- 32. Placa de Rivera\\
lucas\\
3. Placa de Amur 13. Placa Helénica 23, Placa do Mar de Salo- 33. Placa das Sandwich\\
mio\\
4, Placa dos Andes Norte 14. Placa Iraniana 24, Placa das Novas Hébri- 34, Placa das Shetland\\
das\\
5. Placa da Birmania 15, Placa Juan Fernandez 25. Placa de Niuafo’ou 35, Placa da Somilia\\
6. Placa de Bismarck 16. Placade Kermadec 26, Placa de Okhotsk 36, Placa de Sunda\\
Norte\\
7. Placa de Bismarck Sul 17. Placa Manus. 27, Placa de Okinawa 37. Placa de Timor\\
8. Placa da Carolina 18. Placa Maoke 28. Placa do Panama 38. Placa de Tonga\\
9. Placa de Doberai 19. Placa das Marianas —29. Placa de Pascoa 39, Placa de Woodlark\\
10. Placa Futuna 20. Placa do Mar de Banda 30, Placa do Recife de Bal- 40. Placa do Yangtzé\\
moral\\
\\
ee\\
\\
3.1.4 Morfologia dos fundos oceanicos\\
\\
No inicio do século XX, uma tecnologia revolucionaria comegou a se desenvolver, i\\
\\
Pactando tanto o campo bélico quanto o conhecimento cientifico. Essa tecnoll\\
\\
ean que desempenharam um papel significativo durante a Segund\\
\\
logia era a\\
a Guerra\\
\\
a, era necessanio\\
\\
com segurang,\\
8 1920, come:\\
\\
artir de\\
anos, A parur ¢ ;\\
Jo submarine\\
\\
Noentanto, Para que os submarinos pudessem navegat\\
"um conhecimento mais detalhado do fundo dos oce\\
\\
Garay i ‘nvestigac:\\
4 Ser utilizadas sondas eletroacusticas (Figura 3.4) para a invesBas\\
\\
Digitalizada com CamScanner\\

\\section*{Pagina 67}

46 —\\_EXPLORACAO DE HIDROCARBONETOS\\
\\
assonicas que reflete,\\
Essas sondas, conhecidas como sonares, usam ondas ultr: q ™ No fy,\\
\\
refletida reto 1\\
do mar, permitindo registrar 0 tempo que leva para cco 3 finds da ae ie 4,\\
© tempo que as ondas ultrassénicas demoram [sake eorindidade, malop sei POstiygs\\
calcular as profundidades oceanicas. Quanto maior P O temp,\\
\\
My\\
que as ondas levam para retornar ao navio-\\
\\
Figura 3.4: Sondas eletroacisticas para investigacao submarina.\\
\\
Além dos sonares, outros aparelhos sao utilizados, como sismégrafos, magnetémetrose\\
veiculos submersiveis (pequenos submarinos que podem ser tripulados ou controlados\\
remotamente (Figura 3.5)).\\
\\
Figura 3.5: Algumas tecnologi i - .\\
Revista Cientifica), utilizadas para exploracéo do fundo oceanico (Fonte\\
\\
obtidas i ae\\
\\
Pas e cartas detalhadas que Por esses diversos equipamentos, é possivel criat -\\
‘opografia dos fundos Oce4nicos, assim com?\\
\\
3.6 ilustra algumas das estruturas en"\\
\\
dade nao \\& Muito grande (as ioe do continente s0b 0 Oceano. Sua profunt\\
tos transpo © Possui fo suave. B coberta pot areias "\\
\\
tes, =a\\
Por exemplo, Pelos rios desde o interior dos on™\\
\\
Digitalizada com CamScanner\\

\\section*{Pagina 68}

CONCEITO Das ROCHAS 47\\
\\
Figura 3.6: Relevo submarino dos oceanos\\
\\
Otalude continental possui um declive acentuado, estendendo-se desde a plataforma con-\\
tinental até o fundo do oceano (pode chegar a até 1000 metros de profundidade). Tem a\\
aparéncia de uma parede vertical (muro = talude) e marca a transi¢\\&o geolégica entre o\\
continente e 0 oceano.\\
\\
A planicie abissal € a regio do fundo oceanico Propriamente dito. Apresenta uma\\
\\
aparéncia geralmente plana (dai o termo “planicie”) e pode se estender até profundidades\\
de 6000 metros. Essa area é coberta por sedimentos muito finos ou até mesmo rochas\\
nuas.\\
A dorsal médio-oceanica € uma vasta cordilheira de montanhas submarinas, com o\\
vale do rifte no seu centro. A dorsal médio-atlantica (localizada no meio do oceano Atlan-\\
tico, entre as Américas e a Europa/ Africa) é a maior cordilheira submarina do planeta,\\
com cerca de 65.000 quilémetros de extenso. Se nao fosse o oceano, seria visivel do\\
espago.\\
\\
O rifte situa-se no centro da dorsal. Trata-se de uma fissura através da qual o magma\\
€ expelido do interior da Terra.\\
\\
As fossas oceanicas, como o nome sugere, s4o depresses profundas e estreitas, com\\
Paredes ingremes, geralmente atingindo profundidades entre 3000 e 5000 metros.\\
\\
3.2 Conceito das rochas\\
\\
a mais externa\\
ompostas por\\
10 é 0 caso do\\
sta terrestre\\
\\
Astochas sio agregados de minerais sdlidos que constituem a camada solid\\
4a Terra, conhecida como litosfera. Elas podem ser formacées naturals ce\\
‘eis ou mais minerais agregados, ou até mesmo por um tinico mineral, com\\
“lcirio€ do arenito, Existem também rochas que cobrem vastas areas da cro’ reso\\
¢ imbora mais raras, hd aquelas compostas por um unico mineral. A parrriettd i”\\
‘tra que caracterizam uma rocha so determinadas pela sua origem geolos\\
\\
*S pela maneira ¢ pelo local onde se formou.\\
\\
Digitalizada com CamScanner\\

\\section*{Pagina 69}

48 EXPLORAGAO DE HIDROCARBONETOS\\
\\
3.2.1 Tipos de rochas\\
nos considerar\\
rficas ¢ rochas\\
\\
s de rochas: rochas igne.\\
a\\
\\
trés tipos basicos der\\
mo ilustramos na Tabe|\\
\\
sedimentares, co\\
\\
Quanto A sua origem, poder Soy\\
amo! aay\\
\\
magmiticas, rochas met\\
‘o das rochas.\\
\\
Tabela 3.3: Classificaga\\
‘Sedimentares Metambtficas\\
\\
Origem ¢ Processo Magmitica\\
spacso ¢ erosao das Rochas sob alta\\
Fusio de rochas na crosta Meteorizaga o S presses\\
Origem do Material pai ‘ada ¢ quente ¢ 10 rochas da superficie ¢ altas temperaturas na\\
manto superior crosta. profunda ¢ no\\
\\
manto superior\\
Processo geologico Cristalizagio \\_(Solidifica- Deposicio e compactagzo Recristalizagio de novos\\
ao do magma) minerais no estado s6lido\\
\\
le transformacao das rochas\\
\\
As rochas estao em constante transformacao, devido aos processos orogénicos internos e\\
externos da Terra. Os trés grupos de rochas— magmiaticas, sedimentares e metamérficas\\
—transformam-se continuamente por meio de um conjunto de processos geolégicos de.\\
\\
ela primeira vez em 1785 pelo escocés\\
\\
nominado Ciclo das Rochas. Este ciclo foi descrito p\\
James Hutton, em uma apresenta¢ao oral perante a Royal Society of! ‘Edimburgo (Figura 3.7)\\
\\
3.2.2 Ciclo d\\
\\
meteoitzagio ¢ erosto\\
\\
nocnas\\
SEDIMENTARES\\
Pec de\\
prasad mp eratura\\
ics\\
| IETAMORFICAS:\\
—.\\
tuslo\\
\\
Uaeua\\
\\
Figura 3,7: Ci\\
: Ciclo de transformacao de rochas.\\
\\
fecimen!? ©\\
1\\
\\
Digitalizada com CamScanner\\

\\section*{Pagina 70}

CONCEITO Das ROCHAS 49\\
\\
om complexa de material rochoso em estado de fusio,\\
ums tos quimicos, incluindo silicatos, silicio e elementos volateis, como P\\
\\
poo hidrogénio, fluor e outros. ‘ "Por de gua,\\
“ae varios tipos de magmas que, de acordo com sua com\\
pos da quantidade de silica (SiOz) e de gases dissolvidos, sio de:\\
\\
B formado Por diversos\\
\\
POSi¢40 quimica em ter.\\
‘Signados por:\\
, basiicos: poles: en silica ane 7 gases dissolvidos; sao expelidos ao longo\\
= riftes, na dorsal média oceanica, e dos pontos quentes, dando origem a rochas\\
como 0 basalto e o gabro;\\
\\
Magmas andesiticos: tém composicao intermediaria em silica (SiO2) e bastantes\\
dissolvidos; sao expelidos ao longo das zonas de subduccao, dando origem a rochas\\
como 0 andesito e 0 diorito.\\
\\
Magmas rioliticos: ricos em silica (SiOz) e gases dissolvidos; estao associados a fusio\\
parcial das rochas constituintes da crosta continental, que ocorre aquando da colisio\\
entre duas placas continentais, dando origem a rochas como 0 riolito eo granito.\\
\\
As rochas magmiaticas, também conhecidas como igneas (Figura 3.8), séo formadas pela\\
\\
solidificagao do magma expelido por vulcées e podem ser subdivididas em dois tipos:\\
\\
Rochas intrusivas: sao as rochas formadas pelo magma que se solidificou em grandes\\
profundidades. O granito é uma das variedades desse tipo de rocha.\\
\\
Rochas extrusivas: sao as rochas formadas pelo magma que se solidificou na superficie.\\
O basalto é uma das variedades desse tipo de rocha.\\
\\
Figura 3.8: Rochas igneas ou magmiticas.\\
\\
3.2.3.1 Cristalizagéo do magma\\
\\
Pr magma rre: ade 4 aco\\
da fe que o vai arrefecendo, outro fendmeno pode ocorrer dev ido \\& wie\\
de ase dual —a diferenciagao gravitica. Nesse processo, os eersd se b pron\\
1 ling a+ Bravidade — a dife:\\
i istai Jam de acordo ¢\\
* ordem de base em sua densidade. Os cristais se acumul\\
\\
P di\\
Assim, formagio ¢ também tendem a se acumular de acordo com sua ines os\\
Meno, Grstais mais densos que 0 magma residual se deslocam para 0 fundo, enquan\\
08 densos se deslocam para 0 topo.\\
\\
Digitalizada com CamScanner\\

\\section*{Pagina 71}

50 EXPLORACAO DE HIDROCARBONETOS\\
\\
Ss\\
has metamérfica .\\
3.2.4 Roc “fico” tem su origem na lingua ‘Brega, derivan\\
lo vocabulo “metam ‘ combinaga0 de “meta Aque significa \\% além de\\
morphoo Bn forma”). Portanto, metam6rfico” refere se\\
he” (que sign\\
\\
mudanga ou transformagao o ode pe fire “ \\textasciitilde{}\\
a algo que passou por uma interior da Terra no estado 8 » 2 alteracgg\\
chas metamoérficas formam-se no re textura de rochas preexistentes (magmiaticas, seq.\\
da composigao mineralogica f ae esas alteracdes, conhecidas como metamorfismo,\\
mentares ou rane i din micas diferentes de ala €/ou pressio em\\
ocorrem devido a con ae e formadas.\\
relacao aquelas em que as rochas foram Le ithaaei0 de importantes eventos ge.\\
O estudo das rochas metamérficas per! is para 0 entendimento da configuracig\\
\\
otecténicos ocorridos no passado, fundamentai:\\
-ontinentes. Os processos metamorficos ocorrem em wu x: 4\\
[< S. proces: ma faixa de condigi\\
atual dos if\\
\\
.amicas situada entre o ambiente magmatico, onde a fusao dos materiais rocho.\\
ee ade magmiticas, € 0 ambiente sedimentar, onde a diagénese leva a forma.\\
ao de rochas sedimentares. ; stn ‘ :\\
\\
Nas rochas argilosas, os constituintes minerals s4o os produtos mais finos do intem-\\
perismo. Essas rochas passam por reacdes e mudancas bem caracterizadas e so particu-\\
larmente adequadas para identificar os sucessivos graus de metamorfismo. As rochas sio\\
compostas principalmente de quartzo e feldspatos, que so minerais estaveis em condi-\\
ges de temperatura e pressao mais elevadas. Essas rochas sofrern mudancas que podem\\
ser dificeis de acompanhar em detalhes.\\
\\
Calcirios e dolomites impuros, quando sujeitos a variagdes de temperatura e pressio\\
em condicées de equilibrio instavel, podem converter-se em novos minerais. Durante\\
© processo metamérfico, o didxido de carbono liberado por essas rochas facilita as mv"\\
dangas mineral6gicas. Um exemplo de rocha da quarta série é 0 basalto. Os principas\\
\\
minerais presentes nas rochas magmaticas basilticas so:\\
\\
A etimologia d\\
do termo grego “metal\\
ou “mudanga”) e “morp!\\
\\
* Os feldspatos do tipo plagioclasio sédico-calcicos;\\
\\
* Os piroxénios ¢ olivinas (minerai\\
gas metamérficas, (minerais de ferro), que sao facilmente suscetiveis a muda"\\
\\
3.2.4.1 Principals fatores de metamorfismo\\
\\
tensa\\
(0, temperatura, Pressao, fluidos ¢ tempo;\\
\\
\\textasciitilde{} Litostatico: atua ; ,\\
das rochas por pee ani Uniforme, alterando principalmente a mineralod?\\
de minerais com phat arses dos dtomos e fons, resultando ma forms\\
\\
\\textasciitilde{} Nilo litostético: atu\\
\\
408 movimentos\\
\\
de man mais compactas ¢ aumento da de" é\\
ira direcionada, como no caso das forsas “os\\
ticos (compressio, distensio e cisalhament®)\\
\\
Digitalizada com CamScanner\\

\\section*{Pagina 72}

CONCEITO DAS ROCHAS 51\\
\\
tipo de tensao altera principalmente a textura das rochas ¢ Pode levar ao d\\
volvimento de foliacgao. cadaeaal\\
\\
ci ae Mas a alteragdes mineraldgicas, como crescimento de cristais e re-\\
crisalizagao, ¢ também afeta as caracteristicas texturais das rochas metamérficas, A\\
temperatura nao € suficiente Para fundir a: .\\
\\
'S tochas, mas favorece 0) imi\\
: = y Teacoes juimicas\\
entre minerais, aumentando a vulnerabilidade das rochas as Pressdes\\
\\
* Fluidos: os fluidos, devido a sua circulacio e Pressio, cat\\
gGes na composi¢ao quimica e mineralégica das rocha\\
recristalizacdo. Esses fluidos tém a capacidade de diss\\
cos dos minerais necessarios para reacées quimicas. Em sua auséncia, a velocidade\\
das reagGes é reduzida.\\
\\
‘usam principalmente altera-\\
\\
= Tempo: os processos metamoérficos ocorrem lentamente, Portanto, o tempo tam-\\
bém é um fator relevante. O tempo influencia a ob\\
\\
tencao de equilibrio nas reacdes\\
quimicas, permitindo a reorganizacao mineraldégica ea formacio de associacg6es de\\
fases em equilibrio.\\
\\
3.2.4.2 Tipos de metamorfismo\\
\\
Com base na prevaléncia e na intensidade de um ou ma:\\
\\
‘is fatores de metamorfismo, é\\
possivel definir varios tipos de metamorfismo (Figura 3.9)\\
\\
* Extensao da area afetada: segundo este critério, existem dois tipos de metamor-\\
fismo:\\
\\
\\textasciitilde{} Metamorfismo local: associado a intrusdes magmiaticas locais que ocorrem no\\
limite entre 0 magma e as rochas encaixantes. Os fatores predominantes sao a\\
temperatura e a circulacao de fluidos. Isso ocorre na zona de contato entre o\\
magma e as rochas encaixantes, formando as rochas metamorficas que consti-\\
tuem a chamada aureola metamérfica. O grau de metamorfismo dessas rochas\\
diminui com a distancia da intrusio.\\
\\
\\textasciitilde{} Metamorfismo regional: asse tipo de metamorfismo afeta areas extensas da\\
Terra, desencadeado pela aco predominante de tensdes litostaticas e nao litos-\\
\\
taticas, de moderadas a altas, bem como temperaturas e circulagio de fluidos,\\
na sequéncia de movimentos tecténicos.\\
\\
3.2.4.3 Minerais formados a partir de rochas metamorficas\\
\\
Os minerais metamérficos originam-se principalmente devido a aco da temaperarara ¢\\
\\
ene que geram uma grande quantidade de minerais, muitos dos quais dificlimente\\
— formados Por outros processos. As transformagées minerais dependem, oe oe\\
\\
net lugar, da composi¢io da rocha original e, em seguida, da natureza ou a de mete\\
\\
Ser agruy 4 que foi submetida. Do ponto de vista da composigao inicial, as rochas p\\
\\
€M quatro séries diferentes, conforme listado a seguir:\\
\\
—\\
Digitalizada com CamScanner\\

\\section*{Pagina 73}

52 EXPLORACAO DE HIDROCARBONETOS\\
\\
Mérmore Ardésla\\
\\
Figura 3.9: Exemplo de rochas metamorficas.\\
\\
= Rochas argilosas;\\
\\
* Rochas arenosas, igneas dcidas e tufos;\\
\\
= Xistos acidos e gnaisses;\\
\\
» Calcérios e outras rochas carbonatadas.\\
\\
3.2.5 Rochas sedimentares\\
\\
‘As rochas sedimentares sao frequentemente chamadas de rochas secundiarias, pois ge-\\
ralmente resultam da acumulaco de pequenos fragmentos quebrados de rochas pré-\\
existentes, Essa acumulacao é compactada e cimentada. Existem trés principais tipos\\
de rochas sedimentares: rochas sedimentares clasticas, quimicas e organicas.\\
\\
Rochas sedimentares clasticas: sao formadas a partir do actumulo de clastos — peque-\\
nos fragmentos de rocha quebrada que foram compactados e cimentados. Exem:\\
plos de rochas sedimentares clasticas incluem arenito, xisto, siltito e conglomerado\\
(Figura 3.10).\\
\\
Rochas sedimentares quimicas;\\
\\
yepa Para tras os mine:\\
‘™ se formar devido a\\
\\
sh delas se formam quando a agua parada evapo\\
inundacses Depésitos espessos de sal ¢ gel\\
© ¢vaporacdo repetidas por um long? P™\\
\\
Digitalizada com CamScanner\\

\\section*{Pagina 74}

GEOLOGIA DE p\\
” ETROLEOS 5\\
3\\
\\
Gesso (Gypsum) Dolomita\\
\\
Figura 3.11: Rochas sedimentares quimicas.\\
\\
de tempo. Exemplos de rochas sedimentares quimicas incluem\\
\\
gesso e dolomita\\
(Figura 3.1 1).\\
\\
Rochas sedimentares organicas: resultam do actmulo de detritos sedimentares causa-\\
dos por processos organicos. Muitos animais utilizam calcio para conchas,\\
edentes. Esses fragmentos de cAlcio podem se acumular no fundo do mai\\
compactar em uma camada espessa 0 suficiente para formar uma rocha sedi\\
“organica”. Exemplos de rochas sedimentares organicas incluem carvio e\\
(Figura 3.12).\\
\\
Ossos\\
re se\\
imentar\\
coquina\\
\\
Figura 3.12: Rochas sedimentares organicas.\\
\\
3.3 Geologia de petrdleos\\
\\
Ageologia do petréleo éa aplicagdo conjugada da ciéncia e da técnica da geologia, devida-\\
mente assistida por outras ciéncias e técnicas conexas, na busca de hidrocarbonetos. Isso\\
abrange aspectos de prospecc4o, explorac4o, uso e também estudos e pesquisas sobre a\\
natureza, a origem e a composi¢ao dos diversos tipos de hidrocarbonetos (Figura 3.13).\\
\\
3.3.1 Origem do petréleo\\
\\
tiem ea acumulagio do petréleo revestem-se de grande importincia pritica, pois\\
Pt ivamente para orientar os trabalhos de pesquisa e prospec¢io ne —\\
\\
der PrOdusio 0u jazigos com interesse econémico. A refinacdo aborda as uae\\
\\
lerado to do petrdleo e os produtos finais. Finalmente, ¢ abordado o beret .\\
mda Petréleo. Ha um alerta para o perigo da escassez ou esgotamento do petroleo,\\
4 ” Tecursos nao renovaveis da Terra, pelo menos em “tempo util Ap\\
in 8 lati Petrus (pedra) e oleum (6leo), significando “bleo do Pate ue\\
\\
* gua, com NO estado liquido ¢ uma substancia oleosa, iniaradral, mae et\\
\\
cheiro Caracteristico e cor que varia entre O negro e\\
\\
A palavra\\
\\
Digitalizada com CamScanner\\

\\section*{Pagina 75}

54 —\\_ exPLORACAO DE HIDROCARBONETOS\\
\\
petréleo: 1 - Aerofotogrametria; 2 - Terramotg\\
\\
Figura 3.13: Técnicas para 4 busca de\\
artificial; 3 - Satélite.\\
: sair do subsolo. E considerado “con.\\
i ni » quando é negro e viscoso a0 ‘\\
peste We’ as e volatil (tendéncia a evaporar-se). A Figura 3.14 demonstra 9\\
lensado” qui\\
posicionamento dos fluidos no subsolo.\\
\\
Figura 3.14: Posicionamento dos fluidos no subsolo.\\
\\
O petréleo pode ser encontrado em diferentes estados: liquido (6leo), gasoso (gas natu-\\
ral), sdlido (asfalto) e semissélido (betume). Ele se origina da decomposi¢ao nao oxidante\\
de matéria organica armazenada em sedimentos, que migra através de aquiferos e fica\\
\\
aprisionado em reservatérios.\\
\\
3.3.1.1 Teorias sobre a formacao de petréleo\\
\\
nie — teorias sobre as origens dos combustiveis fosseis, Neste trabalho, vamos\\
Bern: Processos diferentes: Processo de metano termogénico, biogénico e abiog\\
\\
nico. A teoria mais aceita é a teoria do metano termogénico.\\
\\
Metano termogénico\\
\\
Esta teoria é ms a\\
¢ plantas) é ere oman de dleo, onde a matéria organica (restos de animals\\
aferraaalta pressio ao longo de milhées de anos. O metant?\\
\\
y\\
\\
Digitalizada com CamScanner\\

\\section*{Pagina 76}

GEOLOGIA DE PETROLEOS,\\
\\
oor terrestre, que quebram ou decomps\\
a - ‘ i, t\\
\\
Quanto maior a profundidade abaixo da superficie terrest\\
maiores volumes de gas natural sao formados emc\\
\\
Te, maior é a temperatura, ¢\\
maiores volumes a baixas temperaturas e profund\\
\\
Omparacio com o éleo,\\
idades (Figura 3.15), rane sPresenta\\
\\
Figura 3.15: Maturacao da matéria organica.\\
\\
O processo de craqueamento do querogénio (cozinha de\\
com o aumento da temperatura e da profundidade.\\
hidrocarbonetos, eliminando as moléculas longas de o:\\
nas moléculas de hidrogénio e carbono. Esse process\\
pode ser dividido em trés fases:\\
\\
Seracdo ou cook) é mais intenso\\
O craqueamento térmico produz\\
xigénio e nitrogénio, deixando ape-\\
‘0 de conversio da matéria organica\\
\\
Diagénese: é um processo que ocorre em baixas temperaturas, até aproximadamente\\
65 °C, quando a matéria organica se transforma em querogénio devido a acio de\\
\\
bactérias. Apenas o gas bioquimico é formado nesse estagio (metano biogénico ou\\
“gas de pantano”).\\
\\
Catagénese: é considerado o primeiro estagio termoquimico, no qual ocorre a degra-\\
dacao do querogénio, dando origem a hidrocarbonetos liquidos e também a “gas\\
umido”. A média deste estagio varia de 65 °C a 165 °C.\\
\\
Metagénese: neste estagio termoquimico, ocorre a quebra de todas as moléculas dos\\
hidrocarbonetos liquidos presentes, e grande parte do querogénio remanescente se\\
transforma em gas metano. Ainda nessa fase, ocorre a geracio de “gas seco”. A\\
média de temperatura desse estagio varia de 165 °C a 210 °C.\\
\\
| Metano biogénico ou biogenético\\
\\
O gis natural pode ser formado a partir da transformagio da matéria organica por micro-\\
Organismos. Metanogénicos, pequenos micro-organismos, produzem metano através da\\
decomposico quimica da matéria organica. Esses micro-organismos sao frequentemente\\
encontrados em dreas proximas a superficie terrestre, onde nao ha oxigénio, como nos >\\
‘estinos de muitos animais, incluindo seres humanos. O metano pr‘ oduzido geralmente®\\
liberado na atmosfera. No entanto, esse metano pode ser aprisionado abaixo da supertici\\
terrestre ¢ Tecuperado como gas natural.\\
\\
Um exemplo de metano biogenético é o gis de aterro: os res\\
“anitétios produzem um volume considerdvel de gis natural d\\
\\
\\{duos presentes em aterros\\
evido 4 decomposicio dos\\
\\
Digitalizada com CamScanner\\

\\section*{Pagina 77}

“cmmen \\_,\\
\\
56 EXPLORAGAO DE HIDROCARBONETOS\\
\\
coletado usando novas tecnologias\\
\\
iais orginicos. Esse gas pode ser qUe pe\\
materiais orginicos, Esse gas f a ds pas neural, Pet\\
\\
5 Mey,\\
seu aproveitamento e inclusaio na ofert: mh\\
Metano abiogénico\\
\\
Nas profundezas da subsuperficie, existem gases Sy Peat de hidro, nj tm\\
b A medida que esses gases migram para a superticie, i ‘agem com minersj\\
\\
nies fa cabeanece i: igénio. Essa interacao pode resy| S Dre.\\
sentes na subsuperficie, na auséncia de oxigénio. f | u tar em res\\
que formam elementos e compostos encontrados na atmos a incluindo Nitrogen,\\
oxigénio, didxido de carbono, argénio € agua. Se es ae ie submet dos ke\\
pressGes durante sua migracao para a superficie, podem formar lepdsitos de Metang, a\\
maneira semelhante ao processo termogénico.\\
\\
3.3.1.2 Origem do petréleo no Pré-Sal\\
\\
O pré-sal é uma sequéncia de rochas sedimentares formadas ha mais de 10\\
anos no espago geografico criado pela separacao do antigo continente Gon\\
especificamente, pela separacao dos atuais continentes Americano e Afric:\\
gou hd cerca de 150 milhGes de anos. Entre os dois continentes formaram-\\
grandes depressdes, que deram origem a grandes lagos, Ali foram depos\\
de milhdes de anos, as rochas geradoras de petréleo do pré-sal (Figura 3.\\
\\
Como todos os rios dos continentes que se separavam corriam para\\
baixas, grandes volumes de matéria Organica foram ali depositados. A\\
\\
0 milhoes de\\
idwana, Mais\\
‘aNO, que come.\\
Se, inicialmente,\\
itadas, ao longo\\
16).\\
\\
as regiGes mais\\
medida que os\\
» OS materiais organicos entao acumulados ne:\\
\\
Digitalizada com CamScanner\\

\\section*{Pagina 78}

GEOLOGIA DE PETROLEOS 57\\
rada, altos gastos com investimentos em tecnologias, Pressio na coluna de 4\\
side? eas baixas tee peretinee que podem danificar componentes, Prejudicando e xe\\
wee reracho, O8 FI5CO8 ambientais destacamos a grande concentracio de didxide\\
wes +> wel pode contetliuir para 0 aquecimento global do mundo. A forma\\
de ultrapassar uma lamina de agua de mais de dois quilometros, uma ween\\
deum quilmetro de sedimentos ¢ outra camada superior a dois quilometros de sal para\\
\\
retirar 0 petrdleo com seguranga sem acidentes ambientais.\\
\\
3.3.1.3 Acumulagao do petréleo\\
\\
fe) petréleo é formado pela decomposi¢ao de matérias organicas e minerais atacados por\\
\\
bactérias. Recoberta por sedimentos, enterrada no fundo dos mares e lagoas dos terrenos\\
\\
sedimentares, esta matéria, apés um longo tempo, se transforma em hidrocarbonetos\\
mpostos de hidrogénio e carbono).\\
\\
Quando a geologia do terreno é favoravel, o dleo fica preso entre as camadas de rochas\\
impermeaveis, ¢ assim se formam as jazidas de dleo e gs natural. Nas jazidas, a camada\\
de éleo é coberta por uma camada de gis e flutua sobre outra de agua salgada.\\
\\
As condices necessdrias para acumulacao do petréleo sao: rocha geradora, rocha re-\\
servatorio, rocha selante ou capeadora, trapas ou armadilhas, migracio e sincronismo.\\
Aauséncia de qualquer um desses elementos impossibilita a existéncia de uma acumula-\\
cio petrolifera. Ao conjunto de todos os elementos e processos geoldgicos necessarios 4\\
existéncia de acumulacées de leo e gas natural é chamado de sistema petrolifero.\\
\\
Importa destacar algumas definigSes relacionadas com bacia sedimentar. Uma bacia\\
sedimentar é a depressAo na crosta terrestre onde ha o actimulo de rochas sedimentares\\
que possam ser portadoras de éleo ou gas natural. A parte de uma bacia sedimentar,\\
formada por um prisma onde se desenvolvem atividades de exploracao e produgao de\\
6leo e gas natural, é chamada de bloco. Um bloco pode ser composto por um ou mais\\
campos. Esses campos referem-se a Area produtora de dleo e gas natural, que contém um\\
ou mais reservatorios, cuja profundidade varia e que contém instala¢es e equipamentos\\
destinados a producio.\\
\\
A existéncia de uma bacia sedimentar nao garante, por si s6, a presenca de jazidas\\
de petréleo (Figura 3.17). Existem alguns elementos necessarios para a acumulagio do\\
\\
Petréleo, que passamos a descrever.\\
\\
(col\\
\\
Figura 3.17: Jazida de Petrdleo.\\
\\
Digitalizada com CamScanner\\

\\section*{Pagina 79}

5B \\_EXPLORACAO DE HIDROCARBONETOS\\
originalmente rico em matérjg ;\\
‘ Opi\\
jente natural, geraram tA\\
nica, Sedimentos de gro fino qu Sate cammareial as pelt nae\\
hidrocarbonetos suficientes para formar i a achonate, ongAtsny , eo ¢ yi,\\
i icas em argila f » Cepositag.\\
As rochas geradoras sao lamas ricas ¢™ ia morta is\\
em baixa unr reduzindo as condigées. O alae os : inte na Beracio\\
, tC) a i: Vln\\
petréleo na rocha geradora éa temperatura. 4 aga pay shed ie organic\\
insolavel (querogénio) contida nas rochas gera' loras G40 de petrélen ,\\
\\
gas.\\
\\
ra é 0 sedimento\\
\\
: a rocha gerado\\
Rocha Geradora: a g! sem seu am\\
\\
nte acumula-se nas rochas em que foi gerado, Natura\\
sedimentos, e tende a subir devido 4 dif\\
\\
dimentos provoca 0 movimento dos fyi.\\
\\
maior porosidade ¢ pressdo mais baixa\\
\\
ficilme:\\
da 4gua existente nos\\
pactagao dos se\\
\\
Migracao: o petréleo di\\
mente, ele se separa\\
renga de densidade. A com\\
dos em diregao as rochas permedaveis, com\\
\\
Sincronismo: é o fenémeno que faz com que todos os fatores presentes na formacig\\
de um sistema petrolifero ocorra em uma escala de tempo geolégico adequada, oy\\
seja, nao basta a ocorréncia das condices de geracao de hidrocarbonetos, de rotas\\
de migracao, de rochas reservatorio ou de trapeamento, sem que elas ocorram de\\
forma favoravel ao longo do tempo. A auséncia de sincronismo tem sido uma causa\\
comum do insucesso nas pesquisas exploratérias no mundo inteiro.\\
\\
Rocha reservatério: é onde o petréleo se acumula. Ela deve possuir boa porosidade para\\
armazenar os fluidos e boa permeabilidade para permitir 0 fluxo dos fluidos no\\
meio poroso. Dois tipos de rochas sedimentares sao considerados importantes como\\
rochas reservatério: arenitos e carbonatos (calcario e dolomita).\\
\\
* Arenitos: so uma das rochas reservatorio mais comuns em campos de petrdleo\\
F abasad bewve Essas rochas podem ser espessas, atingindo varias centenas\\
rosidade nos sine Be Le exibir uma grande continuidade lateral. A po\\
fiddade lokeegia a. le classificada em intergranular e por fraturas. A po-\\
gros. A porosidade inical de i ag aos vazios deixados apés a cimentacio dos\\
\\
¢ tamanho dos graos de arei ‘ P fs nde principalmente do grau de arredondament0\\
¢ de tamanho semelhante, ee ee forem moderadamente arredondados\\
grande variagio de tamanho Porosidade pode variar de 35\\% a 40\\%. Se house\\
\\
30\\%. A cimentacio e as pine Nar sera menor, muitas vezes abaixo de\\
\\
arenitos, conferindo a rocha\\
uma alta permeabilidade. A porosidade em rocks\\
\\
carbonaticas\\
geralme;\\
Solugio ¢ fretting € secundaria, devido a processos como dol omitizas®”\\
\\
* Outras Rochas: f\\
: Além das\\
sideradas reservaté, eg a\\
\\
Digitalizada com CamScanner\\

\\section*{Pagina 80}

GEOLOGIA DE PETROLEOS 59\\
\\
Rocha capeadora: também conhecida co\\}\\
tém o petréleo, impedindo a Migrag:\\
uma armadilha com diferentes elem\\
\\
mo selante, é uma roch:\\
40 até a superfi\\
lentos que per:\\
\\
r ‘a impermeavel que re-\\
Cie. Na Figura 3.18, \\& ilustrada\\
mutem a acumulacao do petréleo.\\
\\
Figura 3.18: Acumulagao do petréleo,\\
\\
Maturacao: € a conversao de matéria organica sedimentar em petréleo.\\
resultantes sio amplamente controlados pela composi¢ao do material\\
ginada da decomposicao nao-oxidante da matéria organica armazenad:\\
tos, que migra através de aquiferos e fica aprisionada em reservatérios.\\
ilustra as temperaturas da decomposi¢o e formacio de petréleo no\\
\\
Armadilhas ou trapas: sio situagdes geologicas estruturais, estratigraficas e mistas que\\
Provocam condicées para a existéncia de acumulac6es petroliferas, nao permitindo\\
que o dleo migre para a superficie. A formacao de quase todos os tipos de armadilhas\\
envolve deformacées da rocha reservatério. Armadilhas ou trapas sao conjuntos de\\
rochas reservatérios e selantes, arranjados em estruturas ou outras situagGes geolé-\\
gicas que permitem a acumulacao de dleo ou gas. Armadilhas ou trapas sao capazes\\
\\
de aprisionar o petréleo aps sua formacao e migracio, evitando que ele escape para\\
a superficie. Elas podem ser:\\
\\
Os produtos\\
original. Ori-\\
laem sedimen-\\
A Figura 3.15\\
reservatério,\\
\\
* Trapas estruturais (Figura 3.19): aquelas cuja geometria é resultado da atividade\\
tect6nica, estando relacionadas a falhas, dobras ou diapiros. Elas sao as mais\\
evidentes nos mapeamentos geoldgicos de superficie e também podem ser ra-\\
pidamente localizadas em subsuperficie. Uma trapa estrutural pode ser identi-\\
ficada por meio de geologia de superficie, perfuracdes estruturais, geologia :\\
subsuperficie, métodos geofisicos ou por uma combinacao destes métodos. .\\
armadilhas estruturais podem estar associadas a dobras ou anticlinais, ee : mo\\
vimentos saliferos com perfuragao das camadas superiores pela massa diapirica,\\
0u podem ser uma combinacio de dobras ¢ falhas.\\
\\
Digitalizada com CamScanner\\

\\section*{Pagina 81}

wos\\
\\
PXPLORACAD DE HIDROCARBON\\
\\
a\\
\\
Pigura 3.195 ‘Trapa Fstrutural.\\
\\
Jas resultantes de variag6es litolégica\\
\\
xemplo: recifes, lentes de atenttones\\
exemplo: truncamentos, barreiras diagenéticas). A exis\\
fio estruturais \\& conhecida desde o final do século xix\\
finidas como aquelas em que as restrigGes geo.\\
ério/rocha de cobertura s4o formadas por va.\\
te de uma deformacio estrutural.\\
\\
\\{ficas (igura 3.20): aque’\\
\\
* ‘Trapas estratigrs\\
rigem deposicional (por ¢\\
\\
podendo ser de or\\
pos deposicional (por\\
téncia de armadilhas n\\
Rssas armadilhas podem ser de!\\
métricas ¢ a combinagao reserval\\
riagdes na est ratigrafia, independentemen\\
\\
Figura 3.20: Trapa Estratigrafica.\\
\\
Rittenhouse 1972 divide\\
as armadilhas estratigraficas\\
Barre sje cine secundarias Stratigraficas em: primarias ou de deposica0 ¢\\
adilhas estratigraficas ‘..\\
Satneite Sto es aad deposi¢io sao produtos diretos do ambiente\\
— exemplos, podemos citar as lent mt trapas deposicionais ou trapas diagenttics\\
PsA ra eet Daten da eelte em saqtion espeesas ce folhelhos *\\
comuns rome ney dio origem a stinallins € preenchimentos de antigos canais\\
estratigr: : eae\\
armadilhas deste am Zonas porosas locais em seas, Campos petroliferos 9°\\
Os recifes . carbonatos dio origem a excelentes\\
cas Beralmente sio Stim,\\
slo aquelas que se\\
\\
reserva 08 reservatérios d\\
trio. Hstas desenvol le petrleo, Armadilhas estratigr\\#\\
‘ vera\\
¢ diflcdl identificar én adilhas esto frequenteme, Aapds a deposigio ¢ diagénese da roc”\\
Mitre de Mergulho, 4 partir de ian associadas a discordincias. As ve\\
’Juntamente com as ie em pocos. Os perfis ele\\
igdes litolégicas ¢ as secdes sismic®\\
\\
Digitalizada com CamScanner\\

\\section*{Pagina 82}

EXPLORACAG PETROLIFERA\\
\\
61\\
emitem seu Mapeamento com alguma precisio, Outro ti\\
\\
importantes resulta de alteracSes sedimentares d\\
cas\\
pode criar:\\
* Rochas com qualidade de reservatério a Partir de sedimentos Original,\\
essas qualidades; Mente sem\\
\\
Po de armadi\\
\\
has est\\
T\\
as camadas. Este ti\\
\\
‘atigrag.\\
po de alteracdes\\
\\
» Rochas de cobertura a partir de rochas que, inicialment\\
\\
e, tinham c;\\
reservatorio.\\
\\
‘aTacteristicas de\\
As rapas hidrodinamicas formam-se em reas onde 0 fluxo descendente de Agua retém\\
petréleo sem nenhum tipo de fechamento estrutural ou barreira estratigrafica. As ar-\\
madilhas mistas sao 0 resultado da combinacio de fatores estruturais e estratigraficos em\\
proporrao aproximadamente igual. Armadilhas combinadas tipicas sao formadas quando\\
uma falha corta um arenito préximo a sua mudanga de facies Para folhelho, ou quando\\
este mesmo arenito é dobrado. Quase todos os grandes campos de petréleo do mundo\\
estao associados a anticlinais, que sao dobras na crosta terrestre com varias origens, como\\
efeito de um relevo soterrado, diapiros de sal, esforcos de\\
\\
compressao, entre outros. Es-\\
sas dobras podem dar origem a falhas, que na maioria das vezes constituem barreiras 4\\
migracao, gerando armadilhas petroliferas. Em alguns casos, essas falhas podem funcio-\\
\\
nar como condutos, permitindo a migracio para reservatérios mais rasos ou até mesmo\\
para a superficie, originando exsudagGes.\\
\\
Structural Traps\\
\\
et\\
\\
Stratigraphic Traps\\
\\
Unoonformity\\
\\
Figura 3.21: Trapa mista.\\
\\
3.4 Exploracdo Petrolifera\\
\\
ivisd ara a peri-\\
\\
ATTetra é formada por trés camadas que apresentam subdivis6es. Do ee E 7 oe\\
\\
ria temos: Nicleo, o Manto e a Crosta (Figura 3.22). As acumulagoes cancedese\\
\\
Scorrem no nivel da crosta terrestre, que compreende os materials Peni\\
\\
wy Petficie até uma profundidade de 30 a 50 quilémetros. A exploragio © ¢ wa\\
Pela equipe de sismi i les utilizam varios me!\\
\\
sismica, geofisica e gedlogos. Eles u erent\\
* descoberta ¢ comprovacio da seaatenidets da existéncia de petroleo.\\
\\
todos e téenicas\\
( importante\\
\\
Digitalizada com CamScanner\\

\\section*{Pagina 83}

62 EXPLORAGAD DE HIDROCARBONETOS\\
\\
obrir 0 reservatOnioy é necessirio pine se ha Vebting\\
\\
campo descoberto- Para realizar a exploragio do oon\\
\\
08 importantes:\\
\\
mentares por meio de anilise detalhada g,\\
0\\
\\
destacar que no basta dese\\
\\
econdmica na produgio do\\
rios basicamente dois passe\\
\\
sho necessd\\
lizaciio de bacias sedi\\
\\
* Prospecgao: éa local\\
solo ¢ do subsolo; ;\\
na vez descobertas as jazidas de petréleo, realiza-se a Marcacig con,\\
rcadoras sobre a Agua do mar. Se for em terra, Tealiza-se\\
Imente existir 0 petréleo, outros\\
\\
primeiro pogo Se rea i\\
cio é viavel economicamente.\\
\\
* Perfuragio: um\\
coordenadas GPS e boias ma!\\
a perfuragio do solo para 0\\
pocos sao perfurados eana\\
\\
lisa-se se a extra\\
\\
Figura 3.22: Deformag6es da Crosta Terrestre.\\
\\
3.4.1 Prospeccao do Petréleo\\
\\
Para encontrar 0 petréleo, sao utilizados métodos e técnicas especificos que permitem\\
— partie! asua formagao. Essas sao as Ae sda: sscnicas a roa\\
of fe) ie ee pogo, a éa etapa que exige a maior parte dos investimentos\\
campy at a Pee e geofisicos estudam detalhadamente os dados das\\
pine hp nice nh () - les buscam parametros que indiquem a possibilidade de\\
cnluaetnantt lana i if locais mais provaveis para sua existéncia. Essas anilises so\\
ag pte, ia ii Ke e relacao custo/beneficio para a exploracio.\\
onsite de latasune ap vido durante a fase de prospeccao fornece uma grande\\
falas siteunae patisectt tt com um investimento relativamente pequeno com\\
probabilidade de haver ite) Fa exploratério. O gedlogo que determina 4\\
weerene oa. a Ys vatorio com petréleo aprisionado pode fazé-lo de di\\
manciras, uso de imagens de satélite.\\
\\
3.4.1.1 Métodos Geoldégicos\\
\\
Aidentificagio de\\
tod uma area fi ‘4\\
propicio econ e siidkleionen acumulacio do Petréleo é realizada através deme\\
Para aperfuracio, Todo 0 ando em conjunto, conseguem indicar o local ™\\
Programa desenvolvido durante a fase de prospect?\\
\\
uma quantidade\\
relativament muito grande de j\\
eaplorseésien Pequeno quando com * informaces técnicas, com um investimen®\\
: parado ao custo de perfuragio de um tinico pos?\\
\\
Digitalizada com CamScanner\\

\\section*{Pagina 84}

EXPLORACAO PETROLIFERA 63\\
\\
Os métodos geoldgicos sao aqueles que utilizam in!\\
ae ouem subsuperficie como forma de analisar a Pp\\
da. Os tipos de rocha que afl 3\\
na drea estudat q ‘oram na superficie for indicacd\\
opassado geoldgico de determinada rea, desta fo necem indicagdes sobre\\
\\
rma  possivel inferir i\\
SS sobre a capaci\\
desta rea apresentar uma estrutura no subsolo capaz de conter 0 petréleo =—\\
\\
" que se identificam as bacias sedi\\
sedi -\\
res, praticamente excluido de estudos posteriore. oe\\
\\
ou metamérficas, e ainda bacias sedimentares\\
\\
desta geologia de superficie pode ser realizada através da aerofotogrametria, onde sio\\
\\
realizadas fotografias da superficie. Esta técnica é bastante utilizada em levantarientoa\\
\\
topogrificos. Outra forma de conhecer a geologia superficial é através do auxilio de infor-\\
\\
magoes de satélite. Em termos de informagiio de subsuperficie, estas podem ser obtidas\\
\\
através do estudo dos dados fornecidos pela perfuracio de Pocos.\\
\\
Durante a perfura¢ao de pocos sao estudadas as amostras dos cascalhos levados a su-\\
\\
perficie, chamadas de amostras de calha. Essas amostras indicam os tipos de rochas que\\
\\
f estao sendo atravessadas e também podem indicar a presenga de hidrocarbonetos por\\
\\
meio de testes especificos realizados ainda no local de perfuracao do poco (locagao). Apds\\
\\
a perfuracao, obtém-se a distribuicao das litologias atravessadas em sua profundidade.\\
\\
Quando correlacionadas com as informacGes de varios Po¢os perfurados na area, pode-se\\
obter o mapa estrutural da area.\\
\\
formagdes sobre as rochas em superfi-\\
ossibilidade de ocorréncia de petréleo\\
\\
3.4.1.2 Métodos Potenciais\\
\\
A gravimetria e a magnetometria, também chamadas de métodos potenciais, sao muito\\
importantes no inicio da exploracao de petréleo por métodos indiretos, permitindo o\\
reconhecimento e mapeamento das grandes estruturas geolégicas que nao aparecem na\\
\\
superficie.\\
Gravimetria\\
\\
Atualmente sabe-se que o campo gravitacional depende de cinco fatores: latitude, eleva-\\
\\$40, topografia, marés e variacdes de densidade em subsuperficie. As variages de densi-\\
dade em subsuperficie so, na verdade, o tinico item de interesse na exploracao gravimé-\\
trica para petréleo, pois permite fazer estimativas de espessura de sedimentos em uma\\
\\
ia sedimentar, presenca de rochas com densidades andmalas como as rochas igneas e\\
domos de sal, e prever a existéncia de altos e baixos estruturais pela distribui¢ao lateral\\
desigual de densidades em subsuperficie. O mapa gravimétrico obtido é denominado\\
mapa Bouguer, em homenagem ao matemitico francés Pierre Bouguer (1698-1728). i"\\
interpretacio do mapa Bouguer (Figura 3.23) é ambigua, pois diferentes ausieiae geole\\
Bicas podem produzir perfis gravimétricos semelhantes. Portanto, a utilizacao indy . ,\\
° no consegue diagnosticar com confiabilidade a real estrutura do interior\\
Terra, apesar de mostrar a existéncia de algum tipo de anomalia.\\
\\
Digitalizada com CamScanner\\

\\section*{Pagina 85}

OCARBONETOS\\
\\
64 —\\_ expronacAo DE HIDR\\
\\
Figura 3.23: Mapa de Bouguer.\\
\\
Magnetometria\\
\\
A prospeccio magnética para petrdleo tem como piles soenie pequanee Variacies ny\\
intensidade do campo magnético terrestre, consequéncia da distri al akg irregular de j,\\
chas magnetizadas em subsuperficie. Nos levantarnentos aeromagnéticos, aS medidas\\
obtidas pelos magnetémetros dependem de varios fatores, dos quais se destacam: |g»,\\
tude, altitude de voo ou elevagio, direcao de voo, variagGes diurnas e presenga localizady\\
de rochas com diferentes susceptibilidades magnéticas (Figura 3.24). As rochas sediments\\
res apresentam, em geral, valores de suscetibilidade magnética muito baixos. Da mesma\\
forma como os mapas Bouguer, os mapas magnéticos obtidos apés as devidas corregies\\
das medidas de campo podem apresentar interpretag6es ambiguas e devem ser utiliza-\\
dos em conjunto com outros métodos. O exame cuidadoso destes mapas pode fornecer\\
tivas da profundidade do embasamento e presenga de rochas reservatorios.\\
\\
Figura 3.24: Ma\\
\\
pa aerom: i 1\\
eton, 1971) senna da campo de petrdleo Puckett, ‘Texas (Extraido de Net\\
\\
3.4.1.3 Métodos Sismicos\\
\\
Os métodos sismicos comuns\\
(Figura 3.25). Este método\\
das rochas da subsuperficie\\
\\
ex?\\
\\
na inddstri\\
ria petrolifera so: sismica de refragio e de rell\\
\\
Cconsiste em analisa; ont\\
Aue se deseja ponhecs snnuneto de ondas sismicas a\\
\\
Digitalizada com CamScanner\\

\\section*{Pagina 86}

EXPLORACAQ PETROUFERA 65\\
\\
sismica de Refragao\\
ométodo sismico de refracao registra somente Ondas refratadas com angulo eri\\
ars B i i : Fi critic\\
rande aplicagao na area de sismologia. Foi através deste método que a estrutura ‘Oetem\\
. 4 ra a 4 Interi\\
da Terra foi desvendada. Na area de petrdleo, sua aplicacao é bastante festtita atta terior\\
embora este método tenha sido largamente utilizado almente,\\
\\
na década de 1950\\
‘ como a\\
refinamento dos resultados obtidos pelos métodos potenciais. A sismologia é one\\
2 i Ht Tt a\\
geologia € da geofisica que se dedica a estudar e Prever os terremotos e as ondas sismicas\\
\\
artificiais, ¢ associadamente determinar a estrutura da Terra,\\
Sismica de Reflexdo\\
\\
O método sismico de reflexao é o método de exploragao mais utilizad\\
distria do petrdleo, destacando-se pelo alto grau de eficiéncia a um\\
baixo. Mais de 90\\% dos investimentos em Prospec¢ao sao aplicados em sismica de Te-\\
flexio. Por este método, obtém-se uma excelente definigio da formacio geoldgica da\\
subsuperficie, permitindo a analise da probabilidade do acimulo de hidrocarbonetos,\\
\\
O levantamento sismico baseia-se nas reflexdes de ondas elasticas geradas artificial-\\
mente, por exemplo, por explosées de cargas de dinamite ou de ar comprimido que se\\
propagam pelo interior da Terra, onde sao refletidas pelas interfaces das diversas forma-\\
ges rochosas. As reflexGes sao captadas por equipamentos especiais denominados geo-\\
fones (para os registros em terra) ou hidrofones (para os registros no mar), os quais con-\\
vertem as vibragGes mecAnicas em oscilagées elétricas que sao transmitidas e registradas\\
nos sismégrafos.\\
\\
A sismica convencional é chamada 2D (duas dimensées). Ja a sismica tridimensional,\\
3D, permite uma melhor defini¢ao, pois determina a imagem tridimensional das feigdes\\
geolégicas de subsuperficie. Com o avanco tecnoldgico j4 € possivel a utilizacdo de sis-\\
mica 4D, onde a quarta dimensao é representada pelo fator tempo. Trata-se da repeticio\\
da sismica 3D em intervalos periddicos (entre 6 a 12 meses), com 0 objetivo de monitorar\\
4 movimentacao de fluidos (extracao, injecdo de Agua etc.) em um campo de petréleo,\\
sendo, por isso, mais empregada em um campo em producio, onde a extracio continua\\
acaba por acarretar uma queda de pressio no reservatério.\\
\\
lo atualmente na in-\\
custo relativamente\\
\\
Figura 3.25: Métodos Sismicos.\\
\\
Digitalizada com CamScanner\\

\\section*{Pagina 87}

66 —\\_EXPLORACAO DE HIDROCARBONETOS\\
\\
Conclusao\\
\\
jdentificar inicialmente as condicdes he\\
‘ 4 preparado para ! .\\
Neste capitulo, ° pod ade pie ulagdo de hidrocarbonetos. aa a hemes, te\\
cessarias para a exist — 4 ra adquirir conhecimentos bisicos le geociéncia relag;.\\
contetidos foram direcionados pa! lacas tectonicas € morfologia dos fundos oce;.\\
\\
ia das pl:\\
\\
dos com a crosta terrestre, teoria dai i i le rochas ¢ a\\
\\
wes Gonetins relacionados para identificar 0s diferentes tipos d quelas que\\
mi a\\
\\
Jo. As técnicas geolégi\\
idrocarbon jlustrados nesta forma¢ao. ve\\
pene rom aen caer depen de petréleo e gas natural foram parte dos\\
¢ geofisic:\\
contetdos de aprendizado deste médulo.\\
\\
4\\
\\
Digitalizada com CamScanner\\

\\section*{Pagina 88}

CAPITULO 4\\
\\
PERFURACAO E COMPLETACAO DE\\
POCOS\\
\\
SS eee\\
\\
As operacées de perfuracao e completacao de\\
\\
; Po¢os tém como objetivo certificar a exis-\\
téncia de hidrocarbonetos,\\
\\
comunicar 0 reservatério a superficie (plataforma de produ-\\
640) € preparar 0 poco para produzir os hidrocarbonetos, Essas operacGes sao realizadas\\
de forma a otimizar o tempo e, consequentemente, os custos operacionais, garantindo a\\
seguranga inerente 4 operac4o e adequando-se as leis governamentais.\\
\\
As fungdes de um engenheiro de perfuracio envolvem a elaborac4o de programas\\
de perfuracao, estimativas de custos modificados do programa de perfuraco (incluindo\\
Authorization for Expenditure/ Autorizagdes para Despesas (AFE)), especificacao de licita-\\
s4o para empreiteiros de perfuracio e contratos de perfuracao. Esse engenheiro trabalha\\
€m conjunto com um gedlogo (no programa de registro) e outros especialistas, como\\
engenheiros direcionais, engenheiros de lama, pessoal de servicos e suprimentos.\\
\\
Mesmo com os recursos tecnolégicos provenientes dos métodos sismicos, somente\\
com a perfuracéo de um poco é possivel comprovar ou refutar a tese de acumulacao\\
Proposta nas anilises geolégicas e geofisicas. Tecnicamente, a perfuracao consiste no\\
conjunto de varias operac6es e atividades necessérias para atravessar as formacées geolé-\\
gicas que compéem a por¢ao superficial da crosta terrestre, com objetivos predetermina-\\
dos, até atingir 0 objetivo principal, que é a prospecco de hidrocarbonetos.\\
\\
O objetivo geral deste capitulo é entender as operagées de perfuracao e completacao.\\
Nesse contexto, os objetivos especificos sao:\\
\\
* Descrever os equipamentos de perfuragao e completacao;\\
* Conhecer os problemas comuns durante a perfura¢4o e completacao;\\
* Planejar a execucio do projeto de poco;\\
\\
* Explicar os processos de perfuracio, completagao, intervengao ¢ abandono, por\\
meio do fluxograma da cadeia de valor da industria petrolifera,\\
* Conhecer os riscos associados 4 seguranga na industria petrolifera, com eficiéncia.\\
\\
67\\
Indistria do petréleo para ndo-engenheiros,1* edigdo.\\
By Geraldo ne tp acca Copyright © 2023 Editora Dois03.\\
\\
Digitalizada com CamScanner\\

\\section*{Pagina 89}

Ss\\
\\
68 PERFURACAO E COMPLETAGAO DE POCGOS\\
\\
4.1 Perfuracdo de pocos de petrdleo\\
\\
O petréleo se encontra na natureza ocupando os espagos vazios de uma tocha\\
chamada rocha reservatério. O pogo é 0 elo de ligacdo entre tal rocha e a Superfi, he\\
perfuracao de pogos é feita através de uma sonda cuja sele¢ao depende da lamina d vi\\
A profundidade dos pocos é programada em fun¢ao dos estudos sismicos, Bristen,\\
tipos de métodos de perfuragao de pocos: método percussivo ¢ rotativo. dois\\
No método percussivo, a perfuragéo é feita golpeando a rocha com uma broca\\
sando a sua fragmentagao por esmagamento. Os cascalhos gerados no interior ng\\
apés varios golpes sio retirados posteriormente através de uma ferramenta cha:\\
\\
cacamba (drilling bucket), como ilustrado na Figura 4.1.\\
\\
© Poge\\
Mada de\\
\\
Figura 4.1: Cacamba de perfuracao.\\
\\
Foi o primeiro método utilizado para a perfuracao de pogos na industria petrolifera. Co-\\
mecou a ser utilizado em 1859, quando o Coronel Drake (Figura 1.13) descobriu petrd-\\
leo em Titusville, Pensilvania, utilizando o processo de percussao movido a vapor. Nesst\\
momento, foi iniciada a exploragéo comercial do petréleo nos Estados Unidos. Esse pogo\\
produziu 2m3/dia de dleo com uma profundidade de 21 metros.\\
\\
No método rotativo, a perfuracao é realizada através do movimento de rotacao de\\
uma broca, comprimindo a rocha e causando seu esmirilhamento. A retirada dos casce\\
Ihos gerados é realizada através do bombeio de fluidos de perfuracao através de tubos\\
de perfuracio que retornam pelo espaco anular entre a coluna de perfuracao ¢ © poco\\
perfurado.\\
\\
Este método foi utilizado pela primeira vez no século XX na cidade de Spindletop. Te\\
xas, quando o americano Anthony Lucas encontrou éleo a uma profundidade de 354 al\\
tros. Nos anos seguintes, a perfuracdo rotativa progressivamente substituiu a perfuragi?\\
pelo método de percussio, tornando-se 0 método padrao na industria petrolifera. Ae\\
\\
Hoje, para fazer o petréleo chegar a superficie, é necessrio perfurar um pose ba\\
atinge o reservatério € o eleva até a superficie. A tecnologia envolvendo a aie\\
de pogos desenvolveu-se bastante nos ultimos anos, permitindo alcangar profundi re\\
antes nunca imaginadas, acima de 6.000 metros. A perfuragao de poses pode 00°\\
tanto em terra (onshore) quanto no mar (offshore).\\
\\
\\_ at\\
\\
Digitalizada com CamScanner\\

\\section*{Pagina 90}

PER\\
FURACAO DE POCos pg PETROLEO «gg\\
\\
4.1.1 Vantagens e desvantagens do\\
\\
vencional Método rotativo con-\\
\\
jo transmitida 4 broca provém da mesa rotativa ici\\
rain pen da coluna de perfuragio, existe uma haste a Meee ah Las\\
gue deslza livremente dentro da mesa rotativa (kelly). Assim,\\
haste quadrada que, por sua vez, transmite a coluna de perfur:\\
\\
A coluna de perfuragao também exerce duas funcées imp\\
sobre a broca através dos comandos de perfuragao e circula\\
racio, que por sua vez tem a fungao de reftigerar a broca e\\
superficie.\\
\\
Todo o conjunto constituido pela broca, coluna de perfuracio, haste quadrada e ca-\\
beca de injecdo fica suspenso pelo gancho, que é movido verticalmente por um sistema\\
de polias, chamado de bloco-catarina ou simplesmente catarina. O bloco se apoia na torre,\\
que por sua vez transmite toda a carga a subestrutura e esta a fundacio ou base.\\
\\
‘ tficie. Na\\
“10 hexagonal ou quadrada\\
a mesa transmite Totacdo a\\
‘agao, fazendo girar a broca.\\
ortantes: aplicacdo de Peso\\
a0 de um fluido de Perfu-\\
remover os cascalhos até a\\
\\
« As principais vantagens do método rotativo convencional sao:\\
\\
— Permite rapidez nas operagées, propiciando baixo custo por metro perfurado;\\
- Possibilita a perfuragdo de pocos de grandes diametros;\\
\\
- Permite a retirada continua dos cascalhos;\\
\\
Possibilita a perfuracao de pocos direcionais;\\
\\textasciitilde{} Permite o controle dos fluidos das formagées, impedindo “blowouts”;\\
— Proporciona pocos de boa calibragem;\\
— Possibilita alcangar grandes profundidades devido ao emprego do fluido de per-\\
furagao.\\
" As principais desvantagens do método rotativo sao:\\
— Determinacao imprecisa das caracteristicas das formagdes penetradas — devido\\
a formacao do reboco nas paredes do pogo e penetracao do filtrado da lama,\\
junto as paredes do pogo;\\
- Pogo tortuoso, dependendo dos parametros empregados;\\
- Custo elevado da sonda de perfuracao;\\
\\textasciitilde{} Custo operacional elevado;\\
\\textasciitilde{} Tempo de manobra — bastante elevado.\\
\\
4.1.2 Classificagaéo de pocos de petroleo\\
\\
Os pogos de petréleo podem ser clasificados quanto 4 trajetéria ou geometria em pocos\\
Yerticais, direccionais, horizontais e multilaterais (Figura 4.2). ue\\
A escolha do percurso dependera de critérios técnicos especificos, sendo certo q .\\
todo poco apresenta desvios e variag6es, normalmente relacionados 4 velocidade da per\\
\\$40. Se a perfuraco for excessivamente rapida, o poco tende a ficar mais tortuoso,\\
\\
Digitalizada com CamScanner\\

\\section*{Pagina 91}

; OS\\
7O —\\_penruracad E COMPLETACAO DE POC\\
\\
Jo de alguns posos quanto A finalidade e trajetéria,\\
2: Exemplo\\
\\
Figura 4.\\
\\
F Jo feitas em seu interior. Quanto a profy,\\
\\
i acdes posteriores que scrao .\\
\\
bie ies as Serpodett ser clasificados como rasos (< 1 500 m), médios (1500 m,\\
\\
sere , is wos (> 2.500 m). Quanto A finalidade, os pocos classificam-se €M™ pogog\\
de cakancts desenvolvimento ou especiais, conforme descrito a seguir:\\
\\
® Pocos de exploragao:\\
\\
\\textasciitilde{} Pioneiro: descobrir petréleo baseado em indicadores obtidos por métodos geo-\\
\\
légicos e/ou geofisicos;\\
\\
— Estratigréfico: obter dados sobre a disposig\\&o sequencial das rochas de subsu.\\
perficie;\\
\\
— Extensdo ou delimitacao: fora dos limites provados de uma jazida com o obje-\\
tivo de amplid-la ou delimita-la;\\
\\
\\textasciitilde{} Pioneiro adjacente: descobrir novas jazidas adjacentes;\\
\\
\\textasciitilde{} Jazida mais rasa: quando queremos testar se existem jazidas mais rasas do que\\
as jé descobertas numa determinada 4rea;\\
\\
\\textasciitilde{} Jazida mais Profunda: quando queremos testar se existem jazidas mais profun-\\
das do que as j4 descobertas numa determinada area,\\
\\
* Pocos de desenvolvimento:\\
\\
— Produgao: drenar racionalmente o\\
mento entre pocos);\\
\\
\\textasciitilde{} Injegao: pocos para aumentar\\
tural de um reservatério,\\
comuns, como apagar\\
\\
petrdleo (preceitos econémicos e de espa\\
\\
entar ou melhorar a recuperacao de petréleo ¢ gis ™\\
Myetamos fluidos como Agua e gés. Outros fins mends\\
um incéndio em poco em erupcio.\\
\\
" Pocos especiais:\\
\\
\\textasciitilde{} Para quaisquer outros tipos d :\\
sie redugo de espasanienc (inl th exemplo pogo de refinamento de malhs\\
\\
jecao do campo,\\
\\
Para descobrirmos novos cam i\\
tar dados ¢ avaliar a éxtensi Ones de Petrileo, por exemplo, precisamos cole\\
\\
io das rese: : bs.\\
vas por Meio dos chamados pogos exploraron™\\
\\
J\\
\\
Digitalizada com CamScanner\\

\\section*{Pagina 92}

PERFURACAO DE POCOs DE PETROLEO 71\\
\\_— efetuadas todas as pesquisas ¢ testes exploratérios Nos pocos pionei\\
\\
ciros ja-\\
\\
es, quando nos certificamos de que uma descoberta tem viabilidade Rbab i oe\\
\\
companhias podem decidir pela perfuracao de pocos de desenvolvimento oi\\
\\
gxistem outras formas de perfuracao de pogos que podemos classificar como:\\
\\
+ Pogo repetido: quando, por algum motivo, precisamos perfurar novamente\\
pogo, com os mesmos objetivos; um\\
\\
= Pogo partilhado ou multilateral: nesse tipo de pogo, aproveitamos um poco ja per-\\
furado, ou parte dele, para objetivos diferentes; P\\
\\
* Poco desviado: quando precisamos desviar a trajetéria da perfuracao devido a um\\
obstaculo.\\
\\
Os pogos podem ser clasificados como surgentes ou nao surgentes (artificiais). Define-\\
se como pogo surgente qualquer poco que tenha pressao suficiente no reservatério para\\
promover 0 fluxo de fluidos até a superficie, sem a necessidade de um sistema de elevacio\\
artificial. Ou seja, o petréleo ou gas natural flui naturalmente até a superficie. Este tipo\\
de poco exige pouco equipamento ou despesa para fazer fluir o petréleo e g4s natural até\\
a superficie. Entre os equipamentos comuns usados esto os tubings (tubagem), a cabega\\
de poco ea arvore de natal.\\
\\
Os custos da perfuracao de pocos sao significativos, sendo ainda mais elevados para\\
0s pocos offshore. A perfuracao de pocos tem diversas finalidades e pode ocorrer em va-\\
rias fases da exploracao e producao de petrdleo. Os pogos estratigraficos sao utilizados\\
na fase de producao, e na avaliacao de descobertas, tém vez os pocos de extensio e de\\
delimitacao. Os pogos de producao e de injecdo podem ser perfurados tanto na fase de\\
desenvolvimento como na fase de produgao de um campo.\\
\\
Durante a atividade de perfuracao, nem todos os pogos sao perfurados com sucesso.\\
Segundo os estudos realizados sobre a probabilidade de sucesso:\\
\\
* Média histérica mundial: entre 8\\% e 12\\%;\\
\\
* Varia de acordo com o conhecimento do potencial do campo;\\
\\
* Para o cdlculo da recupera¢ao de uma jazida, deve-se levar em conta:\\
\\textasciitilde{} Fator de saturacao do dleo, porosidade (volume de vazios) e espessura real(NTG\\
- Net to gross) para obter o volume de dleo in place;\\
\\textasciitilde{} Unidade usada: barril (159 L);\\
\\textasciitilde{} Fator de recuperacao: FR: 25\\%;\\
- Apés a retirada do dleo, deve-se levar em consideragao outras perdas,\\
fuga de volateis para a atmosfera: +/- 12\\%.\\
\\
como a\\
\\
4.1.3 Planeamento de pocos\\
ncionamos nas segdes an\\
\\
ariando de poucos ™!\\
para\\
\\
hPerfuracio de um pogo depende da finalidade, conforme me\\
“ores. B um proceso que envolve um investimento relevante, variando ¢\\& PO\\
Bes de délares para um poco em terra (onshore), até 100 milhdes de délares ou ma"\\
\\
Digitalizada com CamScanner\\

\\section*{Pagina 93}

cAo pe POGOS\\
\\
fas. Nessa atividade, estao envolvidos engen)\\
\\
denominagio assumida por companhia “etry\\
\\
lor desse investimento através da aplica On\\
\\
erfurar 0 pogo com custo mot\\
Mim,\\
\\
72 —\\_-renruracdo F COMPLETA\\
rofund\\
dendo da\\
ro va\\
dequados para P\\
e ambientais.\\
\\
pogos de exploragio em aguas Pp’\\
de perfuragio ou de pogos, depen ‘a\\
jetivo desses engenheiros é maximiza\\
recnologias € processos de aoe aah\\
dentro das normas de igs gust a Seen diaveaticc, iomiens:\\
taml s disciplinas especificas de georitns”\\
\\
Vay,\\
\\
O sucesso do processo a tas OUNTE\\
fo e ii indo t\\
\\
nheiros de perfuracao € integra a eecfurtd\\
\\
i 4o perfurados com os segui \\textbackslash{}\\
\\
de servigos. OS posos sao P unt objet\\
\\
Engenharia ¢ prestagao\\
* Coleta de informagdes;\\
\\
* Produgio de hidrocarbonetos;\\
ter a pressao do reservatério ou melhorara pro\\
WU:\\
\\
* Injecio de gas ou 4gua para man|\\
tividade do petréleo;\\
perfuragao ou sequestro de CO).\\
\\
* Descarte de agua, cascalhos de\\
4.1.3.1 Etapas de perfuracéo de pocos\\
\\
As principais etapas necessarias para a perfuracao de um poco sao:\\
1. Verificar os dados sismicos e outras informagées;\\
\\
2. Alugar/comprar ou obter a concessao do terreno;\\
3. Calcular/estimar da forma mais precisa possivel o volume dos hidrocarbonetos exs-\\
\\
tentes no reservatério;\\
\\
4, Se as estimativas se mostrarem promissoras, prosseguir;\\
5. Obter as licencas necessarias;\\
\\
6. Preparar o programa de perfura¢io e completacao;\\
\\
7. Realizar editais/leil ick,\\
na¢do por meio ba Para a medi¢fo perfurada, dias de trabalho ou uma combi\\
presas contratadas baseadas no programa de perfuracao;\\
\\
8. Se necessirio, modificar o pro;\\
5 icar 0 pI a pa:\\
rama para se adequar aos equipamentos;\\
\\
9. Construir estradas, \\}\\
; , locacdes/plata! .\\
sdrios para acessar 0 local: Plataformas e outros equipamentos (maritimos) nec\\
\\
10. Reunir todo o pessoal e: i\\
i, bead ‘Avolvido antes de iniciar a perfuragio (pre-spud mectins)\\
rio, modificar novamente 9\\
12. Perfurar o poco; pid aed\\
13,\\
Contratar a empresa que Tealizard a co;\\
14, Completar o pogo; Sone\\
15. i\\
Instalar os equipamentos de Superfici\\
ie;\\
\\
16. Realizar a anilise das o\\
PeracSes com\\
Specialistas,\\
\\
y\\
\\
Digitalizada com CamScanner\\

\\section*{Pagina 94}

TIPOS E SELEG, FUR\\
AO DE SONDAS ou PLATAFORMAS DE Pen\\
URAC AG\\
\\
73\\
41.32 Ftapas da planificacéo da perturacao ¢ complet.\\
“ etagao\\
\\
gos\\
acdiferentes etapas de planeamento do poco de Petrdleoe\\
\\
jos da drea até o desenho ou projeto do pogo. Nesse cont\\
dos\\
\\
gasegio 4.13.1, sobre as etapas de perfuragio de Po¢o, ou\\
de perfuracao e completagao de pogos, deve ser usadas\\
\\
1, Estudo da area;\\
\\
de Po-\\
\\
gas comecam a partir dos estu-\\
EXtO, para atender ao Tequisito\\
1 Seja, preparacao do Programa\\
aS Seguintes etapas:\\
\\
2. Dados da locagao;\\
3, Trajetoria do pogo;\\
4. Geopress6es;\\
5. Assentamento de sapatas;\\
6. Fluido de perfuracao;\\
7. Coluna de perfuracao e brocas;\\
8. Revestimento e cimenta¢ao;\\
9, BOP;\\
10. Previs6es e otimizagdes de tempo e custo;\\
11. Programa de perfura¢ao;\\
12, Programa de completacao;\\
13. Projeto de poco.\\
\\
Muitos termos usados nesta secdo serao tratados nas se¢Ges posteriores. Esta introdu-\\
¢io nio tem como objetivo aprofundar a matéria da disciplina em questao. Para aqueles\\
interessados em se aprofundar, recomendamos adquirir livros sobre perfuragao ¢ com-\\
Pletacio de pocos de petréleo e gas natural.\\
\\
4.2. Tipos e selecdo de sondas ou plataformas de perfuracao\\
\\
Uma sonda de perfuracio é uma estrutura de ago em formato de torre, composta por um\\
Conjunto de equipamentos utilizados para realizar operagGes de perfuracao, completagio\\
*interven¢io de pogos, Pode ser montada em uma unidade terrestre ou maritima. he\\
\\
as sondas terrestres (onshore) tem como principais recursos 0 design por :\\
* Profundidade operacional maxima, onde as torres s40 construidas no local. As = s\\
jatimas (offshore) tém como principais recursos 0 design portatil ¢ a Water Depth! Pre\\
mane D'égua (WD) de operagio. Muitos especialistas faze ee eat\\
- ha Plataformas com base nas atividades realizadas — pet! it Ape nant\\
\\textasciitilde{} “SS¢ contexto, sondas so usadas para atividades de perfuragio de P «« * furaclo)\\
\\
Plataformas sig usadas para atividades de produgio ou ambas (produce ¢F\\
\\
m\\_ uma distingdo entre\\
fio ou am\\
\\
enquanto\\
\\
Digitalizada com CamScanner\\

\\section*{Pagina 95}

74 = PerFURACAOE COMPLETAGAO DE POCOS\\
\\
plataformas\\
\\
ao de sondas ou\\
4.2.1 Selecao a dependera de varios parametros, Ee.\\
\\
cionad:\\
O tipo de sonda ou plataforma a ser sele'\\
\\
esses, destacamos:\\
\\
* Custo e disponibilidade;\\
\\
S to);\\
« Profundidade da Agua de locacao (em mar aber\\
\\
= Mobilidade / transporte (regio costeira);\\
vo e presses de formacao esperadas;\\
\\
* Profundidade da zona-al\\
climaticas predominantes na area de operacio,\\
\\
= Condicdes meteoceanograficas €\\
\\
da equipe de perfuracao (especialmente seus recordes de seguranca),\\
\\
= Experiéncia\\
4.2.2 Tipos de sondas ou plataformas de perfuracao\\
Nas atividades de perfuracao de pocos de petréleo, sao utilizadas sondas de perfuracio,\\
\\
que consistem em um conjunto complexo de equipamentos. Existem diversos tipos de\\
\\
sondas, e elas podem ser terrestres ou maritimas, dependendo do local de operacao. Um\\
\\
elemento caracteristico das sondas de perfuragao € a presenca de uma torre (torre de\\
\\
perfuracio ou derrick), que facilita o manuseio dos tubos de perfurac4o em secées de tis\\
\\
tubos, tornando a operacao mais agil.\\
\\
As técnicas usadas para explorar 6leo no mar sao semelhantes as utilizadas em terra. As\\
primeiras sondas adaptadas para perfuracao maritima eram basicamente as mesmas son-\\
das terrestres, porém adaptadas com estruturas que permitissem a perfuracao em aguas\\
rasas. No entanto, a medida que a perfuracdo em Aguas profundas se tornou mais comum,\\
surgiram novas técnicas especificas para atender a essas demandas,\\
\\
Mesmo apés 0 abandono do poco, 4\\
'o de um poco envolve varias etapas, coe\\
"na, instalar o tubo condutor, preparar os!\\
das, cercase escavacies, e, as vezes, pel furl\\
\\
Digitalizada com CamScanner\\

\\section*{Pagina 96}

TIPOS E SELECAO DE SONDAS ou PLATAFORMAS DE PERFURAGAO 75\\
\\
Tabela 4.1: Tipos de sondas ou Platafor:\\
\\
™as de perfuragio,\\
\\
ondas ou plataformas de perfuragio\\
ve) Mével\\
\\
Com suporte na parte infe. Flutuantes\\
ior\\
\\
Mastro portétil (Portable Jackups\\
= SS\\
\\
Navio-sonda\\
\\
Alta tecnologia (High- “Jacknife” (Perfuragao Plataformas\\
\\
Semissubmersfvel\\
Tech) mais profundo)\\
\\
‘Alta poténcia (High horse- Submerstvel\\
\\
Barcaga de perfuragio\\
power)\\
\\
=\\
\\
um poco de 4gua. Depois disso, a plataforma é movida para o local, montada e preparada\\
- mee awe enn as da categoria de torre convencional e as méveis jacknife,\\
4 ‘eal ao solo com estacas (pins) e ent4o elevadas como uma unidade usando o equi-\\
wan de elevacao da sonda. O mastro portatil (Figura 4.3), adequado para pocos de\\
me "i oderada, geralmente é montado em caminhées ou reboques que incor-\\
oanec “4 ees de clevacio, Os motores € a torre como uma unica unidade. O mastro\\
crs este ico é elevado a posic¢ao vertical e depois estendido até a altura maxima\\
fore niet 4 istao hidraulico na unidade. A profundidade maxima de perfuracdo das\\
ae aber de cerca de 300 metros, e a operac4o dos pogos é controlada a partir\\
soni\\
da superficie.\\
\\
: de mastro\\
Ao (a rda) e sonda terrestre\\
i rfuragao (a esque:\\
Figura 4.3; Sondas terrestres de alta pe:\\
portatil (a direita).\\
\\
jo maritima\\
4.2.2.2 Sondas ou plataformas de perfuracao\\
\\
‘ti utilizadas\\
i das ou plataformas maritimas, Juir: bar-\\
orias das son ae dae rfuragao podem incluir:\\
\\
io-sonda e\\
missubmersivel, fixa, navio-son\\
\\
Na Tabela 4.1, ilustramos as categi\\
Para perfuracao offshore. Essas sondas ene\\
“sa de pantano, jaqueta, autoelevavel (jack up).\\
FPSO Monocoluna.\\
\\
Digitalizada com CamScanner\\

\\section*{Pagina 97}

1S\\
76 PERFURAGAO E COMPLETAGAO DE pogo’\\
\\
Barcaca de pantano\\
ndas que\\
\\
no (Figura 4.4) sao SO! ree\\
acerca de 6 metros ¢\\
\\
em seguida,\\
\\
operam em dguas muito £8835 (eon,\\
o yane se imo 50 metros). Elas podem \\$y\\
i inferio!\\
profundidade in\\
bocadas até o local ¢,\\
A unidade de perfuragao é\\
em Areas pantanosas, com!\\
petrolifera dos Estados Unidos\\
\\
Jastreadas pare que fiquem apoiadas no fundo dg mar\\
da sobre a barcaga. AS barcagas de pantano sio Usad,\\
\\
“Nas\\
ria, venezuela\\
\\
e na Costa do Golfo, sob juris\\
\\textasciitilde{} ‘0\\
\\
Figura 4.4: Barcaca de pantano de perfura¢ao.\\
\\
Jaquetes de perfuragao\\
\\
‘0 sao plataformas fixas de pequenas estruturas de ago utilizadas\\
\\
As jaquetas de perfuraca\\
tir de uma jaqueta,\\
\\
em areas de Aguas rasas e calmas (aproximadamente 300 metros). A part\\
muitos pocos podem ser perfurados. Uma vez que 0 desenvolvimento do campo seja vit-\\
vel, é extremamente eficaz em termos de custo construir e operar jaquetas em ambientes\\
de mar raso. Em particular, elas possibilitam uma progressao flexivel e passo a passo das\\
atividades desenvolvidas no campo. A construgao da jaqueta é realizada nos estaleiros ma\\
posigao ae Apés a conclusao da construgao, a jaqueta é levada para sua locali-\\
nm parhipa a Ld woe ao mar na posig\\&o vertical, no ponto onde\\
superiores da plataforma, que abaeat aa oe fundo do mar, a jaqueta recebe as part?\\
de materiais, alojamento de pessoal, gens tains = ae Serer\\
producao dos pogos. A producio é escoada a Le aps ap aaa ve a\\
vez que as jaquetas nao possuem tanques de pad pariiipa a on nee ™\\
mazenamento.\\
\\
Plataformas autoelevéveis\\
\\
As plataformas autoelevaveis (Fi\\
\\
de balsa ou plataforma, hh 4.5) sio constituidas basicamente por uma espéc\\
como sonda de perfura¢io, rindtoge 3s todas as facilidades de operacio ¢ om\\
Porto etc, Blas possuem trés ou quatro poe rios, laboratérios, salas de controle. helt\\
ue, acionad,\\
\\
A idea\\
las mecanicamente ou hier\\
\\
Digitalizada com CamScanner\\

\\section*{Pagina 98}

TIPOS E SELECAO DE SONDAS OU PLATAFORMAS DE PERFURACAC\\
CAD, 77\\
\\
avimentam-se para baixo até atingirem o fundo do mar\\
\\
jeamente, me\\
da plataforma acima do nivel da égua, a uma altura se;\\
\\
Em seguid\\
7 ‘ a, inicia\\
ee a clevagio\\
\\
gura e fora da acic\\
saan, ESAS plataformas sio méveis ¢ podem ser transportadas por rebocadores’ ;\\
| ou\\
\\
por propulsio propria. Sao Cestinadas 4 perfuragio de pocos exploratérios na plataforma\\
continental, em Aguas consideradas rasas para a industria offshore, em profundidades que\\
yariam de 5 a 150 metros. Apés a conclusao da perfuragio de um determinado nevhs o\\
convés da plataforma desce até o nivel do mar e a unidade pode ser rebocada para outra\\
\\
Jocalizagao.\\
\\
Plataformas semissubmersiveis\\
\\
As plataformas semissubmersiveis (Figura 4.5) sio compostas por uma estrutura de um\\
ou mais convés, apoiada por colunas em flutuadores submersos. Uma unidade flutuante\\
sofre movimentagées devido a acao das ondas, correntes ¢ ventos, com a possibilidade\\
de danificar os equipamentos a serem descidos no pogo. Portanto, é necessario que ela\\
fique posicionada na superficie, opera¢ao que ocorre em lamina de dgua superior a 2.000\\
\\
metros.\\
Dois tipos de sistemas sao responsaveis pelo posicionamento da unidade flutuante: 0\\
\\
sistema de ancoragem e¢ 0 sistema de posicionamento dinamico. O sistema de ancora-\\
gem é constitufdo por 8 a 12 ancoras e cabos e/ou correntes, atuando como molas que\\
produzem esforgos capazes de restaurar a posi¢ao do flutuante quando ela é modificada\\
pela acdo das ondas, ventos e correntes.\\
\\
No sistema de posicionamento dinamico, nao ha ligacao fisica da plataforma com 0\\
fundo do mar, exceto a dos equipamentos de perfuracio. Sensores actisticos determi-\\
nam a deriva, e propulsores no casco, acionados por computador, restauram a posi¢ao da\\
plataforma. As plataformas semissubmersiveis podem ou nio ter propulsio prépria. De\\
qualquer forma, apresentam grande mobilidade, sendo preferidas para a perfuracio de\\
pocos exploratérios.\\
\\
Plataformas fixas\\
\\
As plataformas fixas (Figura 4.5) so estruturas apoiadas no fundo do mar por meio de\\
estacas cravadas no solo, com o objetivo de permanecerem no local de operagao por um\\
longo periodo. Foram as primeiras a serem utilizadas em campos localizados em laminas\\
dégua de até 300 metros, e esse também é o limite de sua utilizagio. Devido ao custo\\
elevado, compreendendo desde o projeto até a montagem ¢ instalagao, sua aplicagdo ¢\\
Festrita aos campos que jé tiveram sua exploracdo comercial comprovada. As platator-\\
mas fixas sio projetadas para receber todos os equipamentos de perfuraga\\
\\& materiais, alojamento de pessoal, bem como todas as instalagdes nece\\
\\
Producio dos pocos.\\
\\
estocagem\\
\\
rias para a\\
\\
Digitalizada com CamScanner\\

\\section*{Pagina 99}

E POGOS\\
\\
78 PERFURACAO E COMPLETACAO D!\\
\\
rfuragao maritimas: fixa, autoelevaveis e Semissubmersy\\
\\
Figura 4.5: Plataformas de pe!\\
\\
Navio-sonda\\
\\
um navio equipado com sonda de perfuracao € demais egy,\\
izagio de perfuragao rotativa. Ele é projetado Para 3\\
\\
=. ‘1 erfuracado esta localizada no centro\\
\\
perfuragio de pocos submarinos. A torre de Pp ¢ Bieatins dine do m.\\
vio, com uma abertura no casco que permite a passagem le perfuracio, 9\\
sistema de posicionamento do navio-sonda, composto por sensores acusticos, propulso.\\
res e computadores, compensa os efeitos do vento, ondas e correntes que podem deslocar\\
o navio de sua posi¢ao.\\
\\
O sistema de cabeca de poco é posicionado no solo marinho usando um BOP sub-\\
merso. A operacao do navio-sonda é dependente das condi¢des maritimas, principd-\\
mente do movimento de afundamento (heave), em que o navio sobe e desce de acordo\\
com 0 movimento das ondas. Navios-sonda possuem sistemas de compensacio de hea.\\
€ sio autopropulsantes. Eles sao capazes de operar em aguas ultraprofundas, ancorados\\
em laminas de agua de mais de 2.000 m. Também possuem grande capacidade de arm:\\
zenamento de suprimentos para perfuracao.\\
\\
O navio-sonda (Figura 4.6)\\
pamentos necessirios para a reali\\
\\
FPSO monocoluna\\
\\
A plataforma FPSO monocoluna (Figura 4.6) é uma plataforma destinada 4 producia\\
\\
et ak arate eae de dleo bruto, sendo usada em 4guas ultraprofundas. Seu\\
\\
Boe 3 Agia entie Bi take 2 ae estabilidade. Sua tecnologia inovadora permit\\
\\
plataforma desse tipo no mu ae ao casco, reduzindo os efeitos das ondas. A primel\\
\\
Ser realizadas em laminas dete Pe As operacées com essas plataformas pode\\
As plataformas do ti ts '€ 2.000 m. ae\\
\\
de um casco com paar Wrst t€m como caracteristica marcante 4 ueiza\\
na, diferente das semisubmersiveis. Nesse s¢™\\
\\
© conceito assemelha-se as pla\\
ta! ‘\\
relaco aos conceitos ‘hy spa do tipo SPAR, porém com calado menor\\
\\
+ ic ative®\\
Entre elas, podemos citar: tiene oe Monocoluna apresenta vantagens significa\\
flexibilidade operacional e boa rel movimentos, maior reserva de estabilidade aril\\
eS Foe ce convés por Deslocamento.\\
\\
Digitalizada com CamScanner\\

\\section*{Pagina 100}

TIPOS E SELECAO DE SONDAS oF\\
1U PLATAFORMAS DE Pp\\
ERFURACAO = 7g\\
\\
piataforma SP, AR\\
\\
Aplataforma SPAR (Figura 4.6) é uma estrutura que contém um\\
\\
oco, semelhante a uma grande boia, podendo também ter uma\\
\\
plementar. O projeto inicial era utiliz4-la para pesquisa Oceanog:\\
\\
de petréleo como um grande tanque. Atualmente, ela é usa\\
roducao, principalmente no Golfo do México.\\
\\
Com 90\\% de seu casco submerso, ela oferece vantagens sobre as plataformas semissub-\\
mersiveis devido ao seu menor movimento de pitch, especialmente devido a od grande\\
inércia, o que é uma vantagem em Aguas profundas. Essa caracteristica a aproxima das\\
Tension Leg Platform/ Plataforma de Pernas Atirantadas (TLP), permitindo o uso de tie\\
sers rigidos, mas com um sistema de amarragio simplificado em comparacio com uma\\
plataforma semissubmersivel tradicional.\\
\\
unico corpo cilindrico €\\
estrutura de trelica com-\\
Tafica e armazenamento\\
ida para perfuragio e/ou\\
\\
Figura 4.6: Plataformas maritimas de perfuragao: FPSO Monocoluna, Navio-Sonda e\\
SPAR.\\
\\
Nas tabelas 4.2 e 4.3 estao resumidas as sondas ou plataformas de perfuragdo tendo em\\
conta as atividades, profundidade, controle de pocos, escoamento de produgio e vanta-\\
gem de cada um.\\
\\
Tabela 4.2: Comparac¢ao sondas terrestres e barcaca de pantanos.\\
\\
—\\
\\
‘Terrestre (Land Rig) Barcaca (Swamp-Barge)\\
Atividade Perfuracio Perfuracao\\
Profundidade Até 300 m At 50m\\
Controle de pogos Superficie Superficie\\
Escoamento da produgio Nio Nio\\
Vantagem Mobilidade Mobilidade ¢ facil fixagio\\
\\
4.2.3 Navios-Sonda da Sonangol\\
\\
A Sonangol, E.P., no ambito do programa de dinamizacio da industria petrolifera, pos-\\
sui atualmente dois Navios-Sonda cujas especificagées permitem superar os desafios das\\
Novas fronteiras de exploracao e produgio, tornando assim seus servi¢os integrados na\\
cadeia Primaria e reduzindo a dependéncia de disponibilidade externa para cumprir os\\
Programas de trabalho de exploracio.\\
\\
Estes equipamentos sao o “Sonangol Libongos” ¢ 0 “S\\
de Bera¢ao, com capacidade de perfuracgao em Aguas ultraprofundas.\\
\\
onangol Quenguela”, ambos\\
Eles esto em\\
\\
Digitalizada com CamScanner\\
\\
|\\

\\section*{Pagina 101}

MMPLETAGAO DE POGOS\\
\\
80 —\\_ PERFURACAO E CO\\
plataformas de perfuragao de petrdleo ¢ is\\
\\
Tabela 4.3: Comparacio das sondas ou\\
\\
vel ‘Semissubmersive! FPSO monocoluna — Navig mds\\
a Auto-clevive\\
fms (Jack up) ics\\
Perfurasio ¢ produ: Perfuragio, produ: Perfuracig\\
jo Perfuragio io, armazenamento\\
Atividade ac’ ° me e transferéncia\\
produgio ‘ji\\
Mais de 2.000 m Mais de 2.000 m Mais de 2 049)\\
Hateaiade ‘até 300 m. eX) aa Fundo do mar Fundo do mar Pundo do may\\
Controle de posos Superfice SU Olcodutos ou fundo ‘Transferéncias para Ni\\
Escoamento da pro- Oleodutos Nio do mar petrolciros\\
; je- Movimentos meno. Maior\\
dugio Joca- Especialmente proje auto,\\
Vantagem Instalagdes—‘Facilidade ee ta. tada para ter pouco res que FPSO nomia par,\\
mals simples © \\$10 a estrutura movimento perfurar em,\\
controle de po- mento 2 Brandes distin\\
\\
cos A superficie fixa as da costa\\
\\
ja ae a\\
\\
construcio no estaleiro naval “Daewoo Shipbuilding and Marine Engineering Co., Lid\\
\\
DSME”. ard htt 5\\
A gestao das Sondas sera assegurada pela Seadrill Limited, que firmou parceria coma\\
\\
Sociedade Nacional de Combustiveis de Angola, conhecida como Sonadrill. 7\\
\\
O projeto de construgao dos dois navios foi concebido para seguir a estratégia do pro.\\
grama de exploracio petrolifera, visando novas descobertas e consolidacao de reservas\\
para o futuro, considerando principalmente as necessidades do pais. Esses navios sio\\
capazes de operar em dguas profundas, ultraprofundas e na camada pré-sal, e estario\\
disponiveis para trabalhos no ambito nacional como no internacional.\\
\\
Vale ressaltar que o nome “Libongos”, escolhido para o novo Navio-Sonda, é inspirado\\
em uma localidade na regiao de Bengo, onde ha afloramentos de petréleo a superficie. A\\
Figura 4.7 ilustra a cerimGnia tradicional de batismo do navio Libongos em 21 de margo\\
\\
na Coreia do Sul, enquanto o “Quenguela” foi entregue em 17 de maio de 2019 em Okpo,\\
também na Coreia do Sul.\\
\\
Nesta seco, vamos Tesumir como\\
longo do ciclo de vida dos campos\\
\\
8 Sondas oy\\
de petréleo ¢ Bias.\\
\\
Plataformas de perfuracao sao usadas 20\\
\\
Digitalizada com CamScanner\\

\\section*{Pagina 102}

TIPOS E SELECAO DE SONDaS OU PLATAFORMAS DE PERFURACA\\
o 681\\
\\
4.2.41 Fase de explorac¢ao\\
\\
Durante a fase de exploracao, as sondas ou plataformas sio usada:\\
de exploracao ¢ pogos pioneiros em locais com potencial para estruturas geolégicas con-\\
tendo hidrocarbonetos, depois de serem identificados pela andlise dos resultados de varios\\
estudos seclOeices ¢ Pes Saar Geralmente, pocos verticais sao perfurados para\\
garantir a seguranga € a estabilidade do pogo, além de adquirir dado:\\
suficientes sobre a subsuperficie.\\
\\
'S para perfurar pocos\\
s € conhecimentos\\
\\
4.2.4.2 Fase de avaliagao\\
\\
Durante a fase de avaliacao, as sondas ou plataformas sao usadas para perfurar varios\\
pocos a fim de compreender as taxas de fluxo, a dinamica do reservatorio e 0 tamanho e\\
limites do reservatério. Isso é feito para confirmar a suposicao de que hidrocarbonetos\\
podem ser produzidos de maneira econémica.\\
\\
4.2.4.3 Fase de desenvolvimento\\
\\
Durante a fase de desenvolvimento, as sondas ou plataformas sao usadas para perfurar\\
pocos até a profundidade de uma zona produtiva do reservatério. Os P0¢os nesse ponto\\
podem ser verticais, horizontais ou desviados, e podem ser perfurados numa grade ou\\
bloco.\\
\\
4.2.4.4 Fase de producaéo\\
\\
Durante fase de produco, as sondas ou plataformas sao usadas para perfurar mais pogos,\\
também conhecidos como pogos de refinamento de malha ou restauracio/trabalho em\\
Posos existentes. Dependendo da complexidade, uma sonda ou plataforma de menor\\
atividade de restauracao pode ser usada para um programa de restauracao, melhorias de\\
manuten¢ao ou produc4o de pocos, ou para outros tratamentos de pocos.\\
\\
4.2.5 Equipamentos e sistemas de perfuracao\\
\\
Todas as Plataformas tém os mesmos componentes basicos. De forma geral, a sonda é\\
Composta pelos seguintes equipamentos: (1) Tanque de lama; (2) agitadores de argila; @)\\
infa de succo de lama; (4) bomba do sistema de lama; (5) motor; (6) mangueira vibrat6-\\
"3; (7) draw-works; (8) standpipe; (9) mangueira de perfura¢ao; (10) goose-neck (Pescoco de\\
Banso); (11) traveling block; (12) linha de perfuragao; (13) crown block; (14) derrick (torre .\\
Perfuracio) (15) monkey board; (16) stand do duto de perfuracao; (17) pipe rack; (18) a\\
re kelly drive; (20) mesa rotat6ria; (21) superficie de perfuracdo; (22) bell nipple: is\\
Ps OP; (24) dutos do BOP; (25) linha ou coluna de perfuraco; (26) broca de pa 3 ,\\
)cabeca do Casing; (28) duto de retorno da lama, como ilustrado na Figura 4.8.\\
\\
Digitalizada com CamScanner\\

\\section*{Pagina 103}

0S\\
82 PERFURACAO E COMPLETAGAO DE POG\\
\\
Figura 4.8: Principais equipamentos que compoem uma sonda de perfuragao de petréleo\\
\\
Essas partes ilustradas na Figura 4.8, que fazem a sonda de perfurac4o rotativa funcionar,\\
destacam trés fungées basicas executadas durante as operagGes:\\
\\
* Otorque é transmitido para a broca a partir da energia na superficie através da co-\\
luna de perfuracao;\\
\\
* O fluido de perfuracao é bombeado para baixo da unidade de armazenamento atra-\\
vés da coluna de perfuracao ¢ sobe pelo espaco anular. Os cascalhos de pogo (frag-\\
mentos cortados ou triturados pela broca) sao transportados para a superficie pelo\\
fluido de perfuracao. Assim, além de remover os cascalhos para a superficie, 0 fluido\\
contribui Para a limpeza do poco, resfriamento da broca e lubrificagao da coluna de\\
perfuracao;\\
\\
, pies se empertcias acima € dentro das camadas contendo hidrocarbonetos\\
a aebe 0 Peso do fluido de perfuracao e por grandes conjuntos de selas\\
na superficie BOP,\\
Os equipamentos da sonda de 5 :\\
tas no pardgrafo anterior, bara 5 na Figura 4.8, assim como as trés fungoes descr\\
. e i x\\
ser clasificados em sete principais sited poaaie de sondas de perfuracio, que podem\\
* Sistema de Sustentacao de Cargas;\\
* Sistema de Movimentacio de Cargas;\\
* Sistema de Rotacio;\\
\\
* Sistema de Geracio e Transmissio de nal\\
\\
Digitalizada com CamScanner\\

\\section*{Pagina 104}

TIPOS E SELECAO DE SONDAS\\
OU PLATAFORMAS Dy\\
EPERFURACAO = 3\\
« Sistema de Circulagao de Fluidos;\\
\\
« Sistema de Seguranga de Pogo;\\
\\
« Sistema de Monitoramento.\\
4.2.5.1 Sistema de Sustentacao de Cargas\\
\\
O sistema de sustentagao de cargas (Figura 4.9) \\& com\\
\\
subestrutura e pela fundacao. A carga corresponde ao peso da coluna de perfuracao ou\\
revestimento no pogo, ¢ é transmitida para o mastro ou torre, que, por sua vez, a transfere\\
para a subestrutura e, finalmente, para a fundacao ou base.\\
\\
Posto pelo mastro ou torre, pela\\
\\
Figura 4.9: Esquema de uma sonda rotativa com uma torre de perfuracio (a esquerda) e\\
0 Mastro de perfuracao (a direita)\\
\\
Torre ou Mastro\\
\\
A torre ou mastro é uma estrutura de aco especial, de forma piramidal, que fornece um\\
€spaco vertical livre acima da plataforma de trabalho para permitir a realizagdo das mano-\\
bras. Uma torre é composta por muitas pecas, que sao montadas uma a uma. Por outro\\
lado, o mastro é uma estrutura trelicada ou tubular que, apés ser baixada pelo guincho\\
da sonda, é subdividida em trés ou quatro se¢Ges, que sao transportadas para a locacao\\
do novo poco, onde sao montadas na posic¢ao horizontal e entao elevadas para a vertical.\\
\\
O mastro tem sido preferido devido a facilidade e economia de tempo na montagem em\\
Perfuracdes terrestres,\\
\\
Subestruturas\\
\\
Asubestrutura € constituida por vigas de aco especial montadas sobre a fundacgao ae\\
i . criando um espaco de trabalho sob a plataforma, onde sao instalados os es +\\
lentos de Seguranca do poco. As fundacdes ou bases sao estruturas rigidas cons' mu us\\
foe » 8GO OU madeira que, apoiadas sobre um solo resistente, suportam com s\\
as deflexdes, vibracdes e deslocamentos causados pela sonda.\\
\\
Digitalizada com CamScanner\\

\\section*{Pagina 105}

84 PERFURACAO E COMPLETACAO DE pocos\\
\\
Estaleiros\\
flica composta por varias vigas apoiadas acima tt a9)\\
‘ Solo\\
\\
frente da sonda e permite manter todas t\\
“aS ty.\\
\\
O estaleiro é uma estrutura met\\
revestimentos etc.) dispostas paral\\
Paralelamente g\\
\\
por pilares. O estaleiro ¢ posicionado n\\
bulagdes (comandos, tubos de perfuragao, ces\\
uma passarela, facilitando seu manuscio \\& transpor'’\\
\\
4.2.5.2. Sistema de Movimentagao de Cargas\\
mite mover as colunas de perfuragao, reves\\
). Os principais componentes desse sistema\\
bo de perfuragao, gancho\\
\\
cio de cargas pert\\
(Figura 4.10\\
bloco viajante (catarina), cal\\
\\
O sistema de movimenta\\
timento € outros equipamentos\\
sio: guincho, bloco de coroamento,\\
e elevador.\\
\\
Bloco do\\
coroamento\\
\\
Bloco viajante\\
(Traveling Blo\\
Guincho\\
\\
de ago novo\\
/\\
\\
Figura 4.10: Sistema de movimentagAo de cargas.\\
\\
Guincho\\
\\
O guincho recebe\\
aenergia Ani P\\
da transmissfo principal re sea ge necessdria para a movimentacao de cargas através\\
, le sondas di oi\\
acoplado a ele, nas las diesel, ou di:\\
, Nas sondas elétric: A ’ retamente de um motor elétrico\\
bor auxiliar ou de lim frei as. O guincho é constituido : inci\\
ipeza, freios, molinetes e embreagy aaa ll ali\\
ens.\\
\\
O tambor princi\\
ipal tem a funcii =\\
g40 de acionar 0 cabo de perfuracéo movimentando as\\
\\
O tambor auxilia;\\
i r ou de lim :\\
acima do tam! peza é\\
0 tambor principal. Tem a dentine No eixo secundario do guincho, fican do\\
le m incho,\\
\\
ovimen| i\\
tar equipamentos leves no pos\\
\\
Digitalizada com CamScanner\\

\\section*{Pagina 106}

TIPOS E SELECAO DE SONDAS OU PLATAFORMAS DEPERFURAGAO gg\\
como registradores de inclinacio e diregio do Pogo,\\
tos de completagio ¢ teste de Pogo etc.\\
\\
O molinete é um mecanisma semelhante a uma embreagem que permite tracionar\\
cabos ou cordas. Existem dois tipos de molinetes na sonda: © molinete das chaves flutu-\\
antes, para apertar ou desapertar as conexdes da coluna de perfuracio ou revestimentos,\\
\\
eo molinete giratério, cathead, que permite 0 igamento de Pequenas cargas quando uma\\
corda, chamada catline, é enrolada nele,\\
\\
, amostradores de fundo, equipamen-\\
\\
Bloco de coroamento\\
\\
Obloco de coroamento é um conjunto estacionério de 4 a7 polias montadas em linhas\\
num eixo suportado por dois mancais de deslizamento. Ele esta localizado Na parte su-\\
perior do mastro ou torre. O bloco de coroamento suporta todas as cargas transmitidas\\
pelo cabo de perfuracao.\\
\\
Bloco viajante\\
\\
O bloco viajante (travelling block, também conhecido como “catarina” no Brasil) é um\\
conjunto de 3 a 6 polias méveis montadas em um eixo que se apoia nas paredes exter-\\
nas da propria estrutura do bloco viajante. O bloco viajante fica suspenso pelo cabo de\\
perfuracao, que passa alternadamente pelas polias do bloco de coroamento e pelas polias\\
do bloco viajante, formando um sistema com 8 a 12 linhas. Na parte inferior do bloco\\
viajante, ha uma alga onde o gancho é preso. O gancho consiste de um corpo cilindrico\\
que internamente contém garras que engatam no cabo de perfuracao, impedindo que as\\
cargas se propaguem para o bloco viajante.\\
\\
Cabo de perfuracéo\\
\\
O cabo de perfuracio é feito de aco trancado em torno de um niicleo ou alma. Cada\\
tanca ¢ formada por diversos fios de aco especial de pequeno diametro. O cabo é passado\\
do carretel para uma ancora préxima 4 torre, onde um sensor mede a tensio no cabo, rela-\\
Sionada ao peso total sustentado pelo guincho. O cabo passa pelo sistema bloco-catarina\\
€ €enrolado e fixado no tambor do guincho.\\
\\
Elevador\\
\\
O elevador é um equipamento em forma de anel bipartido com partes ligadas por uma\\
adica resistente. Ele contém um trinco especial para fechamento. O elevador é usado\\
Para Movimentar elementos tubulares, como tubos de perfuracao e comandos.\\
\\
Digitalizada com CamScanner\\

\\section*{Pagina 107}

S\\
86 PERFURAGAO E COMPLETAGAO DE pogo:\\
\\
5, Ss\\
4.2.5.3 Sistema de rotaca° de carga\\
\\
luna de perfuragao é realizada pela mesg j\\
forma da sonda. A rotagao é transmitida a um tubo boligong,\\
é enroscado no topo da coluna de perfuragao. Nas sop\\
equipadas com top drive, a rotagao é transmitida a ion ies de perfy.\\
racio por um motor acoplado 4 ea ah es ‘ados a torre,\\
i acao da colu!\\
\\
bart rh prayer usando um motor de re oe logo acima da tio,\\
O torque necessario é gerado pela passagem do fluido de ae ee em seu interior,\\
Esse motor pode ser de deslocamento positivo ou uma turbina. sistema de rotacio\\
convencional é composto por equipamentos que promovem ou permitem a rotacio da\\
coluna de perfuragao, incluindo a mesa rotativa, 0 Kelly e a cabega de circulacao ou swivel,\\
\\
—\\_ A t\\
Nas sondas convencionais, a rotacio da\\
\\
tativa localizada na plata\\
externo chamado Kelly, que\\
\\
Mesa rotativa\\
\\
A mesa rotativa (Figura 4.11) é 0 equipamento que transmite a rotagao a coluna de per-\\
furacao e permite o livre deslizamento do Kelly em seu interior. Em algumas operagées,\\
a mesa rotativa também deve suportar o peso da coluna de perfuracao.\\
\\
MESA ROTATIVA\\
\\
LLY.\\
\\
KELLY\\
mar se ft MASTER BEEING\\
\\
Figura 4.11: Mesa rotativa, Kelly, bucha de Kelly,\\
\\
bucha da mesa rotativa (master bushing)\\
\\
do a , €nquanto\\
devido a sua maior resistén, iq em soni\\
\\
cia A tracio, tore’ das maritimas é utilizada a seca hexagonal\\
\\
10 € flexo,\\
\\
Digitalizada com CamScanner\\

\\section*{Pagina 108}

TIPOS E SEU\\
SELECAO DE SONDaS oy PLATAFORMAS DE PERFURACAO 87\\
\\
cabeca de injecao\\
\\
Acabeca de injecao, também conhecida como swivel (Figura 4.12),\\
separa 0S elementos rotativos dos estacionarios Na sonda de pectin\\
de perfuragio é injetado no interior da coluna através da cabe. a de inj\\
sistemas alternativos para a aplicacao de rotagao a broca: top drive e motor de fundo.\\
\\
éo equipamento que\\
a¢a0. A parte superior\\
\\
Figura 4.12: Cabeca de inje¢ao (swivel).\\
\\
Top drive\\
\\
O top drive (Figura 4.13) é um motor suspenso que entra em contato direto com os tubos\\
de perfuracao, transmitindo rotacao e permitindo o manuseio dos tubos. A perfuracio\\
com um motor no topo da coluna (top drive) elimina a necessidade da mesa rotativa e do\\
Kelly. O sistema Top drive permite a perfuragao do pogo de trés em trés tubos, em vez de\\
um a um como na mesa rotativa. Esse sistema também permite a retirada ou descida da\\
coluna tanto com rota¢4o quanto com circulagao de fluido de perfuracao pelo seu interior.\\
\\{sso ¢ especialmente importante para pocos de alta inclinagao ou horizontais.\\
\\
Motor de fundo\\
\\
Neste caso, um motor hidrdulico tipo turbina ou de deslocamento positivo é colocado\\
acima da broca, O giro s6 ocorre na parte inferior do motor de fundo, que esta solidario\\
* broca. Assim, esse tipo de equipamento é amplamente empregado na perfuragao de\\
Posos direcionais, onde o objetivo a ser atingido nao se encontra necessariamente sob a\\
Mesma vertical que passa pela sonda de perfuracao. Como a coluna de perfuragdo nio\\
82, o torque imposto a ela é nulo e seu desgaste é reduzido.\\
\\
Digitalizada com CamScanner\\

\\section*{Pagina 109}

POGOS N\\
BB —\\_—PERFURACAO E COMPLETAGAO DE\\
\\
Figura 4.13: Top drive do sistema de rotagao da sonda de Perfuracio,\\
\\
4.2.5.4 Sistema de Geracao e Transmissao de Energia\\
\\
A energia necesséria para 0 acionamento dos equipamentos de uma sonda de Perfurags,\\
énormalmente fornecida por motores diesel. Nas sondas maritimas em que existe Prody.\\
g40 de gas, é comum e econémico o uso de turbinas a gas para a geracao de energia par,\\
toda a plataforma. Quando disponivel, a utilizagao de energia elétrica de redes Piblicas\\
pode ser vantajosa, principalmente quando o tempo de permanéncia da sonda em cads\\
localizacao for elevado.\\
\\
Uma caracteristica importante dos equipamentos de uma sonda, que afetao Process\\
de transmissao de energia, é a necessidade deles oOperarem com velocidadee torques varis-\\
veis. Dependendo do modo de transmissa0 de energia para os equipamentos, as sonds\\
de perfuracao sao classificadas em sondas mec4nicas ou diesel-elétricas.\\
\\
Sondas mecénicas\\
\\
pesleneeebancinnrie (Figura 4.14), a energia gerada pelos motores diesel é transmitidz\\
Por meio de acoplamentos hidrdulicos (convers\\
ound é constituido de diversos eixos, rodas dentadas\\
uem a energia para todos os sistemas da sonda. As embreagenspe\\
Pplados ou desacoplados do compound, proporcionando\\
motores a diesel,\\
\\
2 uma transmissao principal (compound)\\
res de torque) e embreagens, O compi\\
€ correntes que distrity\\
mitem que os motores sejam aco\\
maior eficiéncia na utilizacio dos\\
\\
Pequenos motores AC\\
\\
Figura 4,14; uem:; ;\\
a Hersems ta uma sonda MecAnica com cinco motores diesel.\\
\\
4\\
Digitalizada com CamScanner\\

\\section*{Pagina 110}

S E SEL\\
\\
TIPO: ECAO DE SONDAS OU PLATAFORMAS DE PERFURACAG. 89\\
sondas piesel-elétricas\\
js sondas diesel-elétricas geralmente sio do tipo AC/DC (Figura 4.15), no qual\\
\\
s ita em corrente alternada ¢ a utilizacdo é corrente continua. Os atlas y inn\\
juas turbinas @ gas acionam geradores de corrente alternada (AC) que alltttcarn an\\
parramento trifisico de 600V. Este barramento, alternativamente, também pode acalice\\
energia da rede publica, Pontes de Sillicon Controlled Rect ifier/ Retificador Controlado de\\
silicio (SCR) recebem a energia do barramento e a transformam em corrente continua.\\
que alimenta os equipamentos da sonda. Os equipamentos auxiliares da sonda ou plata.\\
forma, como iluminagao e hotelaria, que utilizam corrente alternada, recebem a\\
go barramento apés passar por um transformador.\\
\\
‘As sondas diesel-elétricas com sistemas tipo AC/ AC (geragio ¢ utilizacao ocorrem em\\
corrente alternada) tém uso incipiente, mas com tendéncia de aumentar no futuro. A\\
energia é fornecida por motores diesel, turbinas a gas ou através da rede publica. Por\\
utilizar motores AC, nfo ha necessidade de retificacio da corrente, mas do controle da\\
frequéncia aplicada aos motores.\\
\\
gio é fe\\
\\
energia\\
\\
Bombas de lama\\
Mesa rotativa\\
\\
BS\\} Guincho\\
\\
Transformador Motores AC\\
\\
Figura 4.15: Esquema de uma sonda AC/DC, tipica de sondas maritimas.\\
\\
4.2.5.5 Sistema de circulagaéo de Fluidos\\
\\
O sistema de circulagao de fluidos consiste em equipamentos que permitem a circulagio\\
€0 tratamento do fluido de perfuracdo. Na circulacdo normal, o fluido de perfuracao €\\
lo através da coluna de perfura¢o até a broca, retornando pelo espaco anular até\\
a superficie, trazendo consigo os cascalhos cortados pela broca. Na superficie, o fluido é\\
Tatado e armazenado em tanques (Figura 4.16).\\
Fase de injegao\\
0 fluida de perfuracSo é succionado dos tanques pelas bombas de lama e injetado ae\\
‘ina de perfuracdo até passar para o espaco anular entre 0 pogo ¢ a coluna Oe :\\
ities da broca. Durante a perfuracao, as vaz6es e pressdes de bombeio se ‘teal\\
da ate a geometria do poco. As bombas sao associadas em pone nie an\\
‘“ragio, quando sio requeridas grandes vaz6es. Conforme a pe\\
\\
és\\
\\
Digitalizada com CamScanner\\

\\section*{Pagina 111}

90 —\\_ PERFURACAO E COMPLETAGAO DE POGOS\\
\\
Figura 4.16: Sistema de circulacao de lama,\\
\\
altas pressdes, mas baixas vazbes sio necessarias, entao apenas uma bomba é\\
\\
Pistdes ¢ camisas de menor diametro sao empregados pa\\
\\
Fase de retorno\\
\\
Essa fase comeca quando o fluido de perfuracio sai pelos\\
\\
ele chega a peneira vibratéria, percorrendo o €spaco anular entre a\\
\\
a parede do poco ou o revestimento,\\
\\
Fase de tratamento\\
\\
Utilizada ¢\\
\\
ra atender as demandas do Poco,\\
\\
jatos da brocae termina quando\\
coluna de Perfuracio\\
\\
A fase de tratamento ou condicionamento do fluido de perfuracio envolve a remocio de\\
Sdlidos e gis i Y : x\\
\\
que se in\\
\\
corporam durante a perfuracao, bem como a adicao de produtos\\
quimicos Para aj SI ‘ ani\\
\\
uas propried:\\
\\
Apéso desareiador, o fluido\\
\\
Para remover particulas de dinmenegee siltador, que contém hidrociclones meno"\\
\\
es quivalentes\\
© mud cleaner, que é um dessiltador bi\\
do material é descartada en uma Peneira\\
\\
e\\
Algumas sondas utilizam aindy O° \\# fluido, red\\
\\
‘ i inte\\
silte. © equipamento one\\
Para recuperar te\\
luzindo a necessidade de adi"\\
\\
ainda nid\\
capturadas pelos hidrociclones. “entrifugadora para remover particulas menor’\\
\\
Um equipamento essencial na sonda\\
de perfuracio, Durante a Perfuracig éo despaseific\\
\\
. ido\\
ador, que remove o gas 40 fui .\\
\\
5 ui\\
de uma formacio com gis, ou quando hi U™\\
\\
Digitalizada com CamScanner\\

\\section*{Pagina 112}

Ti\\
POS E SELECAO DE SONDAS oy PLATAFORMAS DE PERFURACAO «= 9.\\
\\
fluxo de gas da formacio para o poco, as particulas di i m\\
mdreulacio no poco \\& pete, le gas se incorporam ao fluido e sua\\
\\
4,2.5.6 Sistema de seguranca do Poco\\
\\
O sistema de seguranca é constituido pelos Equipamentos de Seguranga de Cabeca de\\
Poco (ESCP) ¢ equipamentos complementares que permitem o fechamento e soar\\
do poco. Um dos elementos mais importantes Nesse sistema é o BOP, que consiste em\\
um conjunto de valvulas que pode fechar o PO¢o 4.17, a\\
\\
Os preventores sao acionados sem\\
\\
de fluido de uma formacao Para dentro do Pogo. Se esse fluxo nao for controlado ade-\\
quadamente, pode resultar em sérias Consequéncias, como danos aos equipamentos da\\
\\
sonda, riscos para a seguranga das pessoas, perda parcial ou total do Teservatorio, poluigao\\
e danos ao meio ambiente, entre outros,\\
\\
Cabeca de poco\\
\\
A cabeca de poco  composta por diversos equipamentos que permitem a ancoragem e\\
vedacao das colunas de revestimento na superficie, Esses equipamentos incluem a cabeca\\
de revestimento, carretel de perfuracio, adaptadores, carretel espacador e seus acessérios.\\
\\
A cabeca de revestimento é 0 primeiro equipamento a ser adaptado no topo do reves-\\
timento de superficie. Suas funges incluem sustentar os revestimentos intermedidrios e\\
de producao por meio de seus suspens6rios, proporcionar a vedacao do anular do reves-\\
timento intermediario ou de produc3o coma propria cabeca, permitindo o acesso a esse\\
anular, e servir como base para a instalac\\&o dos demais elementos da cabeca de poco e\\
dos preventores.\\
\\
O suspensor de revestimento é 0 componente responsavel por ancorar e vedar 0 reves-\\
timento, proporcionando a vedacao do anular do revestimento com 0 corpo da cabeca\\
onde foi ancorado. A vedacio é realizada automaticamente quando o peso do revesti-\\
mento é aplicado, comprimindo um elemento de borracha. O carretel de revestimento\\
€ semelhante a cabeca de revestimento, mas apresenta um flange adicional na parte in-\\
ferior. Ele também possui saidas laterais para acesso ao espaco anular e um alojamento\\
Para 0 suspensor do revestimento a ser descido posteriormente. Elementos de borracha\\
na parte inferior interna fornecem veda¢ao secundaria no topo do revestimento anterior-\\
\\
A cabeca de producao também é um carretel que possui uma sede na parte inferior\\
interna Para receber os elementos de vedacao secundaria que atuam no topo do Tevesti-\\
mento de producio, impedindo a passagem de pressGes elevadas, geralmente is ee\\
“Tesisténcia do flange inferior. Na parte superior interna, possui uma sede para ° tu \\textasciitilde{}\\
hanger, que sustenta a coluna de producao. Além disso, possui saidas laterais para acess'\\
nat? anlar et ena base\\
\\
Ocarretel de perfuragdo é um equipamento com flanges de conexdo no ha ee :\\
© @Presenta duas saidas laterais flangeadas que recebem as linhas de controle do pos:\\
linha de mata (kill line) ea linha de estrangulamento (choke line).\\
\\
Digitalizada com CamScanner\\

\\section*{Pagina 113}

GAO DE pocgos\\
\\
92 PERFURAGAO E COMPLETA\\
\\
preventores (BOP) ve\\
as de Brupsao Obturadores de Seguranga (Bops,\\
\\
ntores odem ser de dois tipos: preventor any, a\\
\\
Pp a funcao basica de fechar 0 espacg Fe €\\
\\
0 ser deslocado dentro de um corpo tii\\
\\
yout Preventers/ Preve!\\
do esp’\\
reventol\\
\\
Os Blow\\
mitem 0 fechamento\\
gaveta. Op\\
\\
ago anular €\\
r anular tem\\
\\
tao que, a\\
or um pista 4 ajusta contra a tubulacdo\\
rracha que se Jt C40 dentro q,\\
\\
preventor de\\
de um pogo  € compo:\\
drico, comprime um ¢!\\
poco. A Figura 4.17 ilustra um ¢\\
\\
Figura 4.17; Conjunto tipico de um BOP.\\
\\
alquer diametro de tubulacao e pode até fechar um poco\\
\\
sem coluna, embora este procedimento possa causar danos ao elemento de borracha. 0\\
preventor de gavetas tem a funcao de fechar o espaco anular do poco através de dois\\
pistées que, ao serem acionados hidraulicamente, deslocam duas gavetas, pressionando\\
uma contra a outra transversalmente ao eixo do pogo.\\
\\
Quanto ao arranjo dos preventores, geralmente, em terra sao utilizados trés: um pre-\\
ventor anular e dois de gavetas, No mar, existem duas possibilidades: nas plataforms\\
fixas ou apoiadas no fundo do mar, em que os equipamentos operam na superficie, si0\\
ee ¢ trés ou quatro de gavetas. Nas plataformas funn,\\
do mar, geralmente sao utilizados de ape = Seer operam © =\\
Resumindo, o BOP tem as sranlites fi ge entores anulares e trés ou quatro de gee\\
\\
in¢Oes:\\
\\
O preventor anular atua em qu:\\
\\
= Impedir que i is gies ;\\
Pedir que os fluidos da formacio atinjam a superficie de maneira descontrolada,\\
= Bombeament i tis,\\
‘0 de fluido para o interior do poco para promover 0 ¢ ontrole;\\
\\
* Permitir 0 controle das \\textasciitilde{}\\
ressdes san ; do\\
poco; R enquanto o fluido invasor é expulso para for\\
\\
Osinal\\
inal de comando pode ser hidrdulico, elétrico ou ético;\\
\\
* Em algumas so .\\
ndas, como Semissubmersiyeis € Navios-sondas, o BOP fica n° fund\\
\\
do mar. Em outras\\
sivel, plataformas ew ot ping plataformas maritimas com? subme!\\
+ *LP ¢ Spar, o BOP fica na superfici\\
i perficie.\\
\\
4\\
Digitalizada com CamScanner\\

\\section*{Pagina 114}

2 RPT\\
\\
TIPOS E SELECAO DE SONDAS OY PLATAFORMAS DE pER\\
FURACAO 93\\
\\
Kick\\
\\
Ui “kick” \\& a invasao de fluidos de formacao Para dentro do POCO. Isso oco;\\
-essio dos fluidos de formaco se torna maior que a pressa rre quando\\
a\\
\\
‘0 hidrostatica i\\
lo fluido\\
erfurag¢ao no pogo. Esse evento pode ser desencadeado Por diversas razées: mode\\
\\
» Perfuracao nao planejada de zonas com Pressdo anormalmente alta;\\
» Gas misturado 4 lama de perfuracio;\\
» Falha em manter o abastecimento adequado do Po¢o durante manobras (trip tank);\\
» Pistoneio (movimentos de pistao na coluna de perfuracao).\\
Os sinais de ocorréncia de “kick” incluem:\\
* O poco entrando em fluxo mesmo com as bombas desligadas;\\
* Aumento do volume de fluido de perfuracao nos tanques;\\
» Aumento na taxa de penetracao da perfuracao;\\
= Aumento na velocidade das bombas de lama.\\
\\
Para controlar um “kick”, é necessario:\\
\\
* Fechar 0 poco imediatamente utilizando 0 BOP ao primeiro sinal de suspeita;\\
\\
« Medir as pressdes na cabe¢a do poco (Shut-in Drill Pipe Pressure/ Pressio de Fecha-\\
mento no Tubo de Perfuracao (SIDPP) e Shut-In Casing Pressure/ Pressio de Fecha-\\
mento do Revestimento (SICP));\\
\\
* Expulsar o fluido invasor do pogo enquanto mantém a pressio constante no fundo\\
do poco;\\
\\
* Substituir o fluido de perfuracao por um fluido mais denso, mantendo a presséo\\
Constante no fundo;\\
\\
* Remover Possiveis bolsdes de gas aprisionados abaixo do BOP;\\
\\
* Ajustar a pressio no fundo controlando a vazio através do estrangulador (choke).\\
\\
4.2.5.7 Sistema de monitoramento\\
rte a arios para O controle da\\
roca, indicadores de tor\\
u claro que a eficiéncia\\
e varios para:\\
ara registrar\\
\\
rf monitoramento consiste nos equipamentos nccess:\\
orto. Isso inclui manémetros, indicadores de peso sobre f »\\
ea cone entre outros. Com o avango da perfuracao, ci aca\\
‘Onomia maximas poderiam ser alcancadas ao combinar per! fe a ‘\\
ens Perfuracio, Isso levou a necessidade de utilizar equipamentos P\\
\\
4\\
\\
Digitalizada com CamScanner\\

\\section*{Pagina 115}

94 ppnrunacad F \\& oMPLeTAcAd DE PE 08\\
\\
pain 00 BONDADON -”\\
= rn uaen votes\\
Liv ell oes Lend\\
"enn\\
vauenn\\
vor\\
torcem\\
nie\\
racer ronan OA\\
evan pemena CUA PPO omcarn WA ean ROTATION\\
\\
Vesna WAH py pam CANE LUTVANTE\\
\\
le monitoramento de uma sonda de perfuragic\\
agao,\\
\\
Figura 4.18 Paine! do sistem@ d\\
\\
ntos podem ser classificados em indicad\\
Ica\\
rimetro em questi, ¢ registradores, 0\\
na 2 DE + qu\\
\\{lustra o painel do sistema de me in\\
» MOnito:\\
\\
Hsses equipame\\
o valor do pa\\
9, A Higura 4.18\\
\\
pardmetros.\\
fe mostram\\
lores medide\\
\\
¢ controlar esses\\
res, que simplesment\\
\\
tragam curvas dos val\\
ramento de uma sonda de perfuragio.\\
Je peso no gancho € sobre a broca,\\
4,0\\
\\
Os principais indicadores slo 08 indicadores ¢\\
nandmetro que indica a pressiio de bombelo, o torquimetro para o torque na col\\
a luna de\\
\\
perfuracio, 0 torquimetro instalado nas chaves flutuantes com a fungio de medir ot\\
to\\
aplicado nas conexbes da coluna de perfuragio ou de revestimentos, ¢ 08 tacémet \\_\\
cele TOS para\\
\\
medir a velocidade da mesa rotativa ¢ da bomba de lama. O registrador mais impo\\
: ais or\\
0 que mostra a taxa de penetragio da broca, que é uma informagio crucial A ort\\
avaliar as\\
\\
mudangas nas formagGes perful das sf\\
rr 4 perfuradas, 0 desgaste da broca ¢ a adequacio dos parametros\\
Principals componentes da coluna de perfuracao\\
\\
Para realizar a\\
de Shen caciien um conjunto de ferramentas que compéem a col\\
sagfio, é neceseizio eplicar asia a extremidade inferior da coluna, e aaah wf\\
rochosas. A coluna de pefaracto be rotacio ¢ peso suficiente para cortar as tomate\\
bre a broca, transmitir rotacio a br. sempenha varias fungées, incluindo aplica Sas\\
estavel, e garantir a inclinacio ¢ skin conduzir 0 fluido de perfuracgao ina - iam\\
perfuragio é composta pelos seguints jo desejadas, Entre outros acess6! jp Mmanter Ss pogo\\
49 imentaa: essOrios, a coluna de\\
\\
* Comando de perfuraga\\
bi lo (também\\
re a broca ¢ proporciona rigidez conhecido como drill colla 2c\\
[4 Acoluna (Figura 4,19) r), que exerce peso sO”\\
\\
* Tubos de perfur:\\
‘ago pesado:\\
mitem rigidez 8, feitos de\\
para Material ,\\
0s tubos de perfuracio lous ane 4 fadiga, que trans\\
\\
* Tubos de perfu\\
mente com pirisichad dphAwl\\
Outros acessérios ir\\
conforme “ine, , também faze,\\
™ parte do aparato, permitindo arranjos\\
\\
res, bombas, mesa\\
\\} Mesa rotativa 1 além do apare\\
tC. Relativamente as br Ihamento de suporte, como me\\
‘Ocas empregadas, existem diversos\\
\\
que sao tub\\
01\\
© desgaste ¢ 8 de aco sem costura tratados inter™\\
@ corrosio (Figura 4.21)\\
\\
Digitalizada com CamScanner\\

\\section*{Pagina 116}

Figura 4.19: Comando de Perfuracao (Drill collar)\\
\\
ae Ae a\\
\\
Figura 4.20: Tubos de perfuracao pesados,\\
\\
Figura 4.21: Tubos de perfuracao.\\
\\
tipos, variando em termos de aplica¢ao, diametro e material, como as de aco-liga e as de\\
diamantes naturais ou artificiais (Figura 4.22).\\
\\
ee\\
\\
Figura 4.22: Brocas de perfuracao de pogos: (1) brocas sem partes méveis, (2) brocas de\\
cone com partes méveis, (3) brocas trit6nicas com dentes de ago, (4) broca de Pressure\\
Controlled Drilling/ Perfuracdo com Pressao Controlada (PDC).\\
\\
Asbrocas so ferramentas conectadas na extremidade inferior da coluna de perfuragio uti-\\
lizada para cortar mecanicamente a rocha na perfuracao de pocos de dleo e gas, compos-\\
‘as por elementos cortantes (dentes de aco, insertos de carbeto de tungsténio, compactos\\
de diamante sintético ou aletas de diamante sinterizado) e por um sistema de passagem de\\
Aluido, o qual permite que o fluido de perfuragao seja injetado através de orificios (jatos)\\
erando uma perda de carga, a qual é convertida em poténcia hidraulica proporcionando\\
\\
x no fundo do poco.\\
Normalmente sio classificadas em brocas sem partes méveis (nao possuem rolamen-\\
15 € partes miéveis) e brocas com partes méveis, que possuem de um a quatro cones\\
a estrutura cortante e os rolamentos, desta forma apresentando maior eficién-\\
\\
ia em relacdo as primeiras,\\
\\
Digitalizada com CamScanner\\

\\section*{Pagina 117}

os\\
96 peRFURACAO E ComPLETACAO DE poe\\
\\
-o si0 OS tipos mais comuns de equipam\\
tungsténio pode ser inserido y, eNtog\\
108\\
nS.\\
\\
novimento através de materia;\\
\\
Mais m,\\
, iy\\
alinos sint\\
\\
ermitir 0 r\\
\\
ee éticos também podem ser adicionados js brog\\
fortes. Os diamantes a ara rorndlas mais fortes 20 perfurar formagies rach, >\\
“ sess set de relevo € extremamente resistente, um diamante incr,\\
mais duras. =\\
\\
lo essas brocas de 40 a 50 vezes mais fc\\
ilizado, tornand' ites\\
tado na broca pode ser util fo)\\
\\
aco. A inici ‘i\\
\\
a pina a selegao da broca de perfuragao na ae inicial \\% operacio, levandy\\
reco! : ito difici\\
\\
1m conta a rocha € 0S sedimentos, P Ko efi we 9 roca de perfuracig\\
\\
per sso ter sido iniciado. Cada broca esté ligada a uma coluna de tubos\\
\\
apos todo o proce: m 30 pes, equivalentes a 9,1 metros de compr.\\
\\
a a e possue!\\
\\
de perfuracao que s20 montados 3 5\\
pa ers alterar uma broca danificada ou inadequada para 0 tipo de rocha encontrads,\\
todo o aparato, incluindo os tubos e demais componentes da coluna, deve ser retirado dp\\
\\
poco.\\
\\
As brocas de perfuracao com ¢\\
utilizados na extracio de gee. i.\\
positivos de ago para adicionar orga\\
\\
4.2.5.8 Acessorios de coluna de producao\\
\\
Substitutos (Subs) sao pequenos tubos que desempenham varias fungédes de acordo com\\
suas caracteristicas. Eles s40 usados para movimentar comandos, para a conexao de bro-\\
cas ¢ para conexao de tubos de diferentes roscas e diametros (Figura 4.23).\\
\\
\\}\\
\\
Figura 4.23: Substitutos (Subs).\\
\\
Estabilizadores sa\\
oer phe? hey resiad que conferem maior rigidez 4 coluna e ausiliam nam\\
se ee (Figura 4.24).\\
one frente cai resi amesma fungao dos estabilizadores e s40 usados\\
ron eran i a incluem a diminuigao do torque sobre @ oom\\
phe J Partes méveis que podem teachers ¢ na possibilidade de operagoes de pes\\
ae dope Figura 4.26) sio littanoeitas i ahaa a a\\
oi perfurado. Eles se dividem fel x rmitem aumentar o diametro 4¢ ™\\
: is Categorias:\\
izado quando se perfura ny desde a superficie, possui braces fixos °°\\
a descida do condutor;\\
\\
)\\
\\
Digitalizada com CamScanner\\

\\section*{Pagina 118}

TIPOS E SELECAO DE SONDAS OU PLATAFORMAS DE PERFURAGAO, 9\\
'7\\
\\
|\\
\\
Figura 4.24: Estabilizadores,\\
\\
Figura 4.25: Escareadores.\\
\\
* Under reamer: usado para alargar um trecho do poco a partir de um ponto abaixo\\
da superficie.\\
\\
Hole opener\\
\\
——S\\_—\\
——E\\_\\
\\
Under reamer\\
\\
Figura 4.26: Alargadores.\\
\\
4.2.5.9 Ferramentas de manuseio da coluna de perfuracao\\
As ferramentas de manuseio da coluna sao usadas para conectar e desconectar os varios\\
elementos da coluna de perfuracao. Algumas destas ferramentas incluem:\\
\\
Chaves flummantes: possuem a fungao de fornecer o torque necessario para a conexdo\\
€ desconexio das uniées cénicas dos tubos de perfuracdo. Um exemplo é 0 iron\\
roughneck, que desempenha a mesma funcio (Figura 4.27).\\
\\
Cunhas: mantém a coluna de perfuracdo constantemente suspensa na mesa rotativa (F\\
gura 4.28).\\
\\
a de\\
Colar de Seguranca: é um equipamento de seguranca colocado no topo da ens d\\
comandos quando ela esta suspensa pela cunha na mesa rotativa (Figura 4.28).\\
\\
Digitalizada com CamScanner\\

\\section*{Pagina 119}

s\\
pueTacko DE POCO\\
\\
98 rERFURAcAo £ COM\\
\\
fron roughneck\\
tes e iron roughneck.\\
\\
ante\\
\\
Chave flutu'\\
Figura 4.27: Chaves flutuan|\\
\\
Colar de seguranga\\
\\
Cunhas\\
\\
Figura 4.28: Cunhas e colar de segurangca.\\
\\
4.2.6 Flufdos de perfuracao\\
\\
Durante a fase de perfuracao de um poco, sao utilizados fluidos de perfuragao, também\\
conhecidos como lamas de perfuracao. Esses fluidos sao composi¢des complexas de pro-\\
dutos quimicos, liquidos, s6lidos e, por vezes, até gases, cujo principal objetivo é lubrifiar\\
a broca e garantir uma perfuragio eficiente e segura. A lama é injetada dentro da coluna\\
de perfuracio e retorna pelo espaco anular entre a coluna de perfuracio e as paredes do\\
\\
Pogo ou do revestimento. Os fluidos de perfuracdo devem desempenhar varias funcoes\\
essenciais:\\
\\
i: .\\
pc abt © do pogo, removendo e transportando para a superficie os detritos\\
tae ou s6lidos) cortados pela broca durante a perfuracao, e separar ess\\
\\
©quipamentos do sistema de extracdo de sélidos;\\
\\
* Lubrificar ¢ resfriar a coluna de perfuracio;\\
\\
* Bxerceruma Pressio hidrost4tj ‘\\
a entrada de fluidos de stitica adequada contra as zonas permedveis para prevent\\
\\
40 no Poco;\\
* Mantera estabilizacao das Paredes do poc\\
0;\\
* Manter sélidos em SUspensio;\\
\\
* Inibir a reatividade de formacdes argil\\
losas;\\
Resfriar ¢ lubrificar a broca;\\
\\
Digitalizada com CamScanner\\

\\section*{Pagina 120}

| eli\\
\\
TIPOS E SELECAO DE SONDAS OU PLATAFORMAS DE PERFURACAO 99\\
\\
+ Reduzir 0 atrito entre a coluna e bes paredes do POG¢o e permitir a coleta de informa-\\
ges sobre © poco através dos detritos.\\
\\
a selecio do tipo de fluido deve ser criteriosa, uma vez que um fluido de ma qualidade\\
\\
e causat problemas na perfuragao, resultando em aumento de custos. Caracteristicas\\
como estabilidade quimica, fluidez e custo/beneficio comparavel com a fase operacional\\
também devem ser consideradas.\\
\\
4.2.6.1 Classificagaéo de fluidos de perfuragao\\
\\
0s fluidos de perfuracao podem ser classificados em fluidos de base nfo aquosa, fluidos\\
de base aquosa € fluidos de base gasosa (ou pneumiticos), como ilustrado na Figura 4.29.\\
\\
Figura 4.29: Classificacao dos fluidos de perfuragio.\\
\\
| Fluidos de perfuragao de base aquosa\\
\\
| A grande maioria dos fluidos de perfuracao utilizados no mundo é formada por liquidos\\
\\
| 4 base de agua. Um fluido de base aquosa é constituido de Agua e diversos componen-\\
tes como argilas e coloides organicos, que sao adicionados para conferir as propriedades\\
viscosas ¢ de filtracio necessdrias. Esses fluidos consistem na mistura de sélidos, liquidos\\
€ aditivos quimicos, tendo a Agua como a fase continua. O Iiquido base pode ser a agua\\
salgada, Agua doce ou agua salgada saturada (salmoura), dependendo da disponibilidade\\
€ das necessidades relativas ao fluido de perfuracio.\\
\\
Fluidos de perfuracéo a base de dleo\\
\\
Os fluidos de perfuracio a base de dleo sao similares em composicio as a base de agua,\\
Sxceto pela fase continua que passa a ser o dleo. A dgua esta presente na lama a base de\\
leo sob a forma de uma emulsio. Os fluidos 4 base de dleo ganharam destaque, apesar\\
\\
Custarem de 2 a 4 vezes mais do que os de base aquosa. Os fluidos a base de dleo sio\\
"nuito utilizados ¢ indicados para a perfuracio maritima. As vantagens de desempenho\\
4 perfuracio com lamas a base de dleo em comparacaio com as de base aquosa \\$0:\\
\\
Digitalizada com CamScanner\\

\\section*{Pagina 121}

¢ POGOS\\
100 PERFURACAO E COMPLETAGAO pe POG\\
\\
* Com pilid ym as forma OCS sensivels 4 Agua,\\
ibi Je cor as formagor’’\\
it patibili jad\\
\\
agao da corrosdo;\\
\\
* Minimiz; :\\
utural na perfuragao de pogos profundos @ ;\\
Om\\
\\
i ica ¢ est"\\
* Maior estabilidade térmica ec\\
\\
altas temperaturas; ae\\
\\
6a ‘os direcionais;\\
* Melhor lubrificacao, facilitando a perfuragao de pogos direcionais,\\
o adequado;\\
\\
ento apos tratament\\
proporcionando um aumento da Rate of Peney\\
Ta.\\
\\
* Reaproveitam\\
is rapidamente,\\
\\
* Perfuracio feita ma\\
Penetragao (ROP).\\
\\
tion/ Velocidade de\\
\\
Fluidos de base sintética\\
\\
"Trata-se de um fluido cuja base principal é um 6leo sintético. B o tipo de fluido mais\\
utilizado em plataformas de perfuracao ocednica (offshore), j4 que possuem propriedades\\
semelhantes as dos fluidos 4 base de dleo, porém, seus vapores apresentam grau de to.\\
xicidade inferiores. Esse fato é de grande relevancia no caso de manuseio em espacos\\
fechados, como normalmente ocorre em plataformas de perfuracao oceanica. Os mes.\\
mos problemas de contaminagao para andlises quimicas de rochas ou Gleo original da\\
formaciio ocorrem com 0 uso de fluido de base sintética.\\
\\
Fluidos de perfuracao a base de ar\\
\\
pas ne ar a é mee oar ou gas como fluido principal durante a perfu-\\
sinnagbes. 0 fluido a nd ae lensidade e seu uso é recomendado somente em algumas\\
veras, formagies pro aii ar pode sen aplicado em zonas com perda de circulagio se-\\
danos, formagées muito inca pressio muito baixa ou com grande susceptibilidade\\
¢ regides glaciais com como 0 basalto ou diabasio, regides com escassez de agua\\
\\
€spessas camadas de gelo. Quando se utiliza o ar puro ou outro\\
\\
gas.\\
\\
Propriedades de fluidos de perfuracdo\\
\\
Para maior eficiénci.\\
ficiéncia das atiyj\\
mantidas: tividades de perfuraca : . o\\
C40, as seguintes propriedades devem set\\
\\
Densidade:\\
: este parametro ji\\
\\
i TO influenci:\\
\\
fluido ao longo do pos a diretamente na coluna hidrostatica exercida pelo\\
\\
d ‘0. O peso\\
aos “specifico, ou densidade, é determinado na balang!\\
\\
Pardmetros reolégj,\\
oldgicos: es\\
; estas\\
nhecido por viscosimetra, @ of! sio medidas com o Redmetro, também\\
lente transporte dos cascalhos a superficie bet\\
\\
d\\
\\
Digitalizada com CamScanner\\

\\section*{Pagina 122}

TIPOS E SELECAO DE SONDAS OU PLATAFORM\\
AS DE PERFURACAG\\
101\\
» sua sustentagio durante as parad.\\
\\
come :\\
co do fuido.\\
\\
reoldgi\\
gocosidade marsh: mesmo nao fazendo parte de qualquer calculo de\\
Vis amento, @ viscosidade marsh é utilizada como um Posstvel indicativo de alteracio\\
no comportamento do fluido. Essa viscosidade é medida pelo funil marsh e leva oa\\
conta o tempo gasto para 946 ml de fluido escoar pelo citado funil,\\
\\
las de circulacio, fica Por conta do controle\\
\\
reologia ou carre\\
\\
parimetros de filtragao: (API e High Temperature High Pressure/ Alta ‘Temperatura Alta\\
Pressio (HTHP)) determinagao do volume de agua livre que por agio da pressio\\
hidrostatica, estatica e dindmica, forma um reboco ao longo das paredes do poco,\\
Essa determinacao é efetuada em condigdes normais, a 100 psi (689,5 kPa), e em\\
condigées especiais quando adotamos uma pressio de 500 psi associada a uma tem-\\
peratura de 250 °F (121,11 °C). Estas medidas so efetuadas pelos equipamentos de\\
nome Filtro Prensa.\\
\\
Teor de sdlidos: esta medida é efetuada pelo equipamento chamado de retorta a qual\\
fornece os percentuais de Agua e sdlidos existentes no fluido, lidos diretamente numa\\
proveta de 10 ml.\\
\\
Alcalinidade: importante para inibir corrosdo da coluna de perfuracio. Este deve ser\\
mantido entre 7 e 10.\\
\\
Teor de cloretos ou salinidade: a determinagio do teor de cloretos é efetuada pelo mé-\\
todo volumétrico, onde a precipitago do mesmo na forma de AgCl, com uso de\\
AgNO3 em presenga de dicromato de sédio quantifica a presenca do fon no fluido.\\
Usado para identificar o teor salino da gua de preparacao do fluido de perfuragio;\\
controlara salinidade de fluidos inibidos com sal; identificar influxos de gua salgada;\\
identificar a perfuracao de uma rocha ou domo salino.\\
\\
Lubricidade: determina um melhor desempenho da broca com consequente melhora na\\
taxa de penetracao.\\
\\
Toxicidade: quantifica o grau de agressao do fluido ao meio ambiente.\\
\\
Biodegradabilidade: a quebra facil das cadeias dos produtos utilizados determinard o\\
impacto temporal ao meio ambiente.\\
\\
4.2.7 Operacées normais de perfuracao\\
\\
As operacées normais que envolvem a atividade de perfuragao sao de rotina. A conexio\\
40s tubos de Perfuracdo é um exemplo bem tipico de tais operagées. Cabe a equipe da\\
Sonda executé-las, acrescentando secées de trés tubos 4 coluna de perfuracio, deste modo\\
é 408 poucos as formacGes, Quando se percebe o término da vida itil da broca,\\
necessiria a sua substituicio, operaco conhecida como manobra da coluna. -\\
Tal °peracdo consiste em retirar toda a coluna do pogo, a fim de que uma nova a\\
pos tad, Tanto na descida quanto na retirada da coluna, ass ecdes de tubos, for = a ;\\
pert Unidas, sio devidamente posicionadas na torre, na posigdo ver tical, de modo ¢\\
rows maior agilidade ¢ racionalidade no manuseio das ferramentas\\
\\
\\_f\\
\\
Digitalizada com CamScanner\\

\\section*{Pagina 123}

102 peRFURACAO E COMPLETACAO pe pocos\\
ais de perfuragao\\
\\
goes especi\\
\\
4.2.8 Opera\\
das, indispensévei\\
\\
i casos especificos a\\
Sio operacdes diferenc as pecificos. A seguir, alguns ¢,\\
em.\\
\\
plos:\\
\\
4.2.8.1 Perfilagem\\
m seu interior alguns equipamentos espe,\\
ais.\\
\\
rfurado 0 pogo, sao descidos\\
las formagées que fario parte dy\\
\\
avaliar algummas propriedades d\\
mica do mesmo.\\
\\
tamento de caracteristicas \\& propriedades das rochas\\
camente em fun¢ao da profundidade, mediante 9 re\\
\\
locamento de um sensor dentro do poco. As principais caracteristicas registradas si\\
\\
resistividade elétrica, Raios gama (GR), potencial eletroquimico (porosidade), velo nH\\
sismica (sOnico). A partir da anélise dos perfis, é possivel identificar, por exe ‘i =\\
formagoes rochosas atravessadas, calcular suas espessuras € porosidades, ¢ id 3 she\\
tipos de fluidos presentes nos poros das rochas. nee\\
\\
Uma vez pe!\\
cuja finalidade é medir €\\
caracterizacio € avaliagao econ\\
\\
A operacio consiste no levant\\
\\
furadas, que s40 registradas graf\\
\\
: 4.2.8.2 Revestimento de poco\\
\\
A principal necessidade de revesti\\
tir F\\
grax paredes. Os sees de pblerhers pogo, total ou parcialmente, é devida a protesio de\\
agus Sasa Ge 4 nnamento so consideraveis, havendo também diversos\\
: bre aa arama do revestimento. :\\
Sendo perfurado z\\
Jocados uns por dentro dos eh vao sendo revestidos com tubos de ago especial,\\
ear 0 tubo inicial tem Ce care colnas de revestimento. No comm\\
mando um ajuste tipo telescépic extensio e diametro maior do que os posterior\\
que o diametro diminui, 0 Ratatat para formar a coluna de revestimento. A medi\\
mado de intermediirio ento inicial, antes dito di i\\
J ¢; depois, d ca ito de superficie, passa 2 ser cht\\
revestimento de , de revestimento d 5 ‘\\
aie poco, le producao. A Figura 4.30 ilustra°\\
Protecio das paredes, as princi\\
4 cipai 2 .\\
pais funcGes da coluna de revestimento sao:\\
\\
* Nao permitir a\\
perda de flui\\
uido de Perfuracao para as formacoes;\\
\\
* Permitir o reto;\\
mo do fluido de\\
Perfuraco a superficie, para o devido rratamen'®!\\
\\
* Evitar a contami:\\
taminaca,\\
cdo da\\
agua de possiveis lencéis fresticos:\\
\\
* Dar suporte\\
al 4S\\
Pata os equipamentos de cabeca d\\
; e¢a do poco etc.\\
\\
Digitalizada com CamScanner\\

\\section*{Pagina 124}

J TIPOS E SELEGAO DE SONDAS OU PLATAFORMAS DE PERFURAGAO 103\\
\\
Figura 4.30: Revestimento de pogo.\\
\\
‘o da coluna e isolando as zonas porosas € permeaveis atravessadas pelo\\
Esta operacao é realizada por meio de tubos condutores auxiliares.\\
superficie, toda a extensio é cimentada, e nos demais, normalmente\\
s6 a parte inferior em intervalos predefinidos. O objetivo da cimentacio é fixar o reves-\\
\\
de fluidos entre as diversas zonas permedveis atravessadas\\
\\
timento e evitar a migracao\\
pelo poco por tras do revestimento. A cimenta¢ao do espaco anular é realizada por meio\\
\\
do bombeamento de pasta de cimento e agua.\\
\\
uma melhor fixaca\\
\\
poco (Figura 4.31).\\
No revestimento de\\
\\
Figura 4.31: Cimenta¢gao de revestimento.\\
\\
aditivos\\
mento\\
\\
pois \\&\\
s mate-\\
\\
“i ne usado na industria do petréleo consiste principalmente de cimento, ;\\
fete varios tipos de cimento e aditivos dependem da aplicacio final. od\\
Pront 0 material escolhido em 99\\% das operacoes de cimentagao primaria,\\
Tiais, tamente obtido em todo o mundo e é comparativamente mais barato. Outros m\\
reto, pee idiia. pozolana, resinas epOxi ou cimentos especiais de magnésio ou oxiclo-\\
Py f ser utilizados quando necessarios, Os principais componentes do cimento\\
como a bxido de cAlcio, silica, alumina ¢ ferro, que combinados formam composts\\
aluminate, to tricdlcico (C3\\$), silicato diclcico (C2S), aluminato tricalcico (C 3A) e fo\\
Mpegs oe (C4AF). Essa proporco de compostos no cimento determina suas\\
» Como resisténcia inicial, calor de hidrata¢ao, resisténcia aos sulfatos etc.\\
\\
Digitalizada com CamScanner\\

\\section*{Pagina 125}

0S\\
104 —\\_PERFURACAO E oMPLETACAO DE oF\\
\\
‘0 segundo API\\
\\
Classificaga\\
de AaJ,com base em sua composic3,\\
S680 Quin.\\
\\
rtland ;\\
fundidade e temperatura dos pocos, ey\\
\\
uy os cimentos Po!\\
mo pro!\\
\\
A API classifico\\
dig6es de us co\\
\\
adequada as con!\\
Classe A: Até 6.000 pés quando nao so requeridas propriedades especiais (Cimentp\\
3 : : ‘.\\
\\
tland comum);\\
\\
Classe B: Até 6.000 pés, quando é requerida moderada a alta resisténcia a0s sulfatos\\
\\
Classe C: Até 6.000 pés, quando é requerida alta resisténcia inicial. Apresenta alta wis\\
\\
téncia aos sulfatos;\\
\\
Classe D: 6.000 a 10.000 pés. Temperaturas moderadamente elevadas ¢ altas pressie,\\
Alta resisténcia 20S sulfatos; ;\\
\\
Classe E: 6.000 a 14.000 pés, pressdo ¢ tem\\
aos sulfatos;\\
\\
peratura elevadas. Apresenta alta resisténcig\\
\\
Classe F: 10.000 a 16.000 pés, sob condig6es extremamente altas de pressao ¢ tempera\\
tura. Alta resisténcia aos sulfatos;\\
\\
ditivos até profundidades de 8.000 pés. Compativeis\\
rdadores de pega, podem ser usados praticamente\\
te utilizados na indistria\\
\\
Classes Ge H: Para utilizacao sem a\\
com aditivos aceleradores ou reta’\\
em todas as condicées das classes A até B, sendo amplamen\\
petrolifera;\\
\\
Classe J: Para uso em profundi\\
\\
profundidades de 12.000 a 16.000 pés, sob condic¢o\\
P i x ; ‘ondi,\\
elevadas de pressdo e temperatura. P eo Se\\
\\
Principais aditivos para cimenta¢géo\\
\\
Os principais aditivos i\\
para cimentacio ii\\
filtrado ¢ outros que eaten, ote redutores de friccdo, controladores de\\
\\
Ges com\\
\\
Redutores de fric¢a\\
ic¢4o (ou dispe: .\\
tsantes): Afinam a pasta e permitem maiores va2!\\
duzindo\\
\\
menores perdas de car; or.\\
> ga, facilitand 5 H\\
o risco de fratura das toile Bia lo. a remocio do fluido de perfuragao € te\\
\\
Controladores de fil\\
; Itrado: Evita 4\\
Head mantendo a Mn lel pp prematura da pasta em zonas permel\\
os €m conjunto com os dis Le danos a “pay zone”. Normalment¢ ¥\\
antes;\\
\\
s de pew\\
detect"?\\
cont”\\
\\
Outros: Anti\\
antes (evita\\
AN CepUT shy 7\\
pele n> , descontaminanes, ta da pasta), adensantes, controladore\\
lada rr wae to. Também sahie ‘adores radioativos e corantes para\\
— a0 do itinaees' — areias de granulometria\\
2 altas temperaturas.\\
\\
4\\
\\
Digitalizada com CamScanner\\

\\section*{Pagina 126}

TIPOS E SELECAO DE SONDAS oy PLATAFORMAS DE PERFURA\\
\\
- CAO 105\\
ipo de cimenta¢ao\\
R cimentagio primaria é realizada imediatamente apés a desci\\
simentO no pogo. A cimentacdo secundaria, por outro lado,\\
“a cimentacao primdria quando necessario. Isso pode ser fe\\
ou tamponamento do pogo. O canhoncio envolve perfurar\\
cimento na falha existente no espaco anular entre o furo d\\
tamponamento do pogo é usado quando se deseja abandon:\\
\\
ida de cada coluna de Teves-\\
é usada para corrigir falhas\\
Ito por meio de canhoneio\\
© revestimento Para injetar\\
10 poco eo Tevestimento, ©\\
at O poco ou isolar zonas es-\\
\\
pecificas.\\
\\
Objetivos/Fun¢des da cimenta¢ao primdria e secundaria\\
\\
A cimenta¢ao primaria tem como objetivo posicionar uma pasta de cimento (Agua, ci-\\
mento e aditivos especiais) de forma integra no €spaco anular entre o revestimento e a\\
formacio. A pasta de cimento é bombeada até sua Pposi¢ao final e, apés repouso, adquire\\
resisténcia compressiva suficiente para restringir 0 movimento de fluidos entre diferentes\\
formages (por exemplo, zonas de diferentes pressées) e isolar aquiferos. As funcées da\\
cimentac4o priméria incluem:\\
\\
* Prover aderéncia entre o cimento e a formacio;\\
* Dar suporte mecAnico ao revestimento;\\
* Isolar zonas com diferentes pressdes;\\
\\
* Dar suporte ao revestimento em casos de kick off.\\
Cimenta¢do secundaria\\
\\
A cimentacao secundaria é realizada para corrigir problemas na cimentagao primaria,\\
quando necessério. Essas so operagGes emergenciais que visam permitir a continuidade\\
das operacées, conforme descrito a seguir:\\
\\
Tampées de cimento: consistem no bombeamento de determinado volume de pasta, co-\\
brindo uma secdo especifica do pogo. Essa técnica é empregada em situagdes como\\
perda de circulacio, abandono total ou parcial do poco, base para desvios, entre\\
Outras;\\
\\
Recimentacio: tealizada com a finalidade de corrigir a cimentagao primaria quando o ci-\\
mento nao atinge a altura desejada no espaco anular entre o revestimento e a forma-\\
S40. Nesse procedimento, o revestimento é perfurado em dois pontos para permuar\\
injeco de cimento;\\
\\
ada de cimento sob pres\\
\\
Compressio de cimento ou “squeeze”: envolvea inje¢io forg in\\
£m 0 objetivo de corrigir problemas na cimenta¢ao priméria. Essa téc\\
Para selar yazamentos no revestimento ou para imped a produs\\
\\
‘Me comecaram a produzir Agua.\\
\\
Jio de zo!\\
\\
4\\
\\
Digitalizada com CamScanner\\

\\section*{Pagina 127}

p POGOS ]\\
preTacad 0!\\
\\
acAo E COM!\\
\\
106 rearul\\
\\
Go\\
4.2.8.4 Testemunhage™ dep\\
\\
so de uma amostra da formagio rox hoa ¢\\
10\\
\\
in\\
onc: 7 Sy\\
vere na obtengé tant : coletar informacées 4 pe\\
A testemunhagem consisce 0 0, com objetivo de cole NAGGES Lites Par\\
\\
hi ipes dee iF)\\
+ estemun «. analigadas por equipes de engenh,\\
ficie, conhecida como (es formagoes 840 analisadas | July genharia de fete\\
5s aff co, Essas ir\\
avaliagio do pos\\
\\
-os especialis his a broca vazad,\\
vatdrios, gedlogos ¢ 0UtTO' cect é realizada utilizando uma broca vazada ¢ dois\\
Jo de testemunna\\
A operagao de\\
\\
-oluna de perfuracd bany\\
BS ante com a coluna de perfuragio, enqug\\
ilete externo gi\\
letes. Um barrilete\\
\\
a juntame Pls eter ba\\
testemunho. Conforme a broca avanga, 0 testemunho 6 algj,4,\\
ilete i Joja o tes\\
barrilete interno al\\
no interior do barrilete interno.\\
\\
4.3 Completacao de pogos\\
\\
Jo de um pogo, segue-se 4 fase de completacio, que envolve uma shrie\\
Apés a perfuracao pjetivo de permitir a produgio econémica ¢ segura de hidrocarhy\\
peg nome ean cara fulbs no reservatorio quando necessario. Algumas dis\\
fiefos, bem a nesse proceso incluem a descida do revestimento de producio,\\
Oaepeate “canhoncio” (uso de uma carga explosiva para perfurar ° revestimento ¢\\
colocar o reservatério produtor em comunica¢io com 0 pogo), ¢ a instalagdo da cabeca\\
de pogo. , . ;\\
A completagio de um po¢o envolve um conjunto de operag6es para equipar o poco\\
visando a producio de éleo ou gis (ou ainda a injecao de fluidos nos reservatérios).\\
\\
4.3.1 Objetivos da completacao\\
Os objetivos da completacao de um poco sao:\\
f Preparar © poco para producio, utilizando técnicas especificas de acordo com 0 po-\\
tencial produtor e as caracteristicas da formagao do reservatorio;\\
* Evitar danos ao reservatério;\\
* Maximizar a producio ou injegao de fluidos;\\
* Minimizar os custos;\\
\\
* Realizar uma completacao\\
\\
Permanente que nio exija inte Bes frequentes pot\\
problemas mecinicos dura: que nfo exija intervengées freq\\
\\
inte toda a vida produtiva do pogo;\\
* Realizar a completacao de forma simples e Segura;\\
\\
= Cumprir com as leis overnamentais aplicdveis,\\
3.1.1 Categorias basicas de completaca\\
[-)\\
\\
Existem diversos méto,\\
\\
dos d -cifi\\
cados em duas Categorias basic Pletacio utilizados na industria, que podem ser css\\
como ilustrado na Figura 4.32, hole” (sem revestimento) ¢ “cased hole” (revest as\\
\\
Digitalizada com CamScanner\\

\\section*{Pagina 128}

COMPLETACAO DE POCOS 107\\
\\
Figura 4.32: Categorias basicas de completacio.\\
\\
4.3.2 Equipamento de cabeca de poco\\
\\
Naparte superior do poco, é instalado um equipamento conhecido como cabeca do poco.\\
Aconfiguragao desse equipamento varia de acordo com a fase em que 0 pogo esté, seja\\
durante a perfuragao ou durante a producao. A funcio primordial da cabeca do poco é\\
vedar as colunas de revestimento e servir como ancoragem para as mesmas.\\
\\
Durante a producao, é instalado sobre a cabeca do poco um conjunto de valvulas co-\\
nhecido como arvore de natal. Essa arvore de natal contém dispositivos de seguranca e\\
controle da producao, bem como outros componentes necessarios.\\
\\
Em pogos terrestres, a 4rvore de natal fica na superficie. Ja em pocos submarinos,\\
os equipamentos podem ser mais complexos, podendo ser instalados na superficie (em\\
plataformas) ou no fundo do mar (submarinos). As arvores de natal submarinas podem\\
ser do tipo “arvore de natal seca”, protegidas da agua e pressao externa, ou “Arvore de\\
natal molhada”, expostas 4 4gua. A Figura 4.33 ilustra uma drvore de natal convencional.\\
\\
Figura 4.33: Arvore de natal convencional.\\
\\
4.3.2.1 Etapas de uma completacao\\
\\
4 completacao de um poco envolve uma série de operagoes realizadas apés a perfuragao.\\
\\
™a completacio tipica de um poco maritimo, com arvore de natal convencional e pel\\
Pamentos de gas lift, segue as seguintes fases em sequéncia cronolégica. Embora possa\\
td Pequenas variacées, essas fases sao essencialmente as mesmas para Posos\\
\\
Digitalizada com CamScanner\\

\\section*{Pagina 129}

E POGOS\\
108 pPERFURAGAO E COMPLETAGAO ie)\\
ie superficie\\
=0 dos equipamentos di ip\\
stalagao\\
\\
4.3.2.2 In ‘0 instalados, incluindo a cabega de ;\\
\\
seguro a0 interior do pogo para a execusig lu\\
Aguas rasas, é possivel trazer a cabeca de\\
\\
2 ; "\\
‘os que esto ancorados nos equipamentos »\\
\\
ie SA\\
superficie\\
Inicialmente, OS equips oa ee,\\
; ibilitar\\
de possibil it\\
30 co BOP, a fim nes\\
préximas fases. Em pogos pou cbt\\
até a superficie, estendendo os\\
\\
talados no fundo do mar (tie-back)-\\
\\
lo pogo\\
4.3.2.3 Condicionamento do p'\\
jpamentos de superficie, ocorre a fase de condicionamento dy\\
Apés a instalagao dos ove substituigao do fluido de perfuracao (lama) presente no inn,\\
pe noone aie. Para 0 condicionamento, uma coluna coma bro “aes\\
rior do poco por um ae mostrado na Figura 4.34, de modo a limpar o interior dore\\
é baixada no pogo, com e prepard-lo para receber os equipamentos necessarios. A broca\\
imento de produgao ¢ P . tampées mecanicos, se existentes no interior\\
€ usada para cortar tampoes de cimento € p'\\
\\
do poco, assim como quaisquer residuos de cimenta¢ao.\\
\\
Orit pipe\\
\\
Condiclonador\\
do topo Ener\\
\\
Figura 4.34: Condicionamento do pogo.\\
\\
4.3.2.4 Avaliagao da qualidade da cimentacao\\
\\
A cimentaao desempenha um papel crucial ao promover a vedacdo hidraulica ente 0s\\
diferentes intervalos permedveis, inclusive dentro de um mesmo intervalo, prevenindo a\\
\\
migracao de fluidos atras do revestimento, Além di P ecanico\\
park paosan pak ecalit . Além disso, ela proporciona suporte m\\
\\
tos de medi¢ao baseados em aid qualidade da cimentacdo, so empregados instrumet\\
\\
; riedades acisstica: F oe io cient\\
a0 revestimento ¢ a formacao rochosa circundant s, que avaliam a aderéncia do\\
e.\\
\\
4.3.2.5 Canhoneio\\
\\
Digitalizada com CamScanner\\

\\section*{Pagina 130}

OPERACAO START UP \\& FIRST on 10\\
9\\
jive o revestimento, bem como suas Paredes metalic,\\
ee ser produzido. A Figura 4.35 ilustra o resultado\\
por pee produtora. Geralmente, diversos disparos expl\\
. oan toda a espessura da camada produtora. Os diferentes tipos de canh i\\
tet sio demonstrados na Figura 4.36, ose wa\\
\\
a8, ¢ chepue ao interior do\\
do processo de canhoneio\\
lOSiVOs sig Necessirios para\\
\\
@)\\
(Maing Comveyed Partorster) Cohuna de predugse\\
\\
Figura 4.36: Tipos de canhoneio utilizado.\\
\\
4.3.2.6 Instalagdo da coluna de produgao\\
\\
Dentro do revestimento de producio, desce a coluna de producio, um tubo de diametro\\
reduzido, geralmente com cerca de 3,5 polegadas, através do qual o petréleo é produzido.\\
A producao pode ocorrer de forma natural ou artificial, envolvendo a elevagio do petré-\\
leo por bombeamento ou a injecao de gas no pogo. A coluna de producio € composta\\
Principalmente por tubos metilicos, nos quais os componentes sao conectados. Ela é\\
inserida no interior do revestimento de produgio e desempenha as seguintes fungées:\\
\\
* Transportar os fluidos produzidos até a superficie e proteger o revestimento contra\\
fluidos agressivos e press6es elevadas;\\
\\
* Permitir a instalacdo de equipamentos para a elevagio artificial;\\
\\
* Facilitara circulagao de fluidos para o amortecimento do poco durante intervengo:\\
futuras,\\
\\
“4 Operacéo Start-up \\& First Oil\\
\\
‘ beracio de start-up marca o inicio da producio de um poco apos\\
Talmente,\\
\\
Tealizada injetando nitrogénio na tubulacio para alivia\\
\\
a completagio. Ge\\
ro peso da coluna\\
\\
Digitalizada com CamScanner\\

\\section*{Pagina 131}

OMPLETAGAO DE pogos\\
\\
110 penruracAad \\& C\\
\\
atta natural do pet réleo no pogo. O termo\\
ye um pogo, campo OU bloco com Firs 1\\
F (\\
ro milhar de barri: ecaa ny\\
= Darris pr Pr\\
Pp Oduzidg, Od,\\
€\\
\\
Ne,\\
\\
estimular a surge\\
0 momen!\\
te, é ma\\
\\
roem q\\
\\
hidrostatica ¢\\
rcado pelo primei\\
\\
meiro Oleo) simboliza\\
trdleo. Frequentemen\\
perados na superficie.\\
\\
4.5 Abandono de pogos\\
a série de operagGes que culminam no dest\\
esli\\
\\
os envolve um\\
instalagdes de produgao e recuperacig ¢, Men,\\
S areas.\\
\\
O abandono de po\\
Ss\\
\\
definitivo do poco, desativacao da:\\
padas por essas instalac6es. O procedimento de abandono deve estar em ¢\\
com os regulamentos e padroes espectficos de cada pafs. Um campo é con, pas ida\\
3 9 Si\\
donado quando os pogos nao produzem mais petrdleo, ou quando a prody Cradg aban\\
principalmente de 4gua ou gas. O abandono pode ser definitivo ou tem, 1ga0 COnsistg\\
cesso de abandono, as estruturas metalicas das plataformas de pocos e ee No 0,\\
a ee e, em alguns casos, recicladas (por exemplo, siderurgia sa Ocessamentg 5,\\
cas). Em plataform: i Fa estrui\\
PI as offshore, as estruturas podem ser deixadas no local, cri a Meti.\\
» “TIAN recif\\
es\\
\\
artificiais que beneficiam o meio ambiente marinho.\\
\\
: seo,\\
4.6 Risco na atividade de perfuracgao e completaca\\
cao\\
\\
Aatividade de z\\
perfuracaoe G\\
‘iptianes un i pete de pocos de petréleo envolve riscos sien;\\
a grande hhicieeo de eh a io ambiente € a economia. Esses riscos 08 Significativos\\
ipamentos de alta pressao e temperatura, 4 esto relacionados\\
» 4 presenca de fluidos\\
\\
Perfuracao de pocos incluem:\\
\\
* Risco de perfurar um pod\\
i ‘0 exploratori\\
\\
netos; exploratorio ou de desenvolvimento que na hi\\
\\
jue no contenha\\
\\
. Risco de descoberta de uma\\
\\
cobrir os quantidade i\\
cus' lade \\_\\
tos de producao; insuficiente de 6leo em uma reserva pit\\
\\
" Risco politi\\
tico associ:\\
Ociado as ;\\
Tecursos petroliferos, 4S incertezas legais ¢ j\\
institucionais de paises detentores*\\
\\
" Riscos di\\
S de explosdes e acid np\\
‘aves,\\
\\
Digitalizada com CamScanner\\

\\section*{Pagina 132}

CUSTOS E CONTRATOS DE PERFURACAG \\& ComPLer\\
SCAO DE POCOS 111\\
\\
7 Custos e contratos de perfuracéo e completa:\\
gos\\
\\
4 Gao de po-\\
\\
os reais envolvidos na perfuracao e completacao de\\
\\
Os cust as\\
egorias distintas:\\
\\
tres call\\
\\
Custos fixos: relacionados a revestimentos e tubulagées, perfis ¢ registros,\\
brocas de perfuracao, mobilizacao de equipamentos, movimentacio de\\
ou sondas;\\
\\
Po¢os podem ser divididos em\\
\\
cimentacio,\\
plataformas\\
\\
Custos didrios: incluem servicos de contratantes, tempo de uso de plataformas insu\\
mos;\\
\\
Despesas gerais: abrangem custos com escritérios, salarios, beneficios, satide e viagens.\\
\\
Além disso, existem outros custos relacionados 4 mobilizacio ¢ desmobilizacao de plata-\\
formas ou sondas. Esses custos podem ser significativos, variando entre 5 ¢ 10 milhées\\
de délares ou mais, dependendo do programa de perfuragao especifico. Os custos de um\\
poco podem variar consideravelmente devido a diversos fatores, incluindo:\\
\\
* Tipo de pogo (exploratério, avaliativo, desenvolvimento);\\
\\
« Trajet6ria do poco (vertical, desviada, horizontal, multilateral);\\
\\
* Profundidade total do poco;\\
\\
* Condicdes subsuperficiais (temperaturas, pressdes, agressividade dos fluidos);\\
* Tipo de operac4o (em terra, maritima);\\
\\
* Infraestrutura disponivel, transporte e logistica;\\
\\
* Clima e geografia (tropical, artico, localizagao geografica).\\
\\
Na Figura 4.37, ilustramos o resumo de custos de atividades de perfuragao e completagao\\
de pocos de petréleo e gas natural, onde o maior investimento é para o custo das sondas\\
\\
ou plataformas de perfuracao e suas equipes.\\
\\
Figura 4,37: Resumo de custos de atividades de perfuragio € completagio de poses.\\
\\
Digitalizada com CamScanner\\

\\section*{Pagina 133}

112 PERFURACAO E COMPLETAGAO DE POCOS\\
\\
Principais custos na implantagio de um projeto de desenvolvinen od\\
Petréleo offshore sio ilustrados na Tabela 4.4, onde 6 demonstrag *\\
\\
10 Gg i Gn\\
a ot a\\
furacdo representam quase 50\\% de todas as despesas de desenvolvinenn te\\
‘ 0 de\\
de petréleo, de y,\\
\\
Tabela 4.4: Principais custos na implementagao de\\
\\
Segio Equipamentos\\
\\
Inve ‘Stimenty\\
\\
1, Casco de plataforma\\
Unidade de Exploracio e Produgio 2, Médulo\\
\\
35-45\\% do IVestimentg \\textasciitilde{}\\
Obs: no cas de\\
\\
afretany\\
3. Ancoragem este custo @ considered dy way\\
Pera\\
. 1, Linhas de coleta/ escoamento oral\\
\\
‘elses Saburo) 2. Bquipamentos submacinos (manifold, 15-25\\% dp investimeny\\
\\
PLEM, PLET etc.) 0\\
\\
1, Perfuragio € completagio.\\
Perfuracio e Completagio " ieee)\\
\\
2. Arvore de Natal Molhada (ANM) 40° 50\\% dovin\\
\\
* Plataformas e equipes precisam ser m\\
\\
antidas e pagas independentemente das atv:\\
dades Operacionais da empresa;\\
\\
wide - letehon Prestadoras de Servicos, as empresas participam de concurs0s\\
\\
shes ponte feltesbes Para atividades €specificas, Por exemplo, em atividades de petfu\\
\\
Satoh ae existem empresas Para fluidos de Perfuracio, cimentacio etc. Bntre 6\\
Part aady es de perfuracio e completacio estao:\\
\\
‘ ’ — ry\\
* Nesse tipo de contrato de perfuragio, a\\
tante (empreiteiro) baseada na perfurai\\
mig, “anit\\
Pendentemente as operacées de perfuracio. Gat\\
\\
‘ desafio nesse ti iso contratante pole\\
visar perfurar de forma répida e seangaaie de contrato, pois 0 co\\
H\\
\\
Contrat ri -\\
; pelo piss ita ado (Footage): esse contrato se baseia no ant\\
 medido em Pés. Ele inc ‘aa completagdo MPP"\\
\\
r centiva a come vat\\
\\
No entanto, carrega os MesMi04 riscos do contrato de obra completa. Geralmenl®\\
\\
Digitalizada com CamScanner\\

\\section*{Pagina 134}

CUSTOS E CONTRATOS DE PERFURACAO\\
E COMPLETACAg\\
DE Pocos 1\\
13\\
\\
éusado para a segdo acima do reservatério Prospectivo, onde as co;\\
\\
Fes, nak ndicé\\
<io menos criticas em termos de avaliacio ou producao; 'g0es do furo\\
\\
Contrato de Incentivo: esse contrato, frequentemente usado nas atividades de perf\\
cioe completagao, envolve bénus por desempenho acima da média. As onesies\\
\\
goes do pogo S40 acordadas entre a empresa operadora e 0 contratante, com base\\
\\
nos custos hist6ricos de pocos semelhantes. Isso permite estimar os custos para 3\\
\\
novo poco;\\
\\
Contrato de Taxa Diaria: a empresa operadora aluga a sonda ou plataforma e a equipe\\
em uma base diaria, geralmente com controle total da operacao de perfuracio. O\\
contratante recebe uma quantia especificada por cada dia no poco. Nao ha bénus\\
de desempenho, mas 0 contratante pode ser penalizado por negligéncia;\\
\\
Contrato de Parcerias e Aliancas: esse contrato envolve a definicao e fusao de objetivos\\
comuns de negocios, compartilhamento de riscos financeiros e recompensas. Visa\\
melhorar eficiéncia e reduzir custos. Equipes conjuntas sao criadas, e a empresa de\\
perfuracao pode formar aliancas com subcontratantes para ampliar o espectro de\\
\\
atividades.\\
\\
Conclusao\\
\\
Nesse capitulo, os leitores adquiriram conhecimentos para compreender operacées de\\
perfuracdo e completacao. Estao aptos a descrever equipamentos, explicar processos de\\
perfuraco, completacio, intervencao € abandono, bem como compreender riscos e efici-\\
éncias da industria petrolifera. Contratos de perfuracao foram apresentados, assim como\\
operacSes normais e especiais de perfuracao € completa¢ao.\\
\\
Digitalizada com CamScanner\\

\\section*{Pagina 135}

CAPITULO 5\\
\\
AVALIACAO DE FORMACGES\\
ee,\\
\\
Avaliacao de formacées sao estudos e atividades que visam avaliar qualitativamente e\\
quantitativamente o potencial de uma jazida petrolifera, ou seja, a sua capacidade pro-\\
dutiva ¢ a valoracdo das suas reservas de petrdleo e gas natural. Ou seja, a avaliacio de\\
formacao é usada para:\\
\\
« Avaliar reservat6rios de hidrocarbonetos e Prever a recuperacao de petréleo;\\
\\
= Fornecer aos engenheiros do reservatério os parametros geolégicos e fisicos da for-\\
macao necessarios para a construgao de um modelo de fluxo de fluido do reservaté-\\
Tio;\\
\\
* Medir a pressao do fluido de formacao in situ e adquirir amostras de fluido de for-\\
macao;\\
\\
* Naexploracao e desenvolvimento de petréleo, a avaliagao de formacao é usada para\\
determinar a capacidade de poco para produzir petréleo.\\
\\
As atividades referentes a avaliacio de formago sao o teste de formacao ea perfilagem.\\
Além do teste de formacio e a perfilagem, todas as informacées anteriores 4 perfilagem\\
do intervalo de interesse que podem ser obtidas através das etapas do estudo geolégico\\
€ geofisico ou na etapa de perfuracdo de poco fazem também parte das atividades de\\
avaliac4o de formacées. Uma avaliacao de formacao efetiva de reservatério exige uma\\
integrac4o efetiva de todos os dados disponiveis.\\
\\
5.1 Avaliacéo de macroestruturas\\
\\
Muitos dos Conceitos de avaliacdo da macroestrutura foram abordados nos capitulos 3\\
€4. Como mencionamos no prefacio deste livro, muitos dos conceitos serao repetidos\\
0s capitulos Para facilitar que o leitor estude o capitulo de seu interesse, ou seja, estudo\\
\\
A avaliacio de Macroestruturas antecede a etapa de perfuracao de pogos € a fase =\\
da formacio, visto que esta ultima somente ocorre apés a perfuracio do pogo.\\
115\\
\\
8 Cer rd Petrileo para ndo-engenheiros, 1* edigio.\\
André Raposo Ramos Copyright © 2023 Editora Dois03.\\
\\
Digitalizada com CamScanner\\

\\section*{Pagina 136}

116 AVALIACAO DE FORMACOES\\
\\
+6 ia de S\\
ras € compost4 pela Geologia de Superficie, cin\\
ade superficie.\\
\\
ic macroestrutu\\
\\
A avaliagdo di n J\\
a, Sismica € Geologi\\
\\
Magnetometn\\
\\
5.1.1 Geologia de superficie\\
\\
erficie realiza 0\\
\\
de rochas que afloram na superfi\\
delimitar as bacias sedimentares cid\\
\\
onetos. Os mapas geoldgicos, gue ent\\
ontinuamente construidos e atvalnde\\
rficas sao praticamente eliminadas os\\
a sedimentar muito reduzida oy he\\
\\
mapeamento\\
Cie. Ag,\\
\\
é possivel conhecer €\\
\\
de acumular hidrocarb\\
eressantes, sao Cc\\
\\
A geologia de sup\\
vés desse mapeamento,\\
ficar estruturas capazes\\
\\
cam as areas potenci® mente int\\
As areas CO! as igneas ¢ metamo’\\
compostas pot rochas ign\\
\\
mapas, assim como pequenas bacias com espessur\\
estruturas favordveis aacumula¢ao.\\
\\
5.1.2 Gravimetria\\
alores da intensidade do campo magnético na Tera,\\
etémetros, onde a ordem de grandeza do valor do\\
de 50.000 gamas (unidade de medida da intensidade\\
do campo magnético). As anomalias produzidas pelas estruturas geoldgicas de interesse\\
tém ordem de grandeza de cerca de 10 gamas. Os magnetémetros também devem ser\\
bastante precisos. Enquanto 0 valor da gravidade esta associado a densidade das rochas\\
de superficie, o valor da intensidade do campo magnético esté associado a suscetibilidade\\
magnética das rochas.\\
Como i i\\
= pena hae Wee apresentam valores baixos de suscetibilidade magnética,\\
. las magnéticas esto relacionados a presenga de rochas intrusivas ou\\
embasamento, Através do mapa magnético obtid 6 6 i\\
\\& eee iofoencan Bee en obtido apés as corregoes devido a outros\\
deed . spas lo campo magnético, é possivel estimar a profundidade\\
embasam lo pacote sedimentar e fazer estimativas quanto 4 estrutue\\
\\
geol6gica da area.\\
\\
A gravimetria utiliza a medicao dos v:\\
Os aparelhos utilizados sio os magn\\
campo magnético terrestre € de cerca\\
\\
5.1.3 Magnetometria\\
\\
S Pee Betuenes variag6es na intensidade do campo\\
vem ser utilizadas idas de ca ae pate de rochas na subsuperficie. ‘i\\
pode fornecer esti €M conjunto com pieced lathe podem ser ambiguse :\\
dos sedi estimativas da profundida i todos, A andlise cuidadosa desses map :\\
cas. A oman tOS estruturais do pr ei aetpahi magnético ou da Cee\\
Variages do Sadenen Mais complexa Py € presenga de rochas intrusive\\
5‘ rae també, lo que no método gravimétrico, devi\\
™ porque diferentes instruments!\\
\\
eg vind métod ae No método gravimétrico 0 velo"\\
ue ndo ime © vetor é horizontal no ¢4\\
£m com os polos geogrificos:\\
\\
apon?\\
ul dot ¢\\
\\
Digitalizada com CamScanner\\

\\section*{Pagina 137}

| — |\\
\\
AVALIACAO DE MACROESTRUTURAS 117\\
\\
Ssismica\\
\\
5\\
\\
jsismica tem como objetivo conhecer informagées da subsuperficie e é um dos métodos\\
mais empregados na industria do petréleo, Este método envolve a analise da Propagacio\\
de ondas sismicas através das rochas da subsuperficie que se deseja conhecer. Embora as\\
condas sismicas possam ter origem natural, como nos terremotos, para fins de exploracio\\
clas sio geradas artificialmente por meio de explosdes na superficie. Em levantamentos\\
sismicos terrestres, as explosdes sio provocadas por dinamites, enquanto em levantamen-\\
tos maritimos sao geradas por canhées de ar comprimido. Cada fonte gera ondas sismicas\\
caracteristicas, formando um pulso conhecido, chamado de assinatura da fonte.\\
\\
Apés a emissio da onda, ela se propaga para o interior da Terra e retorna Para a super-\\
ficie através do fendmeno da reflexao, A perturbacdo gerada na superficie pelo retorno\\
da onda é captada por receptores, que convertem essa perturbagao (energia cinética ou\\
de pressio) em sinais elétricos para posterior processamento adequado. Em levantamen-\\
tos terrestres, esses receptores sio chamados de geofones (eletromagnéticos), enquanto\\
em levantamentos maritimos sao chamados de hidrofones (piezoelétricos). As etapas da\\
investigacao do subsolo pelo método sismico ocorrem em trés etapas: aquisicado de dados\\
sismicos, processamento dos dados adquiridos e interpretagao.\\
\\
5.1.4.1 Aquisigao de dados sismicos\\
\\
O processo de aquisi¢éio de dados sismicos é semelhante em terra ¢ no mar, variando\\
0s equipamentos utilizados. Como 0 método é 0 mesmo, ou seja, envolve a geracio\\
de uma perturbaco na superficie através de fontes de energia e a recepcao das reflexdes\\
dessa onda através de receptores, outra mudanga ser a forma de dispor esses equipamen-\\
tos. Em terra, os geofones sao enterrados no chfio, enquanto no mar os hidrofones sio\\
dispostos em cabos sismicos ¢ puxados por uma embarcacao. Os receptores sao afasta-\\
dos equidistantemente (20 a 50 metros), € 0 comprimento do cabo sismico pode variar de\\
alguns a varios quilémetros (de 5 a 15 quilémetros).\\
\\
Apés 0 disparo, a onda sismica se Propaga no subsolo. Ao encontrar uma interface\\
entre dois tipos de rocha, parte da onda sofre refracfio e continua se propagando para\\
0 interior do subsolo, enquanto outra parte da onda sofre reflexao e retorna para a su-\\
perficie. As partes da onda que retornam A superficie sio captadas pelos receptores, que\\
Tegistram o tempo de chegada da onda ea quantidade de energia retornada. A primeira\\
figura mostra um exemplo de aquisicfo no mar e em terra. Apés o disparo e€ 0 registro\\
dos dados, tanto a fonte quanto os receptores sio movimentados para a frente para que\\
um novo disparo e registro sejam realizados. Apés varios disparos, toda uma linha na area\\
4 Ser estudada é coberta, e o mesmo processo é repetido para outras linhas, até que toda\\
4 drea seja abrangida, Analisando esse processo de aquisicio, podemos perceber que um\\
Ponto no subsolo é registrado varias vezes por receptores diferentes 4 medida que a fonte\\
\\
8¢ desloca,\\
\\
5.1.4.2 Processamento\\
\\
Na etapa de Processamento dos dados sismicos, busca-se eliminar alguns erros inerentes\\
49 processo ¢ agrupar os dados dos receptores, de forma a obter, para cada ponto na\\
\\
Digitalizada com CamScanner\\

\\section*{Pagina 138}

AVALIACAO DE FORMAGOES\\
\\
118\\
a de acordo com o temp\\
0 €\\
\\
‘ a a recebid\\
superficie, amplitude da seit oe sismico. A juncio de v4 qu\\
esta foi percebida. A este registro da-se ono! ) eo conjunto ieee le virion\\
a sismica resulta na linha sismica (figura 5 mecteaie wat linhas Sismj :\\
e sismico. Para a obtencao deste vo ume sismico € necessaria a eategt\\
rita posteriorment®- - . ; Ce\\
\\
nsistia em uma unica linha de re <eptores\\
Co.\\
\\
ramento sfsmico CO’\\
micos. Esse tipo de levantamento é chamado de sismica an\\
sigao dos dados so os mesmos, mas, em Seas Na\\
temos cinco OU seis, o que faz com que um ie ter.\\
mo disparo em um cabo e também seja iain\\
os 3D sao mais caros que os tevaiacags\\
Ja quantidade € qualidade dos equipamentos utilizados quanto ‘0s\\
exo dos dados obtidos. No entanto, os resultados sao deh lo\\
ntamentos sismicos 4D envolvem uma quarta aifiéiesy\\
tempo, ou seja, consistem em fazer levantamentos 3D em intervalos regulares de tig °\\
(de 6 a 12 meses). Os levantamentos sismicos 4D, por serem cada vez mais Gilly\\
com o tempo, podem fornecer informac6es sobre a movimenta¢ao do éleo produzi .\\
\\
em uma Area.\\
\\
um registro da\\
\\
na linh\\
forma o volum'\\
de sismica 3D, desc\\
\\
Até agora, 0 levan\\
nectados por cabos sisi 1\\
sismica 3D, os procedimentos de aqui\\
mos apenas uma linha de receptores,\\
\\
o seja varrido varias vezes pelo mes\\
\\
pont \\textbackslash{}\\
5 levantamentos sismic\\
\\
por outros cabos. O\\}\\
sismicos 2D, tanto pe!\\
processamento mais comp!\\
tivamente melhores. Os leva\\
\\
5.1.4.3 Interpretagao\\
\\
Ai A —\\_\\
\\
rie ee dos dados sismicos processados busca identificar indicios que possibili\\
\\
sm estimar ta! Ogi a a\\
nto a estrutura geolégica da Area quanto a ocorréncia de acumulacdes\\
\\
hidrocarbonetos.\\
\\
5.2 Técnicas de avaliagao de formagées\\
\\
As técnicas de avaliaca\\
iagao de Fa F\\
pitulo, sio seecinee a ae considerando a informacio introdutéria deste\\
fo: 5 em acom, Boe "\\
rma¢4o, que passamos a descrever a a eee i geoldégico, perfilagem e teste de\\
\\
5.2.1 A\\
icompanhamento egolégico de poco:\\
Ss\\
\\
Durante a operaci\\
a¢ao de perfuraca\\
necidas aa ‘acao de\\
de aeccenthae e formacoes. fae mas informacées importantes S40 for-\\
tempo 640, como a inclinacx ‘maces incluem 6 Alises\\
de set poy epee fodeck ee do pogo, dados paler\\
abordaremos os tré: 7 testemunhagem e Pee Além disso, sio  calisadas desci@t\\
lacionadas aos cb ii Pardmetros, beta de gas. Para o propésito deste livro.\\
etros de perfuracao, uanto o lei\\
‘a¢4o em fon itor pode buscar informagoes ©\\
tes adicionais\\
\\
5.2.\\
2.1.1 Amostras de calha\\
\\
Calha, ou cascalh\\
pela Jama de 08, Sao fragmento;\\
perfura 8 de rocha tri\\
co, sendo capturados enbitrados pela broca ¢ trazidos 3 uperfice\\
™aPencira. A chegada das amostrasas”\\
\\
Digitalizada com CamScanner\\

\\section*{Pagina 139}

Tl\\
ECNICAS DE AVALIACAO DE FORMACOES 119\\
\\
sficie depende do tempo de retorno em Telagdo a\\
de retorno é 0 tem on\\
da coleta. O tempo PO necessario Para que a lama di.\\
Ps ‘ 4 le\\
\\
qirculagio alcance a superficie. As amostras sao coletadas em intervalos ne em\\
mente a cada 3x3 metros. Ao descrever as amostras de calha, as Porce ti ares, geral-\\
fragmento de rocha sao estimadas. Os fragmento: ntagens de cada\\
utilizando uma lupa binocular.\\
\\
Indicios em amostras de calha\\
\\
A presenca de petrdleo mével ou residual pode ser revelada através de fragmentos de\\
rocha nas amostras de calha. Essa presenca pode se manifestar na forma de manchas,\\
fluorescéncia e corte. °\\
\\
As manchas de dleo nos fragmentos de rocha geralmente esto associadas ao 6leo re-\\
sidual, enquanto a fluorescéncia indica a presenca de Petréleo ou gas nas amostras. E\\
importante ressaltar que a fluorescéncia é a luminosidade emitida pelo fragmento de ro-\\
cha quando exposto 4 luz ultravioleta. A fluorescéncia pode ser classificada como total,\\
esparsa ou pontual. Ela é considerada total quando todos os grdos da amostra apresentam\\
fluorescéncia, esparsa quando 10 a 90\\% dos gr4os a apresentam, e pontual quando menos\\
de 10\\% dos gros apresentam fluorescéncia. No entanto, nem sempre a fluorescéncia\\
indica a presenca de dleo. Alguns minerais, como a calcita, também sio fluorescentes,\\
fenémeno conhecido como fluorescéncia mineral.\\
\\
O teste de corte ocorre quando, sob a luz ultravioleta, é aplicado um solvente (triclo-\\
roetano) na amostra, € 0 liquido se torna fluorescente. Esse teste pode ser classificado\\
como imediato, moderado ou provocado. O corte imediato sugere boa permeabilidade,\\
enquanto o corte provocado indica baixa permeabilidade ou presenca de dleo residual.\\
\\
Vale notar que a andlise de fluorescéncia e corte sé é possivel quando a lama utilizada\\
€ 4 base de agua. Em lama a base de 6leo, a presenca natural de petréleo na rocha fica\\
mascarada.\\
\\
5.2.1.2 Testemunhagem e amostragem de poco\\
\\
Durante as operaces de perfuracio de pocos exploratorios, sio colhidas amostras das\\
formagées atravessadas. Em certas situac6es, é necess4ria a obten¢io de uma amostra\\
integral da rocha perfurada. A selecdo das camadas onde as amostras serao coletadas ¢\\
tesponsabilidade do gedlogo.\\
\\
As amostras podem ser pequenos fragmentos retirados das paredes do pogo, conhe-\\
cida como amostragem lateral, ou um cilindro de alguns metros obtido na perfuracao do\\
trecho de poco, chamado testemunhagem. A testemunhagem é 0 método mais frequen-\\
temente utilizado. Normalmente, os testemunhos tém o comprimento de um tubo de\\
Perfuracio (ou multiplos): 9m, 18m, 27m.\\
\\
tascceecee sin thei m: éutilizado um equipamento chamado barrilete. os\\
testemunhos sio coletados, descritos na sonda, embalados e enviados ao bearer i 4\\
Beologia, onde so serrados e detalhadamente descritos. Além da descrigio, sio arn\\
informacdes sobre o tipo de rocha, porosidade, permeabilidade e outras proprieds ss\\
Petrofisicas. Também podem ser coletados microfésseis para datagdo e amostras para 8\\
\\
Digitalizada com CamScanner\\

\\section*{Pagina 140}

120 AVALIACAO DE FORMACOES\\
\\
otencial gerador, tipo © maturagao da Matérig\\
Pp dem ser usadas como padrao de caliby,\\
0, selegio de fluidos de completacy,\\
\\
ao do\\
ara a obten¢@\\
munhagem PO\\
testes de forma\\
\\
oquimica, utilizadas p:\\
As operagoes de teste!\\
\\
organica. io ds perfis,\\
\\
ao para a interpreta\\
entre outros.\\
\\
5.2.1.3 detetor de gas\\
nhagem, utiliza-se um detetor de gas para ands ;\\
Durante a operacio de testemu) i calha, a lemma Seal ati pees dear\\
\\
amostras : A F\\
amostras. Juntamente com as ido proximo 4 pencira de lama, 10 gu\\
\\
4s é recoll\\
' erfurada. O gas\\
provenienns Ed oe alsa no detetor de gases. Alguns detetores podem ser bas,\\
trap, e continuame\\
\\
a As\\
icados, realizando a cromatografia do gas. A presene ri gas pesado sugerea\\
tante ane ee eservardrio, enquanto 2 ocorréncia apenas de gases leves sugere g\\
pecans 86 q eee\\
page le endo associado. Geralmente, ‘os indicios observados nas amostras de calha\\
presenga de g iS\\
\\
i de gas.\\
estiio associados aos indicios detetados pelo detetor de g'\\
\\
5.2.2 Perfilagem de poco\\
\\
A operagio de perfilagem é o registro continuo ou discreto, em escala, de parametros\\
fisicos (resistivos, actisticos e radioativos) ou quimicos. Devido ao seu custo relativamente\\
baixo em comparacao com outras operacSes mais dispendiosas e 4 obtencao ao longo\\
de todo 0 poco, essa opera¢ao pode ser considerada a principal ferramenta da geologia\\
de reservatérios. Essa técnica foi desenvolvida a partir da década de 1920 pelos irmios\\
Schlumberger. A perfilagem de poco, juntamente com o teste de formagao, é um dos\\
dois pilares basicos da avaliacao de formagées.\\
\\
Otermo “perfil” refere-se a qualquer dado que seja apresentado em forma grafica em\\
funcao da profundidade do pogo ou do tempo, acompanhado por algumas notas escritas.\\
Isso engloba registros feitos em pocos abertos e revestidos, registros de sondador, regis\\
ban Measure While Drilling/Medicao Durante a Perfuracao (MWD) e outros.\\
\\
\\} OS principais objetivos da operacao de perfilagem sao:\\
\\
= No curto prazo: i \\} Ay.\\
Prazo: obter informacées necessArias para a tomada de decisao sobre 4\\
\\
conclusa inci\\
: eng ideas do poco. As principais informacées incluem profundidade\\
espessura dos reservatérios, porosidade e saturacio de fluidos;\\
\\
de reservas e modelagem ge\\
\\
de fluidos e inclinacao das camadas,\\
\\
Além dos objetivos mencionad i\\
nat 0 calibre do poco, a variseg, son \\# OPEraGaO de perfilagem é usada para deter\\
\\
: a variacio\\
geolégica das formacées fies da temperatura com a profundidade, a identificas*\\
XIStO, arenito, calcdrio ete.), 0 mergulho das camadis\\
\\
a qualidade da cimentaci:\\
640 (aderénci;\\
mento) e para indicar danos na Fala Pasta de cimento 4 parede externa do reves"\\
\\
Digitalizada com CamScanner\\

\\section*{Pagina 141}

TECNICAS DE AVALIAi\\
GAO DE FORMA:\\
GOES, 121\\
\\
22d Tipos de perfilagem\\
rfilagem pode ser classificada quanto ao estado do\\
\\
Po¢o ou quanto g\\
Quanto ao estado do pogo, a perfilagem poi —\\
\\
de ser feita em:\\
\\
Ape\\
perfilage™-\\
\\
poco aberto: antes de descer o revestimento. Essa é a mais\\
5 ; ; 5\\
6 maces, € seus resultados podem influenciar a decisao d\\
\\
momento da\\
\\
‘omum na avaliacdo de for-\\
\\
¢ revestir ou nio o Pogo;\\
\\
oats revestido: no caso de pocos j4 em producio, Para verificar as mudangas nas pro-\\
priedades da formacio, principalmente a saturacao de agua, coma produgao. Nem\\
todos os perfis podem ser realizados em Po¢os revestidos,\\
\\
Quanto ao momento da perfilagem, ela pode ser feita:\\
\\
Durante a perfuracio: quando é necessario obter informacgées em tempo real enquanto\\
perfura (Logging While Drilling/ Perfilagem Continua (LWD)). Essa operacao é mais\\
cara e é geralmente realizada em pocos horizontais e offshore;\\
\\
Ap6s a perfuracéo: € menos dispendiosa e mais comum na avaliagao de formacgoes.\\
Perfilagem de produ¢éo\\
\\
A perfilagem de producao é realizada por meio de perfis feitos apds a descida do reves-\\
timento de producao e a completacio inicial do Poco, com o objetivo de determinar a\\
efetividade da completagéo ou as condicdes de produtividade (ou injetividade) de um\\
poco. Algumas ferramentas usadas na perfilagem de producio sio:\\
\\
* Production Logging Tool/ Ferramenta da Base Adaptadora de Produgio (PLT): essa\\
ferramenta de perfilagem possui perfis como continous flowmeter, gradiomanémetro,\\
perfil de densidade, hidrolog e perfil de temperatura.\\
\\
\\textasciitilde{} Medidor de vazao continuo: mede a velocidade do fluxo nos trechos de inte-\\
resse do poco aberto, determinando a contribuicdo de cada trecho na vazio total\\
de producao do pogo;\\
\\
\\textasciitilde{} Gradiomanémetro: registra a densidade média dos fluidos contidos em um tre-\\
cho do pogo em funcio da profundidade, através da medicao da pressao em dois\\
Pontos distantes;\\
\\
\\textasciitilde{} Perfil de densidade: registra a densidade do fluido utilizando uma fonte de raios\\
gama e um detetor que mede os raios gama dispersos no fluido apés colidirem\\
com os elétrons;\\
\\
\\textasciitilde{} Hidrolog: usado quando hé fluxo trifasico (gas, dleo e 4gua); mede a constante\\
dielétrica e determina a porcentagem dessa mistura;\\
\\
\\textasciitilde{} Perfil de temperatura: registra a temperatura do fluido do pogo, indicando os\\
trechos em producio, vazamentos e outros elementos.\\
\\
DT é usado para\\
\\
. Decay Time/ Tempo de Decaimento Térmico (TDT): 0 TE emia\\
Um perfil qualitativo da saturacio de fluidos no reservatorio,\\
\\
Digitalizada com CamScanner\\

\\section*{Pagina 142}

=\\_"”""—\\textasciitilde{}—\\textasciitilde{}S—\\
\\
122 —avauiacAo DE FORMAGOES\\
\\
do, que emitem raios gama com intensidades dif\\
da rocha em 6leo, gas ou agua. Em outras Palavr\\
\\
6leo-4gua-\\
\\
s contra a forma:\\
da saturagao\\
ntatos bleo-gas ¢\\
\\
emite néutron:\\
rentes dependendo\\
ele determina os CO!\\
\\
5.2.2.2 Tipos de pertis\\
anos de perfis que podem ser classificados como mecanicos, elético,\\
Existem pape ed sjerromagnéticos. Vou descrever alguns dos tipos de perfis maig\\
radioativos, acustico:\\
comuns. 7 t Formation Test/ Teste de Formac\\
anicos i caliper e 0 Repea ; aco\\
Os perfis mecanicos incluem o calip DS Te ee\\
\\
elétricos sao 0 SP (Pote: v\\
aa incluem Raios Gama, densidade e neutrao. Os perfis\\
\\
enquanto o perfil eletromagnético ¢a\\
\\
Repetido (RFT). Os pe!\\
gem elétrica. Os perfis ra\\
acusticos so 0 perfil s6nico\\
ressonancia magnética.\\
\\
£ importante notar que cada emy|\\
curva gerada nos perfis. Por exemplo,\\
\\
ea imagem acustica,\\
\\
presa pode ter sua propria nomenclatura para cada\\
densidade é chamada de RHOB na Schlumbergere\\
RHOZ na Halliburton. Para evitar confusées, ja que as empresas de petréleo muitas vezes\\
contratam diferentes empresas de servicos, sao utilizadas nomenclaturas padronizadas\\
nos bancos de dados. Por exemplo, a curva de densidade é identificada como RHOB, a\\
curva de neutrao como PHIN e a curva de resistividade como ILD.\\
A resolucao dos perfis varia dependendo do tipo de ferramenta e perfil. Geralmente,\\
a resoluco esta em torno de 1 metro. Ou seja, camadas com espessura menor do que\\
1 metro no sao detetadas. Existem alguns perfis especiais, como a microrresistividade\\
¢ perfis de imagem, que possuem resolugdo centimétrica e conseguem detetar camadas\\
muito finas € estruturas como fraturas e estratificagdes. Apesar das limitaces, a resolu-\\
ee oa émuito maior do que a da sismica (que é de 10-50 metros), embora\\
: loners a eer pode ™mapear entre os pocos). Na Figura 5.1, ilus:\\
resolucdo (preto) ¢ um perfil rit : eeelae da s\\{smica (tracejado), um perfil de baixa\\
ta resolucdo (vermelho). A camada de arenito com 30\\
\\
cm de espessura é detetada so:\\
mente pelo perfil de alta resolucio.\\
\\
quali licagée: eae,\\
itativas incluem: PicagOes tanto qualitativas quanto quantitativas. As aplicagoes\\
\\
* Ident 0 litolégica (ti\\
tificagao litolégica (tipo de Tocha) a partir dos valores das curvas;\\
. 1-40 do amb; ici rv:\\
identifica ambiente deposicional a Partir da forma das curvas;\\
bd entre com) )\\
Correlacio PoG0s comparando a Similaridade entre as curvas;\\
\\
* Identificacio dos fluid\\
OS que\\
* Controle da profund; que saturam a rocha com base nos valores das curva\\
Profundidade ¢ do diametro do Pi\\
O¢o,\\
\\
Aplicacdes quantitativas dos\\
\\
: cdlculo Pertfis 33\\
oa da saturacio dos fluidos, pate an Obtencao da porosidade e per™ abil\\
sp + Alguns dos reservatérios; densidade das roc\\
\\
imativa de reseryas, No hetantn\\
\\
Digitalizada com CamScanner\\

\\section*{Pagina 143}

TECNICAS DE AVALIACAO DE FORMACOES, 123\\
\\
\\textbackslash{} tonne\\
\\textbackslash{}\\
\\
SES, fil de bal\\
\\textasciitilde{}-..\\_\\_\\_\\_ Perfil de balka resolusto be rtil de atta resolucdo\\
\\
trago sismics\\textasciitilde{}\\textasciitilde{}\\textasciitilde{}------\\textasciitilde{}\\
\\
Figura 5.1: Compara¢ao entre a baixa resolugao da sismica e um perfil de baixa e alta\\
resolugao.\\
\\
solugao requer a combinagio de varios perfis. Na Tabela 5.1, fornecemos um resumo das\\
aplicagées dos principais perfis abordados.\\
\\
Tabela 5.1: Resumo da aplicacao dos principais perfis.\\
\\
Perfil Medigio Unidade Valor alto Valor baixo\\
RaiosGama(GR) —\\_Radioatividade naru- API Xistos, argilas, saisde Arenito, calcirio, ha-\\
ral - identificagio do potiissio Tita\\
tipo de rocha\\
Resistividade (ILD) Resistividade 4 cor- Ohm.m Oleo, Gés, sal, areni- Agua salgada\\
rente elétrica - identi- tos cimentados, 4gua\\
ficagio de fluido doce\\
Densidade (RHOB) Densidade da rocha g/cm? Porosidade baixa Porosidade alta\\
\\
baseada no contetido\\
de matéria - determi-\\
\\
nagio da porosidade\\
Neutrio (NPHI) Densidade da rocha Porosidade (\\%) Porosidade alta Porosidade \\_baixa,\\
baseada no conteudo presenga de gis\\
de hidrogénio - poro-\\
sidade e presenga de\\
Lad :\\
Porosidade baixa Porosidade alta\\
\\
Sénico (DT) Velocidade de propa- jis/ft\\
gacdo de onda acis-\\
tica - porosidade ¢ ca\\
\\
librago com sismica\\
\\
5.2.3 Teste de pressao em pogos\\
\\
Presséo uma is i a dugio de\\
Yi iedades mais importantes na exploracao e produgac\\
pap ine ; na produtividade dos reservatérios e\\
\\
9 € gis natural, com implicag6es nas reservas,\\
\\
Digitalizada com CamScanner\\

\\section*{Pagina 144}

124 AVALIACAO DE FORMAGOES\\
\\
na seguran¢a operacional, é importante revisarMos 0 seu Conceito, assim co,\\
de gradiente.\\
\\
m\\
9 “Once,\\
\\
5.2.3.1 Press\\&o e gradiente de pressao\\
\\
A pressao é definida como a forca exercida por unidade de area. Este conce;\\
\\
; ea ito \\& do\\
minio do leitor nas primeiras lic6es de Fisica. A pressao € expressa como iustramos\\
Equacio 5.1. Suas unidades de medi¢ao sio o Pascal (Pa), a atmosfera (atm), o keflem?\\
ea libra por polegada quadrada (psi). :\\
\\
iE\\
we A (5.1)\\
\\
O gradiente de pressao é definido como a taxa de aumento de\\
devido ao peso exercido pela coluna de rocha ou fluido. I\\
ilustrado na Equagio 5.2.\\
\\
presséo coma Profundidade\\
sso pode ser XPresso como\\
\\
dP\\
O gradiente de pressio depende fundamentalmente da densidade (P) média da coluna de\\
\\
Tocha ou fluido, da Profundidade (h) e da acelerac3o da gravidade ( © 10 m/s?), Isso\\
pode ser expresso na Equacio 5.3:\\
\\
A=\\
\\
P=pgh (53)\\
pa seig as) a densidade do fluido, € possivel obter o gradiente de pressio e, com isso,\\
estim: Pressdo em uma determinada profundidade. Na Tabela 5.2, ilustramos as den-\\
Sidades e gradientes de Pressao tipicos de alguns fluidos,\\
\\
Tabela 5.2: Exemplos de\\
\\
gradientes e massas especificas tipicas de fluidos.\\
Fluido Densidade A kgflem? my\\
\\
(g/m) A@si/m) 4 (lb/gal) Obs\\
\\
Aguadoce 1,9 01 ta\\
\\
Agua salgada 1,01 -1,2 0,101 - : cH\\
\\
Cleo 97-095 gor. Ha 3 PU MRCSFOS. oe eRe\\
\\
Gs 03-05 0.03 - 0,05 0, \\% 135 58-79 Depende da densidade (AP!)\\
S WR O2 bs Bs ctag Depende da composigio e pres-\\
\\
Lama de per tants ay, — sio (fungio da profundidade)\\
\\
1,56 -\\
213 92-425 Depende da composi¢io\\
\\
Digitalizada com CamScanner\\

\\section*{Pagina 145}

TECNICAS DE AVALIACAO DE FORMACOES 125\\
\\
5.2.32 Tipos de Pressao\\
o, vamos abordar os tipos de pressio, ou geopressdes,\\
\\
esta 9064\\
e roliferas (Figura 5.2);\\
\\
vidades pet\\
\\
Predominantes nas ati-\\
\\
pressio litostitica: é a pressio exercida pelo peso das camadas de rocha. Também co-\\
nhecida como pressio de sobrecarga ou overburden;\\
\\
Pressio estatica: éa pressio exercida pela coluna de fluidos de formacio, da superficie\\
até o reservatorio. £ a mais importante para os estudos de reservatério. Também\\
conhecida como pressao de poros, pressio do reservatério ou da formagio. Na se-\\
cao 5.2.3.2, vamos descrever as variages da pressiio estatica;\\
\\
Pressao estatica original: é a pressio do reservatério virgem, antes de entrar em produ-\\
ao;\\
\\
Pressao estatica atual: é a pressio do reservatério em uma determinada data, apés a en-\\
trada em produgao. A diferenga entre a pressao estatica original e a atual é chamada\\
de deplegao;\\
\\
Presso hidrostatica: é a pressdo exercida no anular do poco pela altura da lama de per-\\
furacao (ou fluido de completagao). Para manter a estabilidade e seguranca do poco,\\
deve ser maior do que a pressio estatica da formacio. Ao mesmo tempo, nao pode\\
ser excessiva para nao danificar o reservatério;\\
\\
Pressao de fratura: é a pressdo acima da qual a rocha reservatorio se fratura. Seu conhe-\\
cimento é importante para planejar operagGes de estimulacio e para evitar a perda\\
de circulacao durante a perfuracdo de pocos. Normalmente, a pressao de fratura é\\
maior do que a pressao estatica e menor do que a litostatica.\\
\\
* Pressdo\\
\\
Gradiente de\\
Pressio litostatica\\
\\
Gradiente de\\
Gradiente de Press4o hidrostatica\\
Pressao Estatica\\
\\
Figura 5.2: Grafico de gradientes de press\\
\\
\\_A\\
\\
Digitalizada com CamScanner\\

\\section*{Pagina 146}

‘AVALIAGAO DE FORMAGOES\\
\\
50 estatica\\
\\
126\\
\\
Variagdes na pres:\\
\\
de variar em difere\\
pre\\
\\
ntes cenarios, sendo classificada coma\\
Nor,\\
\\
5 Se 01 A i\\
\\
A pressao estatica P' ssurizado, conforme descrito a seguir:\\
\\
rio que segue o gradiente normal q,\\
© Cres.\\
\\
anormal, sobrecargas deplecao 04 sub]\\
\\
go do reservat6\\
\\
press\\
entar;\\
\\
Pressio Normal: é a 3\\
da bacia sedim'\\
\\
mento de press40\\
lo reservatorio acima da obtida através q\\
lo\\
ae\\
\\
rios limitados, em reservatérios com\\
ry if\\
\\
dos e soterrados muito rapidamente : ‘ande\\
\\
™ sug\\
\\
a: éa pressao de\\
\\
Pressao anormalmente alt\\
em reservato\\
\\
diente normal. Ocorre 2\\
coluna de gas e em sedimentos depositat\\
histéria geoldgica;\\
quando 0 pogo foi perfurado com lama excessivam,\\
lente\\
\\
da pressao hidrostatica para o reservatério\\
Nas\\
\\
Sobrecarga: em alguns casos,\\
feréncia\\
‘o alta é restrita 4s proximidades do pog\\
Oe\\
\\
densa, pode ocorrer trans\\
imediagées do pogo. Nesse caso, a pressai\\
pode ser confundida com presséo anormalmente alta;\\
\\
Pressdo anormalmente baixa: € 2 pressao do rese:\\
gradiente normal. Ocorre em reservatorios que\\
\\
Deplecao: éa perda de pressio devido a producao de fluidos;\\
\\
rvatorio abaixo da obtida através q\\
sofreram perda de fluidos; °\\
\\
Reservatorio depletado: € um reservatério que teve sua pressdo reduzida devido a\\
a pro-\\
\\
dugo de fluidos;\\
Reservatério sobre i\\
pressurizado: é um reservatori\\
vatociG 200) : -Orio que teve sua a\\
vido a injecao de fluido maior do que a produco. ae\\
\\
5.2.3.3 Teste de pressdo\\
\\
O teste de pressai\\
do envolve a operaco de bai\\
arame para a obtenca a0 de baixar registrado Et\\
cao de 5 res de pressa\\
teste de producio, pois pt aoe de fluxo ou estatica de fundo do en aes ado\\
inclui operacées de gradi Ive medicdes de vazao na superfici jae tien deus\\
Ma Wes ente estitico ou dinamico, e, perficie. Este tipo de teste\\
setae , uma ferram: » entre outros.\\
varios pontos do poco durante =, de perfilagem realiza um “mini-teste” (pré\\
de pigs a ysl mesma operacao de perfil este” (pré-teste) em\\
disso, possibilita a Pontos do reservatério em u: ee rap ro berate 0 mek\\
amos ma uni i\\
tragem do fluido, embora cee descida da ferramenta. Além\\
penas uma amostra por operaca0-\\
\\
5.2.4 Teste de formacado\\
\\
Os testes de form\\
F ‘agdo\\
\\
de identificar os Saiich Wake completacio ;\\
\\
tividade e avaliar a extensio da ried cnt Provisé:\\
\\
Tia do bjetivo\\
\\
10s de pogo com o obje\\
\\
teste Seki relacionados 4 produ:\\
lita a obtencao de propriedades\\
\\
Digitalizada com CamScanner\\

\\section*{Pagina 147}

TECNI\\
CAS DE AVALIACAO DE FORMACOES 127\\
\\
<titica, permeabilidade, transmissibilidade, danoe skin /.c, \\} bressdo\\
erie deteccao de descontinuidades Cathe ae ed Med ae\\
avaliagao da extensao do reservatério e volume de hidrocarbon oe eee\\
\\
Oteste de formacao envolve o isolamento do intervaloaser cad\\
\\
mais obturadores (packers). Em seguida, é estabelecido um\\
a formagao € 0 interior do pogo, forcando os fluidos da fo\\
Através da valvula de fundo, periodos de fluxo e estatica\\
presses de fundo sao registradas continuamente em fung:\\
‘A anilise dos dados coletados durante 0 teste de pressio\\
produtivo da formacao testada.\\
\\
testado por meio de um ou\\
diferencial de Pressao entre\\
rma¢ao a serem produzidos.\\
sao alternados, enquanto as\\
40 do tempo durante o teste.\\
Possibilita avaliar o potencial\\
\\
5.2.4.1 Principais tipos de teste de formacao\\
\\
Existem diversos tipos de teste de formacdo, cada um com suas caracteristicas e objetivos\\
especificos. Alguns dos principais tipos sao:\\
\\
» Repeat Formation Test/ Teste de Formagao Repetido (RFT): este teste é realizado\\
durante a operacdo de perfuracao, ao término da ultima fase, antes da descida do\\
revestimento. Uma ferramenta com uma pequena 4rea de investigacao é utilizada\\
para coletar amostras de fluido e registrar dados de pressao estatica da formacio em\\
diferentes profundidades. Isso permite determinar o gradiente de pressio do poco\\
e estimar os contatos entre os fluidos;\\
\\
* Teste de Formacao a Poco Aberto (TE): realizado durante a perfuracao antes da\\
descida do revestimento de producio, ou durante a completacao de um pogo aberto\\
a formacao. Similar ao Teste de Formacao em Poco Revestido (TFR), mas nao é\\
recomendado em zonas de perda de circulacao;\\
\\
* Teste de Formacao em Poco Revestido (TFR): realizado em um poco revestido\\
nsiste em um Periodo de Fluxo e um Periodo\\
\\
conectado a formacio para fluxo. Co Fai tte.\\
\\
de Estatica, onde o reservatorio alimenta e pressuriza 0S fluidos no\\
\\
Dependendo do objetivo do teste de formacao, ele pode ser classificado como:\\
\\
pressdo do reservatério\\
\\
é fechado, a\\
ae da de pressio resultante\\
\\
* Teste de Fluxo (Drawdown): nesse teste,\\
medida e depois é aberta para uma vaza0 constante. Aaa i\\
da descompressao dos fluidos do reservatério nas proximi sil\\
\\
Esse teste é titil para fluxos de longa duracao ¢ permite in\\
\\
+ aati me;\\
Servatério, fornecendo uma estimativa do seu volu\\
\\
limites do re-\\
\\
opogoé colocado em fluxo\\
\\
teste, -\\
: atorio. A pressio\\
\\
* Teste de Crescimento de Pressao (Buildup): nes ress do reser\\
\\
© depois fechado, medindo-se a recuperagao da\\
\\
Digitalizada com CamScanner\\
\\
les do pogo é medida.\\

\\section*{Pagina 148}

128 = AvaLiacAo De FORMAGOES\\
\\
1 yr médio aps um certo tempo de fechamenty, F\\
vale\\
\\
a aremum > te durante a medica we\\
ende a se estabilizar ev azao constante @ \\$40 da proces\\
tende ase manter a vazdo Pressig,\\
\\
teste tem a vantagem de\\
\\
ste de queda de pressao, mas a pre: >€ med\\
: h\\
\\
ade: similar ao te: : =\\
jade: simile Jmente Agua) no reserva 6rio com vazao ae\\
fluido (normal :\\
buildup, mas medindo a queda de Btesefy '\\
\\
“Pos\\
\\
stabilizar em um valor médio do reservaty\\
\\
Teste de Injetivid\\
enquanto se injeta\\
\\
similar ao teste de\\
o tendera a se \\&\\
\\
te\\
\\
Teste de Falloff:\\
cessar a injegao. A pre\\
\\
ro; :\\
\\
de Interferéncia: Nesse teste, a variagio de pressdo no poco Observado é me.\\
\\
Fig) ido a produgio ou injegdo em outros pogos que drenam ou injetam no te\\
lida devi\\
\\
servatério;\\
\\
Teste a Poco Aberto Drill Stem Test/ Teste de Formagao (ST): esse teste é real.\\
zado durante ou apés a perfuracao do pogo, utilizando a prépria coluna de perf.\\
ragao como coluna de producao. Sao testes curtos, geralmente com dois fluxos ¢\\
dois registros de pressio do reservatério, medindo a recuperacao da pressao apés o5\\
\\
fluxos.\\
\\
Quanto a qualidade do teste realizado, ele pode ser classificado como conclusivo, parcial\\
mente conclusivo ou falhado, conforme descrito a seguir. Um teste conclusivo é aquele\\
que é realizado sem problemas e obtém as informagées planejadas. Um teste parcial-\\
mente conclusivo é aquele em que os resultados sao parciais, como conseguir medir va-\\
z6es, mas nao as pressdes. Um teste falhado é aquele que nao é bem-sucedido devidoa\\
problemas operacionais, como falta de assentamento dos packers, entupimentos, ou falha\\
na abertura das valvulas. Um teste falhado deve ser repetido.\\
\\
5.2.4.2 Ferramentas da coluna de teste de formacgao\\
\\
Acoluna de teste d 5\\
\\
patos = psi € composta por um conjunto de ferramentas selecionadas\\
\\
principais co) ee reads Condigdes mecanicas do PO¢o € nos objetivos do teste. Os\\
mponentes de uma coluna de teste de formaco incluem:\\
\\
* Registrador de Pressio Externo:\\
sdo e uma unidade de registro par:\\
de teste ao longo do tempo;\\
\\
. Tubos de Perfuracao: ferram\\
\\$40 para dentro da tubulacao;\\
\\
* Obturador (P; :\\
(Packer): quando ssentado, suas borrachas ve\\
\\
lando a formacio da 7\\
anular; s40 da pressio hidrostatica do fluido de a\\
\\
* Registrador de Pressio In\\
“ terno: seme\\
localizado dentro da coluna de ie Ihante ao registrador de pressao exte\\
\\
essa ferramenta consiste em uma unidade de pres\\
a medir continuamente a pressfo externa a coluna\\
\\
lentas 2 .\\
dN¢ Permitem a passagem dos fluidos da form\\
\\
dam o espaco anula\\
: x ¢0\\
Mortecimento no ¢sps\\
\\
rno, ms\\
\\
d\\
Digitalizada com CamScanner\\

\\section*{Pagina 149}

ESTUDOS DE RESERVATORIO. 129\\
\\
+ Conjunto de Valvulas: essas valvulas sio operadas na superficie ¢\\
tura ou fechamento da coluna de teste, A valvula de fu\\
fluidos na coluna de teste durante a descida da colun,\\
\\
Permitem a aber.\\
\\
indo impede a entrada de\\
\\
‘4 NO pogo;\\
\\
+ Valvula de Circulagao Reversa:\\
espaco anular para o interior da\\
dos durante o teste.\\
\\
essa ferramenta permite a circulagao no sentido do\\
coluna, permitindo a remogio dos fluidos produzi-\\
\\
5.2.5 Teste de producao\\
\\
Durante o ciclo de vida de um campo, os testes de producao sio usados para coletar dados\\
de produgao, como vazio e pressdo. No teste de produgio, o poco é fechado na cabeca.\\
A pressio é medida na cabeca do pogo, o que gera incerteza na estimativa da pressio\\
estatica devido ao efeito da compressio dos fluidos dentro do poco (estocagem). Apesar\\
disso, cle tem a vantagem de, as vezes, nao necessitar de uma sonda.\\
\\
Oteste de producao no utiliza uma valvula de fundo, O controle do poco (abertura\\
e fechamento) é realizado na superficie. A coluna de teste pode ser descida da mesma\\
maneira que em um teste de formagao, com registradores acoplados, ou estes podem\\
ser descidos posteriormente. O teste de producao pode envolver uma coluna de teste\\
especifica ou a propria coluna de produ¢ao pode ser usada Para esse propésito. A principal\\
caracteristica de um teste de produgao é 0 tegistro de pressdo e temperatura no fundo,\\
juntamente com as medicées e o controle do po¢o na superficie.\\
\\
Os testes de produc3o apresentam algumas desvantagens em relacdo aos testes de for-\\
macao. Por exemplo, eles envolvem um fechamento na superficie, o que os torna mais\\
suscetiveis a vazamentos. Além disso, ha um periodo excessivo de estocagem e segrega-\\
sao de gas, ea auséncia de uma valvula de fundo reduza seguran¢a do teste. Muitas vezes,\\
€ necessario usar um colchao até a superficie, o que implica em operagées demoradas e\\
Custosas de inducao de surgéncia.\\
\\
No entanto, o teste de producio oferece uma avaliacao mais precisa dos fluidos produ-\\
zidos, permitindo a medicao das vazées de dleo, gas e agua, além de possibilitar 0 acom-\\
Panhamento do comportamento da pressao do poco. Quando o teste é de longa dura¢ao,\\
48 condices sio ideais para avaliar a situacao do reservatério e melhorar as estratégias\\
de desenvolvimento € producao do campo. A producio durante o desenvolvimento do\\
‘ampo é chamada de producio antecipada. Os testes de longa dura¢io permitem obter\\
Pardmetros mais precisos do reservat6rio apés um longo periodo de produgao e também\\
antecipam a producio do reservatorio. Geralmente, esses testes so realizados em dreas\\
“ploratérias antes de realizar investimentos significativos.\\
\\
5.3 Estudos de reservatério\\
\\
. . ala tai reservatorio,\\
\\
Ms nna companhia petrolifera, grande parte do foco esti volta ee o aharia de Re\\
\\
© Capitulo 6, abordaremos temas mais aprofundados relacionados a nge : " ae\\
“rVatérigs, No entanto, a avaliagdo das formagées é uma operagao cruct\\
\\
; ¥. de maneira precisa.\\
*AUisicao de informacées detalhadas para caracterizar 0 reservatorio de maneira pre\\
\\
Digitalizada com CamScanner\\

\\section*{Pagina 150}

|\\
\\
130 AVALIACAO DE FORMACOES\\
\\& essencial qualificar e quantificar as Propriedade\\
s Ne : ;\\
preender as condigées existentes no reservatdri ,\\
produgao de hidrocarbonetos (mecanisme,\\
0\\
\\
5 de reservatorio,\\
dos, além de com|\\
lo reservatorio durante 4\\
\\
Durante os estudo:\\
da rocha ¢ dos flui\\
\\
© comportamento d\\
de produgio). cha, destacam-se a porosidade, per\\
: i riedades da rocna, ted s » Permeabili.\\
No que diz AEE OA ata 3s propriedades dos fluidos, incluem-se a densidade\\
dade e composigao q \\textbackslash{}. ‘mica. Além disso, € importante considera;\\
\\
viscosidade, poder calorifico e composi¢a0 qui\\
\\
as condicdes presentes NO reservatorio, como t\\
\\
fluido e molhabilidade.\\
O sucesso das opera¢'\\
\\
da formagao produtora, ¢\\
\\
emperatura, pressao, saturacao de ca, di\\
\\
envolve analises de laboratério em amostras\\
élulas de Pressao, Volume ¢ Temperatura (PVT), medicées de\\
porosidade, permeabilidade e pressao capilar, bem como modelos fisicos reduzidos e 5.\\
mulacoes teéricas. As ferramentas tedricas permitem simular diferentes configuracées ds\\
estruturas produtivas, localizagao ¢ espacamento entre pogos, vaza0 Por Poco € outros\\
parametros. Um estudo teérico bem elaborado facilita o desenvolvimento de um plano\\
de campo adequado.\\
\\
© entendimento aprofundado das propriedades da rocha e dos fluidos, juntamente\\
com as condicGes do reservatorio, € essencial para a tomada de decisGes informadas ao\\
planejar e otimizar a produgao de hidrocarbonetos. Com base nessas informacées, os\\
engenheiros de reservatério podem determinar a estratégia mais eficaz para maximizar\\
a recuperacio de éleo e gis, otimizando o fluxo de producao e minimizando os desafios\\
\\
cionais.\\
\\
des frequentemente\\
\\
5.3.1 Cut-offs\\
\\
aca os estudos de reservatorio, determinam-se alguns valores limites das propried-\\
Para que 0 reservatério seja produtor de forma economicamente viavel. Esses limites\\
\\
sio conhecidos como “Cut-offs” 4 =\\
desses parimetros: ‘ut-offs” ou parametros de corte. A seguir, descreveremos alguns\\
\\
* Cut-offs de porosidade: utiliza-\\
Para reservatérios de Fy Se 0,08 (8\\%) para reservatérios de dleo e 0,06 (6\\%)\\
\\
* Cut-offs de saruracio: normalm\\
Para valores de Sw (sai ear\\
vatério produza 6leo,\\
produgao de agua;\\
\\
cena é utilizado 0,5 (50\\%) para a saturacao de agua.\\
pois a eee acima de 50\\%, é improvavel que 0 rese\\
eabilidade relativa favorecera completamente @\\
\\
Digitalizada com CamScanner\\

\\section*{Pagina 151}

=\\
\\
ESTUDOS DE RESERVATORIO 131\\
\\
conclusao\\
\\
Neste capitulo, © leitor teve a oportunidade de aprender conceitos basicos relacionados\\
go processo de avaliacao de formacGes, uma operacao crucial Para o desenvolvimento de\\
campos de petrdleo e gas natural. Ao final desse Processo, sio obtidas as informacdes\\
para a tomada de decisées sobre a viabilidade do prosseguimento de um pro-\\
jeto, seja para completar um poco, iniciar a produgo ou até mesmo perfurar novos pocos.\\
Entre os muitos parametros obtidos no processo de avaliacao de formacées, destacam-se\\
aqueles necessarios para estimativas de volumes e reservas.\\
\\
O conhecimento adquirido através da avaliago de formagées é essencial para o pla-\\
nejamento estratégico das atividades de exploracao e producao de hidrocarbonetos. As\\
informacGes obtidas contribuem para uma compreensi4o detalhada das caracteristicas do\\
reservatorio, permitindo a formulacio de estratégias eficazes para a otimizacao da recupe-\\
racho de petréleo e gas. Dessa forma, os estudos de reservatério desempenham um papel\\
fundamental na industria do petréleo e gas, auxiliando na maximizacao da produgao e na\\
tomada de decisdes informadas para o sucesso dos empreendimentos.\\
\\
necessarias\\
\\
Digitalizada com CamScanner\\

\\section*{Pagina 152}

CAPITULO 6\\
\\
RESERVATORIOS\\
a\\
\\
Aengenharia de reservatérios é um ramo da engenharia de petréleo no qual se aplicam\\
os principios cientificos aos problemas de drenagem que surgem durante o desenvolvi-\\
mento e producao de reservatérios de petrdleo e gas natural, de modo a obter uma alta\\
recupera¢do econémica.\\
\\
O objetivo principal da engenharia de reservatérios é maximizar a producdo de hidro-\\
carbonetos com o menor custo possivel. Neste contexto, os estudos na engenharia de\\
reservatorios baseiam-se na caracterizacio das jazidas, nas propriedades das rochas e dos\\
fluidos de reservatérios, na intera¢ao dos fluidos dentro das rochas e nas leis fisicas que re-\\
gem o movimento dos fluidos no seu interior. As ferramentas de trabalho do engenheiro\\
de reservat6rios incluem geologia, matemitica aplicada e modelos computacionais, além\\
das leis basicas da fisica e da quimica que regem 0 comportamento das fases liquida e va-\\
por de petrdleo bruto, gés natural e Agua na rocha-reservatorio.\\
\\
O objetivo geral é a aquisicao de conhecimentos basicos das atividades de um Enge-\\
nheiro de reservatorios e a compreensio das propriedades de rocha-reservatério de hidro-\\
\\
carboneto e das diferentes fases de Tecuperacao de hidrocarbonetos, enquanto os objeti-\\
Vos especificos pretendem:\\
\\
* Descrever as principais atividades de um Engenheiro de reservatdrios;\\
\\
* Conhecer as propriedades basicas das rochas e dos fluidos de reservatérios e calcular\\
0S seus parémetros para estimar o volume original e as reservas de hidrocarbonetos;\\
\\
* Explicar a importancia do envelope de fase para definir 0 tipo de reservatorio;\\
* Identificar e entender os mecanismos de recupera¢ao de hidrocarbonetos;\\
\\
" Compreender e aprender a calcular as reservas e diferenciar as principais definicdes,\\
como fator de Tecuperacao, fracao recuperada, volume original, volume recuperavel\\
€ producdo acumulada.\\
\\
6.1 Atividades basicas de um engenheiro de reservatérios\\
\\
De forma detalhada, dentro das atividades de engenheiro de reservatérios destacamos:\\
\\
: 133\\
ce Pees Para ndo-engenheiros, 1* edigdo.\\
André Raposo Ramos Copyright © 2023 Editora Dois03.\\
\\
Digitalizada com CamScanner\\

\\section*{Pagina 153}

134 RESERVATORIOS\\
sedades do fluido contido no reservatorio: viscosidade .\\
prie to\\
\\
= Determinar as pro) :\\
densidade etc.\\
\\
posigao da mistura, .\\
\\
-reservat6rio: porosidade, permeabj;\\
i iedades da rocha-r" hing\\
* Determinar as proprieda se ‘\\
\\
capilaridade, saturagao etc.;\\
nto do fluido no i\\
ervatorio durante\\
rtamento de reservat6rios: baseando-se no hits\\
muladores de fluxo;\\
\\
interior da rocha-reservatério; compo\\
= Estudar o comportame s paodlooten 7\\
mento da pressao do res\\
\\
ao de compo i\\
lizando-se de modernos si\\
\\
to, planejamento ¢ desenvolvimento de\\
\\
» Simulagio e previs\\
rico de produgao ut\\
\\
. 1\\}\\
«= Estimativa de reservas: acompanhamen'\\
\\
campos;\\
» Interpretacio de resultados de testes de pressao;\\
\\
= Métodos de recuperacao.\\
\\
6.2 Propriedades basicas das rochas e fluidos de reservat6-\\
rios\\
\\
s basicas das rochas e fluidos de reservatério é crucial\\
para o desenvolvimento dos estudos no reservatorio de petrdleo. Essas propriedades sio\\
usadas para determinar 0 volume ou a quantidade dos fluidos existentes no reservatorio\\
(meio poroso), sua distribuicao, a capacidade desses fluidos fluirem e a quantidade de\\
fluidos que pode ser extraida. As propriedades basicas das rochas de reservat6rio incluem\\
porosidade, satura¢ao, molhabilidade, pressao capilar, permeabilidade, compressibilidade\\
¢ mobilidade. As propriedades bisicas dos fluidos incluem massa especifica, densidade do\\
ae | e es °API), fator de volume de formacio de fluidos e razao de solubilidade\\
pa — a cumprir com os objetivos deste curso, vamos nos concentrar no estudo\\
\\
porosidade, permeabilidade e saturagao, bem como na interagao rochas-fluidos.\\
\\
© conhecimento das propriedade:\\
\\
6.2.1 Porosidade\\
\\
mera rae ensleo » ma piieeta de se acumula na rocha reservatorio. Essa\\
0s vazios estejam interconectados te (porosidade) e é essencial que esses ae\\
volume total ocupado pela rocha pie i ah o-lhe a caracteristica de permeabilidade.\\
igi eis © canto) ¢ dos espai ae a soma do volume de materiais sOlidos\\
ISOS vazios existentes entre eles. O termo “cimento\\
\\
refere-se aos grios ligad\\
i los uns\\
\\
nial muito fino sio chamados ne outros por um material, enquanto os graos 0U mate:\\
reservatorio. atriz. A Figura 6.1 ilustra a composigao de uma 0\\
\\
Digitalizada com CamScanner\\

\\section*{Pagina 154}

PROPRIEDADES BASICAS DAS ROCHAS E FLUIDOS DE RESERVATORIOS 135\\
\\
Figura 6.1: Composicfo da rocha-reservatério,\\
\\
palavras, porosidade (q) é definida como uma porcentagem ou fracao do volume poroso\\
\\
‘al da rocha (bulk volume). A porosidade dos\\
reservatérios comerciais pode variar de cerca\\
\\
de 5\\% a cerca de 30\\% do volume aparente\\
total da rocha (bulk volume). A porosidade de uma rocha é definida pela Equacio 6.1:\\
\\
V, = V,\\
g=—= MENA = =k (6.1)\\
V; V, Vp +V;\\
€o volume total da rocha (bulk volume) \\& definido como:\\
\\
Vi =Vp+V; (6.2)\\
\\
€0 volume aparente total da rocha, Vp € o volume poroso, e\\
Na Figura 6.3 ilustramos a classificacao de porosidade.\\
\\
onde: * € a porosidade, V;\\
V, € 0 volume de sélidos.\\
\\
Figura 6.2: Classificagao de porosidade.\\
\\
Exemplo 6.1\\
\\
ig ic\\
Calcular a porosidade de uma amostra com o volume poroso de 10 cm” e um volum\\
total de 60 cm?.\\
\\
Solucao 6.1\\
\\
Digitalizada com CamScanner\\

\\section*{Pagina 155}

136 —\\_RESERVATORIOS\\
\\
ou\\
\\
a\\
V, 10em 498\\
p - OSB. x 100 17\\%\\
ory, x 100= Goem?\\
dade\\
\\
6.2.1.1 Fatores que afetam a4 porosl'\\
Tamanho do grio: © tamanho do gre nao afetaa porosidade. Sedimentos bem arredon\\
dados que sio embalados 10 mesmo arranjo geralmente apresentam Porosidade de\\
\\
26\\% a 48\\%, dependendo da embalage™:\\
\\
cio: sedimentos bem clasificados (sy\\
netro, geralmente tém porosidades\\
ly sorted). Seo sedimento tiver uma\\
es poderao preencher os\\
\\
(sorting). arranjo ou sele\\
os esféricos, do mesmo did\\
\\
mal classificados (poor\\
as particulas menor’\\
\\
Classificagao das rochas\\
sorted), ou seja, gra\\
mais altas que sedimentos F\\
variedade de tamanhos de particulas,\\
espagos vazios entre as particulas maiores;\\
\\
‘om formas irregulares t\\
\\
‘o as particulas arredondadas,\\
\\
Forma dos grios: particulas ¢ endem a nao se empacotar de ma\\
neira tio ordenada quant resultando em proporcées\\
\\
mais altas de espago vazio.\\
\\
6.2.1.2  Porosidade absoluta\\
\\
A porosidade absoluta é a relagdo entre 0 volume total de vazios (vi\\
Id\\
s (volume poroso) de uma\\
\\
6.2.1.3 Porosidade efetiva\\
\\
A porosidade pode ser efetiva ou al rr ovolume\\
le ser efeti i\\
bsoluta. A porosidade absoluta é a razao entre |\\
\\
4\\
\\
4 de todos os poros interconectados ou nio e o vol\\
\\
4 porosidade efetiva é a razio entre o volume d une total da rocha. Por outro lado, 2\\
\\
bs rocha, A Figura 6.3 ilustra e de poros interconectados e o volume total da\\
a porosidade de uma rocha.\\
\\
rosidade priméria fluid\\
‘ los\\
secundaria, A AH hae ser extraido da rocha, Bxis\\
Primaria é gerada junto ao sediment?\\
\\
ou\\
\\
Digitalizada com CamScanner\\

\\section*{Pagina 156}

PROPRIEDADES BASICAS DAS ROCHAS E FLUIDO:\\
S DE RESERVATOR|\\
rl 10S 137\\
\\
ha. Por outro lado, a porosidade secundaria se desenvolve apésa formagio das roch\\
wn “ee ochas,\\
\\
On exemplo de porosidade secundaria é o aparecimento de fraturas 0U seja, o surpi\\
nais espacos vazios. A medicao da porosidade ¢ feita a partir de perfis elétrico\\
” s\\
\\
n\\
\\
nent de Sees . ra\\
\\
m ecutados nos pogos ou de ensaios de laboratério em amostras da rocha,\\
ex .\\
\\
6.2.1.4 Porosidade primaria e secundaria\\
\\
porosidade primaria é definida como a porosidade em uma rocha devido ao Processo de\\
cedimentagio, ou seja, aquela que se desenvolve durante a deposicao do material sedimen-\\
tar, Exemplos de porosidade priméaria ou original incluem a porosidade intergranular dos\\
arenitos e as porosidades intercristalina e oolitica de alguns calcdrios (Figura 6.4 Ae D).\\
\\
© ©\\
\\
Figura 6.4: Porosidade quanto a sua origem: A, D porosidade priméria; B, C porosidade\\
secundaria.\\
\\
Porosidade secundaria é definida como a porosidade em uma rocha que ocorre apés 0\\
processo de sedimentacao, por exemplo, fraturamento e recristalizagéo. Exemplos de\\
porosidade secundaria ou induzida sao dados pelo desenvolvimento de fraturas, como\\
encontradas em arenitos, xistos e calcarios, e pelas atividades decorrentes da dissolugao\\
de parte da rocha, comumente encontradas em calcarios (Figura 6.4 B e C).\\
\\
6.2.1.5 Medicao da porosidade\\
\\
Para a determinacao da porosidade de rocha-reservatérios, varios métodos sao utiliza-\\
dos, tais como amostragem (medicao em laboratério a partir de pequenas amostras de\\
testemunhos), perfilagem do poco ou anilise de testes de pressao em algumas situacdes\\
“speciais. O método mais comum é 0 que usa pequenas amostras da rocha-reservatorio.\\
O valor da porosidade de grandes porgées da rocha é obtido estatisticamente a partir dos\\
Tesultados de andlises de numerosas amostras. .\\
\\
Na medigio da porosidade em laboratério, é necessario determinar apenas dois dos\\
tés parimetros basicos: volume total, volume de vazios (volume poroso) © Mee q\\
\\
i . Nesta anilise, pode-se usar 0 método de peso timido e seco, 0 método usando 0\\
Petisometro da lei de Boyle e somatério de fluidos. di e 0 attostta seca\\
\\
No método que utiliza a diferenca de peso entre a amostra saturada e Se gan\\
= inar 0 volume poroso da rocha, pesa-s¢ a amostra seca €, em seguida, s2\\
\\
‘¢ com o mesmo peso de fluido por algumas horas ¢ calcula-se a diferenga.\\
\\
Digitalizada com CamScanner\\

\\section*{Pagina 157}

138 RESERVATORIOS\\
\\
Exemplo 6.2\\
.. Sabe-se que a mesma amostra nj,\\
0 sy\\
\\
esa 120 gramas-\\
P de 0,84 g/c1 3\\
especifica do dleo é ,84 g/cm” (ou 840 kg/m’,\\
\\
mostra saturada com 6leo p!\\
3. Determine a porosidade da amostra\\
\\
eamassa\\
de 173,864 cm\\
\\
Uma a\\
turada (seca) pesa 90 gramas,\\
© volume total da amostra é\\
\\
Solucao 6.2\\
V \\_ 1208-908 . 35,714cm?\\
\\textasciitilde{} p 0, 84g/cm’\\
\\
Pp\\
V, 35,714cm?\\
\\
Bea) \\_ 35) TED = 20,54\\%\\
\\
= 7, *100= 173, 864em? °\\
\\
Msat — Miseeo. =\\
\\
o de fluidos\\
\\
netos nas rochas, elas também contém agua. Portanto,\\
volume poroso para estimar as quantidades de ales\\
¢/ou dgua presentes nas formacoes. Para estimar essas quantidades, énecessario determi:\\
\\
nar qual porcentagem do volume poroso é ocupada por cada fluido. Essas porcentagens\\
Quando o reservatério é descoberto, ele apresenta uma\\
\\
6.2.2 Saturaca\\
\\
‘Além da existéncia de hidrocarbo:\\
nio é suficiente conhecer apenas °\\
\\
sio chamadas de saturacoes.\\
certa saturagao de 4gua, chamada de 4gua conata. A saturaca ) a\\
r cao de dleo, agua e gas é\\
porcentagem do volume poroso (V,) ocupada por cada uma dessas fases, ou a .\\
5, = 2\\
V (63)\\
5, =\\
aot\\
V, (6.4)\\
V,\\
Sy=—\\
w Vp (6.5)\\
Sot+Sgt+Sy=1 a)\\
\\
Onde S, é a saturaca\\
sao de dleo, S, é\\
Figura 6.5 : a saturacai 5\\
ilustra a estrutura da g © @ saturaco de gés e S,, é a saturacao de 4gua A\\
\\
6.2.2.1 Métodos di focha coma saturacao de fluidos.\\
le determinagao da satur:\\
‘agao\\
\\
Digitalizada com CamScanner\\

\\section*{Pagina 158}

PROPRIEDADES BASICAS DAS ROCHAS E FLUIDOS DE RESERVATORIOS §=—-1 3.\\
9\\
\\
Oleo\\
\\
Agu\\
\\
Gas\\
\\
Figura 6.5: Saturacio de fluidos.\\
\\
6.2.2.2 Fatores que afetam a saturacao\\
\\
Todos os métodos de medico direta so falhos devido ao modo como é feita a amos-\\
tragem da formacio e ao manuseio do testemunho desde o fundo do poco até ao labo-\\
ratorio. O filtrado da lama de perfuragao normalmente penetra nos poros da formacio\\
¢ consequentemente altera a distribuicdo dos fluidos. Também por ocasifo da retirada\\
do testemunho para a superficie, devido ao abaixamento de pressio, o dleo ird liberar o\\
gis que se encontra em solugdo, bem como havera uma expansio do dleo, da agua e do\\
gis formado, alterando mais uma vez a distribuicio original dos mesmos. Para evitar a\\
contaminacao no trajeto entre o poco € o laboratério, € praxe em certos casos revestir\\
 testemunho com uma fina camada de para. No caso em que se visa somente a medi-\\
gio da saturagao de Agua, os testernunhos podem ser colocados em recipientes fechados\\
\\
contendo leo diesel.\\
\\
Exemplo 6.3\\
\\
Uma amostra de 62,823 cm? de volume total saturada com Agua e éleo tem 9,0 cm? de\\
volume de gua, com uma porosidade de 25\\%.\\
\\
1, Qual é a saturacao de agua?\\
\\
2. Qual é a saturacao de dleo?\\
\\
Solucado 6.3\\
\\
Vp=oxVi = 0,25 x 62,823cm? = 15,708cm?\\
\\
3\\
Vy \\_ 90cm" \\_ \\_ 9/5730\\
\\
Sy= Vp 15,708cm?\\
\\
Digitalizada com CamScanner\\

\\section*{Pagina 159}

140 RESERVATORIOS\\
\\
3\\
\\
V, 9,0cm" \\_ 100 = 57,30\\%\\
\\# x 100= 75, 708 cm?\\
\\
a\\
Pp\\
2.\\
S Bee 1 20,5780= 04070\\
5, a 1 5. = 100\\% — 57, 30\\% = 42,70\\%\\
io = teow\\
\\
6.2.3 Permeabilidade\\
\\
A permeabilidade de uma rocha € uma medi\\
\\
i é de de quao\\
fluir através da rocha. Isso depen t\\
cha esto interconectados. Quanto mais cheios de estrangulamentos, mais estreitos ¢\\
\\
mais tortuosos forem os canais porosos, maior sera o grau de dificuldade para os fluidos\\
ermeabilidade é uma medida da capacidade de um material\\
\\
b um gradiente potencial. Na Figura 6.6 ilustramos 0s tipos\\
\\
da da facilidade com que os fluidos podem\\
bem os espacos de poros dentro dessa ro.\\
\\
percorrerem seu interior. A p'\\
oso de transmitir fluido sol\\
de permeabilidade de uma rocha.\\
\\
Figura 6.6: Permeabilidade d 2 E "\\
\\textasciitilde{} permeabilidade fraca. le de uma rocha; (A) representa uma boa permeabilidade. (8)\\
\\
em um oe si permeabilidade absoluta, que consiste apenas\\
facilidade de movimento dos Hoes meabilidade efetiva, que é a comparacao entre 4\\
€a permeabilidade relativa, que se at ae na rocha em relacao ao fluido considerado;\\
ividii fetes normalizar os dados de permeabilidade, ov seja,\\
ata , geralmente utiliz. le efetiva por um mesmo valor de permeabilidade\\
\\
Ph jainidade de permeabilidade (we ‘a-s€ a permeabilidade absoluta.\\
Gaspard Darcy, que fiver em homenagem ao cientista francés Henti\\
Beige ee Pacha fe) ome ee agua através de filtros de areiaem 1856\\
we perimental é mostrado esquematicament€ na\\
\\
Darcy observou\\
QUE 05 resulta\\
dos da experiéncia seguiam a equacio 6.7:\\
\\
és ‘ 1g =kAtT 2\\
cant Lé 4 a vazdo de apy : ap (67)\\
; a atr: . ;\\
aaltura do avés do cilindro de areia cuja secio transversal €\\
\\
08\\
\\
colocados nas faces de Meio POro0; hy, ¢\\
fe-\\
\\
h a a\\
a ¢ de saida do 2 Sio as alturas da Agua em manome!\\
\\
filtro (medidos a um mesmo nivel de\\
\\
1\\
Digitalizada com CamScanner\\

\\section*{Pagina 160}

PROPRIEDADES BASICAS DAS ROCHAS E FLUIDOS DE RESERVATORIOS L\\
41\\
\\
Figura 6.7: Esquema do experimento de Henry Darcy sobre o fluxo de Agua através de\\
filtro de areia.\\
\\
réncia) e representam o potencial hidraulico nesses dois pontos; e K é uma constante de\\
proporcionalidade caracteristica do meio poroso e do fluido.\\
\\
Mais tarde, outros investigadores realizaram experimentos com outros fluidos, desco-\\
brindo que a constante K podia ser escrita como x, onde pe 7 sao respectivamente a\\
viscosidade e 0 peso especifico do fluido, e k é uma propriedade da rocha somente, de-\\
nominada permeabilidade absoluta. Para o fluxo horizontal, por exemplo, a equac4o da\\
vazio pode ser escrita como ilustrada na equacao 6.8:\\
\\
\\_ KAAP\\
\\
= eit (6.8)\\
\\
Onde q é a vazao de fluido (cm? /s), Aa area da secao transversal (cm?), AP o diferencial\\
de pressao (atm), pt a viscosidade do fluido (cp), L o comprimento do meio poroso (cm).\\
Apermeabilidade tem por simbolo a letra “k” e a sua unidade de medida mais utilizada é\\
© Darcy, em homenagem ao engenheiro francés Henry Darcy (1803 — 1858). 1 Darcy =\\
0.987 x 10-1? m2. A Figura 6.8 é uma ilustracéo do enquadramento destes parametros.\\
\\
RP\\
\\
Figura 6.8: Ilustracdo da definigo do valor de 1 Darcy.\\
\\
te contexto, 1 Darcy é definida como a permeabilidade que permitira que um fluido\\
de viscosidade de 1 centipoise flua a uma taxa de 1 centimetro ciibico por segundo através\\
de uma area de secdo transversal de 1 centimetro quadrado quando o gradiente de pressio\\
€de 1 atmosfera por centimetro. A equacdo de Darcy (equagio 6.8) foi estabelecida sob\\
cettas condicées:\\
\\
* Flaxo isotérmico, laminar e permanente;\\
\\
° i sccosi invaridvel com a pressao;\\
Fluido incompressivel, homogéneo e de viscosidade inva'\\
\\
Digitalizada com CamScanner\\

\\section*{Pagina 161}

142 RESERVATORIOS\\
\\
i ‘om 0 fluido;\\
é nao reage ©\\
* Meio poroso homogenee, que\\
\\
6 ‘ha.\\
» Fluido tnico saturando 100\\% a roc!\\
\\
idades de Darcy, @ unidade de medida de Permesbildas\\
a de unida Ses realizadas no laboratério, esse sistema de unida\\
s ala de campos. Nesses casos, sio utilizados oe\\
o mD (miliDarcy) como unidade de medidsge\\
yiatura md em vez de mD.\\
\\
No chamado sistem :\\
o Darcy, com excegao das medi¢\\
nao é adequado para os estudos em esc\\
tros sistemas de unidades. Por exemplo,\\
permeabilidade. E comum o uso da abre\\
\\
6.2.3.1 A velocidade de Darcy\\
\\
que 0 fluido no meio poroso ocorre 4 “Velocidade de Darcy,” que é dada por:\\
\\
kaP\\
pl\\
\\
Assume-se\\
\\
(69)\\
\\
yet\\
A\\
A velocidade de Darcy é um conceito macroscépico e de facil determinagao. Esta velo.\\
cidade é diferente da velocidade microscépica de particulas do fluido associadas aos seus\\
particulares caminhos de fluxo através dos graos do meio poroso. As velocidades micros-\\
cépicas possuem existéncia real, mas so impossiveis de se medir. A velocidade de Darcy\\
é uma velocidade ficticia, pois em seu cAlculo assume-se que 0 fluxo ocorre através da se-\\
Gio toda do meio poroso, quando na realidade o fluxo se desenvolve apenas através dos\\
poros interconectados. A Figura 6.9 ilustra os caminhos reais ¢ lineares de fluido.\\
\\
Figura 6.9: Velocidade de Da\\
\\
sumidos pela Lei de Darcy, cy: A -caminhos reais; B — caminhos lineares para 0 flux.\\
\\
6.2.3.2 Caracteristicas da Permeabilidade\\
\\
A permeabilidade \\& um d a\\
do reserva tori, Jieskiae done artanies Parametros na definigao do desempenit\\
aia como ilustrada na tes Tepresentada pela permeabilidade, podes\\
\\
Como descrevemo;\\
s anterio\\
rnb pode ser rep re a md na Tabela 6.1 6 abreviatura de millidarcy nA\\
pas a oe le . Esta escala sofre modificaces ao longo do temp?\\
£ra considerado um valor pobre\\
\\
Digitalizada com CamScanner\\

\\section*{Pagina 162}

PROPRIEDADES BASICAS DAS ROCHAS E FLUIDOS DE RESERVATORIOS,\\
143\\
\\
Tabela 6.1: Classificacdo de reservatério com base na permeabilid,\\
\\
lade.\\
\\
valor da permeabilidade Classificagiio\\
eT\\
k<1mD Reservatério Pobre\\
\\
1<k<10mD Reservatério Médio\\
\\
10<k<50mD Reservatério Moderado\\
\\
50<k<250mD Reservatério Bom.\\
\\
k>250mD Reservatério Muito Bom\\
\\
6.2.3.3 Fluxo linear permanente\\
\\
Na Figura 6.10 est4 esquematizado um bloco horizontal de um meio poroso, saturado\\
com um tnico fluido. A equacao geral de Darcy pode ser expressa como:\\
\\
L\\
\\
Py P,\\
\\
Figura 6.10: Fluxo linear.\\
\\
q\\_\\_k(dp\\
Pt eta ue\\
onde q (cm?/sec) é a vazdo de fluxo através da seco transversal A (cm’), k (darcy) é a per-\\
meabilidade, ps (cp) é a viscosidade do fluxo, P (atm) é a pressao, L (cm) é 0 comprimento\\
¢v (cm/s) é a velocidade aparente na dire¢ao L.\\
\\
Para os casos em que a Figura 6.10 toma dire¢des diferentes, a vazdo qna equacao 6.10\\
para o fluxo linear depende também do Angulo de inclina¢ao. A equacao 6.10 da Fi-\\
ura 6.10 pode também ser escrita como:\\
\\
fad\\
\\
Se invertermos a posicao da diferenca de pressdes para AP = P, - Pj, 0 coeficiente da\\
\\
Mesma equacao torna-se negativo, ou seja:\\
\\
kA (24 ) (6.12)\\
L\\
as equac6es de Darcy sao atualizadas com\\
\\
Para casos com i 12\\
inclinago, Figura 6.11 € 6.12, eid Gad\\
\\
@inclusio do Angulo de inclinagio, como jlustrado nas equagoes\\
\\
Digitalizada com CamScanner\\

\\section*{Pagina 163}

144 RESERVATORIOS\\
\\
q NE.\\
ao\\
\\
clinagao com fluxo para cima.\\
\\
Figura 6.11: Fluxo linear com in\\
\\
—pgsin9\\
(acbaee) Rig\\
\\
qe\\
\\
kA\\
ia\\
\\
Figura 6.12: Fluxo linear com inclinagao com fluxo para baixo.\\
\\
= 4(A-* egsin’) (@19\\
\\
PB L\\
Para o sistema de unidades de campo ou unidades de engenharia, a equacao 6.12 \\textasciitilde{}\\textasciitilde{} \\_\\
escrita como ilustrado na equacao 6.15. Para este sistema de unidades, usa-se 0 coehaen!\\
\\
ser\\
\\
de conversio de 1,127 x 10°:\\
g=1ia7 x10 (A=*2) (615)\\
\\# L\\
Onde q (bbI/d) é a vazdo de fluxo através da seco transversal A (ft?), k (md) é 2 perme\\
abilidade, 1 (cp) é a viscosidade do fluxo, P (psi) é a pressao, L (ft) € 0 comprimen' é\\
(ft/s) € a velocidade aparente na direc3o L.\\
\\
6.2.3.4 Aplicabilidade da Lei de Darcy\\
\\
A Lei de Darcy \\& apli a\\
\\
Re, é utilizad ve a apenas quando o fluxo é laminar. O nimero dé ver de\\
\\
Darcy é sclched na She mS é condi¢ao de fluxo laminar ou turbulento- i os \\&\\
: : < 10 (faix, ara es\\
\\
reservatérios), conforme ilustrada na Figs it ete seseaen\\
\\
Digitalizada com CamScanner\\

\\section*{Pagina 164}

PROPRIEDADES BASICAS DAS ROCHAS E FLUIDOS DE RESERVATOR\\
10S 145\\
\\
1010\\
Gradiente Hidréulles, dp/dL\\
\\
Figura 6,13: Intervalo de aplicabilidade da Lei de Darcy.\\
\\
Além da condigao de ser laminar, a equacao de Darcy é aplicavel ou valida Para os seguin-\\
tes casos:\\
\\
* Nio existe acumulacio;\\
* Fluxo do liquido monofasico;\\
* O meio poroso nao é reativo com o fluido que flui.\\
\\
| Na Figura 6.13, o grafico de q/A contra dp/dL deve produzir uma tinica linha reta, onde\\
\\& a inclinac3o ou o gradiente da reta € k / ju, representando a mobilidade do fluido, que\\
vamos tratar nas se¢des seguintes.\\
\\
6.2.3.5 Fluxo radial permanente\\
\\
Aequacio do fluxo radial é usada em cAlculos de engenharia para expressar aproxirnada-\\
mente o fluxo dos fluidos do reservatorio para dentro do pogo. As propriedades de um\\
sistema de fluxo radial esto ilustradas na Figura 6.14, onde r,, € rp representam os raios\\
do poco ¢ externo do sistema, respectivamente, P,, e P, representam as pressdes no poco\\
F no raio externo, respectivamente, e h é a altura do sistema. Por questoes introdutérias,\\
‘S\\%0 vamos ilustrar as derivagdes da equacao neste nivel da engenharia de reservatorio.\\
\\
a\\
\\
Figura 6.14; Fluxo radial permanente.\\
\\
Te Pi ermanente\\
*ndo em conta as derivagdes através da Figura 6.14, a vazio do fluxo radial per\\
Ser expressa como ilustrado na equagio 6.16:\\
\\
Digitalizada com CamScanner\\

\\section*{Pagina 165}

146 RESERVATORIOS\\
\\
gee 2nkh(Py — Pe)\\
\\#ln(te/ty)\\
O sinal negativo na equagao 6.16 apenas indica que 0 fluxo ocorre n, (ie |\\
do crescimento do valor da coordenada r. Como normalmente tabal oe\\
positivos da vazao de produgao q, a equacao 6.16 pode ser escrita sem\\
\\
ou seja:\\
\\
0 Con.\\
\\
°sinal eget\\
\\_ 2nkh(Pe — Py)\\
\\
\\# pln(re/Tw) 6:\\
\\
Y\\
Nas unidades de engenharia ou de unidades de campo, a equacdo 6.17 pod\\
como: POM Ste\\
\\
kh(P. —\\
q=7,08x10°\\textasciitilde{} (Fe Po)\\
\\
pln(re/tw) 613,\\
\\
es alei de Darcy de que 0 fluido que flui é incompressivel e lan.\\
fluidos de reservatérios sio compressiveis, a lei de Darcy é um,\\
de dleo e Agua. No caso do gas, € usada se o gas produzidy\\
e uma forma modificada é usada nos pocos de gis. Esa\\
\\
As suposicées subjacent\\
nar. A rigor, enquanto os\\
boa aproximagao para o fluxo\\
estiver associado ao petréleo,\\
equacao pode ser escrita como:\\
\\
mkh(P? —P5) ‘\\
pat In(r,/ry — 0,5) 6)\\
Onde a vazio do fluxo q (Mscf/d) em condig6es-padrao de temperatura € pressio, per\\
meabilidade k (md) em diregdo do fluxo, z fator de compressibilidade de gas medida em\\
pressdes médias, 4 (cp) viscosidade do fluido avaliado nas press6es médias, T (CR=°F*\\
460) é a temperatura do reservatorio, e h (ft) é a espessura da formacao ou camada.\\
\\
q=703x 10\\%\\
\\
6.2.3.6 Métodos de medicao de permeabilidade\\
. «os de labo\\
A permeabilidade de um meio poroso pode ser determinada atraves de oe ii\\
rarério com amostras extrafdas da formacao ou por meio de testes de Pree ina\\
nos pocos. A Tabela 6.2 demonstra alguns métodos diretos € indiretos usados\\
\\
para medicao da permeabilidade.\\
\\
Tabela 6.2: Métodos de medicao de permeabilidade-\\
\\
Métodos indiretos Métodos diretos\\
\\
Ar da curva granulométrica. \\textasciitilde{}— Permedmetro de carga const?\\
4 do ensaio de adensamento PermeAmetro de carga varidvel\\
\\
os ta\\
\\
ate\\
\\
Digitalizada com CamScanner\\

\\section*{Pagina 166}

PROPRIEDADES BASICAS DAS ROCHAS E\\
FLUIDOS DE Rese\\
RVATORIOS 147\\
\\
metro pode ser utilizado para medicao da permeabili i\\
Oreo para medi¢ao de liquidos: 0 liquido bie oe eet ee de at\\
or forgado a fluir através de uma amostra em um suporte “coreholder” re oto\\
edida, e a permeabilidade € calculada usando a €quacao geral de Diep 1 .\\
aeeetro para Me dicao de gases: um gas nao reativo (normalmente hélio) toad ig\\
oak ds permeabilidade. O gas flui através da amostra, e a taxa de fluxo ae =e\\
vreaida. A Figura 6.15 ilustra o diagrama esquemitico da medicao de permeabilidack ay\\
= Hassler sob condicées de fluxo em estado estacionario. “\\
\\
Figura 6.15: Diagrama esquematico da medicao de permeabilidade do tipo Hassler.\\
\\
A permeabilidade do ensaio experimental pode ser calculada usando a forma modificada\\
da equacao de Darcy, que leva em consideracao a compressibilidade do gas durante o\\
fluxo. Ou seja:\\
\\
2qnP,AL\\
\\
2 2\\
\\
A(P? - P?)\\
Onde q é a vazao de gas (cm?/s), kg é a permeabilidade do gas (darcy), A é a 4rea da seco\\
transversal (cm?), pu € a viscosidade do fluido (cp), P éa pressao de entrada (atm), P) éa\\
pressao de saida (atm), P, é a pressao atmosférica (atm), e L é 0 comprimento (cm). Para\\
95 liquidos, pode-se usar uma das equacdes lineares que apresentamos anteriormente\\
Para calcular a permeabilidade.\\
\\
ig = (6.20)\\
\\
6.2.3.7 Fatores que afetam a permeabilidade\\
\\
A permeabilidade absoluta é uma propriedade do meio poroso e, consequentemente, in-\\
\\
ariavel com 0 fluido que o satura. No entanto, certos fatores podem alterar o valor da\\
Permeabilidade medida. Entre esses fatores esto 0 efeito Klinkenberg, efeito da reagao\\
fuido-rocha e efeito da sobrecarga.\\
\\
Efeito Klinkenberg\\
As Medicdes de permeabilidade realizadas em laboratério usando gas como fluido resul\\
\\
A aredes\\
‘am em valores maiores do que os reais devido ao escorregamento as nents do\\
meio poroso, o que nio ocorre com liquidos. Esse fendmeno de escorreg:\\
\\
Digitalizada com CamScanner\\

\\section*{Pagina 167}

148 RESERVATORIOS\\
proporcionan maior ¢, Consequey\\
pilidade calculada a partir desse tipo de teste. A ey "\\
9 comportamento da permeabilidade se assemelys a\\
a pressao média a diminui até um limite em que, para tng\\
de liquidos, e a perm infinita, 0 gas Se transformaria em liquido ea permeate\\
sio média hipoteticaine  =peatata para corrigir 0 efeito desse fendmeno de escy\\
i ida seria jj ico\\
so ieceees OND propés a relacao:\\
\\
4s, conhecido como efeito Klinkenberg,\\
gas,\\
\\
mente, uma mal\\
Ottege,\\
\\
b\\
k=keo Pt pe (6.21,\\
\\
a ido nos ensaios experimentais, b éuma consta\\
<9 valorda permeabilidade medido no: he Mt,\\
kel : é ae do experiment, € koo € 8 permeabilidade absoluta. O parametry\\
Poa io de fator de Klinkenberg, ¢ fungao do tipo de gas utlizado e da perme\\
lidade do meio poroso. A Figura 6.16 ilustra a influéncia do tipo de gas usado sobre\\
\\
ili erimento.\\
comportamento da permeabilidade observada no exp\\
\\
Hidrogénio\\
\\
Nitrogénio\\
\\
Didxido de carbono\\
\\
Permeabllidade (md)\\
\\
o a2 a4 O6 Os a 12\\
1/P, (1fot)\\
\\
Fi i 16: i i : a: is ae As\\
ie 16: Efeito Klinkenberg no ensaio experimental com hidrogénio, nitrogénioe \\&*\\
\\
ene miso iver da\\
wn a inclinacdo k,b e uma interc\\
A indinacio € uma funco do pes\\
Figura 6.17. -\\
\\
pressio média de fluxo (1/ Py,) produz uma linha?\\
eptacdo de k,,. b é a fungiio de desvio de Klinkenbe\\
\\}0 molecular e tamanho molecular, como ilustrad®®\\
\\
Efeito da feacdo fluido-rocha\\
\\
Este fendmeno >\\
Score com inch?\\
petméabllidede cn aoe quando o meio poroso contém argila hidrativen\\
Seguintes, abardaremos com. sa de salinidade menor que a da formagao- Nass\\
Mas detalhes a interacZo rocha-fluidos.\\
\\
Digitalizada com CamScanner\\

\\section*{Pagina 168}

PROPRIEDADES BASICAS DAS ROCHAS F 1005\\
FLUIDOS DE RESERVA\\
5 OE TORIOS 14\\
9\\
\\
Figura 6.17: Grafico de ky em funcao de 1/P,\\
\\
m:\\
\\
Efeito da sobrecarga\\
\\
Aremogio da amostra da formacao, com consequente alivio da sobrecarga (peso das ca-\\
madas superiores), acarreta alteragSes na rocha e na permeabilidade. A Figura 6.18 ilus-\\
tra a variacao da permeabilidade com a variagiio da sobrecarga para diferentes amostras,\\
\\
onde k,, 6a permeabilidade a uma dada pressio de sobrecarga e ks--g 6a permeabilidade\\
4 pressio de sobrecarga nula.\\
\\
Figura 6.18: Efeito da sobrecarga sobre a permeabilidade.\\
\\
6.2.3.8 Permeabilidade relativa\\
\\
A permeabilidade relativa de uma rocha-reservatério é influenciada por varios fatores,\\
\\
a saturaco de fluidos, molhabilidade, geometria do meio poroso e pressoes\\
. Nas secdes anteriores, aprendemos que a permeabilidade relativa é uma nor-\\
\\
izacio dos valores de permeabilidade efetiva, ou seja:\\
\\
ko (6.22)\\
kyo = =\\
\\
ke (6.23)\\
keg =a\\
\\
ky (6.24)\\
krw =\\
\\
1\\
Digitalizada com CamScanner\\

\\section*{Pagina 169}

150 RESERVATORIOS\\
\\
onde kyy € a permeabilidade relativa ao dleo, Kygéa Permeabilidade Te]\\
\\
fax Ativa ago.\\
€ a permeabilidade relativa 4 Agua. Medi¢des de permeabilidade Telativa gs Sieh\\
\\
iva\\
mente realizadas em amostras de testemunho para determinar a8 quanti\\
\\
de fluidos que fluirao através das rochas quando mais de fea fase Aluida ¢.\\
A permeabilidade relativa em uma rocha sents ie) i. e4\\
sentada em fungio da saturacio de Agua, como pode ser obsery;\\
\\
S40 ro,\\
a\\
des relat,\\
\\
Stiver P\\
re\\
mente ame\\
\\
igura 619,\\
\\
tinein,\\
\\
gua gera|\\
ado na Fi,\\
\\
Figura 6.19: Curva de permeabilidade relativa em funcao da Saturacao de Sgua,\\
\\
Em resumo, a permeabilidade relativa é uma A\\
dade efetiva. Ela varia de 0a 1, sendo 0 quando a Permeabilidade rel,\\
1 quando o meio poroso esta 100\\% saturado. Inicialmente, a perm\\
6leo corresponde a permeabilidade absoluta, e\\
um, ja que o meio poroso esti total\\
Saturacdo de 4gua aumenta gradu:\\
©a saturacao critica de Agua (S\\
No grafico da Figura 6.19, a agua come¢a a fluir quando a Saturacao de agua atinge\\
cercade 20\\%. A medida que as\\
\\
aturacdo de gua aumenta, a permeabilidade efetiva a4gua\\
©, Consequentemente, a permeabilidade relativa\\
\\
de 6leo diminui. 4 Permeabilidade relativa ao éleo também diminui 4 medida que®\\
Saturacao de éleo diminui, até que 0 dleo pare de fluir quando a satura¢ao de égua ating?\\
cerca de 80\\%. Nessa Saturacao, apenas a Agua flui, ea saturacao de dleo corresponte?\\
Saturacdo residual de dleo (Sor), que é de 20\\%, O processo inverso, partindo de um mé?\\
Poroso 100\\% saturado com agua, é andlogo, \\{do mo-\\
O proceso de imbibicao Scorre quando a saturacéo do fluido preferencial (uido ud\\
Ihante) aumenta, enquanto a drenagem ocorre quando a saturagao do fluido nao ie\\
rencial (fluido no molhante) aumenta, A diferenca entre as curvas na Figura 6.20 \\_\\
minaré a saturacio minima que resulta\\
\\
ativa € inexistentee\\
eabilidade efetiva ao\\
a permeabilidade relativa a0 dleo éiguala\\
mente saturado com 6leo (Sy, = 0). A medida quea\\
\\
almente, uma fase continua de gua comegaa se formar,\\
we) € atingida,\\
\\
Digitalizada com CamScanner\\

\\section*{Pagina 170}

PROPRIEDADES BASICAS DAS ROCHAS ¢ FLUIDOS DE Res\\
ERVATORIOS\\
151\\
\\
Figura 6.20: Influéncia da saturacao sobre as curvas de permeabilidade relativa em u\\
sistema dleo/ Agua, preferencialmente molhado por 4gua. mt\\
\\
Medicdo da permeabilidade relativa\\
\\
Ascurvas de permeabilidade relativa podem ser obtidas de duas maneiras distintas. A pri-\\
meira consiste em método experimental onde as medidas sio realizadas em testemunhos\\
de rochas no laboratério. Essas medic6es sao feitas rotineiramente em amostras de tes-\\
temunho, para definir as quantidades relativas de fluidos que fluirao através das rochas\\
quando mais de uma fase fluida estiver fluindo. A segunda é 0 método analitico em que\\
se baseia em correlacGes empiricas disponiveis na literatura.\\
\\
6.2.4 Molhabilidade\\
\\
Molhabilidade é a tendéncia de um fluido de aderir a superficie de um sdlido, em presenca\\
de outros fluidos imisciveis. Esta tendéncia é medida de forma mais conveniente através\\
do angulo de contato (9). A molhabilidade é completa quando o Angulo de contato é\\
igual a zero graus. Ao passo que a molhabilidade nao completa este Angulo é igual a 180°\\
(Figura 6.21).\\
\\
Quando o Angulo de contato é menor que 90° diz-se que 0 liquido mais denso molha\\
preferencialmente 0 sdlido e quando é maior que 90° diz-se que 0 liquido menos denso\\
molha preferencialmente 0 sélido. A distribuicdo dos fluidos no reservatério é fungao da\\
molhabilidade. Geralmente distingue-se:\\
\\
* Fase molhante (aderida a rocha): usualmente é a fase aquosa;\\
\\
* Fase nio-molhante: usualmente a fase organica (leo e gas).\\
\\
Rocha molhada por égua Rocha molhada por dleo\\
\\
: eusiiiaite\\
Figura 6.21: Molhabilidade de fluidos.\\
\\
Digitalizada com CamScanner\\

\\section*{Pagina 171}

RESERVATORIOS\\
\\
152\\
a ocupar OS Poros Menore.\\
fase molhante rende ‘ = Atualmente, a ny . = Mane,\\
yas, a tas ’ be ’ 10) ‘4\\
Devide as forgas atrativass 8° canais mals \\& lidos na indi hab\\
sthante ocupa os P* mente compreendidos na indistria dg yt\\
fase nio-anolha smnenos 140 complet Petriy,.\\
Jes tenor\\
\\
um dos grane\\
\\
6.2.5 Mobilidade\\
\\
capi entre a permeabilidade efetiva des,\\
«dy como a razdo en 88\\
Jefinida com\\
\\
zi ¢ flu, lly\\
\\
A mobilidade do atten’ digdes de reservatorio, Ou seja, a mobilidade de dleo ede ten\\
SI cor 'S\\
¢ sua viscosidade nas ¢ ‘ ties\\
matematicamente S10 determinados ce\\
Ke\\
Ao = Ho (62s\\
kw\\
Aw = Ne los\\
\\
iis ivas, as mobilidades também dependem das\\
. s permeabilidades efetivas, as mo! «| satura\\
pra semper a razio de mobilidade é a razao entre a mobilidade do flutdo des.\\
ee inte w) atras da frente de avango ea mobilidade do fluido deslocado (0), ou seja:\\
\\
M= oa (627)\\
A razio de mobilidade é diretamente relacionada 4 eficiéncia de varrido de um reserva.\\
trio. Razdes de mobilidade favoraveis (M \\textasciitilde{} 1) implicam em melhores eficiéncias de\\
varrido, ji que a interface entre as fases torna-se mais uniforme. Ja razGes de mobilidades\\
muito altas levam a frentes nao uniformes diminuindo a Area varrida do reservatérioe,\\
consequentemente, a eficiéncia do método de recuperagdo. Ou seja, quanto maior fora\\
razio de mobilidades, menor ser a eficiéncia de deslocamento de leo, uma vez que, ¢e\\
vido: 4 sua maior mobilidade, o fluido injetado tender a “furar” o banco de éleo criando\\
caminhos preferenciais entre pocos injetores e os produtores.\\
\\
6.2.6 Pressdo capilar\\
\\
P NOE ha ¢a de pressao existente entre duas fases decorrente das tenso\\
RO na Figura 6.22 e matematicamente escritas como:\\
\\
(6.28)\\
\\
R= gh(py - =\\
Par) = ghAp\\
Onde P. é a presi i\\
\\
Par \\& atid ie 8 €a forca de Sravidade, p,, € a massa especifica Ey i\\
poros de um meio poroes ¢ h €a altura do capilar. O deslocamento de fluidos .\\
meio poroso parcialmente ae ou dificultado pela pressao capilar. Para mane\\
\\
Pressio deste maior do “urado com uma fase ndo-molhante, é necessitio\\
Quando dois Nuidos esa on 8 Olhante. ) 4\\
Monks Na pressio Petia: ees €contidos numa estrutura porosa, ume\\
\\
Pressdo através da interface \\& Perle vieog que separa os dois fluidos. Ess4 ;\\
\\
ecida como pressio capilar e é definida a\\
\\
Digitalizada com CamScanner\\

\\section*{Pagina 172}

PROPRIEDADES BASICAS DAS ROCHAS € FLUIDOS DE Rese AVATORN\\
OS 1B\\
\\
Figura 6.22; Pressio capilar entre Agua e ar,\\
\\
Pe = Pry - Py (6.29)\\
\\
Onde P,,, € a pressio na fase nio-molhante e P,, é a pressio na fase molhante. Conside-\\
rando as condigdes de equilibrio na drea elementar em torno de um ponto na interface,\\
sendo essa interface uma curva, tém1-se dois raios principais Ry e R perpendiculares entre\\
si, relacionados pela Lei de Young-Laplace:\\
\\
L 1\\
Paw — Py = 0 (= - ral (6.30)\\
A pressio capilar nesta equagiio pode ser interpretada como uma medida da tendéncia\\
de um meio poroso de absorver um fluido nio-molhante. Para manter um meio poroso\\
parcialmente saturado com uma fase niéo-molhante, é necessario manter a pressio deste\\
maior do que a da fase molhante. No sistema 6leo-agua, a pressio capilar pode ser tam-\\
bém determinada como:\\
\\
h\\
P.= Taq (Pw \\textasciitilde{} Po) (6.31)\\
2G wo cos O (6.32)\\
\\
1 =\\
78 (Pw \\textasciitilde{} Po)\\
\\
onde h é a altura capilar, r \\& 0 raio capilar, g é a aceleragio gravitacional, \\# é a massa\\
\\
‘specifica e o \\& a tensio interfacial.\\
\\
6.2.6.1 Medi¢ao da pressao capilar\\
\\
aboratério, a partir de expe-\\
a rocha-reservatorio. Os\\
0\\
\\
m :\\
Pied te capilar normalmente sao obtidas em 1\\
tien peque: de testemunhos ¢\\
com nas amostras de testem iene\\
- Mais usados sio o da membrana ¢ o da centrifuga. Usando dados de presio\\
“ mmeadhg medidos no campo, ou seja, nos pogos, também é possivel tragar curva:\\
\\
Digitalizada com CamScanner\\

\\section*{Pagina 173}

154 RESERVATORIOS\\
torio. A partir da pressio capilar (q\\
uma dag\\
\\
jar para um dado reserva\\
\\
de pressio capil\\
sivel determinar:\\
\\
saturacao), € pos\\
» otamanho médio dos poros:\\
Wyo C0SO\\
ia ai\\
: (6.33)\\
\\
= Histerese da pressdo capilar:\\
\\
io nao-molhante desloca fluido molhante;\\
\\
— Drenagem: fluid\\
Ihante desloca 0 nao-molhante.\\
\\
— Imbibigdo: fluido mo\\
\\
6.3 Tipos de reservatorios\\
\\
ras de hidrocarbonetos e das diferentes co\\
Nn\\
\\
Devido 4s diferentes composi¢Ges das mistu\\
\\
dicdes de temperatura € presséo, uma acumulagio de petréleo pode se apresentar\\
\\
mente liquida, totalmente gasosa ou uma parte liquida e uma parte gasosa. D ty\\
. De uma\\
\\
pee ey as acumulagdes podem ser separadas em reservatdrios de dleo\\
rvatorios de gas. Na Figura 6.23 esta ilustrado 0 comportamento de fases de Aides\\
10S.\\
\\
ox\\
\\
Figura 6.23: Di ima de fases de u: OC\\
5 ia d ‘ar\\
gra: le uma mistura de hidr bonetos.\\
\\
6.3.\\
3.1 Reservatérios de éleo\\
\\
d\\
\\
Digitalizada com CamScanner\\

\\section*{Pagina 174}

TIPOS DE RESERVATORIOS 155\\
\\
Tg Temperature Tr\\
\\
Figura 6.24: Diagrama de fases - misturas liquidas.\\
\\
exatamente em cima da curva dos pon-\\
\\
Se 0 ponto representativo da mistura se encontra\\
ou simplesmente saturado. Esse € 0\\
\\
tos de bolha, diz-se que 0 dleo é saturado em gas\\
\\
caso da mistura identificada pelo ponto 1 na Figura 6.24.\\
Sea mistura esta nas condicdes representadas pelo ponto R, por exemplo, diz-se que 0\\
\\
dleo é subsaturado. Para iniciar a vaporizagao da mistura mantendo a temperatura cons-\\
tante, é necessario que a pressao seja reduzida até o valor correspondente ao ponto 2do\\
diagrama, ponto este que se situa exatamente sobre a curva dos pontos de bolha. Para\\
identificar essa pressdo em que comega a vaporizacio da mistura 4 temperatura conside-\\
rada, emprega-se comumente a expressao pressao de saturacao. Assim, os termos ponto\\
debolha e ponto de saturacao, pressao de bolha e pressdo de satura¢ao etc., normalmente\\
sio utilizados de maneira indistinta neste texto.\\
Acurva RS representa 0 comportamento do fluido produzido desde as condicGes ini-\\
dais do reservatorio, ponto R, até as condi¢Ges da superficie, ponto S. No exemplo da\\
figura, nas condic6es de superficie, aproximadamente 60\\% dos hidrocarbonetos produzi-\\
dos estarao na fase liquida e os 40\\% restantes estarao na fase gasosa. De um modo geral,\\
ofluido produzido é submetido a um processo de separacao antes de ser colocado em tan-\\
ques nas condicdes ambientes. A pressdo e a temperatura de separacao sao determinadas\\
através de cdlculos e recebem 0 nome de condigées de separacao.\\
\\_ Ocomportamento do fluido que permanece no reservatorio é representado por uma\\
linha vertical correspondente a temperatura do reservatorio. A pressao vai caindo conti-\\
‘Suamente até as condicdes de abandono, quando ocorre a suspensao da produgao.\\
\\
6.3.2 Reservatério de gas\\
\\
cabo Berar de gis a jazida de petrdleo que contém uma mistura de hidro-\\
\\
Ponto riche eeu no estado gasoso nas condi¢Ges de reservatério, Quando o\\
\\
Tura é Ae nte as condi¢ées de pressao e temperatura a que esta submetida a mis-\\
\\
“emperatura pe hae diagrama de fases, verifica-se que 0 mesmo si\\
, como ilustrado na Figura 6.25.\\
\\
a\\
\\
e localiza a direita da\\
\\
Digitalizada com CamScanner\\

\\section*{Pagina 175}

156 —-ReSE AVATORIOS\\
\\
Figura 6.25: Grafi\\
\\
Os reservatérios de gas sio\\
\\
de gas umido e reservatérios de gas\\
\\
tamento das propriedades di\\
proprio reservatério € tambs\\
cie.\\
\\
separacao dos componentes\\
\\
co de fase de envelopes il\\
\\
ustrando os reservatérios de dleo ¢ gis\\
\\
classificados como reservatorios de gas seco, reservatério,\\
retrogrado. Essa classificagio depende do compor.\\
los fluidos quando sujeitos a reduces de pressio dentro do\\
\\
ém do tipo de fluido resultante nos equipamentos de superf\\
\\
Ao ser levada para a superficie, a mistura gasosa pode ser submetida a processos de\\
\\
mais leves dos mais pesados, resultando dessa separagio duas\\
\\
fases distintas. Os componentes mais leves permanecem no estado gasoso, e 0s mais\\
\\
pesados vao dar origem aos\\
simplesmente por LGN.\\
\\
chamados liquidos de gas natural, normalmente designados\\
\\
6.3.2.1 Reservatérios de gds umido e de gas seco\\
\\
= laut do, \\} a0 ser submetida ao Processo de separacio, produzir uma certa quantidade\\
Oreservatério receberd o nome de rescrvatério de gas timido. Se a quant\\
\\
de liquido for desprezivel, a\\
ser observado que essa clas\\
depende muito fortemente\\
\\
dade\\
\\
jazida receberd o nome de reservatorio de gas seco. Det\\
sificagio depende da composi¢io original da mistura, ™*\\
\\
Z , cer clas\\
dos processos de Separacio. Um mesmo gis poe set clas\\
ara uma determinada condi¢ao de separagio e considerado\\
\\
S€cO para outras condi\\
\\
u S€ 0 processo “iid oe Na verdade, 0 gas s6 deve ser class icado con\\
\\
reserva racao for economicamente vi ce dizer que!\\
tério é considerado de Nomicamente vidvel. Costuma-se (ize a\\
\\
dade economicamente ltencstente Fit quando se consegue obter liquide em\\
\\
izando o equipamento ordinirio de cam?” |\\
\\
prios de\\
nados por alguns autores de reservatones® s\\
\\
Digitalizada com CamScanner\\

\\section*{Pagina 176}

METODOS DE RECUPERACAO DE HIDROCARBONETOS 157\\
\\
22 Reservatorios de gas retrégrado\\
\\
cervatdrio de gas retrogrado recebe esse nome devido a um fenémeno\\
ors vatorios de gas cuja temperatura situa-se entre a temperatura critica\\
enter. A Figura 6.24, a jazida de hidrocarbonetos na qual, nas co\\
Sr emperatura e pressao, toda a pnistira = encontra no estado gasoso. A medida que\\
op fluido vai sendo produzido, a pressao no interior do reservatorio diminui, enquanto a\\
temperatura permanece Constante: Acerta altura da vida produtiva da formacio comecga\\
aocorrer uma condensagao de certos componentes da mistura, ou seja, uma parte do gas\\
se liquefaz. Como prosseguimento da producio, a Pressao continua a cair fazendo com\\
que o gis que tinha se liquefeito volte para o estado gasoso. Diminuindo mais a pressio,\\
todo gas liquefeito eventualmente voltara para o seu estado inicial.\\
\\
O ponto de interesse é 0 fato de uma reducao de pressdo causar condensagao de um\\
gis, quando o esperado é que reduces de pressio causem vaporizacées de liquidos. Este\\
fendmeno retrogrado acontece no interior da rocha-reservatério. O reservatério de gas\\
retrogrado também é conhecido como reservatério de gas condensado.\\
\\
que ocorre em\\
da mistura ¢ a\\
ndi¢6es iniciais\\
\\
6.4 Métodos de recuperacao de hidrocarbonetos\\
\\
A recuperacao de hidrocarbonetos de um reservatério de petréleo é comumente reco-\\
nhecida como ocorrendo em varios estagios de recuperacao. Os fluidos sfo necessarios\\
que estejam submetidos a uma certa pressio para conseguir vencer a resisténcia oferecida\\
pelos canais porosos. A producio ocorre devido a dois efeitos principais:\\
\\
* Adescompressao: que causa a expansao dos fluidos contidos no reservatério e con-\\
tracao do volume poroso;\\
\\
* Odeslocamento do fluido por outro fluido: por exemplo, a invasao da zona de éleo\\
Por um aquifero. Ao conjunto de fatores que fazem desencadear esses efeitos da-se\\
©nome de mecanismos de producao de reservatérios.\\
\\
Os principais mecanismos de producio de hidrocarbonetos sao (Figura 6.26):\\
\\
Recuperacao primaria: recuperacao de hidrocarbonetos do reservatério usando a\\
energia natural do reservatorio como mecanismo de producio;\\
\\
. Recuperacao secundaria: a recuperagao auxiliada ou impulsionada pela injecao de\\
\\#gua ou gis da superficie;\\
\\
* Recuperacdo terciéria ou Enhanced Oil Recovery/ Recuperagio Avangada de Oleo\\
(GOR): uma variedade de técnicas amplamente rotuladas como Recuperagio hort\\
Morada de Petréleo” que sao aplicadas aos reservatérios para melhorar a produga\\
\\
°S campos maduros ou de 6leo pesado;\\
\\
izada quando\\
\\
* Recuperacio em pocos de refinamento de malhas (Injilling Wells): ie\\
\\
ye - rfuragao\\
a a “4 , \\{ sso envolve a per\\
“cuperacio das trés fases anteriores estiver concluida. Is\\
\\
Digitalizada com CamScanner\\

\\section*{Pagina 177}

158 RESERVATORIOS\\
GOs existentes para garantir que Pa\\
T)\\
\\
ado de seu petrdleo.\\
\\
entre 05 P\\
te e588\\
\\
jugio baratos\\
\\
; rod\\
de furos de p tt votalmen\\
\\
reservatorio tenhi\\
\\
cee\\
|\\
\\
||\\
il oe\\
\\
Figura 6.26: Métodos de recuperagao de hidrocarbonetos.\\
\\
6.4.1 Recuperacao primaria (Fluxo natural)\\
\\
ergia natural do reservatério usada para transportar\\
\\
Os mecanismos de produclo sao a en\\
de producao. Existem cinco importantes mecanis-\\
\\
rbonetos para fora dos pocos\\
de producao (ou combinag6es):\\
\\
Mecanismo de gas em solugao: mecanismo de produgao no qual as bolhas de gas dissol-\\
vidas no 6leo se expandem, deslocando o leo para o pogo produtor. Este tipo de |\\
mecanismo de producdo de éleo é causado pela diminuicdo da pressao do reserva-\\
torio, atingindo valores abaixo do ponto de bolha;\\
\\
Mecanismo de capa de gis: mecanismo de producao do reservatério no qual a queda de\\
pressio, acarretada pela produgio da zona de éleo, provoca a expansao do gis. 0\\
7 a ocupar os ee que anteriormente eram ocupados pelo dleo. Como0\\
\\
em uma compressibilidade muito alt: EN j\\
pit Fein o alta, sua expansao ocorre sem que haja queda\\
\\
page per i if fe nenetsie de acionamento ou mecanismo é fornecida\\
vier hipge “as ea com 0 éleo no reseryatério no contato dleo-igY*\\
pire seen ee se esgotam (a produgio continua) € 0 dleo €\\&*\\
Par pe ia se expande levemente. Se o aquifero for grandee\\
suficiente, ree ais grande aumento de volume, que pressionaré\\
\\
is a pressio do reservatério;\\
\\
misto: be ,\\
do qualquer um dos tres ma combinaco ou mecanismo misto ocorte 1")\\
saa anismos opera em conjunto ou quand? que\\
com 0 auxilio da drenagem gravitacional\\
\\
)\\
\\
Digitalizada com CamScanner\\

\\section*{Pagina 178}

METO!\\
DOS DE RECUPERACAO DE HIDROCARBONETOS 159\\
\\
stica, um reservat6rio geralmente incorpora Pelo menos dois mecani i\\
om de produsio; ‘Canismos princi-\\
\\
Mecanismo de segregasao gravitacional ou drenagem gravitacional: mecaniano de\\
ducio de dleo associado A segregacao gravitacional, no qual o gas sai de sol Pro-\\
¢ forma uma capa de gas secundaria, empur gas sai de solucao\\
\\
tando o déleo ee\\
cervatorio onde ele é produzido. Este meca em direcao a base do re-\\
\\
: nismo de producao tem um al\\
de recuperacao do éleo, podendo atingir valores de até 50\\%. O mecaniednd Sie\\
\\
nagem gravitacional geralmente ocorre em campos maduros, que produzem com\\
paixas vazdes de dleo.\\
\\
ich piconet resumidos os principais mecanismos de producao. Embora os mi-\\
meros mostrem percentagens animadoras de recuperacdo (Figura 6.26 e Tabela 6.3), na\\
pritica, a média de recuperacao primaria varia de 5 a 15\\% do volume de éleo original no\\
reservatorio.\\
\\
Tabela 6.3: Resumo de diferentes mecanismos de produsio.\\
\\
ae ‘awed heat emer\\
\\
Pressbo Dedlinio répidoe continuo Permanece elevada\\
\\
Dedinto lentamente\\
continuo\\
RaaBo Gis-Sieo Baba no Inicio, subida Permanece bate Aumenta continuamente\\
abrupta, atinge um méximo os poss estreturas\\
¢ depois cal elevados\\
ProduySo de Agua Desprezivel ou muito baba Comegacedoecresceaté —\\_Desprezivel\\
valores elevados\\
Comportamento de Popo Requerbombelologano —Surgente até a produgo de Longo tempo de surgéncis\\
Inicio Sgua se tome excessiva\\
80 8 60\\% 208 40\\%\\
\\
6.4.2 Recuperacdo secundaria ou convencional\\
\\
A recuperagio secundaria € o resultado da intervencéo humana no reservatério para me-\\
horar a recuperacdo quando a eficiéncia dos mecanismos de producao diminuiram para\\
valores baixos. Duas técnicas sio comumente usadas:\\
\\
* injecio de agua: envolve a injecdo de 4gua na base de um reservatério para manter\\
@ pressio do reservatorio e deslocar 0 dleo para os pocos de produgao;\\
\\
* injecio de gas: esse método é semelhante a injecao hac bee) whl\\
Et j ‘ario O .\\
2 pressio da capa do gas, mesmo que nao s¢ja necess\\
\\
4o utilizam- 4 o gas natural como\\
Nos processos convencionais de recuperacao utilizam-se 2 agua € o ga\\
\\
fluidos de injegio. A Agua injetada pode ter quatro origens:\\
\\
Digitalizada com CamScanner\\

\\section*{Pagina 179}

160 nest RVATORIOS\\
\\
i » subsuperficie por meio de\\
Jetada em manancials de subsuy I - Poop\\
rrinea: coletac\\
\\
+ Agua subterrancat\\
furados para este fim acs\\
E 198, Lagos eles\\
Agua de superficie: coletada em ri\\
2 \\$\\
\\
nai jada a producao de petréleo.\\
A Juzida: Agua que vem associada 4 f\\
« Agua prod uzida: agua\\
Ci la aum tratame:\\
\\
s de ser injetada deve ser submetida 1 < ent, de mg\\
Normalmente, a 4gua as reservatOrio cao fluido nele existente. Eimbora estudog 7\\
atorna-la mais adequa vdemnonstrert numeros animadores, a média mundial d ethonutn\\
pane acho primaria ¢ secundaria geralmente extraem apenas cere,\\
que os métodos de recupera¢: rvatorio. Na Figura 6.27 ilustramos a injegao de Agua\\
\\
6 igi 0 reserva . Na\\
de 35\\% do éleo original n\\
\\
* Pogo do injegtio\\
+ Pogo da produrso\\
\\
Figura 6.27: A - Exemplo de injegio em malhas; B - Exemplo de injegio de base ou\\
\\
periférica.\\
\\
6.4.3 Métodos de Enhanced Oil Recovery/ Recuperacao Avancada\\
de Oleo (EOR)\\
\\
ancados de re-\\
> revisados nesta\\
\\
Método térmico; método fe\\
\\
in situ, que resulta Acer Legs avangada por injegio de vapor ou combustio\\
\\
Seu escoamento através ane na redugao da viscosidade do dleo, facilitand?\\
\\
a abordagem mais eficiente do 4 ltd O método térmico é provavelment\\
Método quimico: esses proce:\\
\\
wnat Para alterar a eficig\\
\\
“eTupados como méx,\\
\\
Serta Interasio quimica\\
\\
ml Produtos quimicos ae\\
‘a de maneira a melhor,\\
‘odos quimicos a]\\
\\
entre 0 fluido inj\\
\\
adicionados 4 agua\\
ar a recuperagio de dle: Bs\\
BUNS processos em que se pressupoe um\\
ftado © fluido do reservatorio. S40 eles\\
\\
i\\
\\
Digitalizada com CamScanner\\

\\section*{Pagina 180}

ESTIMATIVAS DE RESERVAS 161\\
\\
a injegio de solugio de polimeros, injecao de so\\
croemulsio, injegio de solugao alcalina ete, Na\\
como nas outras categorias, sendo que alguns pr\\
tro dos métodos misciveis;\\
\\
lugio de tensoativos, injegao de mi\\
0 existe um PONtO tinico de ataque\\
‘OCeSsOS poderiam se enquadrar den\\
\\
Métodos gasosos: este método usa um fluido\\
tem uma tensao interfacial zero com o dleo e pode, em principio, lib\\
éleo restante no local. Na pratica, é usado um gas, uma vez que os Bases possuem\\
alta mobilidade ¢ podem facilmente entrar em todos os poros da rocha, desde que X\\
\\
gis seja miscivel no dleo, Trés tipos de gis sio comumente usados: CO\\}, Ny gases\\
de hidrocarbonetos; , r\\
\\
que é miscivel com 6 Gleo. Esse fluido\\
\\
erar todo o\\
\\
Outros métodos: existem outros processos que tém sido pesquisados e que nao se en-\\
quadram em nenhuma das categorias acima, como é 0 caso da Tecupera¢ao microbi-\\
olégica ¢ da recuperagao utilizando ondas eletromagnéticas, A recuperacao micro-\\
biolégica é obtida a partir da utilizacgao de diferentes micro-organismos que, quando\\
adequadamente escolhidos ¢ através dos seus Processos bioldgicos no interior do re-\\
servatorio, produzem uma série de substancias que causam os mais diversos efeitos\\
¢ que podem aumentar a recuperacio de petréleo. A recuperacao através de ondas\\
eletromagnéticas é um processo de aquecimento do reservatério por meio de ondas\\
\\
eletromagnéticas ocasionadas pela aplicagio de uma diferenga de potencial entre os\\
pocos do campo.\\
\\
6.5 Estimativas de reservas\\
\\
As estimativas de reservas desempenham um papel fundamental para as empresas de ex-\\
ploracao ¢ produgio de dleo e gis. \\& com base nelas que o planeamento estratégico e\\
tomada de decisées das empresas quanto a implantacao de projetos exploratérios e de\\
Produgao so realizados. Os investimentos necessarios para a implantacio do projeto, as-\\
sim como 05 custos para manté-lo em operacio, devem ser pagos com a receita obtida\\
om a comercializacao dos fluidos a serem produzidos.\\
\\
Para fazer as estimativas, utilizam-se diversas formas de calculo, que serao apresenta-\\
das mais adiante neste capitulo. As estimativas podem ser feitas por meio de duas abor-\\
dagens, a deterministica ¢/ou probabilistica. Atualmente, as companhias tém utilizado a\\
\\
em deterministica, entretanto ha uma tendéncia para o uso da abordagem proba-\\
no futuro. ; al\\
\\
Aestrutura atual de pregos do petréleo, o capital de investimento em valor iat cial\\
€ Oambiente tributario determinarao quanto petréleo pode ser gach cupe\\
\\
. Outros fatores que influenciam as reservas sio localizagio ¢ be hic el\\
tan, 8 produsio € 0 mecanismo de produgao de reservatorio, Antes de (rattt alguns\\
™ativas de reservas propriamente dita, faz-se necessaria uma fwved vive\\
‘ermos intimamente ligados A estimativa de modo a facilitar 0 entene :\\
\\
no reservatorio na época da\\
ado gasoso dase\\
\\
* Volume ori al (N): quantidade de fluido existente ;\\
sua isuoee oe ‘doo acumulagio de hidrocarbonetos no est\\
\\
Digitalizada com CamScanner\\

\\section*{Pagina 181}

162 RESERVATORIOS\\
le gas para a mistura de hidrocarbonetos no estado\\
de volume original ¢ we gintl ‘te 61e0;\\
\\
o nome\\
nome de volume\\
\\
liquido da-se 01\\
\\
* Volume recuperivel (VR): volu\\
acumulagio de petréles . meuido se parle\\
rio, faz-se uma estimativa de quan| wade\\
\\
i -se O\\
A esse volume estimado de fluido da-se\\
de fluido que ja foi produzida de um reser.\\
\\
umulada vai crescendo gradati.\\
\\
u gas que se espera produzir de um,\\
casiao da descoberta de um reservaty,\\
produzir ou recuperar do mesmo\\
lume recuperavel;\\
\\
umulada (Np):\\
ma determinada é\\
eo rese\\
\\
A produgao ac\\
\\
= Produgao ac!\\
4 em produgao;\\
\\
yatorio até u\\
yamente com o tempo S|\\
i F jente\\
* Fator de recuperagao (FR): quocien| :\\
nal, ou seja, fragao do volume origin\\
\\
poca.\\
rvatorio esti\\
\\
re o volume recuperavel e o volume origi.\\
\\
ent ‘i\\
luzir de um reservatério,\\
\\
| que se espera prod'\\
\\
V;\\
FR(\\%) = ml x 100 (6.34)\\
\\
(fr): quociente entre a produgao acumulada e o volume original,\\
fe):\\
\\
« Fracdo recuperada \\} :\\
ae y foi produzida até um determinado instante.\\
\\
ou seja, fragao do fluido original que\\
\\
N,\\
f(\\%) = 7 x 100 (635)\\
\\
6.5.1 Tipos de reservas\\
\\
como recursos descobertos de dleo e gas comercialmente recu-\\
\\
i As reservas so definidas c\\
mo ainda nenhum\\
\\
peraveis a partir de uma determinada data. Na época da descoberta, co!\\
fluido foi produzido, a reserva é numericamente igual ao volume recuperdvel. As reservas\\
podem ser divididas em trés categorias conforme o grau de incerteza: reservas provadas,\\
provaveis e possiveis.\\
ne bd porebaet a) sao aquelas que se estimam com elevado grau de cert\\
Slisicore da engsain - seers descobertos e avaliados, com base dos dados ge-\\
pasate ent <i : nsiderando as condi¢g6es econdémicas existentes, OS métodos\\
2000), A SPE isa @ foqulaecpios da legislacao petrolifera e tributaria (AN,\\
seja igual ou maior do ay bere de pelo menos 90\\% (P90) que a quantidade rel\\
As reservas provaveis, por outro ae ae Eevenneaae ft ene 16-\\
gicos e de engenharia tenham indicad » S20 aquelas em que as anilises dos dados °°\\
reserva em comparacao com a ey, hé uma incerteza maior na recuperagio des\\
dagem probabilistica seja empregada a de reservas provadas (ANP, 2000). Caso"\\
(P50) de que o volume real seja igual . eve-se ter uma probabilidade de pelo menos 50\\%\\
—— € provaveis, igual ou maior do que a soma das estimativas de reservas\\
\\_ “8 Teservas possiveis 5 ;\\
tém demonstrado uma Me ga Cujas andlises dos dados geolégicos ¢ de engenh\\
Provaveis (ANP, 2000). Quando a oe) €m comparaciio com as estimativas de e5°"\\
‘ordagem probabilistica for empregada, devera havel\\
\\
eza recu-\\
\\
ania\\
yas\\
\\
d\\
\\
Digitalizada com CamScanner\\

\\section*{Pagina 182}

ESTIMATIVAS OF Rese wn\\
5 163\\
robabilidade de pelo menos 10\\% (P10) de «\\
e serio jguais ou maiores do que a soma das\\
weld» ri P\\
aa gveis mais possiveis.\\
i\\
\\
Ie as quantidades ar\\
\\
ahh. val\\
estimativas de reseryay nn eure\\
\\
Vas provadas Mais\\
\\
6.5.2: 1 Tipos de reservas quanto ao desenvolvimento\\
\\
asreservas podem ser classificadas quanto ao estagio de desenvoly,\\
dase nio desenvolvidas. As reservas desenvolvidas sio aquelas q\\
dasde poses existentes € ‘quando Os equipamentos necessarios pa\\
instalados. As reservas no desenvolvidas sao aquelas que pode\\
\\
jogos nOvOS em Areas nao perfuradas, re-\\
\\
imento em desenvolyi-\\
ue podem ser recupera-\\
ima Producio j\\& estejam\\
mM vir a ser r\\
\\
cl ; entrada ou recompletacao de Pi shee\\
ou que necessitem a instalacdo de equipamentos para a produgio e transporte previsor\\
ors projet0s de recuperacao convencionais ou melhorados, Na Figura 6.28 esta ther 4\\
as de reservas segundo Securities and Exchange Commission/ Comissao de Va-\\
lores Mobilidrios (SEC).\\
\\
Nio Desenvolvida\\
\\
Desenvolvida\\
\\
Figura 6.28: Tipos de reservas de petrdleo segundo a SEC.\\
\\
6.5.2 Métodos de calculo da estimativa de reservas\\
\\
Para estimar a reserva de dleo e gas, tem-se diversos métodos disponiveis, que variam con-\\
forme as circunstancias. Dentre os métodos existentes, serao enfatizados quatro métodos\\
nesta secao: por analogia, anilise de risco, yolumétrico e desempenho do reservatério.\\
\\
6.5.2.1 Analogia\\
\\
Este é um tipo de procedimento utilizado em uma época que precede 4 perfuracao do\\
Primeiro poco a penetrar na jazida, ou seja, do poco descobridor. Nessa época, as infor-\\
Magées a respeito do reservatério séo praticamente inexistentes. Tem-se uma série de\\
evidéncias, entretanto, ainda nao se tem a comprova¢ao da existéncia de uma acumula-\\
So de petréleo na regido que esta sendo pesquisada. As estimativas sio feitas a partir de\\
dados e resultados de reservatorios localizados nas proximidades, os quais se acredita te-\\
caracteristicas semelhantes as do reservatorio que est sendo estudado. E evidente\\
Que esse tipo de estimativa esta sujeito a erros, uma vez que 0 estudo no se baseia em\\
Teais do reservatério.\\
\\
Digitalizada com CamScanner\\

\\section*{Pagina 183}

164 RESERVATORIOS\\
\\
6.5.2.2  Andlise de risco\\
\\
Como o método anterior, a anilise de risco também ¢ um processo utili\\
\\
perfuracao do poco de descoberta. Da mesma forma, a estimativa é feita oc\\
\\
tados de reservatérios cujas caracteristicas sio semelhantes ds do ferences\\
\\
€ que se localizam nas suas proximidades. A diferenca entre os dois OBL i\\
\\
fato de que na anilise de risco existe uma certa sofisticagio no tratamento est viniee \\textasciitilde{}\\
\\
dados, e os resultados sio apresentados, nio como um valor tinico, mas Siiip tse\\
laixa\\
\\
de resultados possiveis.\\
\\
ANtes\\
ACtir de esy\\}\\
\\
Mestugy\\
\\
6.5.2.3 Método volumétrico\\
\\
O hidrocarboneto no local é uma quantidade fixa que se desenvolveu ao longo do tempo\\
geologico. Pode ser estimado usando métodos volumétricos ou de balango de materiais\\
O cdlculo volumétrico de hidrocarboneto em vigor requer conhecer a extensio real do\\
reservatério, a sua espessura média e porosidade, a saturagio de hidrocarbonetos ¢ o\\
fator volume de formacao de hidrocarbonetos (isto é, 0 volume que uma unidade de\\
volume de hidrocarboneto a pressao € temperatura da superficie ocupa nas condigées do\\
reservatorio). E um método estatico que nao depende do comportamento dinamico do\\
reservatério, ou seja, da resposta da pressao a producao. As equacées para 0 caleulo do\\
volume original de hidrocarbonetos (para reservatérios bifasicos de dleo / agua ¢ gis /\\
Agua, respectivamente) sao:\\
\\
0\\
gac ub Swi) (637)\\
Bui\\
\\
onde C = 7758 €0 fator de com\\
\\
Para unidades de Engenharia ou Unidades de Campos, .\\
é 0 fator de conversio\\
\\
versio de acre-pé para bbl de dleo, e para o gas C = 43560, que : 4\\
de acre-pé para ft?; A é a Area do reservatério (acres) dados do mapa; hé ey ie\\
espessura da zona produtora (ft) do log e/ou dados de amostra; q) é 4 porosid i a\\
cimal) do log e/ou dados de amostra; Syj € a saturacao inicial de 4gua (decimal ¢ ds\\
e/ou dados de amostra; Bo; ¢ 0 fator volume de formacio para dleo nas cond in\\
(reservatério bbl / Stb); G éo Volume original de gas medido em condicoe pad ‘ GNe |\\
Bg; € 0 fator de volume de formacio para gas nas condig6es iniciais (resft™ pi |\\
o Volume original de 6leo medido em condicgées-padrao, Stb. .\\
\\
Em unidades do sistema internacional (SI), N (Sm?) é o volume original\\
dido em condicées-padrao; V, (m?) é o volume total da rocha; ? ¢\\
porosidade; Bg; (m?/Sm?) € 0 fator volume de formagio de dleo\\
reservatério; Bg; (m?/Sm?) é 0 fator volume de formagio de gas nas\\
reservat6rio; G (Sm?) é 0 volume original de gis medido em condi¢\\
constante que é igual a um.\\
\\
de dleo me\\
\\
ig da\\
io média ‘\\
\\
éa ist\\
\\
condi¢! cé ’\\
\\
- a3 ¢\\
des: padre\\
\\
Digitalizada com CamScanner\\

\\section*{Pagina 184}

ESTIMATIVAS py\\
E KES\\
SERVES, 165\\
1 reservatrio de dleo, pode-se calcular também © volum:\\
Nun Fes\\
b\\
\\
Ivido no \\$leo. Esse volume é determinado Pelo produto d\\
disso!\\
\\
ic de pas que se encontra\\
yido em condig6es-padrao, pela razao de Solubilidade inici\\
medid\\
\\
‘9 volume Original de bleo,\\
al, ou Seja, NRg. Entio,\\
8=6 VP (1 \\textasciitilde{} Sui) Ry\\
= CEN = ow) Rs\\
. (6.38)\\
G, 0 volume original de gis em solu\\
aa 29 ilustra como o volume original de 6)\\
1a 6.\\
como deve ser expresso.\\
\\
\\$40 medido em condi¢des-padrao, A Fi.\\
leo é medido no Teservatério e a Maneira\\
\\
VAX 5,9)\\
Figura 6.29: Volume original — reservatério de dleo.\\
Com base nas equagGes 6.36 € 6.37, pode-se obter o volume recy;\\
\\
de recuperacao, onde Rr € 0 fator de recuperacao e\\
de dleo e gas, respectivamente.\\
\\
peravel incluindo o fator\\
Vr e Gr sao 0 volume recuperavel\\
\\
1-Sy;)xF\\
Vg =N x Fp = CMP (= Swi) x Fe\\
\\
‘oF\\
Vr (1— Swi) x Fr\\
Bgi\\
\\
(6.39)\\
\\
Gr=NxFpr=C (6.40)\\
\\
6.5.2.4 Desempenho do reservatério\\
\\
Neste método sio utilizados modelos em que a previsao do comportamento futuro (ou\\
desempenho futuro) do reservatério se baseia em seu comportamento passado. Para\\
tanto, énecessério que o reservatorio ja tenha um histérico de producao. Em alguns casos\\
também sao nécessdrias informagGes sobre o mecanismo de produgao do sear\\
\\
A anilise do declinio de produgio, a utilizagéo da equacao de balanco de —\\
Paraa previsio de comportamento e a simulacdo numérica de scectistenica sem —\\_\\
ese inserem no 8rupo denominado desempenho do reservat6rio. A shar ——\\
SH Outro método depende de fatores tais como a quantidade e ° apo dee Se -\\
\\
Muido disponiveis, a existéncia de recursos de informatica (“software” e “hu )\\
\\
Digitalizada com CamScanner\\

\\section*{Pagina 185}

166 — RESERVATORIOS\\
Andlise do declinio da produ¢a°\\
5 vaz6es de produgao do reserva\\
7\\
\\
mento c\\
ee reservat6rio, a pressio que.”\\
q I ei,\\
\\
-se no compor d :\\
\\
do 0 material retirado de\\
rio tende a diminuir, © que por sua vez acarrei .\\
ervat Normalmente, utiliza uma taxa dy is\\
Prog.\\
\\
= dos poses: rp 5\\
dugao do ea analise do declinio éa taxa didria de Prods\\
a utiliza apenas © histérico de producig dorey,\\
dades fisicas do meio em que o Feservatérigy\\
\\
Essa técnica baseia\\
longo do tempo. Quan\\
inicialmente dentro do res\\
declinio nas vazoes de pro\\
gio diaria com 0 tempo- Qutra forma\\
pela produgao acumulativa. Ess\\# récnic\\
vatério, nao sé importando com prope!\\
\\
encontra.\\
\\
Equacao de palango de materials\\
\\
ue utiliza a relagao entre o balanco de\\
dudes de pressao no seu interior. A ideia cen.\\
da contidas nos reservat6rios com a massace\\
deve ser igual 4 massa de material originalmente existente\\
no meio poroso. Esse método é usado para estimar a quantidade de hidrocarboneto ix\\
situ, Parte-se da hipdtese que © relacionamento pressao-volume de um reservatério per\\
é possivel equacionar a expansio\\
\\
manece constante quando ele produz. Dessa forma, €\\
dos fluidos do reservat6rio com © esvaziamento dos fluidos do mesmo menos a injecio\\
\\
de qualquer fluido. Como essa técnica fornece a relacao entre a producao acumulads\\
versus a queda de pressao, faz-se necessario o uso de outras equagdes que relacionem 2\\
\\
produgées acumuladas com vazées de producdo € tempos.\\
\\
de materiais € uma técnica q\\
\\
reservatorio com as Fe!\\
oma das massas ain\\
\\
A técnica de balango\\
massa dos fluidos do\\
tral dessa técnica é que aS!\\
material ja retirado do mesmo\\
\\
Simulacdo Matematica de reservatérios\\
eae oe ‘ eens, os meio de programas de simulagées numeéricas e el\\
oC RE ae, = rio, © método de simulagdo matematica utiliza Pp\\
do volume das reservas, sao eas ae 2 balance de materiais. Para ® een\\
da rocha, do fluido etc, de f informagies geolégicas e geofisicas, carat\\
histérico de producio, ea adie que se reproduza, com um certo grau\\
culailo tepresente sattatate Ojasaliny possivel prever o futuro, desde que ess\\
inilteicetntanicetonges opin o real, Enquanto o método da equas de bes\\
para determinar o volume ovcien oireagnvatorio nade representado por um nie\\
. \\_ i. fomnece o resultado em fu , cujas caracteristicas © es 4 axod?\\
ulacdo matemitica depend, ngo do tempo. O grau de preciso 1,\\
geologia, do histérico ‘pende de dados (como caracteristicas da rocha, do on,\\
sario que\\
\\
dai ', etc.) «\\
dos sejam de qualidade e ro saga satisfatéria, além disso, é neces\\
\\
Digitalizada com CamScanner\\

\\section*{Pagina 186}

ESTIMATIVAS DE RESERVAS, 167\\
\\
soot dagem determinlstica e probabillstica\\
\\
ado de estimagio é chamado de deterministico se a melhor es\\
eq. com base em dados de engenharia, geoldgicos e econdmicos c\\
feita vnerodo de estimacao é chamado de probabilistico quando dados Beoldgicos, d\\
ue sid e¢ econémicos conhecidos sao usados para gerar uma série de eatidiatvs :\\
enge! robabilidades associadas. Se a abordagem probabilistica for utilizada deve-se oni\\
suas PI probabilidade de 90\\% para as reservas provadas, 50\\% para as reservas Srovados\\
\\
gar uma : 0\\%, f\\
ait reservas proviveis, e 10\\% para as reservas provadas mais reservas provaveis mais\\
\\
reservas possiveis.\\
\\
timativa de reserva é\\
‘onhecidos. Ao Passo\\
\\
oméet\\
\\
Conclusao\\
\\
Este capitulo contribuiu para aquisicao de conhecimentos basicos das atividades de um\\
engenheiro de reservatorio e compreender as propriedades de rocha-reservatério de hi-\\
drocarboneto € as diferentes fases de recuperacao de hidrocarbonetos. Os formandos\\
nessa area estdo preparados para descrever as principais atividades de um engenheiro de\\
reservat6rios, conhecer algumas propriedades basicas das rochas e dos fluidos de reser-\\
vatrios, € calcular seus parametros para estimar o volume original e as reservas de hi-\\
drocarbonetos. Também sao capazes de explicar a importancia do envelope de fase para\\
definir o tipo de reservatério e identificar e entender os mecanismos de recuperacao de\\
\\
hidrocarbonetos.\\
\\
Digitalizada com CamScanner\\

\\section*{Pagina 187}

CAPITULO 7\\
\\
pRODUCAO E OPERACOES\\
SUBMARINAS\\
SS\\
\\
A producio é a fase para extrair o petréleo e gas natural de uma Teserva com 0 intuito\\
de maximizar sua vida util. Neste contexto, existe um conjunto de instalaces destinadas\\
a promover a coleta, a separacao, 0 tratamento, o armazenamento e o escoamento dos\\
fluidos produzidos e movimentados num campo de petréleo ou gas natural, conhecidos\\
como sistema de producao. Este sistema é composto por reservatério, poco, linha de\\
escoamento ou coluna de producio, separadores, bombas, dutos etc.\\
\\
Os fluidos que estao armazenados no reservatério precisam ser elevados e escoados\\
\\
até as instalagdes na superficie (até o separador). De um modo geral, este processo é\\
dividido em trés etapas:\\
\\
" Recuperagio: refere-se ao fluxo de fluidos do reservatério até o fundo do pogo;\\
* Elevacao: refere-se ao fluxo de fluidos do fundo do poco até ao topo do poco;\\
\\
* Escoamento: refere-se ao fluxo de fluidos do topo do poco passando pelo regulador\\
do fluxo e até ao separador.\\
\\
Entre muitos técnicos intervenientes neste processo, destacamos o engenheiro de produ-\\
Gio, cujas principais atividades, incluindo engenheiros de subsuperficie, gerem a interface\\
entre 0 reservatério e 0 poco, incluindo perfuracées, controle de areia, controle de fluxo\\
de fundo do poco e equipamento de monitoramento de fundo do poco; avaliar métodos\\
de elevacio artificial; e também selecionar o equipamento de superficie que separa os\\
fluidos produzidos (éleo, gas e 4gua). De forma resumida, suas atividades envolvem:\\
\\
* Extrair petréleoe gas com as mais diversas técnicas de recuperaco (primaria, secun-\\
daria e avancada);\\
\\
* Manter niveis de producio da reserva otimizados (workover);\\
\\
* Encerrar as atividades de producio (ex.: desativacao de infraestrutura ¢ descarte de\\
tesiduos toxicos).\\
\\
vida de um pogo quando o petrdleo\\
de petréleo usadas para perfurar e\\
\\
169\\
\\
Otstigio de produgio é 0 estagio mais importante da\\
© 0 gis so produzidos. A essa altura, as plataformas\\
\\
Inddstria do petvileo para ndo-engenheiros, 1* edigd.\\
By Geraldo André ieee ees Cogn © 2023 Editora Dois03\\
\\
Digitalizada com CamScanner\\

\\section*{Pagina 188}

RINAS\\
170 PRODUGAO E OPERAGOES SUBMA'\\
\\
so retiradas do pos © o topo normalmente é equipado cg\\
ve ominadas arvore de natal ou arvores de producio, Bssas ys Ung\\
s den\\
\\
Jam fluxos € permitem acesso ao furo de POGO NO caso iN ny\\
ay\\
\\
Ari i valvula de sai P\\
Ihos de conclusao serem necessarios. A partir da dvd aida da arVore de naty\\
OS iaeall |\\
ne ode ser conectado a uma rede de distribuicao de aa e tanques para forme\\
° Le a » vefinarias, estagOes de compressao de gs natural ou terminais de expona,\\
o produto Ss\\
\\
de petrdleo. . ‘ ;\\
fo uanto a pressao no reservatorio permanecer altao stificiente, aarvore de Prods\\
iq! ara produzir 0 pogo. Se a pressao se esgota e é Considerad, on\\
\\
é tudo o que é necessario p: 6 ap 2\\
agile vyiavel, um método de elevacao artificial pode ser aplicado para teltes\\
\\
nomi A ck\\
\\
rvatorio.\\
\\
0 petrdleo remanescente no rese\\
\\
completar 0 pogo\\
colecao de valvula\\
regulam pressoes, contro!\\
\\
7.1 Producado em Angola\\
\\
A produgao em Angola é suportada por 14 Blocos, sendo 3 onshore (em terra) e 11 offshore\\
(mar), Os blocos onshore séo Cabinda Sul, FS (Fina Sonangol), FST (Fina, Sonangol, Te.\\
xaco), €, ao passo que os blocos offshore so Bloco 0, 2, 3, 4, 14, 15, 15/06, 17, 18, 31 ¢32\\
Anomenclatura usada para cada campo esté ilustrada na Tabela 7.1. A produgao atual do\\
pais esta ilustrada na Tabela 7.2.\\
\\
Na Se¢io 7.7, abordamos de forma resumida a OPEP. O petrdleo produzido em An-\\
\\
gola é considerado um dos melhores do mundo, como discutido no Capitulo 8, especial-\\
mente na Tabela 8.1.\\
\\
7.2 Sistema de Producao\\
O conjunto de instalacées destinadas a\\
\\
zenagem e escoamento de flufdo:\\
Jeo ou gas natural é conhecido\\
\\
Promover a coleta, separacio, tratamento, arme\\
's produzidos e movimentados em um campo de ane\\
produca como sistema de producao (Figura 7.1). O sistema ©\\
io Ec Hy 7 ¢ igura 7.1).\\
sbacdats wen, end Partes, incluindo 0 casco, reservatério, poo linha de\\
le produca I os,\\
entre outros, Producao, separadores, bombas, tubagens (tubings) ou dut\\
\\
Digitalizada com CamScanner\\

\\section*{Pagina 189}

SISTEMA DE PRODUCAD =—«-17.\\
\\
LA nomenclatura ou origem de nomes aplicados nos campos petroliferos em\\
\\
er\\
\\
Boos pos petroliferos\\
Nome de érvores tipicas de madeira 21 Nome de parques nacionais ¢ reas\\
\\
o de Cabinda de conservacio\\
Nome de Arvores de fruto 2 Nome de morros ou montanhas de\\
\\
' Angola\\
Nome de peixes ¢ outros animais ma- 23 Nome de cores\\
\\
, rinhos\\
\\
7 Nome de antilopes 2 Nome de dangas tipicas de Angola\\
\\
‘ Nome de legumes tipicos de Angola 25 Nome de animais felinos\\
\\
; Nome de arbustos em linguas nacio- 26 Nome de dancas tradicionais\\
nais\\
\\
P Nomes de aves 3 Nome de planetas ¢ lua\\
\\
? Nome de conchas 32 Nome de especiarias\\
\\
' Nome de pedras semipreciosas 33 Nome de pratos tipicos de Angola\\
\\
9 Nome de frutos 34 Nome de miscaras ou esculturas tra-\\
\\
dicionais de Angola\\
\\
10 Nome de mulher em linguas nacio- 35/11 Nomes de aves em linguas nacionais\\
nais\\
\\
“ Nome de cidades ¢ localidades de An- 36/11 Namero em linguas nacionais kim-\\
gola bundo\\
\\
ry Nome de instrumentos tradicionais 37/11 Numero em linguas nacionais um-\\
de Angola bundo\\
\\
6 Nome de rios de Angola 38/11 Nome de répteis\\
\\
7 Nome de flores 39/11 Nome de lagos\\
\\
“ Nome de minerais-metais 40/11 Nome de linguas nacionais ¢ dialetos\\
\\
LU Nome de insetos Cabinda Sul Nome de sementes\\
\\
» Nome de mamiferos marinhos Cabinda Norte Nome de reis e soberanos de Angola\\
\\
Tabela 7.2: Produgao de Angola de 2017 a 2020. Fonte: Petrangola no programa “Directo\\
a0 Ponto” da TV Zimbo, Janeiro 2021.\\
\\
heed 217 2018 2019 2020\\
Mé Produgo Quota OPEP Producio Quota OPEP Producio Quota OPEP Producio Quota OPEP\\
= 1a 1,673 1,569 1,673 1,466 1,482 1,389 1,482\\
bed ted 1,673 1,503 1,673 1,418 1,482 1,379 1482\\
- 1621673 1,496 1,673 1,373 1,482 1,375 1482\\
vel 1449 1,673 1,498 1,673 1,416 1,482 1351 1.482\\
ie im 1673 1,485 1,673 1,425 1,482 1,180\\
Mu pps 1,673 1,447 1,673 1417 1,482 1,180\\
i‘ rod 1,673 1,455 1,673 1,259 1,482 1,180\\
fe 1673 1481 1,673 1,394 1,482 1,265 1,249\\
Oa = umn 1489 1,673 1,393 1,482 1215 1249\\
Now tbe 1673 1,457 1,673 1,390 1,482 1194 1 me\\
ter ora 1,416 1,673 1,272 1,482 1218 1249\\
1,482 1,198 ase\\
\\
A45 1673 1,366\\
\\
varono\\
\\
Bate sig\\
“ts inmalagten al para elevar e escoar 0s fluidos armazenados no ane\\
na superficie, incluindo o separador. De modo geral, 0 fluxo de\\
\\
produ\\
\\
Digitalizada com CamScanner\\

\\section*{Pagina 190}

172 — prooucho OPERAGO\\
\\
cio é dividido em trés etapas,\\
¢ Escoamento.\\
\\
Figura 7.2: Ilustra\\
\\
7.2.1 Recuperacao\\
\\
A recuperagao se refere a\\
tas vezes, essa atividade est:\\
\\
reservatorio, Os métodos de recuperagao po\\
\\
* Recuperacao primaria: En\\
\\
ES\\
\\
confor\\
\\
BMARINAS\\
\\
me jlustrado na Figura 7.2: Recuperacio, Ble,\\
iy\\
\\
Gis\\
\\
=\\
\\
do das trés etapas do sistema de produgio.\\
\\
0 fluxo dos fluidos do reservatério até o fundo do pogo. Mui-\\
4 intimamente relacionada com o trabalho do engenheiro de\\
\\
dem ser classificados da seguinte forma: |\\
\\
volve o fluxo natural dos fluidos (incluindo gas em solu\\
\\
Gio, influxo de agua, capa de gas, entre outros) do reservatério;\\
\\
* Recuperagao secundaria: |\\
tério;\\
\\
inclui a injegdo de 4gua para manter a pressao do reserva:\\
\\
* Recuy ies\\
oa pera G40 terciéria EOR: Envolve a injecao de fluidos nao produzidos no reset”\\
rio, como métodos térmicos, quimicos, gasosos, entre outros.\\
\\
7.2.2 Elevacaéo\\
\\
Aclevacio se refe 2\\
refere ao fluxo dos fluidos do fundo do poco até a superficie. Esse process?\\
\\
Pode ser natural ou artificial,\\
\\
7.2.2.1 Elevacao natural\\
\\
A facia\\
PresSI0 exis Woe © Petrbleo\\
do reservatérig Mn TY atono,\\
A cleva rio fluem faci\\
Teservatérj\\
m\\
\\
Caso, og\\
\\
‘mente par; fs «\\
\\& cnergia lap também ps peaingl © que chamamos de elevac\\#°\\
\\
ale: ‘| ,\\
or superficie esta diretamente relacionad cae\\
lo a pressio é naturalmente suficient, \\textasciitilde{} ‘|\\
\\
nal\\
, ut\\
iO € aco: ey Produgio por pogos surgentess wi\\
naturale no infcio da vida produtiva do re" er\\
ente devido a energia do reservaton™\\
\\
Digitalizada com CamScanner\\

\\section*{Pagina 191}

SISTEMA DE PRODUCAG 173\\
\\
sem desse sistema de elevagio é que ele apresenta\\
yantag mentos mais simples e possui custos Oper:\\
\\
equip’ =\\
puiliza atari de pogo com elevagao natural.\\
ras\\
\\
menos Problemas Operacionais\\
acionais menores, A Figura 7 3\\
\\
gust\\
\\
Arvore de natai\\
\\
Figura 7.3: Fluidos produzidos com um sistema de elevacao natural.\\
\\
Noentanto, ao longo da vida produtiva, os POGos surgentes podem enfrentar um declinio\\
na pressao do reservatdrio, dificultando a producao econémica. Quando isso acontece,\\
seja no inicio ou durante a vida produtiva, é necessdrio recorrer a métodos de elevacio\\
artificial para manter a producao.\\
\\
7.2.2.2 Elevacao artificial\\
\\
O objetivo da elevacdo artificial em qualquer sistema é adicionar energia aos fluidos pro-\\
duzidos para acelerar ou viabilizar a producio. Isso é alcangado por meio de equipamen-\\
tos que aumentam o diferencial de pressao sobre o reservatério, resultando em um au-\\
mento na vazao. Os métodos de elevaco artificial mais comuns na industria do petréleo\\
\\
incluem:\\
\\
* Gas Lift Continuo/ Continuous Gas Lift (GLC) e Gas Lift Intermitente/ Intermittent Gas\\
Lift (GLI);\\
. BCs;\\
\\
* Bombeamento Mecanico (BM);\\
* BCP.\\
\\
Aescolha do método de elevacio artificial depende de varios fatores, como ° nuimero “\\
Pocos, diametro aa revestimento, producio de areia, razio gas-liquido, vazdo, Lies\\
dade do Teservatério, viscosidade dos fluidos, mecanismo de produ¢io a ae\\
nibilidade de energia, acesso aos pocos, distancia dos pocos As estagoes ou Le pe\\
S, quipamentos disponiveis, pessoal treinado, investimento, custos ¢ segurangs\\
ta i sp 1, pe!\\
\\
Digitalizada com CamScanner\\

\\section*{Pagina 192}

RINAS\\
174 propucdo E opERACOES suBMA!\\
\\
Gas lift ”\\
i sto relativa: ‘\\
\\
étodo bastante utilizado devido ao i i whi evamanite baixo, Mesn,\\
cas lift 6 um Met go artificial utiliza a energia conti »\\
Ogaslift método de eleva¢ Bia contida em ;\\
\\
em pocos profundos. Este idos (6leo e/ou Agua) até a superficie. A injecdo de gig\\
comprimido para aaa eda densidade deste fluido, fazendo com que a razig "\\
fluido promove tm e temente, permitindo que 0 fluido suba com ; La\\
liquido seja mais rg ‘sam dois tipos de gas lift: 0 continuo ¢ 0 intermitente. t\\
facilidade a superficie. Ex" se na injegao continua de gas a alta pressdo na coluna de pro,\\
\\
ij i eia\\textasciitilde{}\\
O gas lift continuo bas de injecdo até a superficie. Por outro lado\\
\\
a i i de o ponto\\
ducdo, gaseificando 0 fluido desde 0 PO! si\\
° = i intermitente baseia-se 14 injegdo de gas a alta pressao na base das golfadas na\\
\\
coluna de produgio, através de valvulas operadoras, em tempos bem definidos, como\\
objetivo de acumular uma certa quantidade de fluido. Bsse processo requer elevadas va.\\
zbes periddicas de gas para imprimir grandes velocidades ascendentes a golfada.\\
\\
O sistema de gas lift € composto por fonte de gas a alta pressao (compressores); contro.\\
ladores de injecao de gasna superficie (choke ou motor valve); controlador de injecao de gis\\
de subsuperficie (valvulas de gas lift); equipamento para separacao e armazenamento dos\\
fluidos produzidos (separadores, tanques etc.). Os tipos de instalagao de gas lift podem\\
ser abertos, semifechados, ou fechados, sendo a escolha dependente das caracteristicas do\\
\\
pogo.\\
\\
Fgura 7.4 Posos equipados\\
\\
com 0 sistema de gas lift.\\
\\
Digitalizada com CamScanner\\

\\section*{Pagina 193}

SISTEMA DE PRODUCAD, 175\\
\\
gombel0 Mecdnico com Hastes (BMH)\\
\\
BMH, jlustrado na Figura 7.5, é purl método de elevacio artificial usado em 0\\
aos pressdes estaticas, onde os fluidos so elevados Para a superficie Wtve as com\\
“atuna de hastes, cujos movimentos sao transmitidos por uma bomba de movi ne\\
rotativo, geralmente acionada por um motor elétrico. vee\\
‘A bomba é projetada para ser inserida na tubulacio do\\
objetivo coletar fluidos abaixo dela e eleva-los 4 superficie.\\
rtantes Sa0 as valvulas (méveis e fixas) e 0 pistao. A bom\\
bombeamento na superficie por uma série de hastes de suc¢do, que séo movimentadas\\
para cima e para baixo na tubulacao, ativando a bomba na parte inferior. Na superficie,\\
um grande dispositivo mecanico chamado unidade de bombeamento de feixe est4 defies:\\
tado.\\
\\
Poco € tem como principal\\
Seus componentes mais im-\\
ba € conectada A unidade de\\
\\
one\\
erpeenie\\
\\
Figura 7.5: Elevacao artificial através do bombeio mecanico com hastes.\\
\\
Dependendo do tamanho da bomba, esse sistema geralmente produz de 5 a 40 litros de\\
liquido a cada curso. Frequentemente, esse liquido consiste em uma emulsio de petréleo\\
¢ gua. O tamanho da bomba também é determinado pela profundidade e peso do dleo\\
aser removido, sendo que extragGes mais profundas exigem mais energia para mover os\\
comprimentos mais pesados das hastes de suc¢ao. ;\\
As pei bombeio de hastes incluem a alta eficiéncia do sistema, a economia\\
€m reparos e manuten¢do, a flexibilidade para ajustar producao por meio do compri-\\
mento € velocidade do Bea, ¢ um alto valor residual para as unidades de wpe:\\
equip : i limitagdo a volumes de\\
desvantagens incluem a ¢ nes d\\
city So pee - Laan) inferiores a 1.000 barris por dia,ea aplicagio\\
ipalmente em terra ee het vez que requer uma grande unidade de superficie.\\
\\
Digitalizada com CamScanner\\

\\section*{Pagina 194}

sUBMARINAS\\
\\
proDucho E OPERAGOES\\
\\
176\\
trifugo SubMerSO/ Electrical Submerspyp,\\
"in\\
\\
‘en\\
Osistema Bombeamento C\\
(BCS) | |\\
rum motor elétrico que recebe o suprimento g\\
2 ee ao o eixo da bomba. Além disso, inclui uma\\
oer seria’ contaminacao do motor pelo fluido do pogo, uma admis,\\
e evita bomba (que pode ser um intake ou um separador de \\_ :\\
na ogo), um cabo elétrico, um transformadore um ti\\
O sistema de BCS esta ilustrado na Figura 7,6, 5\\
\\
; \\_\\
\\
O sistema de BCS € com| r=\\
através d .\\
selo (ou protetor) qu\\
ja qual o fluido entra ;\\
ede das caracteristicas do p'\\
de comando, e a propria bomba.\\
\\
TRANSFORMER\\
page, CLTRCO | IXSLERO ELECTRO\\
» ame\\
\\
c480 CLETRCO] CABLE LE TACO\\
\\
sEson ba RCO | SBsOR OE FORO\\
\\
Figura 7.6: Elevac4o artificial com o sistema de BCS.\\
\\
Mm dest sgh ar Submersible Pump/ Bomba Elétrica Submersa (ESP) incuemum\\
a = re de profundidade, com alta eficiéncia acima de 1.000 barris pot\\
phy i necessi ee panes Tequisitos minimos de equipamento de superfide\\
\\
resisténcia tes corrosivos no fundo do pogo. E adequado para uso em\\
\\
balxa capacidade de bombear ain curvas. A principal desvantagem do ESP \\&\\
\\
de fluzo Por um motor\\
eon Primed ha superficie, geralmente elétrico. 4 A\\
de cavidade proprecae PIO, ha inada pela velocidade rotativa do cont”\\
Progressiy; +4 Pou i 5\\
de * embora haj teen bag Podem dar errado com as ear\\
elevasio artificial Pecialmente po energia devido ao uso da va"\\
Por meio de ave utiliza um wo Ss Profundos e desviados. Este ¢ um mero?\\
\\
YM conjuntg de ‘a de movimento rotativo na superficie, ve\\
\\
208 acessérios dos equipamentos de bombe\\
\\
|\\
\\
=—\\_\\
Digitalizada com CamScanner\\

\\section*{Pagina 195}

SISTEMA DE PRODUCAG. 177\\
\\
10 objetivo de elevar 0 fluido existente no fundo do POGo para a superficie. \\{3\\
con ba de deslocamento positivo que trabalha imersa no Poco de petrdleo rm 2 uma\\
peer de progressiva se destacam em Pocos tasos de baixa produtividade pal oa\\
prutos viscosos € podem lidar com quantidades significativas de sdlidos Produzidi\\
lustra o sistema de elevacao artificial com bombeio de cavidade Progre:\\
\\
6leos\\
los. A\\
\\
Figura 7.7 i Ssiva,\\
\\
Figura 7.7: Sistema de elevacao artificial com BCP.\\
\\
7.2.3, Escoamento do fluido\\
\\
Ocscoamento do fluido se refere ao fluxo de fluidos da superficie (topo do poco) passando\\
pelo regulador de fluxo e indo até o separador.\\
\\
7.2.3.1 Processamento primario de petréleo\\
\\
O processo de processamento primario de petréleo consiste na etapa inicial da refina-\\
\\$40 do petréleo bruto, com foco na produgio de hidrocarbonetos (éleo e gas) para fins\\
ccondmicos. Devido ao interesse exclusivo na producio de hidrocarbonetos (dleo e gas)\\
© 4 influéncia negativa de outras substancias em aspectos como transporte € seguranga\\
peracional, é necess4rio realizar um processamento primirio in loco, ou seja, na propria\\
i de producio, seja ela terrestre ou maritima. A Figura 7.8 ilustra o esquema do\\
Sistema de Processamento primario de petrdleo.\\
Este processo visa a separacao do dleo, do gas ¢ da agua, juntamente com as imp\\
24s em Suspensio, realizada por decantacio. Além disso, inclui o Kappan Sa con\\
dos hidrocarbonetos Para que possam ser transferidos para as ren\\
‘pat xorre 0 propriamente dito, e o tratamento da agua para reinjes v0\\
er sumentar apés a devida compressio e tatamen\\
‘ \\textasciitilde{} producto) na descarte. 0 git phe itilizado na\\
"*MosI0 de COZ ¢ H2S ¢ desidratagao), pode ser exportado como produto,\\
\\
Digitalizada com CamScanner\\

\\section*{Pagina 196}

UBMARINAS\\
UCAO E OPERAGOES 5!\\
178 PRODUC\\
\\
ete\\
7.8: Esquema do processamento primirio de petrdleo.\\
ira 7.8: ES\\
\\
Figu\\
\\
propria plataforma como gis combustivel (fuel gas), reinjetado no pogo ou usado como\\
\\
ift (elevacio artificial). . 245, . a\\
a ee ls prncessnBSe0 depende de critérios de viabilidade técnico-econémica. Pode\\
variar iy unidades de processamento mais simples, baseadas em decantagao e uso de va.\\
\\
sos separadores, a processos fisico-quimicos mais complexos, que incluem tratamento do\\
éleo, compressio do gis ¢ tratamento da 4gua para descarte ou feinjecdo NO poco, a fim\\
de facilitar a surgéncia do éleo. O processamento primario é necessario para promover\\
a retirada de impurezas em suspensio, tratar a 4gua para torné-la livre de impurezas ¢\\
facilitar o transporte para terminais e refinarias, reduzir problemas de corrosio e incrus-\\
taco, aumentar a vida titi] de equipamentos e catalisadores em processos de refino, e\\
reduzir os gastos com produtos quimicos usados para inibir processos corrosivos. Apesar\\
dos cuidados, nem sempre é possivel a separacdo e remogao completas da Agua e do gis.\\
\\
No processamento de gis, parte dele pode ser utilizado como combustivel na propria\\
unidade ou para elevacio artificial do éleo. Se a producao exceder muito o consumo,\\
Pode ser transferida para refinarias ou direcionada para queimadores. A parte liquefeita\\
sae ser adicionada ao dleo para transferéncia e aproveitamento posterior nas unidades\\
\\
proceso.\\
\\
Quanto a agua, considerada o mais indesejavel dos contaminantes, deve ser totalmente\\
removida quando chegar A refinaria, pelos motivos mencionados anteriormente. A trans\\
\\
Tabela 7.3: Composicio Upica de petréleo bruto.\\
\\
Elemento\\
\\
0 eae\\
c. Mea\\
\\
Enzof 83-87\\
\\
Nitrogé 0,06 - @\\
\\
Onigtnio O11. 1,7\\
\\
Metals O12\\
\\
SR A AT ia\\
\\
Digitalizada com CamScanner\\

\\section*{Pagina 197}

SISTEMA DE PRODUCAO 179\\
\\
a 7.4: Composi¢ao tipica do gas natural.\\
\\
Tabel\\
Campo de gas natural Gas natural liberado do\\
parkmetros dleo\\
sagtaio Tragos - 15\\% Tragos - 10\\%\\
odo de carton Tracos - 5\\% Tragos - 4\\%\\
ca sultco ‘Traces - 3\\% Tragos - 6\\%\\
eto Tragos - 5\\% Nio\\
ane 70 - 98\\% 45-92\\%\\
eat 1- 10\\% 4-21\\%\\
ingens Tracos - 5\\% 1- 15\\%\\
| nae Tracos - 2\\% 0,5 - 2\\%\\
| pentanos Tracos - 1\\% Tracos - 3\\%\\
| Hlexsnot Tragos - 0,5\\% ‘Tracos - 2\\%\\
‘Tragos - 1,5\\%\\
\\
Tragos - 0,5\\%\\
\\
Heptanos +\\
\\
Tratamento do dleo\\
\\
|\\
|\\
|\\
\\} o petréleo que sai do pogo é uma mistura de dleo, agua, gas e até areia. Ma-\\
| teriis estranhos, como agua € areia, devem ser separados do petréleo e do gas antes que\\
esso € conhecido como tratamento de éleo ou desidrata-\\
\\
possam ser yendidos. Esse proc\\
cio de dleo, na qual a 4gua é removida do éleo. A quantidade desse material estranho é\\
chamada de sedimento basico € agua (BS\\&W). Normalmente, 0 teor de BS\\&W deve ser\\
\\
| para venda. Os tipos basicos de equipamen-\\
\\
inferior a 1\\% antes que 0 dleo seja aceitave\\
arador, Free Water Knockout (nocaute de 4gua\\
\\
105 utilizados no tratamento de éleo sao sep\\
livre) e aquecedor-tratador.\\
\\
Separador: um separador é um grande vaso de pressao projetado para separar 0s fluidos\\
produzidos em seus componentes constituintes de dleo, gas e agua. O separador\\
utiliza a forca da gravidade para separar misturas de dleo e gas (devido a diferentes\\
densidades dos fluidos). O dleo mais pesado que © gas cai no fundo do vaso e é reti-\\
rado pela linha de fluido. O gas mais leve sobe a0 topo \\& removido para venda. Os\\
componentes principais do separador incluem o dispositivo de separagao primario\\
¢/ou seco secundaria de decantacao (separacao), extrator de névoa para remover\\
Pequenas particulas liquidas do gas e saida de gas. A seco de decantagio (separagao)\\
de liquidos remove o gis € a saida de liquido.\\
\\
acao operacional e fungio. Quanto\\
\\
Pnppnians padon cesdonden men |S n\\
tador verti, operacional, os separadores classificam-se em separador horizontal, sepa-\\
tae ee ou separador esférico. Quanto 4 fungao, 0s separadores classificam-se em\\
eine No separador bifasico, © gas é separado do liquido e o gas ¢ 0 liquido\\
¢m gis, gados separadamente. No separador trifasico, o fluido do pogo é separado\\
Os 6leo ¢ 4gua, com os trés fluidos sendo descarregados separadamente.\\
\\
Separadores horizontais bifasicos sio mais eficientes, pois apresentam U\\
\\
dos em sistemas\\
\\
fea :\\
interfacial que permite uma melhor separacao gis/liquido. Sio usa\\
\\
ma maior\\
\\
Digitalizada com CamScanner\\

\\section*{Pagina 198}

180 —\\_ PRODUGAD E OPERAGOES SUBMARINAS\\
\\
2008 gis/dleo. Desvantagens esto relacionad,,\\
tria dos vasos verticais facilita a remo, “\\
oj e4\\
\\
das (variagOes de fluxo).\\
\\
des ¢ altas ra\\
(a geome!\\
\\
ndes golfa\\
\\
que apresentam emul\\
manuseio dos s6lidos produzidos\\
menor capacidade de absorver gra\\
\\
Figura 7.9: Separador bifasico horizontal.\\
\\
fasicos sao utilizados para separar ¢ remover qualquer 4gua livre\\
\\
que possa estar presente no processo. Mais espago deve ser deixado para a decantacio do\\
liquido, ¢ algum dispositivo deve ser adicionado para a remogao da 4gua livre, Problemas\\
\\
operacionais dos separadores verticais trifasicos incluem:\\
\\
* Espuma: dificulta o controle de nivel do liquido dentro do separador,\\
volume que poderia estar disponivel para a coleta de liquido;\\
\\
Separadores verticais wi\\
\\
pois ocupa um\\
\\
* Obstrucio por parafinas, hidratos ¢ asfaltenos: obstru¢ao das placas coalescedoras\\
na se¢io liquida e os extratores de névoa na secao gasosa;\\
\\
id uns sedimentos: causam a crosio nas valvulas, obstru¢do nos elementos inter\\
nos do vaso ¢ reducio do tempo de residéncia do liquido devido ao acimulo no\\
fundo do separador;\\
\\
* Emulsées: i\\
benny d sim at eras na operacio de um separador e causar proble-\\
; ch sr nivel, 0 acamulo de emulsio diminui 0 tempo de retengao\\
efetivo, resultando em uma reducdo na eficiéncia do processo;\\
\\
* Arraste: o arraste de\\
panlin leo, anadiieie de gis ocorre quando o nivel do liquido est4\\
componentes internos ¢ formagio de espuma.\\
\\
Existem também\\
separadores de ;\\
de produsdo de poros petrolieron, ou ae a 7.11) que servem para medir a capacidade\\
Seja, visam executar somente teste ¢ avaliagi0\\
\\
|\\
4\\
\\
Digitalizada com CamScanner\\

\\section*{Pagina 199}

SISTEMA DE PRODUGAO 181\\
\\
Figura 7.11: Separador de teste.\\
\\
terceira linha. Em principio, o vazamento\\
lor, O vaso é chamado de “Free Water Knoc-\\
nas a Agua livre,\\
\\
que sobe para o topo é extraido por uma\\
de agua livre € quase idéntico ao separadi\\
bout” (nocaute na Agua livre) porque foi projetado para eliminar ape\\
enio a Agua que foi emulsificada.\\
\\
Assaidas de dleo e agua sao controladas por valvulas de nivel controlado, que impedema\\
na parte superior da embarca-\\
\\
drenagem completa da embarcacdo e mantém o gas preso\\
cio. Comumente usado em pocos de baixa pressao € pode ser usado para lidar com altos\\
volumes de fluido, O exemplo deste tipo de vaso horizontal e vertical esta ilustrado nas\\
\\
figuras 7.12 € 7.13,\\
\\
Aquecedor-Tratador: aquecedor-tratador (Heater-treater) também chamado de trata-\\
5 de dleo e agua. Semelhante ao\\
\\
mento por emulsao, usado para separar emulsée: a\\
\\
ento de 4gua livre, mas o tratador tem capacidade de aquecimento com a in-\\
dlusio de tubos de incéndio. A combustao do gas dentro dos tubos de incéndio é\\
usada para aquecer a emulsio de dleo € 4gua.\\
\\
Amedida que a mi\\
forma que a mistura de dleo e Agua fica mais quente, aem :\\
€m 6leo limpo e dgua limpa. A agua é removida do fundo ¢ enviada ao\\
\\
ulsio se quebra € se trans-\\
ema\\
\\
re\\
Digitalizada com CamScanner\\

\\section*{Pagina 200}

182 —\\_ PRODUCAO E OPERACOES SUBMARINAS\\
\\
ons\\
\\
ut horizontal.\\
\\
Figura 7.12: Free water knocko\\
\\
Figura 7.13: Free water knockout vertical.\\
\\
de descarte de 4gua. O dleo limpo é retirado do centro da embarcacao e enviado aos tan-\\
ques de armazenamento para venda. Qualquer gis natural na emulsao de agua € éleo sai\\
\\
na parte superior do aquecedor-tratador. Na figura 7.14 ilustramos aquecedor - tratador\\
\\
horizontal.\\
\\
Figura 7.14: Aquecedor - tratador horizontal.\\
\\
7.3 Vendas e armazenamento do petréleo\\
\\
Nesta secio da nossa formacio, van\\
\\
namento do petréleo em sadais peer ee tratar de forma muito resumida a venda¢ ara\\
\\
snternacional e utilizado nas © offshore, O petréleo de referencia no meres\\
es esto ilustrados na tabela 7.5.\\
\\
Digitalizada com CamScanner\\

\\section*{Pagina 201}

DESCARTE DE AGua SALGADA\\
\\
183\\
feréncia no mercado i i\\
; 75: Petrdleo de re: ‘Cado internacional e utili\\
pele? utilizado nas transacs,\\
Ges.\\
qo de pelea Negécio Caracteristicas\\
Bolsa internacional de Lond\\
Fie ites Alta qualidade, leve ¢ d, 0,3.\\
(PIE) fre, Grau API de 39,40 a emto-\\
WTI(Wesn Texas Intermediate) Mercado de Nova Torque (Ny- Maior qualidade que o Brent. Leve\\
mex) Cotaczo: ficar entre dois quatro digioe\\
: on acima de Brent. Teor de enxofre de 0.3\\%\\
pea ‘0 mercado de matéria-prima Referéncia na Asia. Pesado ¢ sul\\
de Singapura (Simex). Nymex. Grau API de 31; Teor de Babi tone ;\\
Com bi itméti neds\\
cesta de OPEP Gom base ma chamada cesta da Uma média aritmética de sete varieda-\\
\\
des de petréleo: Saharan Blend (Argé-\\
lia), Girassol (Angola), Djeno (Congo), Za-\\
firo (Guiné Equatorial), Rabi Light (Ga-\\
bio), Iran Heavy (Ir3), Basra Light (Ira-\\
que), Kuwait Export (Kuwait), Es Sider\\
(Libia), Bonny Light (Nigéria), Arab Light\\
(Arabia Saudita), Murban (Emirados Ara-\\
bes Unidos) e Merey (Venezuela). Qua-\\
lidade média/baixa. O valor da cesta da\\
OPEP € quatro ou cinco délares inferior ao\\
petréleo WTI.\\
\\
7.3.1 Vendas e armazenamento do petréleo onshore\\
\\
£ uma instalacao tipica, € fornecido armazenamento suficiente para acomodar 0 dleo de\\
dois a trés dias de producao. O dleo pode ser retirado do tanque e vendido por caminhao\\
/ caminhao-tangue ou por oleoduto até a refinaria. A Transferéncia automitica de cust6-\\
dia de locacio mede o volume de éleo que passa para o oleoduto. Ele também mede o\\
conteiido da BS\\&W.\\
\\
eee\\
\\
k 7.3.2. Vendas e armazenamento de petrdleo offshore\\
\\
be As vezes, as plataformas de producdo offshore podem ser equipadas com equipamento de\\
armazenamento de leo se a plataforma estiver localizada em Aguas rasas.\\
\\
Odleo limpo pode ser vendido através de um oleoduto que se estende da plataforma\\
20 longo do fundo do mar até a terra. Na operacéo em 4guas profundas, o petrdleo \\&\\
enviado através de um oleoduto ao longo do fundo do mar para outra plaraforma de\\
‘rmazenamento em Aguas rasas.\\
\\
As vezes, a plataforma de armazenai\\
‘mente para a linha de yendas. Para campos de ee\\
\\
€ pequenos campos de petrdleo, onde os 0! , ,\\
‘candmicos, o uso do FPSO é uma boa opgao. O dleo € armazenado no F\\
no navio-tanque.\\
\\
mento nao é necessiria € © dleo é enviado dire-\\
réleo remotos, locais em aguas pro-\\
eodutos no fundo do mar nao sio\\
PSO e seri\\
\\
tada novamente no solo,\\
da dos equipamentos de\\
\\
|\\
\\
‘3 : .\\
\\
7.4 Descarte de dgua salgada\\
;\\
\\
Para evitar a poli roduzida deve ser inje\\
Poluicdo ambiental, a Agua produz i\\
Stralmente na mesma rocha da qual foi retirada, A Agua é retira\\
\\
Digitalizada com CamScanner\\

\\section*{Pagina 202}

ARINAS\\
184 propucAd E OPERAGOES suBM:\\
jada para um tanque de arma:\\
to de dleo ¢ normalmente enviada Pp: a .\\
\\
e sejaa\\
a jes gua salgada dos pocos pode ser desc,\\
erag «meiro ser enviada através de varios ds ta\\
POsit,\\
\\
processamen|\\
que um yolume suficient\\
Frequentemente, em op!\\
No entanto, 4\\
\\
no proprio mar.\\
vos de limpeza que removem\\
\\
igios de dleo.\\
\\
7.5 Tratamento de gas\\
\\
varios componentes que estado presentes em diferentes concen.\\
O motivo do tratamento é atender s especificacGes de vendis\\
e aos requisitos do cliente. Também é realizado para recuperar outros materiais como 9\\
\\
2S (preso mais alto se comercializado separadamente)¢\\
\\
CO? eo enxofre elementar do H. 4 ;\\
para facilitar a transmissao, impedindo a formagio de hidratos.\\
\\
O gis natural consiste em\\
tracdes (agua € impurezas).\\
\\
7.5.1 Processo de adocamento de gas\\
\\
rocesso de tratamento de gés mais comum. Também é conhe-\\
ou remogao de gas acido. O processo consiste em\\
nerador. No absorvedor, a solugao de amina\\
duzir uma corrente de gis adocado\\
Juco de amina rica em gases éc-\\
da para o regenerador (ume\\
\\
O adocamento de gas € 0 Pp’\\
cido como tratamento de gas amina\\
duas unidades principais: absorvedor e rege!\\
absorve H2S e CO? do gis 4cido ascendente para pro\\
(isto é, um gés livre de H2S) como produto e uma so:\\
dos absorvidos. A amina “rica” resultante é entao encaminha\\
stripper com um reboiler) para produzir amina regenerada ou “magra” que é reciclada pare\\
reutilizacio no absorvedor.\\
\\
e aaah EE Passa por processamento adicional, como desidratacio ¢\\
ppsLare fase sho ee [e\\} Be dcido contém a maioria dos compostos de CO;\\
Claus (Claus Sulfr) enviados para uma unidade de recupera¢a0 de enxofte\\
\\
Send jeactite\\
tenon prestore\\
\\
11 Gi techs prennare\\
\\
Figura 7.15:\\
+15: Prog,\\
de adocamento de gis.\\
\\
Digitalizada com CamScanner\\

\\section*{Pagina 203}

PLATAFORMAS\\
E PRODUGAQ\\
185\\
\\
Figura 7.16: Configuracao da planta de processamento,\\
\\
7.6 Plataformas de producao\\
\\
Asplataformas podem ser de perfuracao, de produco (quando podem extrair o petréleo\\
eseparar dleo, agua € gas) ou ter as duas fungées. As plataformas podem ser classificadas\\
de varias formas, como, por exemplo, pela finalidade (perfuragao de pocos, producao de\\
pocos, sinalizac4o, armazenamento, alojamento etc.), pela modalidade (fixas ou miéveis),\\
pelotipo de ancoragem etc, Sdo muitos os intervenientes nas atividades desenvolvidas nas\\
plataformas de producao. Entre estes destacamos engenheiros, técnicos de varias especi-\\
alidades, profissionais de hotelaria e enfermagem, técnicos de seguranca, mergulhadores,\\
entre outros.\\
\\
Nesta segao, vamos tratar das plataformas de produc\\&o ou as que tém duas funges. As\\
unidades de producao, armazenamento e exportaco de leo, no mar, tornam-se bastante\\
complexas e variadas, para receber 0 dleo produzido dos pocos; fazer a separacao € 0 trata-\\
mento dos fluidos produzidos; em alguns casos, armazené-los; e finalmente distribui-los\\
para terra ou para navios armazenadores. Estas instalacdes, dependendo da profundidade\\
€ da distancia do litoral, podem enviar o dleo/ gs para a terra, através de oleodutos, ou\\
entao para navios armazenadores.\\
\\
As plataformas tém sua utilizacio condicionada a alguns aspectos relevantes como a\\
Profundidade da lamina de Agua, relevo do solo submarino, a finalidade do pogo ea me-\\
thor relacdo custo/beneficio. Em Angola, trabalha-se mais com plataformas de producao,\\
€M sua maioria fixas, As plataformas classificam-se em plataformas fixas, plataformas de\\
Pernas atirantadas, plataformas semissubmersiveis, FPSO e FPSO monocoluna.\\
\\
7.6.1 Plataformas fixas\\
\\
As Plataformas fixas (Figura 7.17) so estruturas apoiadas no fundo do mar por meio de\\
“a8 cravadas no solo com o objetivo de permanecerem no local de operagao Lag\\} oe\\
- Foram as primeiras a serem utilizadas nos campos localizados em laminas ie :\\
pal 300m, © este também é 0 seu limite de utilizacao. Devido ao cust? seven -\\
iat £ntre 0 projeto, montagem e instalagio, a sua aplicagao é restrita aos campos\\
Sua explorac¢io comercial comprovada.\\
\\
—al\\
Digitalizada com CamScanner\\

\\section*{Pagina 204}

RINAS\\
prooucAdo E opERAGOES suBMA'\\
\\
7.6.2 plataformas flutua |\\
\\
=o di ito aos navios sonda, ¢ aS plitafores semissubmersivi\\
\\
fo diz respe resentando yantagens logisticas nas operagies ¢ hoje en,\\
\\
(FPSO vem aP cialmente projetados para operagées de perfuracio, be\\
sao especi@ jal, além de um sistema de posicioname\\
\\
‘Nto\\
\\
Possuem um siste posicao ¢ deste modo nao danificar equipame\\
dinamico qué Ihe pee 40 da a¢a0 dos ventos, ondas e correntes faite\\
Lae -si ‘o estruturas apoiadas por colunas ne\\
rmas sc ulsdo prépria, sendo ead\\
\\
lores submersos,\\
de pogo!\\
\\
186\\
ntes\\
\\
Esta classifica\\
navios sonda\\
de serem adaptados,\\
\\
e prejudicar a\\
\\
Ja as platafo\\
das por flutuad\\
requeridas para perfuracao\\
\\
ubmersiveis sa\\
podendo ou nao ter prop’\\
\\
s explorat6rios.\\
\\
forma de Pernas Atirantadas (TLP)\\
\\
7.6.2.1 Tension Leg Platform/ Plata\\
\\
es utilizadas para a produgao de petréleo. Sua estruturaé\\
forma semissubmersivel. Porém, sua ancoragem ao fundo\\
do mar é diferente: as TLP sao ancoradas por estruturas rubulares, com os tendées fixos\\
ao fundo do mar por estacas € mantidos esticados pelo excesso de flutuaco da plataforma,\\
o que reduz severamente 0s movimentos da mesma. Desta forma, as operacées de perfu-\\
ragio e de completacao sao iguais as das plataformas fixas (Figura 7. 17). As plataformas\\
\\
podem ser utilizadas em laminas de agua até 1.500 m.\\
\\
‘As TLP sio unidades flutuant\\
bastante semelhante a da plata!\\
\\
7.6.2.2 Plataformas semissubmersiveis\\
\\
As platafc pa iter ek\\
pay tipo iy poe eg a 7.17) sio compostas de uma estrutura de um\\
sofre movimentages devido 4 aco Hf m flutuadores submersos. Uma unidade fluruan'é\\
danificar os equipamentos a pie d ve aes correntes e ventos, com possibi\\
ela fique posicionada emoitectina escidos no pogo. Por isso, torna-se necessario\\
de a m. ie, operagao esta a ser realizada em lamina de 4gu@ mais\\
wlan tipos de sistema sio ;\\
é pote cMipery Osleema de ae posicionamento da unidade flucuante: ©\\
812 bneoras.0 cabose/ou near dindmico, O sistema de ancor\\#ge™\\
abs produ:\\
\\
zem esforgos ca\\
das ondas, pazes de restaurar a\\
Ja ago\\
\\
lidade de\\
que\\
\\
Pra Pia neal atuando como molas que\\
sao do flutuante quando é modificada pe\\
\\
© fundo do mar, dinamic\\
"enero a dos equi 0, nao existe ligacio fisica da plataforma com\\
\\
nam a de\\
+ € propulsor; ipamento: :\\
“ ep saeco iguue Perfuragao. Sensores actisticos deter\\
Por computador restauram 4 posigi?\\
\\
Digitalizada com CamScanner\\

\\section*{Pagina 205}

PLATAFORMAS DE PRODUGAO, 187\\
\\
baci ighd bain! Plataforma\\
Atirantadet Semissubersivel\\
[ruwe ou T1P)\\
\\
Plataformas de producio fixas € flutuantes.\\
\\
Figura TAT:\\
\\
16.203 plataformas do tipo FPSO\\
\\
18, so navios com capacidade para\\
do éleo e/ou gas natural. No convés\\
ar e tratar os fluidos produzidos\\
éleo é armazenado nos tanques\\
s em tempos. AS\\
\\
o FPSO, jlustrada na Figura 7.\\
\\
processat earmazenar leo, e prover a transferéncia\\
' de navio, instalada uma planta de processo para separ:\\
| pelos poses Depois de separado na agua € do gas, 0\\
do proprio navio, sendo transferido para um navio aliviador de tempo:\\
operaces usando estas plataformas podem ser realizadas em laminas de 4gua mais de\\
\\
2.000 m.\\
O navio aliviador € um\\
\\
As plataformas do tip’\\
\\
petroleiro que atraca na popa da FPSO para receber dleo que\\
\\
foi armazenado em seus tanques € transport-lo para terra. O gas comprimido é enviado\\
para terra através de gasodutos e/ou reinjetado no reservatorio. Os maiores FPSO tém\\
sua capacidade de processo em torno de 200 mil barris de dleo por dia, com a producao\\
\\
associada de gas aproximadamente 2 milhdes de metros ctibicos por dia.\\
\\
7.6.2.4 Plataformas FPSO monocoluna\\
\\
ma plataforma destinada 4 producio, ar-\\
\\
| Plataforma FPSO monocoluna (Figura 7.18) € w\\
ada em aguas ultraprofundas. Seu\\
\\
mazenamento e transferéncia do dleo bruto, sendo us\\
casco é redondo, promovendo maior estabilidade. Sua tecnologia jnovadora permite que\\
clon entre na parte central do casco diminuindo os efeitos das ondas. A primeira plata-\\
eee nO. mundo é brasileira. As operagoes usando estas plataformas podem\\
Asplataf em laminas de agua mais de 2.000 m.\\
\\
on pang do tipo monocoluna tém como cara\\
celto asseme! bitin uma coluna, diferente das sem!ss\\
205 conceitos ee ts plataformas do tipo SPAR, porém com ca’ ;\\
4 quais Pouals, \\& plataforma monocoluna apresenta vantagens *\\
\\
 podemos citar: menores movimentos, maior reserva de estabi\\
\\
ide operacional, boa relacio peso de convés por deslocamento.\\
\\
marcante a utilizagao de\\
\\
cteristica 1\\
lo, o.con-\\
\\
ubmersiveis. Neste sentid\\
Jado menor. Em relagio\\
ns significativas. Entre\\
ilidade avariada,\\
\\
Digitalizada com CamScanner\\

\\section*{Pagina 206}

5\\
188 — PRODUCAODE OPERACOES suBMARINA\\
\\
mas do tip? SPAR\\
\\
7.6.2.5 Platafor\\
ma estrutura que contém um Gni\\
ICO Cop,\\
\\
rma com U!\\
\\
odendo ainda ter uma eé.\\
‘Strui\\
\\
tura de\\
\\
a platafo’\\
Urely "\\
\\
A Spar (Figura 7.18) é um :\\
\\
cilindrico ¢ 0co, similar a uma grande boia, P' :\\
cial era usa-la para pesquisa oceanografica e arma\\
\\
a ela é usada para perfuracao e/ou Zenagen,\\
\\
Producs,,\\
\\
r. O projeto ini\\
tanque. Agor\\
\\
complementa!\\
do petrdleo como um grande\\
principalmente no Golfo do México.\\
0\\% submerso, ela impée vantagens sobre plataformas semiss\\
mente por sua grande inércia, trazendo a\\
Inde\\
\\
Tendo seu casco 9\\
siveis por ter um pitch\\
vantagem em aguas pro!\\
rigidos, mas tendo seu si\\
semissubm:\\
\\
reduzido, especial\\
fundas. Esta carac!\\
stema de amarragao\\
\\
teristica a aproxima das TLP, permitindo 5;\\
simplificado comparado a uma platafors\\
ma\\
\\
ersivel tradicional.\\
\\
Monocoluna\\
\\
Figura 7.18: Plataformas de produco flutuantes.\\
\\
Na Tabela 7.6 fazem\\
¥ os uma Ci HT\\
na Tabela 27 reaumisibs 6 pe ies ae entre as plataformas de producao, ao passo que\\
tistica de 2012. ro das plataformas existentes em Angola, segundo esta:\\
yee\\
\\
Tabela 7.6: Com a\\
paracao das atividades realizadas pelas plataformas\\
\\
Fira\\
Pemas \\_ atirant\\
2 tad: is\\
Atividade ——\\_Perfuragio e prod Fre aera Seuuesebmertivel| /BPSO mionocoluna FPSO\\
sio lu- Perfuracio\\
Rented\\
Profundidade Até = uragio © produ: Perfuragio, produ: Produsio. |\\
300m ack gio, armazenamento mazenamen?\\
500 m Mais de 2 ¢ transferéncia ¢ transferénc\\
000m Maisde 2.000m — Maisde200\\
\\
Fundo do mt\\
\\
Controlo de po- Si\\
me Superficie\\
Producto 2 Oleodutos\\
\\
Fund\\
lo. do mar Fundo do mar\\
\\
Escoado\\
Vantagem ParaFPSQ O\\},\\
Instalacdes leodur\\
mais sim- Sistema de eae fundo Transferéncias para Transferéno®\\
ancora. petroleiros para petrolei®\\
cap)\\
de\\
\\
Ples econtrolo\\
de po-\\
Bem tigido Especialme:\\
ga ei\\
© movi- te proje- Movimentos meno Maior\\
\\
\\$03.4 superficie\\
Tedurido Para\\
ome Pe SE MNEEPSO cidade\\
oO armazen\\
a esenvolviments deen, mento\\
engenharias cutting —\\
rina su\\
Te NOs anos 80s como parte integr™"\\
\\
Digitalizada com CamScanner\\

\\section*{Pagina 207}

PLATAFORMAS, DE PRODUGAG\\
\\
189\\
‘tabela 7.7: Namero de estruturas de produgao existentes €m Angola\\
Tipo de Plataformas N° de Plataformas Blocos\\
Joquetes 181 0,2,3\\
Compliant Tower \\$ vi\\
Avtoclevatorias t p\\
Gosendionaie 37 0,2,3\\
Acomodacies z oe\\
Pernas Atirantadas 2 5\\
FPSO 12 0, 4/05, 14, 15, 17, 18, 31\\
FSO : 2\\
Boias 8 0, 14,15, 17, 18\\
Barcos de Acomodacio 2 23\\
Toul 251\\
\\
com finalidade de maximizar a produgio principalmente nos campos de aguas profundas\\
| ¢ultraprofundas. Podemos definir um sistema submarino de produ¢ao como sendo um\\
| conjunto de instalagdes submersas destinadas a elevacio, injecao e escoamento de fluidos\\
produzidos e/ou movimentados em um campo de petroleo ou gis natural.\\
\\
O projeto de desenvolvimento de um campo visa a maximizacao da recuperacao de\\
petréleo a um custo minimo operacional e de investimento de capital. A Figura 7.19\\
ilustra o sistema submarino de produgio. O arranjo final de um campo € 0 resultado de\\
um processo de otimiza¢ao que envolve diversas variaveis, tais como:\\
\\
* Numero de pocos e posicionamento dos mesmos;\\
\\
* Comprimento e diametro dos dutos de producio;\\
\\
* Posicionamento da unidade de producio flutuante;\\
* Tipo de ancoragem;\\
\\
* Meios de instalagao;\\
\\
* Perfil de producao desejado;\\
\\
* Necessidade de utilizacao de meios de elevacao artificial etc.\\
\\
7.6,3 Equipamentos submarinos\\
\\
sub ja subm - Jo sao: cabega de pogo.\\
| mm: ino de produgéo si: cab\\
A, dutos I tos do siste jbmarino s ASAP\\
5 inos, manifolde, PLEM (pipeline end manifold), PLET (pipeline end ter\\
nator), :\\
\\
Digitalizada com CamScanner\\

\\section*{Pagina 208}

190 propucho E OPERAGOES SUBMARINAS\\
\\
Figura 7.19: Sistema submarino de producio,\\
\\
7.6.3.1 Cabeca de produgao\\
\\
As cabecas de poco submarinos suportam os revestimentos dos pocos, resistem 0s esfor.\\
Gos do riser e fornecem vedacao para o BOP. Na fase de producio, servem de al\\
‘ravamento ¢ vedacio para o “suspensor” de tubulacao e para a arvore de na\\
namente, as cabecas de poco so preparadas para receber a Base adaptadora di\\
(BAP). Existem dois tipos de Sistemas de cabeca de pocos submarinos (SCPS): para unj.\\
dades flutuantes e para unidades apoiadas no fundo do mar.\\
\\
lojamento,\\
tal. Moder.\\
le producio\\
\\
Figura 7,20: Cabeca de Po¢o submarino.\\
\\
7.6.3.2 Arvore de Natal\\
\\
De uma foi genéric\\
\\
a0 servico wirhogrd de mien. mos (Figura 7.22) classificar as Arvores de natal quanto\\
© horizontal). A ANM é ume 0 © de injecio) ¢ configuracao (convencional ou vertical\\
Posto de um conjunto de cones eet? instalado na cabega do poco submarino, com\\
\\
conectores ¢ valvulas ‘i de fluidos\\
injetados po que permitem controlar o fluxo de |\\
— = 8 Projetado para Suportar elevadas pressdes ¢ te™P*\\
\\
Digitalizada com CamScanner\\

\\section*{Pagina 209}

PLATAFORMAS DE PRODUCAD 1g)\\
\\
s do pogo (além de elevadas pressdes ¢ baixas tem;\\
' ont com suporte de mergulhadores em profundid\\
anes ultraprofundas, com auxilio de um (ROV).\\
\\
peraturas ambien\\
tes). Py\\
lade de até 300 m o aa\\
\\
instal\\
Na figura 7.21 est\\
\\
prof\\
RO.\\
\\
uu, em aguas\\
4 ilustrado °\\
\\
Figura 7.21: Veiculo de Pperacdo Remota (ROV).\\
\\
O nome “arvore de natal” tem origem na década de 1930, quando moradores de pro-\\
vincias petroliferas norte-americanas fizeram a associa¢ao do equipamento eae de\\
peve com um pinheiro natalino. Com a descoberta de petrdleo no fundo do mar, o equi-\\
pamento foi adaptado as novas condicGes e passou a ser chamado de Arvore de Natal\\
Molhada (ANM), muito utilizada em sistemas de producto offshore.\\
\\
Figura 7.22: Arvore de Natal Molhada (ANM).\\
\\
7.6.3.3, Dutos submarinos\\
\\
o dutos submarinos sao responsaveis pela movimentacio dos fluidos produzidos e in-\\
Jetados num campo de petréleo e gas, e também utilizados para escoamento (offloading)\\
dos fuidos processados pela unidade de producao. Na produgao, temos 0 fluxo de dleo\\
gas da ANM e/ou manifolde para a unidade de produgao. Na injecdo, temos 0 fluxo de\\
\\
dos e gis da unidade de produgao para o Manifolde e/ou da ANM. Os dutos podem\\
‘er classificados quanto a estrunura (rigidos ou flexiveis), quanto @ funcao (escoamento,\\
umbilical, misto) e quanto a configuracio (riser ¢ linhas de fluxo).\\
\\
Digitalizada com CamScanner\\

\\section*{Pagina 210}

INAS\\
propuchd E OPERAGOES suBMAn\\
\\
192\\
7.6.3.4 umbilical\\
‘ jras que transportam de;\\
njunto de manguciras , icies\\
9° umbilical (Figu’? 7.23)6 Ly po como cabos elétricos transmissores de Sinais dy\\
quimica, 4 nismos de abertura ¢ fechamento do equipement 7\\
Me de\\
\\
hidraulicos ¢ de injegao\\
ua funcio é acio\\
gas su\\
dem ser hid\\
\\
quimicos.\\
\\
ar os mecal Fund\\
b arinos monitorando as caracteristicas do poco (tempery,\\
m elétricos de sinal, elétricos de poténcig bin\\
ONG 6\\
\\
téncia. S\\
raulicos,\\
\\
extracao de dleo €\\
¢ pressao). Estes po!\\
injecao de produtos\\
\\
Figura 7.23: Umbilical.\\
\\
7.6.3.5 Linhas flexiveis e riser\\
\\
Linhas flexiveis e risers s4o os dutos que conduzem os fluidos produzidos pelo poco para\\
unidades de produgdo. Podem também ser utilizados para interligacdo de uma unidade\\
a outra, para injecao ou descarte de fluidos em reservatérios ou para a exportacio da\\
\\
producdo em terra.\\
\\
As linhas flexiveis (Figura 7.24) so os dutos empregados em todo sistema submarino\\
de coleta e escoamento, que conduzem os fluidos produzidos pelo poco para as unidades\\
de produczo. Elas podem ser utilizadas para a interligacdo de uma unidade a outra, para\\
big ba descarte de fluidos em reservatérios ou para a exportacao da producio para\\
pe a ceare oh com camadas de materiais metlicos e nao metilicos, cada\\
doa Seaacires® 4 ‘e fs Em suas extremidades, possuem acessérios denomin\\
\\
Dancer (flexi Eo diitod as: de natal molhadas a manifoldes ou risers.\\
interligam as linhas sbeasinas dis ah sao os trechos suspensos das tubulagoes que\\
as plataformas. Podem também a il i ae cae cue ae\\
Ieito marinho, como os risers de inj so dos para conduzir fluidos da superficie até °\\
entico de um sistema submaring deg ge ee: Constieul-se num component?\\
de tragao e fadiga, devido ao seu pr Producio, por estar submetido a elevados esfors®®\\
€as movimentacées da unidade vie Peso, 4 aco de correnteza, aos efeitos das om as\\
\\
estacionaria de Producio.\\
\\
7.6.3.6 Manifolde\\
\\
O manifolde (Fj\\
‘ igura 7.25) tem\\
pie hy varios pocos, B cee, principal reunir, numa s6 linha, @ produsi°\\
), Conjunto de Valvulas de bl Porarranjos de tubagens (coleta, injecao. reste\\
queio, valvulas de controle de escoamento (chokes\\
\\
4\\
Digitalizada com CamScanner\\

\\section*{Pagina 211}

PI\\
LATAFORMAS DE PRODUCAO =—-193\\
\\
Figura 7.24: Linhas flexiveis (a esquerda) ¢ risers flexiveis (4 direi )\\
reita),\\
\\
ou neers e subsistemas de monitoramento, controle e interc\\
ente por via elétrica — ‘i = ‘onexio —\\
“ a" me) 6 mal Fa de ered de producao estaciondria. No caso Uietok\\
vindos da unidade estacionaria de seer para os pocos os fluidos de inje a\\
seit coutidas Tat esnd tant produgao. As fungées de produgio e injecio ae\\
o manifolde. yje¢ao podem\\
\\
Figura 7.25: Manifolde.\\
\\
de fee ead a utilizagao de manifoldes sao a redugao do comprimento total\\
Por ais beneficins « Peta de risers conectados a unidade estacionaria de produgio.\\
contribuigio na oh wep oalto custo dessas linhas, esse equipamento é de alta\\
manifoldes pédits Ginkh U gical ¢ econémica da produgio, notadamente no mar. Os\\
them sistemas de j FADED pat utilizados para permitir que um grupo de pogos compar-\\
Defines aus injegio de agua e gas lift.\\
\\
G40 de virios pos fated) e simplificada, eles servem para ©\\
serem bisides 5 unidades de producao e também para d\\
odes ajudam a Re POsoe. Como agrupam os fluidos produzidos p\\
Compriment luzir o numero de linhas conectadas a plataforma,\\
\\
‘0 total das linhas de interligagdo de pocos usados num sist\\
\\
direcionamento da produ-\\
istribuir fluidos destas para\\
jor pogos, OS mani-\\
além de reduzir o\\
ema de produgio.\\
\\
Digitalizada com CamScanner\\

\\section*{Pagina 212}

INAS\\
\\
194 propucAd E operacoes sUBMAR'\\
aGao\\
\\
7.6.3-7 Equipamentos de interligas\\
eo equipamento responsivel por interligar ton .\\
fetes @ boia 0 planta onshore. PLET (pipeline end tering *\\
tagio de leo oe rojetados © desenvolvidos para possibilitar a conexio oa 2\\
Te Fexivel® O objetivo pasico 6 minimizar as ltrveatimentos do si\\
\\
fal los e as .\\
\\
das linhas de escoamento-\\
\\
nd manifold;\\
\\
7.7 Paises da OPEP\\
\\
ntal criada em 14 de setembro de 1960 em\\
ablo Pérez Alfonzo, Abdullah Tariki e como\\
bjetivo é a centralizacao da elaboragao das\\
dos os paises membros, ou seja:\\
\\
A OPEP é uma organizacao intergovername!\\
Bagda, Iraque, ¢ teve como fundadores Juan Pi\\
secretario-geral Mohammed Barkindo. Seu o\\
politicas sobre producio ¢ venda de petréleo de to\\
\\
» Estabelecer uma politica petrolifera comum a todos os grandes produtores de petré-\\
leo do mundo (paises membros);\\
\\
= Definir estratégias de producao;\\
\\
* Controlar pregos de venda de petrdleo no mercado mundial;\\
\\
* Analisare i ff\\
gerar conhecimentos para os paises membros sobre o mercado de petréleo\\
\\
mundial;\\
* Cor eT\\
ntrolar volume de produ¢4o de petréleo da organizac4o.\\
\\
Durante 0 programa “D)\\
que cerca de me ti ae. Fonte Janeiro 2021”, Patricio Quingongo apresentoy\\
grantes da OPEP contra 20,6\\% tre cad de petréleo esto localizadas nos paises inte:\\
ioe nese OPEY cies atic paises nado membros da OPEP. Ou seja, as reservas\\
dos paises nao-OPEP, 40\\% hae. bilhdes de barris contra 308,18 bilhoes de\\
realcar que além de An; HH lucéo mundial e 60\\% das exporta¢6es mundiais\\
gola que é membro da OPEP desde Janeiro de 2007,\\
\\
fazem parte desta\\
Emirados Onganizacio a Ardbia Saudi\\
fecente hes ‘abes Unidos, Argélia ria ry a Kuwait, Ino, Iraque, Venezuela, Libia,\\
Brante em 2018, A co, » Nigeria, Guiné Equatorial, Gabao ¢ 0 Cong? o mais\\
dais\\
\\
est ilustrada mparti\\
na Tabela 7.8, Participacao dos paises da OPEP nas reservas ™"\\
\\
PE\\
1\\
\\
Digitalizada com CamScanner\\

\\section*{Pagina 213}

PAISES DA Onep 195\\
\\
« Reservas provadas de Petréleo Bruto da OPEP em\\
\\
bilhde:\\
qobela , 8 de barris,\\
ge 08 fonte Petrangola no programa Direto ao ponto da TV Zimbo, Janeiro —\\
we neers pes ae Pas Reservas\\
pe seo 20708 24 Lia aad up ap -\\
on\\
- 1860 wa Nigéa 46,9731 Gablo 2.00 02\\
wee 148,02 122 Argtlia 12.20 10 G. Equa. 1.19 a\\
ral\\
Kuwait 101,50 as Equador 8,27 O7\\
\\
a\\
Conclusao\\
\\
Este capitulo serviu para equipar os formandos com conhecimentos basicos sobre a enge-\\
nharia de producio de Petréleo e das Operacdes Submarinas. Com esses conhecimentos\\
bisicos, os leitores estao capacitados para descrever alguns equipamentos utilizados nas\\
operacées da engenharia de producao de petréleo e das operagées submarinas. Foram\\
apresentados contetidos relacionados com os componentes do sistema de produc¢io de\\
Petrdleo, assim como a identificacao das diferencas entre as plataformas de producio pe-\\
trolifera ¢ todas as atividades relacionadas com esta fase da cadeia de valor da industria\\
\\
pevolifera.\\
\\
Digitalizada com CamScanner\\

\\section*{Pagina 214}

cAPiTULO 8\\
\\
MIDSTREAM E DOWNSTREAM\\
Sn rn en ne nn aS\\
\\
Ainddstria petrolifera, como tratamos no capitulo , divide\\
tos: Upstream, Midstream e Downstream. Abordamos 0 segm.\\
anteriores, que é a fase caracterizada pelas atividades de bu:\\
das fontes de dleo, bem como 0 transporte deste dleo extraido até as refinarias, onde sera\\
processado. Resumindp, sao as atividades de exploracio, perfuracio ¢ producio.\\
\\
Neste capitulo, vamos abordar de forma resumida os segmentos de Midstream e Downs-\\
tram. A fase Midstream \\& aquela em que as matérias-primas (hidrocarbonetos) so trans-\\
formadas em produtos prontos para uso especifico, como gasolina,\\
GLP, nafta, leo lubrificante etc. Sao as atividades de refinagdo. O segmento de Downs-\\
tream \\& a fase logistica, ou seja, o transporte dos produtos da refinaria até os locais de\\
) consumo. Isso envolve o transporte, distribuicao e comercializacao dos derivados do pe-\\
\\
trdleo. Portanto, nas proximas sec6es, vamos abordar contetidos essenciais relacionados\\
40s segmentos de Midstream e Downstream.\\
\\
Oobjetivo geral deste médulo é entender as operacGes de transporte, armazenamento,\\
tefinacio e Petroquimica. Os objetivos especificos sao:\\
\\
“se em trés principais segmen-\\
ento de Upstream nos capitulos\\
sca, identificagao e localizacao\\
\\
diesel, querosene,\\
\\
* Aprender sobre a caracterizacdio do petréleo bruto;\\
* Conhecer os meios utilizados para 0 transporte de crude e gas;\\
* Descrever os processos de refinagao;\\
\\
* Explicar os tipos de tanques utilizados para armazenar 0 petréleo bruto € 0 g\\
assim como os seus derivados;\\
\\
* Conhecer os derivados do petréleo;\\
\\
ssos petroquimicos;\\
. Deserever nd diferentes pro dutos derivados dos processos petroquin\\
\\
natural em Angola\\
\\
A eas roleo e gas\\
Comhecer as principais unidades de tratamento do petrdleo ¢ § 197\\
\\
- crag: Pilea para ndo-engenheiroy, 1 liga\\
I = Ade Raposo Ramos Copyright © 2023 Bditora Dots).\\
ee : dacome\\
\\
Digitalizada com CamScanner\\

\\section*{Pagina 215}

powNnsTREAM\\
\\
198 miDSTREAM E\\
to\\
\\
a roéleo bru\\
\\
8.1 caracterizas@° de pet\\
ru (em inglés, crude oil),\\
\\
bém conhecido como troleo ¢ on a on 1 teree,\\
Petréleo Bruto, tam senta na rureza, sem ter pa por processamenns fo\\
no estado em que se Pe pocarbonet0s: contendo ainda compostos years,\\
mistura complexa de a contaminantes. Na Figura 8.1, ilustramos um exem,\\
\\
ere 5 id :\\
\\
nio, oxigénio, metals\\
petrdleo bruto.\\
\\
a 8.1; Exemplo de petréleo bruto.\\
\\
Figur\\
\\
lo no campo de produgao, é chamado petréleo bruto ¢, 2\\
extraido, pode apresentar diversos aspectes\\
ele se apresenta em varias cores, variance\\
inflamavel, menos denso qt\\
por hidrocarbonetos. A as\\
do petroleo bruto define os\\
\\
O petréleo, quando extraid\\
pendendo da rocha reservatério de onde foi\\
visuais ¢ caracteristicas diferentes. Por isso,\\
entre 0 negro ¢ 0 castanho escuro, tendo carater oleoso,\\
a 4gua, com cheiro caracteristico e composto, basicamente,\\
tribuigao de percentuais de hidrocarbonetos na composi¢ao\\
diversos tipos de petréleo.\\
\\
O petréleo, em seu estado natural, também contém proporcdes menores de contr |\\
minantes (enxofe - S, nitrogénio - N, oxigénio - O € metais, como niquel - Ni, ferro -\\
cobre — Cu, sédio - Na e vanadio - V). Entre as principais caracteristicas do Petréleo Brut.\\
posers destacar aquelas que mais influenciam no valor comercial de uma determ=3>\\
\\
* Teor de enxofre;\\
* Densidade ou densidade na escala de °API.\\
\\
Importa realcar AP ii\\
American hadon hilaictAte = densidade de petréleo liquido estabelecida Pe\\
tréleo, como ilustramos na “wo. comercial dos diferentes tipos \\& PF\\
\\
8.1.1 Teor de enxofre\\
\\
O enxofre no petréleo brui\\
to causa\\
Sree ptr tm mn ne\\
\\
4\\
\\
Digitalizada com CamScanner\\

\\section*{Pagina 216}

CARACTERIZACAO DE\\
PETROLEO BRUT;\\
‘0 199\\
+ Manuscio: reducao da eficiéncia dos catalisado:\\
h\\
agentes facilitadores que transformam fracdes\\
através da quebra de moléculas dos compostos\\
\\
Tes nas refinarias,\\
mais pesadas em\\
reagentes;\\
\\
Catalisadores Sio\\
Outras mais leves\\
\\
» Transporte: corrosdo em oleodutos e gasodutos;\\
\\
+ Derivados: causam poluicdo ambiental se presentes\\
\\
€m combustiveis derivad\\
De vados do\\
\\
Com base no teor de enxofre, 0 petréleo bruto pode ser classificado como doce ou acido:\\
\\
+ Oleo doce: apresenta baixo contetido de enxofre (menos de 0,5\\% de sua massa);\\
\\
* Oleo acido: apresenta teor elevado de enxofre (bem acima de 0,5\\% de sua massa)\\
\\
Os petréleos acidos sao aqueles que possuem com\\
tendo cheiro peculiar. Estes contém gis sulfidric\\
doces nao contém gas sulfidrico.\\
\\
Chaminés, filtros e outros dispositivos sio usados\\
pores e poeiras para a atmosfera; unidades de Tecuper;\\
cuja queima produziria diéxido de enxofre, um dos\\
banos.\\
\\
Ha um processo de modernizac¢io das refinarias em todo o mundo para processar\\
petrdleos com baixo teor de enxofre, o que resulta em combustiveis menos poluentes.\\
\\
postos de enxofre em alta Percentagem,\\
©, que € perigosamente téxico, Os dleos\\
\\
Para evitar a emisso de gases, va-\\
‘ago removem o enxofre dos gases,\\
Principais poluentes dos centros ur-\\
\\
8.1.2 Densidade ou Densidade na escala de API\\
\\
Adensidade de um liquido é a razo entre a sua massa especifica na condigao padrao de\\
temperatura € pressdo e a massa especifica da Agua pura sob as mesmas condicées. Sua\\
Tepresentacao matemftica é ilustrada na Equagio 8.1.\\
\\
=at (8.1)\\
Yo Pw\\
\\
Onde y, é a densidade relativa do dleo, Pj € Py S40 a massa especifica do dleo e da agua,\\
"spectivamente. O mesmo conceito é aplicdvel aos gases, mas a referencia é 0 ar seco.\\
Ou seja, a densidade de um gas é a razZo entre a sua massa especifica na condi¢ao padrio\\
\\
temperatura e Pressao e a massa de ar seco também na condic¢ao padrao. A expressio\\
Matemitica esta Tepresentada na Equacio 8.2.\\
\\
\\_ Pg \\_ Mg (8.2)\\
Ye Ou 28,96\\
Onde Ye € a densidade relativa do gas, Pg © Par S40 a massa especifica do ais been\\
"spectivamente, e Mg éa massa molecular do gis. As condicdes padrao de tempera\\
*Pressio referenciadas S40 60 °F (15 °C) € 14,696 psi (101,325 kPa). API, para me-\\
A densidade API é a medida da densidade do petréleo estabelecia pelo APL spitulo\\
\\textbackslash{}dentificaco comercial dos diferentes tipos de petrdleo, como abordac - tow 12\\
idade API de um 6l leo pode ser calculada conforme ilustrado na Equagio 1.\\
\\
Digitalizada com CamScanner\\

\\section*{Pagina 217}

200 MIDSTREAM £ DOWNSTREAM\\
\\
ele flutua na Agua, caso contrario, ele é mais\\
\\
Se adensidade API ou °API for — ma ce aduagio aumenta coma diminunan\\
esado. Nesta escala, a agua apresenta AP" \\textasciitilde{} top, anquanto oléos leves\\
ta densidade do éleo. Oleos pesados possuem API ae aah. paket poly\\
‘API elevados. Em outras palavras, quanto ath ;\\
quanto maior o grau API, maior 0 yalor e ecrdadl\\
intermediarios ¢ pesados aumenta com 0 ¢¢ apr\\
Figura 8.2, podemos observar que 0 dleo oa ‘baer\\
que o dleo com menor API. Essa questo a oe east\\
nas secdes sobre os derivados de petrdleo. A Figura 8.\\
\\
do dleo e¢ os derivados extraidos.\\
\\
maior oO val\\
ono mercado. O teor de componentes\\
\\
mo da densidade API. Por exemplo, na\\
r API apresenta menos residuos do\\
dos sera tratada com mais detalhes\\
aa relacao entre o grau API\\
\\
«\\
2 o 2 oo a\\
\\
mane\\
\\
Figura 8.2: Relacao entre o grau API do dleo e os derivados extraidos.\\
Dessa forma, uma amostra de petréleo pode ser classificada quanto ao grau °API da se-\\
guinte forma:\\
\\
. « Petréleos Leves: acima de 30 °API (< 0,72 g/cm?). Caso do tipo de petréleo para-\\
1 finico;\\
\\
4 * Petréleos Médios: entre 21 ¢ 30 °API. Caso dos tipos de petrdleo nafténico;\\
\\
. ahem Deedes abaixo de 21 °API (> 0,92 g/cm?). Caso do tipo de petrdleo\\
\\
Na\\
poms tenes hoya com detalhes os tipos de petréleo parafinico, nafténico ¢\\
ok acai rs sag tlh r diferentes classificagdes de diferentes organismos.\\
: barril a ser refinado (1 barril = 158,98 li imento em\\
derivados varia de acordo com © tipo de petréleo ere eels\\
\\
timo, com 0 processos utilizados, Por as condicdes operacionais e, PO al\\
\\
quantidade de derivad exemplo, petrdleos mais leves geram uma maior\\
sados geram mais éleo phraiew 10 gases combustiveis, GLP e mere: Petrdleos pe\\
40 menor custo e em “oe -— O objetivo é atender o mercado nacional\\
€ 0 grau API, conforme c isquer circunstncias. Existe uma relagdo entre *\\
A Viscosidade é usada para pte ate na Figura 8.3.\\
Na Figura resisténcia ao fluxo exercida nlp em pesado ¢ leve. A viscosidade ¢ u™\\
APL. 8-4, estd ilustrada a relacs AS atrito interno entre as moléculas de um fluido.\\
tre a viscosidade ¢ a densidade do dleo em 8°\\
\\
Fi\\
\\
Digitalizada com CamScanner\\

\\section*{Pagina 218}

OO\\
\\
CARACTERIZACAO DE PETROLEO BRUTO 201\\
\\
Figura 8. 3; Esquematica da relacao entre a densidade e a densidade API do dleo.\\
\\
Vis. = 10 100 1.000 10.000 100.000\\
\\
wm | (35 20 5 10 5\\
\\
Figura 8.4: Relacao entre a viscosidade e a densidade do éleo em grau API.\\
\\
Opetréleo produzido em Angola é um dos melhores do mundo devido ao seu grau API\\
nforme ilustrado na Tabela 8.1, enquanto na Tabela 8.2 apresen-\\
\\
¢aoteor de enxofre, co\\
‘como podemos ver, a rama Nemba é a melhor\\
\\
tamos as ramas produzidas no mundo. C\\
de Angola, com 39,7 °API, muito proxima do petrdleo comercial de referéncia do nosso\\
\\
crude, que é o Brent do Reino Unido (40,0 © APJ), e ainda mais proxima do petréleo WTI\\
dos Estados Unidos (39,8 °API).\\
\\
Tabela 8.1: Rama de petrdleo bruto produzido em Angola.\\
\\
Rama Densidade API ‘Acidez (TAN, mg KOH/g) Enxofre (\\% Peso)\\
rales 7 0,03 0.18\\
Cabinda 325 0,06 0,12\\
Kens 397 0,1 0,20\\
pe 26 1,39 051\\
G 209 0,36 0,42\\
us 0,49 0,63\\
Mungo 28,3 o72 0,65\\
— 318 03 039\\
ue 30,3 07 039\\
ied 233 18 039\\
‘i esnd 336 0,37 ous\\
uA 077 0.85\\
\\
0,25\\
\\
Sauiy|\\
Batuque 327\\
\\
Digitalizada com CamScanner\\

\\section*{Pagina 219}

202 MIDSTREAM E DOWNSTREAM\\
\\
mas produzidas internacionalmente.\\
\\
Tabela 8.2: Ra\\
Densidade API Enxofre (\\% Peso)\\
\\
Origen\\
Ml\\
\\
Rama =\\
Arabian Extra Lt Ardbia Saudita oun ig\\
Arabia Lt Aribia Saudita 4 FA\\
Brent Reino Unido a te\\
Cano Limon Colémbia FA Z\\
oe o 295 02\\
Forcados Nigéria 50d as\\
Kuwait Blend Kuwait ee pd\\
sans ore 21,3 34\\
ae a 39,8 03\\
West Texas Intermediate EUA\\
(wT)\\
\\
8.1.3 Classificagao dos tipos de petrdleo\\
\\
sificados com base em seus constituintes e de acordo\\
do é relevante tanto para os geoquimicos, que bus-\\
cam caracterizar o Oleo para relaciona-lo 4 rocha-mae e€ medir seu grau de degradacio,\\
quanto para os refinadores, que desejam conhecer a quantidade das diversas fracdes que\\
\\
ser obtidas, bem como sua composi¢ao € propriedades fisicas. De acordo coma\\
predominancia dos hidrocarbonetos encontrados no petréleo bruto, este pode ser class-\\
ficado em parafinicos, nafténicos e aromaticos, como resumido na Tabela 8.3.\\
\\
Os tipos de petroleo podem ser clas\\
com a predominancia. Essa classifica\\
\\
Tabela 8.3: Tipos de petréleo quanto 4 predominancia.\\
\\
ee\\
\\
(1) Gasolina de balxo indice de octanagem\\} (2)\\
Querosene de alta qualidade; (3) Oleo diesel com boss\\
caracteristicas de combustlo; (4) Oleos lubrificantes é¢\\
alto Indice de viscosidade, elevada estabilidade quimics\\
@ alto ponto de fluidez; (5) Residuos de refinago com\\
\\
(2) Gasolina de alto indice de octanagem; (2) Querosene\\
de alta qualidade; (3) Gleos lubrificantes de bso\\
residuo de carbono; (4) Residuos asfélticos na refinasS0;\\
(5) Nafta petroquimica.\\
\\
5\\
|\\
i\\
\\
8:14 Outras defini\\
\\
derivados S6es importantes sobre o petroleo \\& seus\\
\\
Digitalizada com CamScanner\\

\\section*{Pagina 220}

PROCESSOS DE REFINACAO DE a\\
203\\
jo: éuma combustéo ndo controlada que se inicia devid\\
peat na camara de combustao, resultante da pipe\\
soln Nas proximas segdes, abordaremos os diferentes f\\
cies combustiveis;\\
\\
‘0 20 aumento da tem-\\
Pressdo da mistura ar/\\
indices de Octano de dife.\\
\\
ponto de fluidez: é a temperatura abaixo da qual 0 éleo nao fluira;\\
a\\
\\
qoxicidade: éa capacidade inerente de um agente causar efeitos adversos em u:\\
\\
. im orga-\\
nismo VIVO;\\
\\
piodegradagao: € um processo no qual os microorganismos\\
\\
in (bactérias e fungos) presen-\\
tes no meio utilizam os hidrocarbonetos do petréleo co\\
\\
mo fonte de alimento,\\
\\
8.2 Processos de refinacdo de petrédleo bruto\\
\\
Arefinagio de petrOleo é um processo industrial em que o petrdleo bruto é processado\\
erefinado para obter produtos de maior valor agregado. O processo de refinagio pode\\
ocorrer de duas maneiras:\\
\\
* UPGN: como exemplo, a planta de Angola LNG;\\
* Refinaria: como exemplo, a refinaria de Luanda.\\
\\
Assim, todas as unidades de processo da refinaria realizam algum processamento sobre\\
uma ou mais entradas, gerando uma ou mais safdas. Todas as entradas origindrias dire-\\
tamente ou indiretamente do petréleo (gas, petréleo e produtos intermediarios) so cha-\\
madas cargas, entrada do processo ou matéria-prima, conforme ilustramos na Tabela 8.4.\\
\\
Tabela 8.4: Exemplo de uma unidade de processo.\\
\\
Nem todos os derivados si0 gerados de uma s6 vez € no mesmo local na refinaria. Quase\\
\\
ipre, eles sio obtidos apés uma sequéncia de processos, chamados de operagGes uni-\\
\\
tirias e que consistem em transformagées de um ou mais fluidos (gas e/ou liquido), que\\
\\
de entradas do processo, em outros fluidos, produtos da saida do quonsa Ss\\
\\
‘idos numa refinaria, sejam "1 entrada ou de saida de algum processo, sto tambem\\
N idos como correntes.\\
\\
‘4 unidade de processo, por exemplo, na\\
\\
“SS Processos ou etapas: separer, tratar, transformar e conve’\\
\\
ja-prim diferen-\\
ia, a matéria-prima passa por\\
\\textasciitilde{} rter. Além dos equipamentos\\
\\
Digitalizada com CamScanner\\

\\section*{Pagina 221}

204 = MIDSTREAM E DOWNSTREAM\\
\\
funcionamento da fabrica precisa-se “ Mea humanos quali\\
operacionais, manuten¢ae e¢ administratvas. Insumos (utilidades\\
; ento e tratamento da unidade. A es,\\
git, agua /vapor e combustiveis. °\\
\\
destacamos LPG (propano/butano)\\
coque, dissolventes etc.). ee\\
tipo de petréleo a ser pro-\\
de um mercado. Um esquema de refino, exemplo Figura 8.5,\\
cubed d s. Por isso, alguns derivados podem se\\
em algumas Fe\\
\\
ou maquinas, para ©\\
dos para as atividades\\
também sio necessirios pare\\
destacamos os quimicos, catal\\
\\
Os principais produtos da saida do P'\\
Jeo, fuel ail € OU\\
\\
Jisadores, en\\
ocesso\\
tros (asfalto,\\
\\
solina, petréleo /Jet fuel. gaséleo,,\\
\\
Cals refinaria ee jda de acordo com ©\\
cessado e as necessidades\\
define e limita 0 tipo € 34\\
produzidos em todas ou apenas\\
\\
finarias.\\
\\
Figura 8.5: Esquema geral de uma refinaria.\\
\\
Durante a vida de uma refinaria,\\
também podem mudar as espe Awe Oe 0 tipo de petréleo que ela recebe, como\\
bai ados por ela produzidos. Por isso, : edi (qualidade) ou a demanda (quantidade) dos de-\\
ee isto é, uma capaci a ce dizer que toda refinaria tem um certo grav\\
pie refino, que permite reajustar Fie aes dinamica na opera¢ao do seu\\
ngas no tipo de dleo e nas necessi o funcionamento das unidades para se adequat\\
idades do mercado e ambientais.\\
\\
REELED ET LE\\
\\
P,\\
rocesso de decantacao\\
\\
-e,rOr\\
\\
Digitalizada com CamScanner\\

\\section*{Pagina 222}

EERE EP EEE EEE\\
\\
l\\
Figura 8.6: Processo de decantacao do Petréleo.\\
\\
Pré-aquecimento\\
\\
Figura 8.7: Esquemitica do processo de pré-aquecimento do dleo bruto.\\
\\
Processo de dessalgacao/dessalinacao\\
\\
Process de dessalinacio consiste na reducdo do teor de sal no petrdleo. No caso de\\
Porque - ou processamento de dleo, o termo realmente usado é dessalga¢ao.\\
i Agua po-\\
tive i essalinacdo é mais usado para transformar a agua do mar em dgua po\\
\\
Pomas * dessalinagao é realizado num vaso separador como ilustramos na\\
\\
Digitalizada com CamScanner\\

\\section*{Pagina 223}

206 MIDSTREAM \\& DOWNSTREAM\\
\\
Calming patie\\
\\{for FSO")\\
inlet Distributor\\
\\
Water Collector\\
\\
Figura 8.8: Ilustra 0 vaso separador para © processo de dessalinacao.\\
\\
antes ¢ precisa ser tratado para evitar certos inconve-\\
\\
leo bruto apresenta contamin r ce\\
nientes, como incrustagdes @ corrosio dos trocadores de calor, formagao indesejada de\\
\\
coque, diminuigio do desempenho de catalisadores etc.\\
De maneira resumida, 0 dleo bruto misturado e pré-aquecido alimenta o dessalgador\\
\\
onde ocorre a separacio do leo bruto, que sai pelo topo do equipamento, da salmoura\\
que sai pelo fundo do equipamento. O processo de dessalgac4o pode ser feito de duas\\
formas que sio 0 processo elétrico € 0 processo quimico.\\
\\
Proceso elétrico: 0 dleo bruto que escoa na tubagem recebe a adicao de agua para que\\
ocorra a dissolugio dos seus sais na Agua. Em seguida, o dleo passa por uma valvula\\
misturadora que permite o intimo contato entre a 4gua injetada com os sais ¢ os\\
sedimentos. A seguir, a mistura de petréleo, agua e impurezas penetra no vaso de\\
dessalgacio, caminhando através de um campo elétrico de alta voltagem, mantido\\
entre pares de eletrodos metilicos. As forgas elétricas do campo assim criado provo-\\
cam a coalescéncia das goticulas de Agua, formando-se muitas gotas grandes, que,\\
por terem uma maior densidade, caem através do petréleo para o fundo da dessal-\\
gadora, aphid — Os sais e sedimentos. A Figura 8.9 mostra 0 proceso\\
\\
ocorra a alban que escoa na tubagem recebe a adicao de agua para que\\
cador (em Brasil, diate sais na agua. Em seguida, o dleo recebe desemulsifi-\\
revertem 0 cardter 7 cujas particulas, atrafdas J interface agua-dleo.\\
mulsificador ¢ um quimico pee pst 4gua-dleo. Importa realgar que 0 dese:\\
tensoativo. A Figura 8.10 mostr lo para quebrar uma emulsio ja formada. Bu"\\
40 processo quimico de dessalgacio do petrilee.\\
\\
Processo de aquecimento\\
\\
"\\
\\
and ‘\\
ine vaporizagio de parte deste dleo. Na said do\\
Vaporizado e 0 escoamento passa a set dew\\
\\
4\\
Digitalizada com CamScanner\\

\\section*{Pagina 224}

PROCESSOS\\
DE REFINACAO DE PETROLEO BI\\
RUTO 2\\
07\\
\\
Agus\\
‘Oleo bruto\\
\\
Salmoura\\
\\
Figura 8.9: Esquematica do processo eléctrico de dessalgacai\\
0.\\
\\
Agua\\
\\
6leo bruto\\
\\
Salmoura\\
Figura 8.10: Esquemtica do processo quimico de dessalgacao.\\
\\
mistura de\\
\\
salda do phe sesal te pet ea A seraperatiira da mistura de vapor e dleo na\\
TES Ee aia de 300 a 400 °C.\\
\\
ns Ce necessaria para 0 aquecimento 0 calor liberado pelas reagdes de\\
gis sto nl provieg vel que é queimado no interior da fornalha (Figura 8.11). O\\
Ae te cs do processo de destilacao atmosférica ou a vacuo. Os gases\\
de windy ren 0 tubo por onde escoa o petréleo cedem grande parte de sua energia\\
aquecé-lo. Os gases de combustao deixam o forno pela chaminé. Por exemplo:\\
\\
controlada, é possivel fazer a\\
© Alcool ferve primeiro. Ele\\
em ponto de ebulicao mais\\
\\
Se uma mistura de 4gua e 4lcool for aquecida, de forma\\
\\
ree dessas substancias. Aquecendo-se a mistura,\\
\\
por ¢ sai, Depois sai a Agua. A substancia que 8\\
vaporiza primeiro e se transforma em vapors\\
\\
.\\
Sea dgua do mar for aquecida, haverd a separacao dosiled\\
do cloreto de sédio é de 1.440 °C, muito maior que 0 di\\
\\
hi jeito de essas substincias serem destiladas ao mesmo tem\\
\\
ia agua. O ponto de ebu-\\
a agua, de 100 °C. Nio\\
po. A Agua yaporiza.\\
\\
Digitalizada com CamScanner\\

\\section*{Pagina 225}

EAM E DOWNSTREAM\\
\\
208 meeps Tr\\
ay\\
ier\\
yom\\
—-\\
—\\
soa\\
Perdis ow\\
Figura 8.11: Fornalha e torre de destilagao.\\
Destilac¢ao fracionada\\
\\
O petréleo é colocado em um forno, fornalha ou caldeira, e ligado a uma torre de desti-\\
lagdo que possui virios niveis, também chamados de pratos ou bandejas. Conforme vai\\
aumentando a altura da torre, a temperatura de cada bandeja vai diminuindo.\\
\\
Quando o petréleo bruto é aquecido até transformar-se em vapor. O vapor quente é\\
de seguida alimentado pelo destilador (uma torre dividida em intervalos por separadores\\
akg seat os rane Pesadas so retiradas na base da torre e as fracdes mais\\
boride d werden ric na Parte superior da torre de destilagio. Nas se¢des\\
\\% ion A Rigs tad ag Sao retiradas por destilagdo o gasdleo, querosene\\
\\
d-swneche 2 a coluna de destilacdo atmosférica para 0 processo de\\
\\
|\\
\\
|\\
\\
|\\
\\
Resumo da destilagéo atmosférica\\
\\
Em suma, na destilacd, ,\\
de onde so earn por rent © petrdleo é aquecido e fracionado em uma torr\\
solina, nafta, solventes ¢ quer orescente de densidade, gases combustiveis, GLP 8*\\
‘Osenes, dleo diesel € um leo pesado, chamado de residuo |\\
per Este residuo é reaquecido e enviado part\\
| uma pressio abaixo da atmosférica, sendo ent?\\
leo. O residuo de fiundo desta destilacag on Ptoduto chamado genericamente de 6”\\
\\
Gle0 combustivel ov asfalto, oy aré wiaoscharnada 4 vicuo, pode ser especificado com”?\\
° Servir como carga de outras unidades \\#°\\
\\
7\\
\\
Digitalizada com CamScanner\\

\\section*{Pagina 226}

PROCESSOS DE REFINAGAD DE PERG fOoRUTG \\textasciitilde{}—-g\\
a 09\\
\\
Ses Frac horace dy\\
petroleo bruto poe \\&\\
\\
Figura 8,12: Coluna de destilagdo atmosférica.\\
\\
complexas de refinagdo, sempre com 0 objetivo de se produzir produtos mais nobres do\\
que a matéria-prima que os gerou. A Figura 8.13 demonstra o processo de destilacio\\
armosirica,\\
\\
Figura 8.13: Processo de destilacao atmosférica.\\
\\
8.2.0.2 Destilacaéo \\& vacuo\\
Adestilacio A vacuo €um processo que consiste na separa¢do das fragGes mais pesadas do\\
0 (residuo atmosférica) sob pressdes inferiores a pressiio atmostérica. Os produtos\\
' rei \\textasciitilde{}\\textasciitilde{} Proceso incluem 0 gasdleo leve de vacuo, o gasdleo pesado de —\\
A dest rae back. A Figura 8.14 ilustra o pavimento asfaltico como derivado de petréleo.\\
day 4 vacuo faz parte deste esquema justamente para submeter 0 residuo pi\\
ea Processo de separagio, :\\
Pan fe de fundo da coluna de destilacio atmosférica ¢ destinada a um vn ee\\
\\
©m condigdes reduzidas de pressdo, reduzindo des\\
\\
esado\\
\\
o onde é\\
\\
—|\\
\\
Digitalizada com CamScanner\\

\\section*{Pagina 227}

210 MIDSTREAM E powNsTREAM\\
\\
ra 8.14: Asfalto.\\
\\
Figu\\
lo retirar fracdes leves (pelo topo) e médias\\
\\
te de fundo da coluna de destilagao atmosfé.\\
=o 5 fe.\\
ao \\& vacuo com seus produtos de topo,\\
\\
possibilitand\\
corren!\\
de destila\\
\\
rbonetos,\\
da restavam na\\
stra uma coluna\\
\\
ebulicao dos hidroca!\\
(lateralmente) que aint\\
rica. A Figura 8.15 mo\\
intermediarios e de fundo.\\
\\
cy\\
\\
= remy\\
\\
Fi : i\\
igura 8.15: Coluna de destilagéo 4 vacuo e seus produtos.\\
\\
8.2.1\\
Processo de tratamento de nafta\\
Nafta é um termo i\\
genérico ad industri\\
do petrleo, que abrange a me er ae pan es fragoes lev\\
A Saray ot ip de 200 a 2000 °C tilagao da gasolina e do querosene. A fin\\
tida pela destilaco do pet\\
- Gao di\\
pi ig fracionada em peng leo \\& conhecida como nafta Destilaca° Duc?\\
tia naftas, dependendo da faixa de destilacao,\\
\\
* Nafta leve e nafta pesada;\\
= Nafta F ,\\
teve, nafia intermedisria € nafta pesad;\\
a.\\
\\
O fracionamento\\
da nafta,\\
nesses dois ou\\
trés\\
cortes, depende da sua aplicacao final. Des?\\
olin’\\
\\
forma, a nafta\\
. pode ser um\\
Produto final\\
\\}, armazenado em tanques (como nafta,\\
\\
4\\
\\
Digitalizada com CamScanner\\

\\section*{Pagina 228}

PROCESSOS DE REFINAGAO\\
DE PETROLEO ap,\\
WTO aa\\
Hvente), OU produtos intermediarios, indo Para unid,\\
vel . i\\
5 ga para a unidade de reforma catalitica\\
\\
ou 0 "\\
ainda como al\\
\\
).\\
er leve é enviada para tanques, para mais tarde Ser vendida como nafta\\
" -a, ou para ser utilizada na produgao de gasolina automotiva.\\
\\
‘Anafta pesada pode ser enviada para a unidade de\\
ge octanagem (melhoria na qualidade da gasolina) para\\
mente para ser utilizada na mistura de gasolina,\\
\\
lade de tratamento custico, ou\\
(para gerar Basolina de melhor\\
\\
Petrogui-\\
mi\\
\\
reforma catalitica Para aumento\\
Produgio de Sasolina, ou direta-\\
\\
8.2.1.1 Processo de hidrotratamento da nafta\\
\\
O Hidrotratamento (HDT) consiste na elimina\\
depetrileo através de reagGes de hidrogenacio\\
de hidrotratamento de nafta consiste na remo.\\
nafia através da reacao com hidrogénio.\\
\\
cdo de contaminantes de cortes diversos\\
na presenca de um catalisador, O Processo\\
¢40 de compostos sulfurados Presentes na\\
\\
Mercaptanos : RSH + Hy —>» RH +H)S\\
\\
(8.3)\\
Sul fetos : RoS +H, — 2RH + HS (8.4)\\
Dissul fetos : (RS)) +H) —> 2RH + 2H2S (8.5)\\
\\
8.2.1.2 Processo de conversdo — Reforma catalitica\\
\\
A Reforma Catalitica consiste no Trearranjo da estrutura molecular dos hidrocarbonetos\\
contidos em certas fracdes de petréleo, com o intuito de valorizé-las. As gasolinas e as\\
naftas tém, usualmente, um baixo mimero de octanas. Esses produtos sao enviados para\\
areforma catalitica Para que sejam convertidos em naftas ou gasolinas de maior indice de\\
Octanagem.\\
\\
NaReforma, podem ser produzidos, dependendo da faixa de ebulicdo da nafta da carga,\\
uma nafta de alto indice de Octanagem (reformado), para ser utilizada na produgio de\\
Brsolina de alto poder antidetonante, ou um composto rico em hidrocarbonetos aroma-\\
tcos nobres (Benzeno, Tolueno e Xilenos), para serem posteriormente isolados. Neste\\
\\
Processo, também sio produzidas pequenas quantidades de gas combustivel e GLP (Fi-\\
Bura 8.16),\\
\\
8.2.2 Resumo dos processos da refinacao de petrdleo\\
Na Tabela 8.5\\
un li\\
Sig\\
\\
+ apresentamos o resumo dos processos da refinagao do  capeeeenile\\
WTO introdutério com os conceitos bisicos, dessa forma, fem todos A ; v dade os\\
“Plicados neste capitulo, Nesse livro, sio apresentados com maior Pe ds acparesics\\
titos de Tefinacao do petréleo, incluindo o enquadramento das ape “ Coe -\\
40 € tratamento, de acordo com os processos neles Se a Rearranlo de hidro\\
: Craqueamento térmico; Combinagio de hidrocarbonetos:\\
0S; Tratamento e blending, descritos a seguir:\\
\\
T\\
\\
Digitalizada com CamScanner\\

\\section*{Pagina 229}

212 MIDSTREAM E DOWNSTREAM\\
\\
Reformado\\
\\
Figura 8.16: Fluxograma de reforma catalitica.\\
\\
Operagoes topping: podem ser apresentadas como operacées de separacao de hidrocar-\\
bonetos. Os fatores responsaveis por essas operac6es sao fisicos, por acao de energia\\
(alteragdes na temperatura e/ou press4o) ou de massa (relagGes de solubilidade a sol.\\
ventes) sobre o petrdleo ou suas fracdes. Sdo comuns nessa Operacao a destilacio\\
atmosférica, a destilaco a vacuo, a estabilizagao de naftas, extracao de aromiticos,\\
desasfaltacao a solvente, desaromatizacao a furfural, desparafinacao a solvente, de-\\
soleificagao a solvente e absor¢’o de N-parafinas. Podemos resumir, como principais\\
Processos, a destilacao e a desasfaltacao a solventes,\\
\\
Craqueamento térmico: é também conhecido como craqueamento catalitico de hidro-\\
carbonetos, relaciona-se a quebra de moléculas grandes de hidrocarbonetos em mo-\\
Iéculas menores. Para a obtengao dessa quebra de moléculas, os hidrocarbonetos\\
S40 expostos a determinadas fracdes de calor, podendo ter nesse processo a presen¢a\\
\\
de catalisadores, Alguns exemplos de Operacées realizadas nessa etapa sao a Visco”\\
reducao eo termo-craqueamento.\\
\\
oe de hidrocarbonetos: enquanto o craqueamento relaciona-se 4 quebra das\\
moléculas, a combinacio envolve a fusio de duas ou mais moléculas para origina\\
uma molécula maior. Na etapa de combinacio de hidrocarbonetos, encontramo\\
unidades de alquilacio e Polimerizacio,\\
\\
Rearranjo de hidrocarbonetos:\\
\\
: tigi\\
© rearranjo de hi aa estrutura 0\\
nal das moléculas, 0 ‘is 6 jo de hidrocarbonetos alter:\\
\\
! jedades\\
casiona uma nova molécula com diferentes propried em\\
entando o mesmo ntimero de 4tomos de ae\\
€ncontram-se as unidades de reformagio cat\\
\\
4\\
\\
Digitalizada com CamScanner\\

\\section*{Pagina 230}

PROCESSOS DE REFINACAO DE PETROLEO BRUTO 213\\
\\
ngio dos produtos finais, sendo essa a ultima Ctapa do proceso de tefino,\\
obte rtante ressaltar a presenca das unidades de recuperacao de enxofre e as de hi-\\
roeratamento nessa operacao.\\
\\
Tabela 8.5: Resumo dos processos da refinacao de petréleo,\\
\\
Etapas Tipos de pro- Processos\\
cessos\\
\\
Separego Fico Destilagdo Atmosférica; Destilagio A Vacuo; Estabilizacio de\\
\\
naftas; Extragio de aromiaticos; Des;\\
saromatizagio a furfural;\\
ficacdo a solvente; Absoredo de N-parafinas,\\
Visco-Redugio; Craqueamento Térmico; Craqueamento Retar-\\
dado; Craqueamento Catalitico;\\
\\
Hidrocraqueamento; Refor-\\
masao Catalitica; lsomerizagio Catalitica; Alquilagdo Cat\\
\\
Conversio —Quimico\\
\\
talitica;\\
\\
Polimerizagio Catalitica,\\
\\
Tratamento — Quimico Dessalgacao Eletrostatica; Tratamento Caustico; Tratamento\\
Merox; Tratamento Bender; ‘Tratamento DEA/MEA; Hidrotra-\\
tamento,\\
\\
Processos auxi- Quimico Geragio de hidrogénio; Recuperagio de enxofre; Utilidades,\\
\\
liares\\
\\
SS ai et\\
\\
8.2.3 indice de octano\\
O indice de octano é uma medida que:\\
\\
* Avalia a resisténcia de um combustivel a se autoinflamar dentro da camara de com-\\
bustio de um motor;\\
* Quanto maior o indice de octano, maior a resisténcia do combustivel a autodetona-\\
Gao.\\
\\
A maioria dos veiculos utiliza\\
veicul\\
\\
combustiveis com um indice de octano de 95, enquanto\\
los de alta gama usam combustiveis com indice de octano de 98, Em contraste,\\
\\
(barcos de Passagem) requerem um indice de octano de 100, e aeronaves podem\\
Usar indi\\
\\
ices de octano de 120 a 180. Na Tabela 8.6, ilustramos os indices de octano de\\
componentes.\\
\\
Tabela 8.6: \\{indice de octano de alguns componentes.\\
ee a OE REM ce ry\\
\\
Componente \\{ndice de Octano\\
\\
Maioria dos veiculos 95\\
\\
Veiculos de alta gama 98\\
\\
Ferryboat (barcos de passa- 100\\
\\
gem)\\
ites\\
\\
Digitalizada com CamScanner\\

\\section*{Pagina 231}

STREAM E pownsTREAM\\
\\
214 mp\\
rivados do petréleo\\
as refinarias s40 responsaveis por ger\\
dos campos de producao. Os derivado, \\%\\
Alguns desses derivados saem das i\\
s diretamente para distribuidores ¢ a\\
i G\\
em diversas industrias o\\
ra\\
\\
8.2.4 Classificasa° dos de\\
\\
valor da industria petrolifera,\\
\\
do petrdleo recebido\\
Ja gama de aplicag6es-\\
\\
0 comercializado’\\
\\
mo matérias-primas\\
\\
Na cadeia de\\
produtos finais a partir\\
petrdleo tém uma amp.\\
rias prontos para consumo € S\\
sumidores. Outros de!\\
a produgao de outros pro\\
\\
rivados server CO\\} ;\\
dutos finais, como jlustrado na Figura 8.17.\\
\\
Figura 8.17: Esquema do percurso dos derivados do petréleo.\\
‘os obtidos a partir do petréleo incluem solventes, dleos combusti-\\
ene, lubrificantes, asfalto, plasticos, entre outros. Esses\\
le de aplicacdes. Alguns deles sio comercializados dire-\\
enquanto outros server como matéria-prima\\
ilustra uma coluna de destilacao contendo\\
\\
Os principais produt\\
veis, gasolina, dleo diesel, queros\\
derivados tém uma ampla variedad\\
tamente a distribuidores e consumidores,\\
para a indistria petroquimica. A Figura 8.18\\
diferentes derivados do petréleo.\\
\\
cote te destepto <i\\
= 6 i nase\\
Cad 1.4\\
‘ena Nota\\
ry me\\} eo: al .\\
, | oes od\\
ee. Cortes we\\
= wor\\} cece —s\\
| aatsears\\
i —\\
el oe eet\\
| ————\\
Gssce0\\
wre] cuca! CA Tee\\
Geo\\
Sear\\
ao) ae wil tS\\
| os \\_ ;\\
ee\\
i\\
mot\\
\\
Figura 8.18: Coluna\\
tilac¢o com diferentes derivados de petréleo.\\
\\
Petréleo men\\
cionados a, cim:\\
‘a podem ser usados em aplicagdes energe\\
\\
. Os der;\\
\\
€ So classific: rivados ¢\\
\\
lenicic (calor ou luz) mo inne leves, el sJo também conhecidos como co™\\
oucentelha m em co; ties ‘0S ou pesados, Eles geram energ\\
: ss na presenga de ar e uma font de\\
\\
Digitalizada com CamScanner\\

\\section*{Pagina 232}

PROCESSOS DE REFINACAO DE PETROLEO BRUTO\\
215\\
0s destilados leves incluem gasolina eee mouvEs nafta, querosene de aviagio (comb\\
J de jato), querosene € dleos combustiveis, Os destilados intermediarios Cag \\textasciitilde{}\\
asdleo, dleo diesel e éleos combustiveis destilados, Os destilados pesados euglebann\\
éleos combustiveis destilados, dleos minerais pesados, dleos lubrificantes, éleos de flots\\
=, pesados € ceras (parafinas). ta.\\
a também os derivados nao energéticos,\\
gerade de energia. Estes incluem:\\
\\
we\\
\\
que tém aplicagdes nao telacionadas 4\\
\\
« Nafta e gasdleos;\\
\\
+ Lubrificantes, utilizados como matéria-prima na fabricacio de dleos Para veiculos e\\
maquinas industriais;\\
* Solventes domésticos e industriais, como 0 querosene;\\
\\
* Parafinas, utilizadas na industria alimenticia, na fabricagio de velas, ceras, cosméti-\\
cos etc.;\\
\\
* Asfalto, empregado na pavimentacio de ruas e estradas;\\
* Coque, utilizado por industrias na produgio de aluminio e na industria de cimento.\\
\\
Na Figura 8.19, ilustramos alguns derivados do petréleo bruto.\\
\\
Figura 8.19: Derivados do petréleo bruto.\\
\\
B i .\\
*rvados energéticos ou combustiveis\\
a Produsio dos derivados energéticos ou combustiveis é a aplicagio mais comum do pe-\\
él 0. Os combustiveis incluem gas combustivel, GLP, gasolina, querosene, leo diesel,\\
te ‘ombustivel ¢ coque (usado na industria de cimento e aco). Gasolina, diesel ¢ que-\\
ne Fespondem Por cerca de 72\\% do consumo total de petréleo no mundo.\\
\\
* Gasolina: éum dos produtos mais importantes do petréleo. \\& um liquido inflamavel\\
\\
€ volitil Composto por uma mistura de hidrocarbonetos de C5 a C9. A gasolina ¢\\
\\
Digitalizada com CamScanner\\

\\section*{Pagina 233}

MIDSTREAM E DOWNSTREAM\\
\\
216\\
\\
sos na refinaria. Produtos nao deriy 4\\
a\\
\\
8 do\\
\\
Ds process\\
adicionados para aumentar a octan\\
Bem dy\\
\\
fio ¢ outre\\
anol, s40\\
\\
obtida por destilag\\
nol e et\\
\\
petrdleo, como metal\\
\\
gasolina;\\
m motores a diesel ¢ sua caracteristicg\\
Prin,\\
\\
Oleo diesel: € um combustiv\\
\\
cipal é a viscosidade, que Bat\\
ermediaria entre a gasolina ¢ 0 dleo diesel, obtida \\_\\
\\
4leo bruto. B usado como combustivel em turbin\\
i as\\
\\
¢. Tema caracteristica de produzir queim,\\
la\\
\\
el usado \\&\\
nte a Jubrificagao;\\
\\
Querosene: é uma fragio int\\
destilagio fracionada do petr\\
de avido a jato e também como solvent\\
\\
isenta de odor e fumaga;\\
‘o de Petréleo (GLP): ta!\\
\\
mbém conhecido como gas de cozinha, é uma\\
rbonetos obtida a partir do gas natural das reservas sub.\\
de refinagio do petréleo bruto nas refinarias. O GLP \\&\\
stado liquido, com 85\\% em estado liquido e 15\\% em\\
transformar em vapor a medida que os aparelhos\\
\\
Gas Liquefeit\\
mistura gasosa de hidroca\\
\\
terraneas ou do processo\\
armazenado em cilindros em ¢\\
estado vapor. O GLP comega a sé\\
a gis sio utilizados.\\
\\
Derivados nao energéticos\\
\\
Como mencionado anteriormente, os derivados nao energéticos incluem nafta e gasdleos,\\
\\
lubrificantes, asfalto ¢ solventes domésticos € industriais (como aguarras, querosene etc.).\\
\\
* Asfalto: é obtido a partir do residuo das destilagées do petréleo. B amplamente util\\
\\
zado na pavimentagao de estradas, e sua forma oxidada é usada como revestimento\\
impermeabilizante;\\
\\
es ee penal que possui propriedades intermediarias entre a gasolina ¢ 0\\
\\
q ah i” § pri leve, em geral, é destinada a ser misturada com outras naftas\\
\\
—— a re aria pata compor a gasolina. As naftas também sao usadas como\\
\\
ho di sari Ointas, na limpeza a seco e como matéria-prima para a produ-\\
\\
sie er ee aioerr ice Hits Anafta “pesada” pode ter 0 mesmo destino\\
\\
ry i 2 7a\\
\\
peendets gulcee ae para unidades de Reformacio Catalitica, onde sofre\\
\\
pesadas sao usadas para pa ‘ee a Pigs al nobees Aizu\\
\\
? ir a viscosidade do é ‘i i-\\
\\
cado como revestimento em estradas aetalim, Guo \\& posteclonnes *\\
\\
8.2.5 Principal refinaria em Angola\\
\\
Atualmente, Angola\\
\\
imiran ifn possui um .\\
primeiras instalacoes pa cage, conhecida como a Refinaria de Luanda. 4S\\
nagao do petréleo foram inauguradas em malo de\\
\\
erie arewlatnrinn Bafa de Luanda,\\
\\
2007, devido a importancia vcr\\
‘atedepane inngel doen ae desta unidade fabril, que ¢\\# nica\\
as ages da entio proprietiria Total. 18s?\\
\\
A\\
\\
Digitalizada com CamScanner\\

\\section*{Pagina 234}

PROCESSOS DE REFINACAQ DE PETROLEO 8Rur\\
io)\\
\\
217\\
ym 0 objetivo de completar sua cadeia de Negocios,\\
ito 60 7\\
\\
© a refinaria\\
ee ig de Sonango! Refinaria de Luanda, Posteriormente, € Passou a sep\\
mada de 9\\
\\
ss . ™ 2012, ela foi int\\
li de de Negocio de Refinacao e Petroquimica (conforme Figura 8.20), egrada\\
s\\
\\
Figura 8.20: Refinaria de Luanda.\\
\\
Com uma capacidade de tratamento atual de 2.800,\\
\\
equivale a cerca de 57.000,00 barris por dia, os com\\
\\
bustiveis produzidos s\\
apartir da Refinaria por terra e por mar, através de uma rede de Pipelin\\
\\
000,00 toneladas métricas/ ano, o que\\
\\
do transportados\\
es € sealines, paraa\\
\\
ARefinaria de Luanda considera a Seguranca, Higiene e Saiide, a Protegao do Ambi-\\
ente a Qualidade dos produtos fabricados como parte integrante de suas operacées. Por\\
\\
Atualmente, a Refinaria utiliza matérias-primas de origem angolana, como a Palanca,\\
Hungo, Nembae Cabinda, para extrair g4s, combustivel, gas butano, nafta, gasolina, JetB,\\
AL, perrdleo (querosene de iluminacio), gasdleo, fuel-dleo e asfalto. Os excedentes de\\
alguns Produtos, como o fuel-dleo ea nafta, sao exportados para os mercados europeu\\
§ icano. Além disso, esto em andamento outros projetos de construcio de\\
refinarias, inclui um projeto de reforma da Refinaria de Luanda, que aumentaria a\\
\\
Producdo para 300,000 barris Por dia, quatro vezes mais do que a producio atual, como\\
dustrado na Tabela 8,7.\\
\\
8.2.6 Industria Petroquimica\\
Existem Outros derivados que servem como matéria-prima para varias a Ln Pro-\\
0 de diversos Produtos. Para fabricar plasticos, fertilizantes, aia aes oe ™\\
4 lustriais a partir do petréleo, é necessirio, a see\\
a = nner Se doeciagsbs sea transformacdes cue\\
“ de bapueaten sio vineoenlila por exemplo, para produzir ul\\
q ,\\
\\
Digitalizada com CamScanner\\

\\section*{Pagina 235}

AM\\
218 MIDSTREAM £ DOWNSTRE\\
5 para constru¢gao de refinarias em Angola,\\
je to:\\
\\
Tabela 8.7: Proj =\\
apacidade de ‘T\\
Provincia Derivados mento La\\
Proyeto ov\\
ibustivel de avi- 200.000 Barris /d;\\
\\
2 1 chumbo, gasdleo, com! d va\\
ame Bengal Ga lo inant, LPG © quanta\\
\\
acto ie enzfie€ COE |\\
\\
da Gasdleo, gasolina fuel ofl eJet Al Btu att iae\\
\\
Cabinda Cabin: Gaséleo e gasolina To. 0 Bie\\
Soyo Zaire a Lal\\
Namibe Namibe N/A\\
\\
camento. A industria petroquimica 60 ramo da industria quimica que utiliza derivados\\
do petréleo como matéria-prima para a fabricagao de novos materiais (conforme a Fi.\\
\\
| ]\\
j\\
a Separagao dos of;\\
componentes Tinta\\
\\
<—\\_\\_seortos quinios —. LP\\
— “ +e Explosivos\\textbackslash{}\\
\\
|\\
pertumes Corentes tergentes icamentos I fas Feriilizantes\\
\\
Figura 8.21: Importancia e utilidade do petréleo bruto.\\
\\
gura 8.21).\\
\\
= Cosm 7\\
de ear ser to pets ene) some éleos, perfumes e ceras, compoem cerca\\
tém sua origem no petréle ron condicionadores ¢ tinturas para cabelo também\\
de organismos; © bruto, resultado de milhares de anos de decomposi¢a?\\
\\
* Borracha sintética: m :\\
\\
rene sctnente a variacGes intensas de temperatur’ '\\
Pamentos esportivos, ténis ¢ stitui o l4tex em diversos produtos, como eque\\
gal6es (cerca de 30,28 litros) et Em média, um pneu requer o equivalent a8\\
PetrOleo. A borracha sintética também é usada \\&™\\
\\
Prevenindo 1, 0 dle. R\\
produto que. regime Premsciiva, Veilcante Teduz o atrito entre as pegas 40 motoh\\
e-em alguns casos, consiste e ignicio também contém lubrificante: \\&™\\
€M até 90\\% de petrdleo em sua composis3”\\
\\
Digitalizada com CamScanner\\

\\section*{Pagina 236}

PROCESSOS DE REFINACAO DF PETROLEO Bauto 219\\
\\
pode parecer estranho usar uma substan\\
\\
‘ nemdos: Pe muitos remédios, especialmente os\\
sett oe contém benzeno, um derivado do Pp\\
homeop:\\
\\
cia altamente\\
analgésicos e\\
etréleo;\\
de limpeza: a maioria dos produtos de lim,\\
’ Prodan ntes artificiais que podem ser eficazes n;\\
de some para a satide ou 0 meio ambiente;\\
ser pen\\
\\
Poluente em medi.\\
até mesmo alguns\\
\\
peza é composta\\
\\
Por uma série\\
a limpeza, mas ta\\
\\
mbém podem\\
\\
idos sintéticos: geralmente mais econémicos di\\
‘ per icos como nylon, acrilico, spandex ¢ poliéste\\
seit |\\
= produtos para o lar, como cortinas e tapet\\
ie estaria menos cheia;\\
rovavelmente\\
roupas Pp!\\
\\
0 que os tecidos Naturais, tecidos\\
T sdo amplamente usados em rou.\\
es. Sem o petréleo, sua Baveta de\\
\\
« Alimentos:\\
\\
Derivados do petréleo sao usados diretamente (como corantes e conservantes)\\
bad i) =, ay pe) \\textasciitilde{} a\\
indiretamente (em fertilizantes artificiais pesticidas) na Producio de alimen-\\
oui\\
tos;\\
\\
- Chicletes (gomas de mascar): poucas\\
\\
Pessoas sabem, mas os chicletes, consumi-\\
dos no dia a dia, também contém deri\\
\\
vados do petréleo (conformea Figura 8.22).\\
\\
Figura 8.22: Chicletes e refrigerantes,\\
\\
OS iv: 6 é ducdo de\\
* Plisticos; um dos usos mais conhecidos de derivados de petréleo éa pro ih\\
im] tidiano da maioria das\\
i intétic lamente presente no cot\\
Plistico, um composto sintético ai ys)\\
\\
‘a escartaveis\\
lo em garrafas de Agua, caixas de DVD e nos copos descart\\
Pessoas, encontrad NV)\\
\\
usados em festas e no ambiente de trabalho.\\
\\
ral\\
8.2.7 Unidades de Processamento do gas natu\\
A Unidade de\\
\\
cio do\\
Unidades de Recupera¢\\
Processamento do Gas Natural (UPGN) ou Uni\\
\\
Gis Natural (\\
\\
IN) tem com: especificados\\
jung i ionar o gas em produtos ¢ ae\\
ms rl riaabies sherk ear Uiottos desses produtos sao de\\
atender dj aplicagdes requeri hes Biba :\\
8 ara e' Gas Na\\
\\
. 2 om oo o; ma bois toilets de Proeamiins wap he yi\\
\\
arg — ida no Soyo u capuchin de produgio de 6,8 bilh¢\\
\\
), Angola LNG) com uma\\
\\
Digitalizada com CamScanner\\

\\section*{Pagina 237}

220 MIDSTREAM E DOWNSTREAM\\
\\
Jo desta fabrica teve inicio em 2008, asua Conclus,\\
0\\
\\
eladas). A construg neiro LNG em 2013 (Figura 8.23).\\
\\
ilhdes de ton!\\
(5,2 milhGes i produgio do prir\\
\\
em 2012 ¢ 0 inicio\\
\\
Figura 8.23: Planta de Angola LNG no Soyo, Angola.\\
\\
No entanto, antes da construcdo desta planta, Angola ja produzia GLP em uma plata-\\
forma fluruante, a Sanha FPSO. Quando o primeiro dleo foi produzido em agosto de\\
1983, a Sanha LPG FPSO (Bloco 0 - Angola) foi a primeira instalac4o flutuante a combi-\\
nar todas as funcdes de processamento e exportagao de GLP a bordo da mesma unidade\\
(Figura 8.24). Antes de explorarmos a utilidade desta unidade, vamos revisitar alguns\\
conceitos sobre o gas natural.\\
\\
Processo de retirar impurezas, condicionar © gas, eX\\
et como etano, propano, butano e gasolina natural.\\
WWoive enquadrar o gas nas especificacdes necessirias pa\\
\\
transporte, seguranca ou\\
Processamento posterior, Alguns desses processos foram trae\\
\\
dos no Capitulo 7, ou seja:\\
* Remogio de dleo e condensados;\\
* Desidratacio; remocio de dgua;\\
\\
* Adocamento: remogio de CO>. HS ¢ enxofre\\
\\
Digitalizada com CamScanner\\

\\section*{Pagina 238}

PROCESSOS DE REFINACAO DE PETROLEO BRUTO 22\\
1\\
\\
1 perivados do gas natural\\
\\
samento primario é feito na cabeca de P9¢0, en\\
cess . toes\\
\\
‘ealizado nas UPGN, cujo objetivo é a separacao\\
desidratacao, resfriamento e fracionamento\\
\\
quanto 0 processamento com-\\
da parte liquida (agua e hidro-\\
€m produtos especificos para a\\
\\
Metano e etano, que formam o gas residual;\\
.\\
+ Corrente liquida de etano, para fins petroquimicos:\\
\\
» Propano e butano, que formam o GLP ou gas de cozinha e ui\\
\\
'm produto na faixa de\\
gasolina denominado C5+ ou GN.\\
\\
Unidades de Processamento ou Recuperacao de Gases é um\\
parar 0 produto do topo nos seguintes compostos:\\
\\
Processo fisico que visa se-\\
« Gas combustivel;\\
\\
* LPG ou gas de cozinha;\\
\\
* Nafta leve;\\
\\
= Nafta pesada.\\
\\
Nas Figuras 8.25 e 8.26, apresentamos as diferentes configuraces para a obtencao dos\\
derivados do gas natural em uma UPGN.\\
\\
=\\
or\\
a\\
a\\
a\\
a\\
e\\
a\\
a\\
\\
Figura 8.26; Opgées C e D do desenho de uma UPGN.\\
\\
i Je acordo com as\\
‘ male ee @presentamos as configuragdes para 0 design do GLP de acore\\
internacionais de STM.\\
\\
Digitalizada com CamScanner\\

\\section*{Pagina 239}

22 MIDSTR! wnsTREAM\\
2 IDSTREAM E DO! NSTRE | |\\
\\
s de Gas de Petroleo Liquefeito (GPL)\\
ificagdes ae\\
\\
Tabela 8.8: Espe“\\
\\
an\\
\\
Konwseudo\\
\\
Figura 8.28: see\\
Esquema da industria do Petréleo do poco ao produto final.\\
\\
8.2.7.2 Produtos finals de uma refinaria ou UPGN\\
\\
Os produtos finais de uma\\
refinaria ou UPGN esto resumidos nas Figuras 8.27 ¢ ® 28.\\
\\
Digitalizada com CamScanner\\

\\section*{Pagina 240}

TRANS ERE \\& ARMACENAMENTO DO PETRO!\\
LEOBRUTO = -g.\\
23\\
qransferencia e Armazenamento do Petréleo Bruto\\
\\
on atividades de exploragao ¢ producio, estio, em igualdade\\
goles? ede transferéncia ¢ armazenamento. O éleo resultante da P\\
peta ar, precisa ser transportado para as refinarias,\\
\\
“qq ou nom z, ays\\
t on produtos de maior utilidade e valor agregado, co’\\
mado\\
\\
pany petrdleo etc. F\\
oer sdor de transferéncia e armazenamento necesita ter conhecimentos de se,\\
0 wri tubos e acessorios, instrumentos de temperatura, Pressao, nivel e yaz:\\
we enormas de operacao, valvulas, tanques etc., bem como uma s6lida formacao\\
a acetic, fisica equimica para compreender os processos e realizar suas tea\\
com consciéncia € qualidade. 5\\
\\
As operaces de transferéncia € armazenamento sao realizadas via oleodutos, navios\\
eterminais, terminais e refinarias, terminais e terminais, ou s\\
movimentar volumes de leo ou derivados.\\
\\
de importancia, as\\
TOspeccao, seja em\\
onde é processado e transfor.\\
MO gasolina, nafta, querosene,\\
\\
‘guranga\\
40, tipos\\
\\
eja, sempre que se deseje\\
\\
g,3.1 Transporte de petréleo bruto e seus derivados\\
\\
0 dleo produzido nos pogos precisa ser transportado até as refinarias, sendo este trans-\\
porte feito por oleodutos para pocos em terra e por oleodutos ou navios para pocos lo-\\
aalizados no mar. Uma vez refinado e obtidos os derivados, estes precisam chegar ao\\
mercado consumidor, o que é feito através da distribuicdo.\\
\\
83.1.1 Transporte por oleoduto\\
\\
Oleodutos, ou simplesmente dutos (Figura 8.29), so o nome genérico dado aos tubos\\
utlizados para transportar grandes quantidades de petréleo e derivados. Consiste no\\
meio mais econémico e seguro de movimentacao de cargas liquidas derivadas do petré-\\
leo, através de um sistema que interliga as fontes produtoras, refinarias, terminais de\\
aumazenagem, bases distribuidoras e outros consumidores.\\
\\
Os primeiros oleodutos surgiram nos Estados Unidos, ha mais de 100 anos, quase ao\\
mesmo tempo que a indiistria do petréleo. A principio, os oleodutos, ou pipelines, eram\\
"servados ao transporte do éleo bruto, desde o poco até a refinaria ou até 0 porto de em-\\
barque. Posteriormente, tal sistema foi aplicado aos produtos refinados, como gasolina,\\
Pettileo, deo combustivel.\\
\\
© gas natural, existe o chamado gasoduto. Os gasodutos transportam gases sob\\
Pressio. Seu Principio é igual ao dos oleodutos, porém a manutengao da pressao interna é\\
\\
Nada por meio de estagdes de compressio, instaladas a cada 300 km apecsnadamicnls\\
A dstincia entre as estacdes deve ser bastante regular, porque o gas é mais compre\\
\\$60 petréleo e sua pressao diminui rapidamente com a distancia.\\
\\
Digitalizada com CamScanner\\

\\section*{Pagina 241}

‘REAM\\
\\
224 = MIDSTREAM E pownsT\\
\\
Figura 8.29: Oleodutos para transferéncia de petrdleo bruto ¢ seus derivados,\\
\\
8.3.1.2 Transporte hidrovidrio\\
\\
O transporte hidrovidrio compreende os transportes que utilizam o meio aquatico, seja\\
maritimo ou fluvial. Existem varios tipos de embarcagoes, e a escolha do tipo apropriadg\\
de embarcacdo é influenciada pelo tipo de carga, 0 percurso, as condi¢des do porto de\\
origem e destino, entre outros aspectos (Figura 8.30).\\
\\
O transporte de cabotagem é realizado por embarcacées ao longo da faixa costeira, Re.\\
presenta uma parte significativa da movimenta¢ao de cargas pelo modal hidrovidrio, Para\\
© transporte de leo e derivados, é comum a utilizacdo de navios de grande capacidade,\\
com 35 mil, 45 mil, 60 mil e 90 mil toneladas.\\
\\
Figura 8.30: Navio petroleiro Jahre Vicking e Irati.\\
\\
Além disso, grandes navi f\\
transporte, ¢ esses waa aoa como petroleiros, também realizam esse\\
los de superpetroleiros. O transporte dos derivados |\\
\\
@ temperaturas cri i\\
Gigura 31) eM) COM capacidade para a\\
\\
8.3.1.3 Transporte rodovidrio\\
\\
ram relegados i Fi .\\
*segindo plano, pita Privilegiado, enquanto os demais modais !\\
\\
Digitalizada com CamScanner\\

\\section*{Pagina 242}

Oo\\
TRANSFERENCIA E ARMAZENAMENTO Do PETROLEO\\
BRUTO =a.\\
25\\
\\
Figura 8.31: Navio para o transporte de GNL,\\
\\
¢ feito em caminhGes-tanque. Alguns apresentam apenas um tnic\\
outros ja possuem tanques segmentados, possibilitando o transport\\
de produto. A Figura 8.32 ilustra o transporte rodoviario de deriva\\
\\
‘© tanque, enquanto\\
‘e de mais de um tipo\\
dos do petréleo.\\
\\
Figura 8.32: Transporte rodovidrio dos derivados do petréleo.\\
\\
8.3.1.4 Transporte ferrovidrio\\
\\
mi i ‘\\
ding porte de derivados pelo modal ferrovidrio ¢ amplamente utilizado em paises onde\\
Senta de transportes privilegia a intermodalidade. O modal ferroviario repre-\\
\\
deimados alternativa econdmica para o deslocamento de grandes volumes de gasolina e\\
do petréleo,\\
\\
Po Entanto, a velocidade de deslocamento das composicdes deve ser levada em con-\\
a “S40 na anilise de custo/beneficio. Em locais onde seja possivel a integracio =\\
plana? de modal, lise de custo. ee i velado vantagens, quando bem\\
— nes - ny ape antidades. Na Fi\\
\\
» Cxatamente possibilidade de se transportar grandes qu ,\\
Bara 8.33, ilustramos rab ferrovidrios om ° seats de derivados do petrdleo.\\
\\
Digitalizada com CamScanner\\

\\section*{Pagina 243}

jOWNSTREAM\\
\\
226 —MIDSTREAM ED\\
\\
Figura 8.33: Vagao ferrovidrio para o transporte de derivados de petréleo.\\
igura 33:\\
\\
8.3.2 Armazenamento do petroleo bruto e seus derivados\\
\\
ados em tanques apropriados, de acordo com suas\\
\\
Nas bases, os derivados sao armazen d\\
tanques de armazenamento sob pressdo e tanques\\
\\
caracteristicas: tanques atmosféricos,\\
de contencio.\\
\\
8.3.2.1 Tanques atmosféricos\\
\\
Os tanques atmosféricos so equipamentos destinados ao armazenamento de combust-\\
veis liquidos e gasosos e podem ser construidos em dimensées e formas variadas, depen-\\
\\
dendo do tipo de produto e da quantidade a ser estocada. Um exemplo é mostrado na\\
Figura 8.34.\\
\\
los A\\
a reas peliculas protetoras como ph ties) 4 corrosao atmosférica, costuma-se reves™\\
ernamente com tinta de silicato ines tintas especiais, galvanizacdo com zinco etc»\\
\\
Banico ou outros materiais especificos.\\
\\
Digitalizada com CamScanner\\

\\section*{Pagina 244}

TRANSFERENCIA E ARMAZENAMENTO Do\\
PETROLEO BRUT\\
o 227\\
2A rmazenamento sob presséo\\
gat\\
nitro arbonetos no sho liquidos 4 pressio atmosférica @ Necessit\\
eytts sdes superiores para continuarem liquidos, Neste ¢\\
\\
i\\
cagds \\& PRESS i f ‘al\\
\\
pe egaddos €M VASOS de pressiio, que podem ser de dois tipas (\\
gear\\
\\
aM ser arma-\\
\\
480, OS Produtos sio\\
\\
Figura 8,35);\\
\\
+ ciindsicos vol eihures riage formatos, como elipsoidal, hemisférico, cdnico,\\
gorocdnico € rorosférico, Normalmente sio usados para volumes relativame: ,\\
quertas (100 a 200 m”) e destinam-se ao armazenamento de ga\\
diesel, fuel oil e asfalto;\\
\\
inte pe-\\
solina, querosene,\\
\\
« Esfericos: Sto usados para volumes maiores (2000 a 3000 m)e destinam-se\\
\\
a0 ar-\\
mazenamento de LPG ou gis butano.\\
\\
» i es\\
\\
Figura 8.35: Tanques cilindricos e esféricos.\\
\\
8.3.2.3 Bacias de contencgao\\
\\
Os tanques de armazenamento devem estar contidos dentro de uma bacia de contengio\\
que possa conter eventuais derrames em caso de sinistros. A Norma Brasileira (NBR)\\
7505, que regula no Brasil a armazenagem de produtos, prevé a necessidade de bacias de\\
contencio, bem como estabelece os critérios para sua construgio, principalmente no que\\
concerne 4 sua capacidade. A norma estabelece que uma bacia de contengio deve ter\\
capacidade para conter o volume equivalente ao seu maior tanque, acrescido de 10\\% do\\
somatério de todos os demais tanques. Ressalte-se que, para atender as normas vigentes,\\
as bacias devem ser impermeaveis, de modo a nao permitir a contaminagio do solo e de\\
Possiveis lencdis fredticos existentes na regio.\\
\\
Classificagdo dos tanques quanto a finalidade\\
\\
De acordo com as finalidades a que se destinam, os tanques podem ser classificados de\\
\\
‘ Tanques de Armazenamento: onde sio estocados os derivados (gasolina, quero-\\
sene, diesel, GLP etc.) € produtos de alimentagio para as unidades de processo,\\
Quando necessario;\\
\\
a unidade sio arma-\\
nento final,\\
\\
1 Tanques de Recepgdo; onde os prod rovenientes de um\\
" Pp lutos p! : .\\
Zenados antes de serem enviados para outra unidade ou para armazenaly\\
\\
‘50 estejam dentro das especificagdes;\\
\\
Digitalizada com CamScanner\\

\\section*{Pagina 245}

228 MIDSTREAM E DOWNSTREAM\\
nde os produtos fora de especificaco ou provenientes\\
rmazenados € aguardam reprocessamento; \\&\\
Jizadas misturas de produtos ou adigoes de -\\
\\
namento final, quando estiverem dentro da\\
5\\
\\
= Tanques de Residuos: oO\\
operacdes indevidas sao a\\
onde sao rea\\
\\
= Tanques de Mistura:\\
5 io ao armaze\\
\\
tivos, para posterior en\\
especificacoes.\\
g i¢ao\\
\\
Classificagdo dos tanques quanto a posi¢a\\
Quanto 4 posi¢ao, os tanques podem ser clasificados como verticais (para grande porte\\
¢ alta capacidade volumétrica) 0U horizontais (para baixa capacidade volumétrica, nor.\\
malmente usados para produtos especiais, como solventes). Os tanques verticais, consi.\\
derando 0 tipo de teto, podem ser clasificados da seguinte forma:\\
curvos, esferoidais etc., com ou sem selo flutuante.\\
tilizados para produtos nao volateis, como diesel,\\
‘os fixos sao utilizados para arma-\\
\\
* Teto Fixo: podem ser cénicos,\\
Esses tanques sao geralmente u'\\
éleo combustivel e lubrificantes. Tanques cilindric\\
zenar gasolina, querosene, diesel, fuel oil, asfalto (Figura 8.36).\\
\\
Figura 8.36: Tanque cénico e curvo.\\
\\
* Teto Flutuante: sio\\
na superficie do oad a 8.37) em que 0 teto est diretamente apoiado\\
mento ou esvaziamento, eae lo, flutuando durante os movimentos de enchi\\
as perdas por €vaporacio co; nies nals volateis, como gasolina e nafta, reduzem |\\
mM a utilizacio do selo flutuante. Esse tipo de tanque\\
\\
consiste em um selo\\
\\
muito fin, ry\\
\\
oa que ndo produz dearer material especial (espuma de uretano, ne\\
elimina o contato da sy atrito. Ao flutuar sobre o produto, ele pratica-\\
\\
Seguranca perfici i\\
do armazenamento, le liquida com o oxigénio, contribuindo pata 4\\
\\
tivados muito Necessidades\\
vols\\
com teto are ae Petreo brass’, Sio adequados para armazenat\\
faz com 80s de Basolinas ¢ naftas. Existem també™\\
MEO tet0 se mova teto flutuante, mas com uma cima?\\
Verticalmente, orientado por guias. sio\\
\\
Digitalizada com CamScanner\\

\\section*{Pagina 246}

MAZENAMENTO E DISTRIBUIGAO DOS DERIVADOS DO PETROLEO =. 229\\
ARI\\
\\
Figura 8.37: Tanque flutuante.\\
\\
para armazenar gas de rua, propano e aménia. Para\\
sado, podem ser utilizados tanques com teto mé-\\
al é telescépico, acompanhando o movimento das\\
\\
mais frequentemente usados\\
 armazenamento de gas proces\\
yel, cujo funcionamento estrutur\\
paredes da estrutura 4 medida que o produto entra no tanque.\\
\\
agma Flexivel: esses tanques possuem grande capacidade de\\
\\
* Teto Fixo com Diafr:\\
da, alterando o vo-\\
\\
variar o espaco interno, pois a pressao interna pode ser modifica\\
lume do vapor. Essa variagdo é controlada pela deformacao de uma membrana\\
flexivel interna. Normalmente, sao fabricados em plastico rigido para suportar as\\
demandas Iiquidas ou de vapor. Sao frequentemente usados em sistemas fechados\\
para minimizar as perdas causadas pela evaporacao de vapores.\\
\\
84 Armazenamento e distribuicado dos derivados do petro-\\
leo\\
\\
Nesa secao, abordaremos algumas instalagdes pertencentes 4 Sonangol, E.P,, cujas ativi-\\
\\
eas oa com o armazenamento e distribuicao dos derivados do petréleo\\
\\
B € gés natural em Angola. Entre essas instalacdes, destacam-se 0 ICPN, 0 IBV-1, 0\\
V5 ¢ 0 IMUL.\\
\\
8.4, Z\\
1 Instalacdo de enchimento de gas butano\\
\\
Aprincip :\\
Essas j instalacdo para o enchimento dos derivados do gis natural éa instalagao ICPN.\\
\\
Fevereirg de Ses comecaram a ser construidas em 1985 e foram inauguradas em 25 de\\
‘em mrad pelo antigo Chefe de Estado, Eng. José Eduardo dos Santos. A ICPN\\
"Peta 24 horas @ capacidade de enchimento de 40.000 garrafas por dia. A instalagio\\
| Por dia, em trés turnos de trabalho, com cinco carrosséis disponiveis par?\\
\\
Digitalizada com CamScanner\\

\\section*{Pagina 247}

230 MIDSTREAM E DOWNSTREAM\\
iperiee de uma placa especifica Para Barrafy levir\\
a\\
\\
o enchimento de botijas de 5,\\
\\
(Figura 8.38).\\
\\
Figura 8.38: Carrosséis de enchimento de gis butano - ICPN,\\
\\
Cada carrossel possui 18 balangas para 0 enchimento de garrafas de 5 ¢ 12 kg, enquanto\\
dois carrosséis tém 24 balangas para garrafas de 12 kg. As garrafas de 6 kg sao enchidas\\
por um carrossel com 12 balangas, e as garrafas de 51 kg sao tratadas por um carrossel\\
que comporta seis garrafas. Todos os carrosséis funcionam sem interrup¢ées durante os\\
trés turnos, dependendo da demanda dos clientes.\\
\\
A instalagio ICPN recebe o gas butano do porto para duas esferas com um volume de\\
220 m? cada (Figura 40), e a partir dessas esferas, o gas é transferido para seis tanques de\\
armazenamento, cada um com capacidade para 150 m, que alimentam as areas de enchi-\\
mento de gas, incluindo um contentor para o gas levita. A ICPN também possui uma sala\\
de controle de sistemas elétricos, oficinas de manutengio e reparagio de botijas de gis,\\
compressores ¢ bombas. Atualmente, esta em andamento um projeto de “REVAMPING”\\
com o objetivo de melhorar a capacidade e eficiéncia do enchimento de gas butano.\\
\\
A Figura 8.38 mostra o Tanque esférico de 220 m? de armazenamento de gas butano-\\
ICPN. Existem outras instalagdes como a de Lobito, com capacidade de enchimento dié-\\
ria de 20.000 garrafas, e ado Zango, com capacidade de enchimento diério de 10.000 gar-\\
renee bes — ree 1 : instalacdes, embora as provincias do Leste ah\\
traa linha de enchimento de © tenham instalacdes para enchimento. A Figura 8.39 ilus\\
\\
garrafas na fabrica de gas de ICPN.\\
\\
Figura 8,39:\\
Linha de enchimento das garrafas levita - ICPN.\\
\\
4\\
\\
Digitalizada com CamScanner\\

\\section*{Pagina 248}

ARMAZENAMENTO E DISTRIBUICAO DOS DERIVADOS 00 PETROLED\\
\\
231\\
instalagdo de armazenamento de derivados do Petré.\\
8.4 leo bruto\\
\\
is instalagdes para armazenamento e distribuicdo dos derivados do Petréleo\\
pai Ja sio as Instalag6es de Boavista (IBV-1 e IBV-5), Existe um manifold de\\
puto ot Refinaria instalada nas instalages de IBV-1 com as linhas para IBV-1 ¢ IBV.\\
gaibuicao as linhas provenientes dos Navios (Embondeiro e Welwitschia). A instalacio\\
s.bem ak tanques de armazenamento, placa de abastecimento dos cami6es, processo\\
ee eet ‘0 de combustivel contaminado, registro da entrada de camiées, Tegistro de\\
wind camides etc. 7\\
*" ginstalacio de IBV-1 destina-se ao stmazenamento de\\
Bs éusado para o armazenamento e distribuicao de g\\
aula de IBV-1 € composta por oito tanques de gaso\\
jpde otal de 17.900 m?, quatro tanques de gaséleo (Tan\\
de 18.300 m?, um tanque de dgua, um tanque slope par:\\
tanques com capacidade de 10.000 litros cada para o pr\\
tel contaminado. Na Figura 8.40, ilustramos a plac:\\
instalagdes de IBV-1.\\
\\
as pind\\
\\
gasdleo e gasolina, enquanto\\
asdleo e JetAl (querosene), A\\
lina (Tanques 100) com capaci-\\
ques 300) com capacidade total\\
‘a produtos contaminados e dois\\
‘ocesso de tratamento de combus-\\
a de enchimento de camides nas\\
\\
Figura 8.40: Placas de abastecimento de camides em IBV-1.\\
\\
BV recebe diariamente trés mil metros clbicos de gasdleo e mil e quinhentos inet\\
cibicos de gasolina. Amostras dos produtos sao enviadas para o laboratério a fim de\\
\\
€ aprovar sua qualidade. Os analistas levam trés horas para concluir a anilise e\\
ones oe i idores. A\\
Os principais clientes sio a UNDC ¢ Pumangola, que distribuem aos sani —\\
ibuicdo de derivados do petréleo para varios pontos do pais inclui panes ef a\\
\\
; » 4€roportos, portos e caminhos-de-ferro. Esta plataforma tem ae eteda\\
“ender a dez carros a0 mesmo tempo: metade para gasolina e cinco ilhas (design:\\
\\
bomba) para o carregamento de gasdleo.\\
\\
94.3 Fabrica de lubrificantes - IMUL\\
\\
i apacidade nominal de\\
me iada a UNDC, possui uma capaci’ als\\
\\_ ttecades Rei ene restrigGes, sua capacidade atual ¢\\
\\
Digitalizada com CamScanner\\

\\section*{Pagina 249}

232 MIDSTREAM E DOWNSTREAM\\
\\
toneladas por ano. A fabrica é certificada com Iso 9001 € 2015 e sey Process,\\
\\
40 envolve a fabricacio ou formulagio de lubrificantes, com Capacidades de Shey Prod,\\
€embalagem de 1, 4, 5, 20 e 200 litros, além da rotulagem dos Produtos, a, fabrigg\\
vidida em duas zonas de fabricagao: Zona A (dois tanques de fabricacio ¢ Seis de Sta gi,\\
acabado) e Zona B (cinco tanques de fabricagao e quatro de produto acabado), TOduto\\
A fabrica IMUL (Instalagdo da Mulemba) é um empreendimento qe produz 8\\
cante NGOL, herdado da Shell e Mobil pela Sonangol em 1989, e Ocupa uma ex,\\
30 mil metros quadrados com capacidade de producio de 20 mil toneladas por\\
\\
Os lubrificantes da fabrica sio aprovados por API e SAB, duas i\\
nalmente reconhecidas, ea matéria-prima é importada em forma d\\
@ refinaria de Luanda nao produz essa matéria-prima.\\
\\
E importante destacar que a IMUL atualmente produz quatro marc:\\
duas nacionais e duas internacionais: Ngol, Wodex, Lubmarine e TG\\
Oil), que estdo presentes no mercado nacional, demonstrando sua Pi\\
tivo mercado de lubrificantes nacional. Foi realizada uma visita ao\\
Produtos sao testados quanto a qualidade antes e depois da fabrica\\
ilustramos os produtos produzidos na IMUL.\\
\\
lubrig.\\
tensig de\\
; ano,\\
\\
nstituicdes internacig\\
© granel, uma vez ae\\
\\
‘as de lubrificantes\\
‘MO (Toyota Motor\\
Tesenca no Competi-\\
laboratério Onde os\\
Gao. Na Figura 8.41,\\
\\
—\\
\\
Figura 8.41: Lubrificante NGOL - IMUL.\\
\\
Conclusao\\
\\
azenamento de petréleo b 4 Je seus derivados\\
\\
foram a petréleo bruto, gas naturale se\\
\\
po fhe mre longo da formacio, Foram abordados também os conhecimentos\\
Petréleo, a descricio dos diferentes produtos derivados dos procs\\
\\
© as principais unidades de tratamento de petrdleo e gis natural em\\
\\
Digitalizada com CamScanner\\

\\section*{Pagina 250}

CAPITULO 9\\
\\
SEGURANCA, MEIO AMBIENTE E\\
SAUDE NA INDUSTRIA\\
\\
SS\\
\\
E crescente a preocupacZo com o aquecimento global e o impacto das emissdes de gas\\
carbonico (CO2) produzidas pela atividade direta ou indiretamente relacionada com a\\
| indistria petrolifera. O mundo esta empenhado principalmente na reducio das emissdes\\
\\
de CO. Isso exigira a identificacao de fontes de energ'\\
\\
ias alternativas e uma mudanca na\\
forma como os in:\\
\\
dividuos e as grandes indistrias a utilizam.\\
A indistria petrolifera dos Estados Unidos investiu desde 1990, USS 148 trilhdes para\\
melhorar seu desempenho ambiental. Esse valor equivale a US\\$ 504 para cada homem,\\
mulher ou crianga nos Estados Unidos. Esse programa esta resultando na reducao das\\
areas afetadas, com menos geracao de residuos e operagées mais limpas, seguras e de\\
menor impacto ambiental.\\
\\
Nas atividades petroliferas, os gases do efeito estufa, a 4gua de extragao de petréleo e\\
gas natural, e a seguranca das plataformas e sondas maritimas estao entre os elementos\\
\\
que devem ser considerados para melhorar o desempenho ambiental, como ilustramos a\\
seguir:\\
\\
* Efeito estufa: a radiaco solar aquece o solo, que, por sua vez, emite radiacao infra-\\
vermelha para a atmosfera. Grande parte dela se perde no espaco, porém uma parte\\
€ retida na atmosfera por causa de certos gases, como didxido de carbono, vapor de\\
4guaemetano, que atuam como 0 vidro de uma estufa. Esse “efeito estufa” mantém\\
a terra quente o suficiente para favorecer e manter a vida. No entanto, © aumento\\
do diéxido de carbono na atmosfera leva a uma maior “captura” dessa radiagdo infra-\\
vermelha, elevando 0 aquecimento do planeta. O didxido de carbono é emitido por\\
usinas de geracio de energia que queimam combustiveis fésseis, principalmente car-\\
vio, mas também por veiculos e edificios. O desflorestamento é 0 segundo maior\\
\\
Fesponsivel pelo aumento do didxido de carbono na atmosfera. 25\\% de todas as\\
emissSes vérn da queima e do corte anual de cerca de 13 milhdes de drvores. Ome-\\
tano, 0 segundo gis que contribui para o efeito estufa, é produzido en are\\
pela agriculrura, em especial pelo cultivo de arroz ¢ pela flatuléncia do gado ory\\
assim como pela producio de combustivel fossil. A melhoria da eficiéncia energ 5\\
tica, a redugio do consumo, o aumento do uso de fontes alternativas de energia que\\
nio tenham base féssil (edlica, solar, das marés, geotérmica) ¢ a mudanga para com-\\
\\
233\\
Industria do petréleo para ndo-engenheiras, 1* edigdo.\\
By Geraldo André Raposo Ramos Copyright © 2023 Editora Dois03.\\
\\
Digitalizada com CamScanner\\

\\section*{Pagina 251}

USTRIA\\
234 — SEGURANCA, MEIO AMBIENTE E SAUDE NA IND!\\
\\
J, assim como 4 captura € 0 armazenamento de\\
\\
i fe)\\
Stivers: renovaveis como 0 etan\\
ina feito estufa.\\
\\
carbono, podem ajudar a reduzir 0 ¢!\\
m 4gua produzida durante a extracig\\
\\
usal\\
a dyes te COS aah de metano de carvao, essa agua deve ser\\
\\
dugio\\
de petréleo e gas natural ou da produs: de ser usada para irrigacao de culturas\\
\\
ntaminada us: jas antes\\
\\
i usando tecnologias ;\\
descor " i Unidos, novas técnicas de descontaminagao da agua produzida durant\\
Nos Estados S,\\
\\
a extragio de petroleo e gas melhoraram significativamente 2 qualidade da agua,\\
tanto para seu descarte na superficie, como para injecdo ou seu aproveitamento,\\
\\
s; no golfo do México existem mais de 4.000 plataformas,\\
reviveram a essas tormentas histéricas. As plataformas\\
de 1988 foram projetadas para suportar “tormentas de\\
pestades e furacGes de extrema\\
\\
* Segurang¢a contra furacée\\
das quais mais de 97\\% sob\\
maritimas instaladas a partir\\
100 anos”, uma expressao usada para designar tem,\\
intensidade, chegando a categoria 5 (a mais alta).\\
\\
* Captura de carbono: a conversio do carvao para um gas de combustio limpae a\\
captura do diéxido de carbono podem ser usadas como parte do processo de recupe-\\
ragio avancada de hidrocarbonetos ou injetadas nos reservatérios subterraneos de\\
onde foi extraido o petréleo e gas, para reduzir drasticamente a poluicdo do ar. A\\
queima de carvao, na producio de energia, éa maior fonte de emissdes de didxido\\
de carbono. Os reservatorios de petréleo e gas natural nas rochas sao as estruturas\\
geolégicas com maior potencial para confinar o carbono capturado.\\
\\
9.1 Principais impactos ambientais\\
\\
A atividade petrolifera apresenta impactos ambientais tanto positivos quanto negativos.\\
Muitos dos impactos ambientais e grandes desastres ambientais estao relacionados com\\
a haa MP transporte ou refinacao do petréleo. Esses desastres incluem 0\\
rompimento de oleodutos ou pogos maritimos. De forma geral, os im bientai\\
aes ‘ pactos ambientais\\
podem ser subdivididos em impactos atmosféricos,\\
tema e socioeconémicos e culturais,\\
Nem todos 08 ir 4 ivos:\\
iiekiives 66 sei negativos; alguns também sdo positivos. Um dos impactos\\
secnologtas os F ulso oe a traz para toda a cadeia de bens e servicos, trazendo\\
pt cg i snare Bran oportunidades de empregos e renda. Para avaliar ¢s-\\
caracteristicas fisicas, alizar um diagnéstico ambiental detalhado que considere as\\
\\
a € socioeconémi\\
Apés o conhecimento dos némicas da érea,\\
\\&spectos ambientais, a andlise dos impactos ambientais pode\\
\\
, terrestres, aquaticos, sobre 0 ecossis-\\
\\
9.11 Melo flsico\\
\\
na Ti\\
fera. Os principais meant ¢ constantemente degradado pela indastria petroli\\
ais potenciais causadores de impactos negatives 9°\\
\\
4\\
\\
Digitalizada com CamScanner\\

\\section*{Pagina 252}

PRINCIPAIS IMPACTOS, AMBIENTAIS,\\
\\
235\\
\\
Tabela 9.1: Componentes ambientais dos diferentes meios\\
\\
Meio fisico Meio biolégico Meio ange\\
\\
°\\
asd st) Sociedade\\
Componentes - Ecossistemas Cultura\\
:\\
Aguas\\
\\
EE\\
\\
\\textbackslash{}ioo fisico para as diferentes etapas das atividades Petroliferas estao ilustrados na Ta-\\
pela 9.2.\\
qebela 9.2: Os principais aspectos ambientais potenciais causadores de im\\
\\
ipactos negati-\\
vos.\\
\\
© Mecoambinas\\
\\
Gerasio de efluentes domésticos e despcjo no mar.\\
\\
Despejo de residuos oleosos no solo e/ou mar.\\
\\
Mi disposiclo de residuos sélidos contaminados,\\
\\
Emissio de gases.\\
\\
Geragio de ruidos ¢ vibragdes.\\
\\
Vazamentos de dutos ¢/ ou tanques de armazenamento no solo ¢/ou mar.\\
Colisdes entre navios € despejo de dleo bruto no mar.\\
\\
Vazamento ou derramamento de residuos da lavagem dos tanques de navios petroleiros.\\
Despejos de rejeitos do processamento industrial,\\
\\
ssaaeaeeere\\
\\
Variagio da qualidade das 4guas, varia¢ao da qualidade do ar e variagao da qualidade do\\
solo sho os principais impactos ambientais no Meio fisico causados pela produgao de pe\\
tileo.\\
\\
9.1.1.1 Variacéo da qualidade das 4guas\\
\\
\\textbackslash{}s80 ocorre durante a operacao das atividades de descarte na Agua do mar, como agua de\\
drenagem, liquidos de limpeza, efluentes sanitarios e residuos alimentares. O unpanate\\
\\&e efluentes sanitarios ¢ residuos alimentares pode promover o aumento ccapenti ss\\
mate orginica nas 4guas ocednicas. Além disso, o descarte de efluentes et\\
Per substincias quimicas no mar altera as caracteristicas fisico-quimicas, comoa ;\\
\\$e aumenta, ¢ a concentracio de oxigénio dissolvido, que diminui.\\
\\
9.1.1.2 Varlacéo da qualidade do ar\\
\\
nadas com o funcioname! de\\
\\
do aré afetada i das no funcionamento ¢\\
\\
austores pelas emissées de gases relacio! anaber <a =\\
de maquinas ¢ turbinas a diesel, bem como pela que ima de hidrocarbone!\\
\\
Digitalizada com CamScanner\\

\\section*{Pagina 253}

ENTE E SAUDE NA INDUSTRIA\\
\\
236 — SEGURANCA, MEIO AMBI\\
i ralmente consistem em\\
di teste do poco. Essas emissoes atmosféricas ge! NO,, Co,\\
Jurante o tes! .\\
\\
i: rticulado. -\\
SO» Secabe a Ce pleut processo de combustao, as emissdes de dig.\\
importante ressa\\
\\
bi ‘CO;) sio inevitaveis. Embora nao represente riscos para a satide, hy\\
xido = carbono eat seu efeito estufa, que contribui para o aumento da tempera.\\
grande preocupa¢a\\
\\
tura do planeta.\\
\\
9.1.1.3 Variagao da qualidade do solo\\
\\
vegetal para a instalagdo de pocos, ocorre uma variacio na\\
rde sua defesa natural e torna-se propenso a erosio. Ou-\\
de cascalhos contaminados por dleo, que sao dispostos\\
disso, derramamentos de dleo podem causar danos\\
\\
Com a remocio da cobertura\\
qualidade do solo, pois o solo pe!\\
tro motivo é a disposicao no solo\\
ao redor da cabeca dos pogos. Além\\
a vegetacao.\\
\\
9.1.2 Meio bioldgico\\
\\
Os aspectos ambientais listados na Tabela 9.2 interferem na biota, seja marinha ou terres-\\
tre. A biota refere-se ao conjunto de todos os seres vivos de um determinado ambiente\\
yu periodo. Isso pode levar a perda ou extin¢do de organismos. Por exemplo:\\
\\
= Quando os mangais sao atingidos por dleo, o sistema radicular das plantas fica im-\\
permeabilizado, tornando-as incapazes de absorver oxigénio e nutrientes, o que faz\\
com que elas comecem a perder as folhas e no consigam realizar a fotossintese.\\
\\
* Os animais podem morrer em poucos dias por asfixia, enquanto outros podem se\\
intoxicar por ingerirem alimentos contaminados. O déleo também pode afetar o sis-\\
tema de isolamento térmico dos corpos dos animais, principalmente das aves, cuja\\
\\
camada de penas tem a importante funcao de P : duti-\\
vidade térmica. ¢ armazenar ar e gas de baixa condul\\
\\
Para atividades de exploragao maritima,\\
benténicas so os mais afetados,\\
\\
Or ai “, "\\
aplica eh eh Brego Planktos que significa “vagar”) é bastante genérico ese\\
Mo¢ao e sio rilbecrads oe quanto animais, que tem um poder limitado de loco-\\
nica é alramente dedi ee no meio marinho, A comunidade planct6-\\
Embora ¢ inclui representantes de quase todos os filos existentes:\\
\\
esses organismos tenham | i\\
0s verticais na coluna de agua Ocomoco restrita, muitos deles realizam movimen\\
\\
A reducao na intens;\\
pee emporaamente seta Bs na coluna de agua, devido ao aumento da turbide?:\\
fe) to da turb; same . Capacidade fotossintética dos organismos fitoplancténicos.\\
concentracio de fitoplance Zooplancton esta Principalmente relacionado a diminuisi0\\
vel Scorreria sobre os ‘on, \\% Seja, da oferta de alimento. Um impacto direto pos\\
uPidos pelos sélidos em ‘radores, que poderiam ter seus aparatos de filtrage™\\
\\
“SPensio, dificultando a alimentacio.\\
\\
Os organismos das comunidades plancténicas €\\
\\
Digitalizada com CamScanner\\

\\section*{Pagina 254}

PRINCIPAIS IMPACTOS AMBIENTAIS == - 23,\\
Jade bent nica \\& bastante complexa @ inclui uma grande divers\\}\\
coma sbactérias, fungos, plantas e animais, que podem ser classific\\
“0 com etamanho. A fauna invertebrada também pode ser classifi\\
com set OF i ue ocupa no substrato. O depésito de cascalho no fun\\
oma eas organismos, especialmente do macro e megabentos,\\
sar a A\\
et easfixia, : des relacionadas a i i\\
meio abrange todas as ages te lacionadas economia, sociedade e cultura, con-\\
ae astrado na Tabela 9.1. As atividades de exploragao e extracao de petréleo causam\\
forme impacts neste meio, que podem ser positivos ou negativos. Portanto,\\
ene no meio antrépico esto ilustrados na Tabela 9.3,\\
\\
idade de Orga-\\
‘ados de acordo\\
cada de acordo\\
do do mar pode\\
devido ao soter.\\
\\
Os princi-\\
\\
Tabela 9.3: Principais impactos no meio antrépico e suas formas de ocorréncia.\\
\\
ispoctos Ocorréncias\\
\\
Geragio de expectativas As expectativas positivas ocorrem especialmente em relacio aos\\
royalties, A geracdo de empregos c ao estimulo 4 economia, As\\
negativas manifestam-se publicamente através da Preocupacio\\
\\
com as questdes ambientais e com as interferéncias na atividade\\
pesqueira e turistica.\\
\\
Aagio de populacao e aceleragio da ex- Impacto sinérgico vinculado a atragio de trabalhadores de ou.\\
\\
pansio do espaco urbano tros municipios ou até de outros estados, devido a possibilidade\\
de obrer emprego, com consequente aceleragio da expansio ur-\\
bana.\\
\\
Pressio sobre a infraestrutura urbana eso- Aumento da demanda Por infraestrutura re;\\
cal\\
\\
senga dos empreendimentos,\\
Gerasio de renda, dinamizagio da econo-\\
\\
Impacto sinérgico vinculado \\& geracéo de empregos e 4 de-\\
malocale demanda de bens e servigos manda por servigos, promovendo a contratagao de bens ¢ ser.\\
\\
gional devido a pre-\\
\\
vigos.\\
\\
Aomento da demanda por dreas para des- Aumento da demanda por Areas em terra para a disposigio dos\\
\\
nacho final de residuos sélidos residuos a serem gerados pelas atividades de implantagio dos\\
empreendimentos.\\
\\
‘Aumento da producto nacional de hidro- Aumento da producto nacional de dleo ¢ gis natural proveni-\\
carbonetos ente de novas exploragées.\\
\\
Repasse de royalties\\
\\
Incremento na arrecadacdo e contribuigao da receita municipal\\
durante as atividades de produgao e escoamento,\\
——\\_\\_\\_\\_\\_taranee as atividates de produgio eescoamento,\\
Interfertncia na pesca Em geral, relaciona-se \\& influéncia direta da exploragio de pe-\\
tréleo € a alterago no pescado, seja pelo risco de acidentes €\\
vazamentos ou em relacio a érea delimitada para a pesca junto\\
as plataformas de petréleo, que corresponde, segundo determi-\\
\\
naco legal, a um raio de 500 metros em torno da plataforma.\\
\\
we Tbe 9-3, podemos identificar facilmente alguns impactos de cardter positivo, tais\\
0.0\\
\\
A ‘ral 2 i uni-\\
ipl Tepasse de royalties, que é responsavel pelo acréscimo de capital a iste ECE\\
oe *€€ geracio de Tenda, que esta relacionada A geracao de empregos. i” ndimento,\\
\\
7 Para a melhoria da qualidade de vida da sociedade afetada pelo empr\\
Ser considerados uma vantagem para a mesma.\\
\\
Digitalizada com CamScanner\\

\\section*{Pagina 255}

238 SEGURANGA, MEIO AMBIENTE E SAUDE NA INDUSTRIA\\
\\
9.2 Relagdo humana com a natureza\\
\\
A relagio humana com a natureza tem impactos Pe rea nents. Re eseg\\
tacamos 0 crescimento populacional desordenado, a revo] ugao industria, aevo eg dey,\\
industria petrolifera, 0 uso excessivo dos meCUrSOS saree i expansio da agricy ie \\&\\
aglomerados urbanos, o desmatamento, a extingao de esp cles de animais C Vegetais *\\
\\
O capitalismo tornou-se o principal modelo de desenvolvimento Econdmicg atual i\\
maioria dos paises do mundo. Sem respeitar os limites de recuperacio da natureza y\\
voca varios problemas ambientais, com maior destaque Para as alteracdes Climéticas\\
luig\\&o, destruicdo da camada de ozénio, erosio, desertificacao, queimadas, Perda da bin.\\
diversidade e escassez de Agua, além da exclusio social.\\
\\
Este modelo de desenvolvimento é insustentdvel, havendo, Portanto, a\\
de mudangas profundas. E essencial trabalhar o desenvolvimento de forma Sustentive|\\
atendendo as necessidades atuais sem comprometer o futuro das Préximas Reracdes, Bae\\
€um problema crescente, e uma série de medidas de orgdos governamentais, ONGS e da\\
Sociedade civil estao sendo implantadas para contrapor 0 modelo de desen\\
\\
Wolvimento Dio\\
sustentavel. No entanto, neste contexto, vamos nos abster de discutir a gestio ambiental\\
na atividade petrolifera.\\
\\
NCCessidade\\
\\
9.2.1 Residuos sélidos na industria petrolifera\\
\\
Os residuos sdlidos so um fator de agravamento da crise ambiental e\\
\\
tém sido objeto de anilise, discusses,\\
jotivacao para o uso de novas tecnologia:\\
bem como para a aplicacdo mais intensa de tecnologias limpas.\\
\\
A evolugao dos modelos de gestao dos residuos solidos industriais, com énfase nos\\
originados na industria do Petrdleo, pode ser descrita em trés principais fases: a primeira\\
fase, até a década de 70, Priorizou o descarte dos tesiduos; a segunda fase, até o final dos\\
anos 80, teve metas prioritarias que incluiam a recuperac3o e teciclagem; a terceira fase\\
teve inicio na década de 90, com metas de evitar a producio de residuos desde o inicio\\
da cadeia, empregar tecnologias li Pas € adotar os 4 R’s: Repensar, Reduzir, Reutilizare\\
\\
de satide humana,\\
estudos, argumentos Para novas legislacdes e\\
's para maior aproveitamento da matéria-prima,\\
\\
As tecnologias limpas exigem aa\\
biental e tecnolégica integrada aos\\
uso de matérias-pri » Agua e ene:\\
\\
Plicagao continua de uma estratégia econémica, am\\
Processos e produtos para aumentar a leavers\\
Tgia, através da nao eracao, minimizagao ou recic*\\
gem de residuos gerados em 4m Proceso produtivo, Esta abordagem induz a inovaG\\#®\\
eee canbresas, dando um passo em diresao ao desenvolvimento econémico sustentivel¢\\
sap ine p a rin s Mas para toda a regio que abrangem. hierarquia de\\
objetivos e ee S, Segundo a Agenda 21, baseia-se em uma hie\\
\\
duos,\\
©M quatro premissas Principais: reduci inimo dos resi ‘\\
aumento ao maximo da ili pais: redugdo ao mini 0\\
" utilizagao e reciclagem ambientalmente saudaveis dos residu0\\
\\
prom denial a\\
do — re sore es ica to ambientalmente saudaveis dos residuos ¢ aye\\
TVICOS que se Ocupam dos Tesiduos. Na Tabela 9.4, descrevemos oS\\
\\
de residuos Sdlidos, ste\\
— a8 © Seu enquadramento na industria perrolifer\\
\\
Digitalizada com CamScanner\\

\\section*{Pagina 256}

ASPECTOS AMBIENTAIS OE PERFURACAO TERRESTRE DEPOCOS 239\\
\\
pROOUGAG MATA LIMPA\\
\\
9.211 Tratamento dos residuos industriais\\
\\
etrdleo podem causar diversos impactos sobre 0 ambiente,\\
e dimensio do projeto, natureza € sensibilidade do ambi-\\
das tecnologias de prevencao, mitigacao\\
\\
Asatividades na industria do p\\
\\
os quais dependem do estagio\\
\\
|\\
|\\
\\
|\\
\\
\\}\\
\\
j\\
\\
Figura oa: Escopo de atuacao de metodologia da produgao mais limpa.\\
da eficdcia do planejamento €\\
\\
ente circundante \\&\\
\\
t\\
e controle da poluicao. ; ;\\
I O tratamento dos residuos industriais tem suas vantagens e desvantagens. A medida\\
t que a legislagao se torna mais abrangente e restritiva nas quest6es dos residuos sdlidos,\\
principalmente os provenientes da indtistria do petréleo e afins, os aspectos tecnolégicos\\
dos tratamentos e da disposico final desses residuos devem ser abordados e discutidos. O\\
conhecimento das tecnologias existentes, tanto as consolidadas como as novas, suas van-\\
tagens e desvantagens, permitira 20 legislador analisar com mais propriedade os efeitos\\
da implementagao de suas aces € medidas.\\
Existem tecnologias consolidadas que apre'\\
além de assegurarem a completa rastreabilidade de seu pi\\
distria do petrdleo séo as seguintes: aterros, landfarming, tratamento ou processamento\\
dos residuos oleosos, incineracao, dessor¢ao, remediacao biolégica do solo, biopilha, en-\\
capsulamento de residuos industriais, plasma.\\
\\
sentam tamanho e perenidade industrial,\\
rocesso. As mais usadas na in-\\
\\
93 Aspectos ambientais de perfuracgao terrestre de pocos\\
\\
= secao, abordamos as vantagens e desvantagens das técnicas de disposi¢ao final dos\\
dies Solidos provenientes da perfuracao terrestre de pocos de petréleo, em fungao de\\
dos s fatores que tém influéncia direta nessas aplicagoes, de forma a priorizar os méto-\\
mals eficazes para mitigar os impactos causados pelos fluidose residuos da perfuragio\\
ie petréleo ao meio ambiente e a satide publica. ;\\
um si useio e disposicao final dos residuos de forma efetiva e responsive!\\
sistema de gestio de risco ambiental na perfuracao de pogos de petrdleo. O poten\\
\\
| so a chave\\
\\
\\textasciitilde{} — |\\
\\
Digitalizada com CamScanner\\

\\section*{Pagina 257}

U INDUSTRIA\\
240 SEGURANGA, MEIO AMBIENTE E SAUDE NA\\
residuos sOlidos - NBR 10.004,\\
\\
Tabela 9.4: Classificagao de\\
Caracteristicas Indastria do Petré| leo\\
Tipos\\
am risco a satide Borras oleosas ¢ catalj.\\
\\
Ghsse fr Patigos" publics, e meio ambiente sadores contendo me.\\
ou apresentam uma das tais pesados, cloretos\\
seguintes caracteristicas: ou fendis.\\
reatividade, inflamabilidade,\\
toxicidade, corrosividade ¢\\
patogenicidade.\\
Classe Il: Nao Inertes Nao se enquadram nas clas- Borras oleosas e cata-\\
sificagdes de residuos Classe lisadores, dleos lubrifi-\\
I ou II nos termos desta cantes usados.\\
Norma. Podem ter propri-\\
edades tais como: combusti-\\
bilidade, biodegrabilidade ou\\
solubilidade em agua.\\
\\
Classe III Quaisquer residuos  cuja Tijolos, vidros, certos\\
amostra representativa, sub- plasticos e borrachas\\
metida a um contato estatico que nao sao decom-\\
ou dinamico com gua postos prontamente.\\
destilada, conforme teste de\\
solubilizagio (NBR 10.006),\\
nao tiveram nenhum de seus\\
constituintes solubilizados a\\
concentra¢6es superiores aos\\
Padrées de potabilidade da\\
\\
dgua.\\
Caries\\
\\
cial de causar danos a satide PY f\\
industriais pode ser minimizado prose cp dos residuos provenientes das atividades\\
\\
Digitalizada com CamScanner\\

\\section*{Pagina 258}

ig ASPECTOS AMBIENTAIS DE PERFURACAG TERRE\\
\\
STRE DE Pocos 24)\\
\\
31 Disposicao e técnicas de disposicao de rejeitos\\
9,3.\\
\\
erfuracao do pogo, os residuos siio armazenados em dig\\
Durante *e nes recebem os efluentes liquidos oriundos d\\
ee rand, restos de cimentos). Por esse motivo, el\\
igt ai de garantir que nao haja percolaciio dos contamin\\
ae dequada para cada poco deve levar em conta aspec:\\
ae alae do ambiente, a viabilidade técnico-econd:\\
re e materiais necessarios para a disposicao final,\\
mere técnicas de disposicdo dos rejeitos\\
que operam com perfuracao de Pogos de pe\\
pelos mesmos ao meio ambiente e a satide\\
em rés grupos de métodos: fisicos, quimic\\
\\
Wes. Além dos cas.\\
‘aS Operacdes (restos de lama,\\
es devem ser impermeabiliza.\\
‘antes. A técnica de disposicio\\
tos como a legislacao local, a\\
mica ea disponibilidade de Te-\\
\\
\\{\\
|\\
\\
(cascalhos) sao empregadas pelas empresas\\
\\
trdleo, visando minimizar 0 impacto gerado\\
publica. Estas técnicas podem ser divididas\\
‘0s € bioquimicos e termoquimicos.\\
\\
9.3.1.1 Métodos fisicos\\
\\
Os métodos fisicos incluem a impermeabilizacdo de diques de Perfuracao e injecdo de\\
cascalhos em pocos por fraturamento de formacées, Entre varias técnicas usadas, desta-\\
camos as seguintes:\\
\\
 minimizacao do impacto produzido pelos\\
\\
Tejeitos da perfuracdo de pocos petroliferos busca reduzir a quantidade de res\\{duos:\\
(cascalhos e afluentes) gerados no ato da perfuracao.\\
\\
Aterro com diluicao:\\
sem contaminacai\\
Centracdo desses\\
\\
técnica de disposico dos tesiduos sélidos na qual utiliza-\\
io misturados aos res\\{duos sélidos contaminados para reduzir a con-\\
contaminantes a niveis aceitaveis, Nao deve ser aplicada a cargas\\
de hidrocarbonetos nos residuos maiores que 3\\% do seu peso, antes do enterro. A\\
Profiindidade do lencol freatico para a aplicagao desse tecnologia é critica e deve ser\\
Pelo menos 6 metros abaixo da superficie do solo.\\
\\
se solo\\
\\
Nas secdes a Seguir, descreveremos as vantagens e desvantagens dos métodos fisicos.\\
\\
Vantagens @ des\\
\\
vantagens do método de impermeabilizagao de diques\\
© Perfuracdo\\
\\
iF Vantagens:\\
\\
\\{9) Baixo custo: a\\
\\
proximadanieate US\\$ 7,50/m? de cascalhos;\\
(()\\
\\
) Rapida instalagdo da nhs de polietileno de alta densidade: maximo de 2 dia\\
Para diques com drea de 450 m2. » dique,\\
©) Rejeitos sdlidos do poco (cascalhos) sio jogados diretamente nc\\
Sando Temocio e transporte;\\
\\
dispen-\\
\\
Digitalizada com CamScanner\\

\\section*{Pagina 259}

fg \\& SAUDE NA INDUSTRIA\\
\\
242 — SEGURANCA MEIO AMBIENT\\
\\
c que errado com 0s 2 si os, Sendo removida somente a parte |\\
1) Di ti cascalhos, § ndo r | tis\\
\\
com os cas\\
( i é at\\
\\
2. Desvantagens:\\
\\
inagio do subsolo, caso haja problemas com a manta;\\
agi ,\\
\\
(a) Possibilidade de contam\\
amento através de pogos de monitoramento constrys.\\
\\
(b) Necessidade de acompanh\\
dos proximos a0 dique;\\
\\
(c) Nio ha recuperagiio, reciclage!\\
\\
Vantagens e desvantagens do método de injecdo de cascalhos em pogos\\
\\
por fraturamento de formagdes\\
\\
m ou reuso dos contaminantes ou cascalhos,\\
\\
1. Vantagens:\\
(a) Eliminagao dos diques apés concluidos os trabalhos de perfuragao;\\
\\
(b) Disposigio efetiva ¢ final de rejeitos sOlidos e liquidos dos diques em reservatd-\\
rios que nio requerem tratamento prévio de impermeabiliza¢ao;\\
\\
(©) Nao hd necessidade de 4rea na superficie para a disposicao dos cascalhos;\\
\\
(a) Baixo custo operacional: US\\$ 10,00/m? de rejeitos injetados.\\
\\
2. Desvantagens:\\
\\
(a) Necessidade de anilise prévia das formacdes a serem usadas como reservat6rio\\
dos rejeitos quanto 4 sua capacidade de receber os materiais ¢ quanto ao seu\\
isolamento de aquiferos;\\
\\
(b) Disponibilidade de pocos para efetuar a injeco;\\
a* " de preparo prévio dos rejeitos sélidos (redugdo do tamanho dos\\
\\
gros);\\
d) Necessi jei\\
; lecessidade de transporte dos rejeitos até 0 local de injecao;\\
e ‘\\
Disponibilidade de unidade de bombeio para efetuar a inje¢’o;\\
\\
(f) Necessidade de moni\\
\\
injegio ponbey a rer * ntp da inje¢do e do comportamento do pogo apos*\\
\\
(g) Restrigdes i veis canalizacSes das fraturas para formagoes permedvels:\\
"postas pela legislaczo local,\\
\\
Vantagens e a\\
esvantagens do método de perfuracdo de pocos delgados\\
\\
lidos\\
(b) Reducio da érea do diene. (cascalhos) gerados na perfuracio;\\
\\
Digitalizada com CamScanner\\

\\section*{Pagina 260}

ASPECTOS AMBIENTAIS DE PERFURACAG TERRES\\
\\
“TRE DE Pogo 2a3\\
Redugio na geragio de residuos lquidos;\\
\\
(Co) ‘\\
\\
4) Redugio no custo com revestimentos em fungiio dos Menores didmetros\\
\\
( r .\\
\\
2, Desvantagens:\\
\\
(a) Necessidade de uso de equipamentos e ferr\\
\\
amentas ¢\\
cao para manter a verticalidade do Pogo;\\
\\
Speciais durante a Perfura\\
(b) Uso de fluids de perfuragio com\\
\\
des mecanicas do poco;\\
(© Alto custo operacional: US\\$ 98,00 por metro perfur,\\
(d) Necessidade de uma segunda técnica de disposic:\\
\\
(e) Aumento do tempo operacional de\\
trole de diregao do pogo.\\
\\
qualidades especiais para parantir boas condi\\
‘ado;\\
\\
‘io para os reje\\
\\
itos Berados;\\
Perfuragio em raziio d\\
\\
‘4 necessidade de con.\\
\\
vantagens e desvantagens do método de aterro com diluigéo\\
\\
1, Vantagens:\\
\\
(a) Necessidade de area reduzida para a disposicgaio dos rejeitos;\\
\\
(b) Possibilidade de uso do proprio dique de perfuracao para a confecgio das val:\\
\\
(©) Baixo custo de implementacio dessa técnica: US\\$ 11,00/m3;\\
\\
(4) Devido a profundidade das vala:\\
abaixo da superficie),\\
tas ali colocadas;\\
\\
as;\\
\\
S\\$ Com Os rejeitos (topo das va\\
\\
las a 1,5 metro\\
nao ha contato dos contaminantes com as\\
\\
raizes das plan-\\
(e) Monitoramento posterior desnecessfrio,\\
\\
2. Desvantagens:\\
\\
(a) Necessidade de\\
dade;\\
\\
que © lengol freatico esteja a pelo menos 6 metros de profundi-\\
\\
(b) Carga de hidrocarbonetos na mistura solo/ rejeitos contaminados deve ser in-\\
ferior a 3\\% em Peso, uma vez que a biodegradacio é reduzida em fungio do\\
ambiente anéxico criado apés 0 enterro da mistura;\\
\\
(©) Nao ha recuperacdo, reciclagem ou reuso dos contaminantes ou cascalhos.\\
\\
93.1.2 Métodos quimicos e bioquimicos\\
Nestes Métodos, os\\
‘aminantes neles im\\
\\
Ges especiais\\
Compo\\
\\
i "1 5 con\\
cascalhos sdo tratados de forma a reduzir o grau de poluigio dos . ve\\
none\\
pregnados, Em uma das técnicas, utilizam-se elementos quimico va\\
, rias é a prineipa\\
acado de bactérias é a p\\
de temperatura e pressio, ¢ na outra, aa\\
nente de reducdo de contaminantes.\\
\\
mI Agios: licagdo\\
: ulsifi po ida em dois estagios: aap\\
\\
ts 0: esta técnica pode ser resumida em ¢ be: spllcacte\\
ema palo em goticulas menores ¢\\
‘mulsificante, que separa o\\
\\
Digitalizada com CamScanner\\

\\section*{Pagina 261}

A INDUSTRIA\\
SEGURANGA, MEIO AMBIENTE E SAUDE N\\
244 st /\\
xo de um silicato alcalino ao dleo er pater fthe\\
* cat i de sflica inerte ao redor das micr\\
microns, € 4 aplic oe ogee “am\\
roduzindo\\
\\
reagio instantinea, P'\\
\\
de dleo.\\
ros cascalhos sobre 0 solo em cam, ;\\
\\
‘ste em espalha\\
a consiste ome ‘, : erAbica. atra\\
atividade microbiana aerdbica, através da tag\\
\\
e controle da umidade.\\
\\
dos: a técnic\\
‘m, estimulando @\\
nutrientes\\
\\
io método de microencapsulamento\\
\\
* Fazenda de lo :\\
\\
das de até 90 a1 ;\\
e/ou adigao de minerals,\\
\\
Vantagens e desvantagens d\\
\\
1. Vantagens\\
(a) Permite a imobilizagao de rejeitos com qualquer tipo e quantidade de contami-\\
nantes;\\
(b) Possibilidade de reutilizagao em sub-base de estradas e cobertura de aterros sa-\\
\\
nitérios.\\
2. Desvantagens:\\
(a) Alto custo operacional: US\\$ 112,00/m?;\\
(b) Necessidade de transporte até o local da reutilizacao;\\
(c) Existem duvidas sobre os efeitos da lixiviagao a longo prazo (tempo superior a\\
\\
50 anos).\\
\\
Vantagens e desvantagens do método de fazenda de lodos\\
\\
1. Vantagens;\\
\\
(a) Relativamente facil de projetar e implementar;\\
(b) Efeti ituii i\\
. neue Para constituintes organicos com baixas taxas de biodegradacio;\\
custo ;\\
operacional: média de US\\$ 45,00/m?> de material a ser tratado;\\
\\
(d) Tempo de tratame:\\
! nto\\
ny * bioktgico curto: de 6 meses a 2 anos, sob condigGes con-\\
\\
ten\\
atmosfera; dem a evaporar antes da biodegradacio, poluindo 4\\
\\
(b) Requer extensas éreas:\\
\\
(c) Pode nio ser\\
eficaz caso\\
carbonetos a serem haja alta concentracio de metais pesados nos hiro”\\
\\
tados\\
— at PPM), © que inibe o desenvolvimento os\\
\\
Digitalizada com CamScanner\\

\\section*{Pagina 262}

ASPECTOS AMBIENTAIS DE PERFURACAG TERRESTRE DE Poco\\
IS\\
\\
245\\
1.3 Métodos termo-quimicos\\
9,31.\\
s termo-quimicos empregam técnicas €m que os Cascalhos sio aqueci\\
Os yaorey i contaminantes. Apds 0 aquecimento, qecidos para\\
qextrag\\
\\
OS contaminantes sa\\
or exemplo, na Pavii\\
Ordo térmica indire\\
\\
dos, ¢ os cascalhos podem entao ser teciclados, p\\
\\textasciitilde{}\\textasciitilde{} Tipos de extracio com CO supercritico, des\\
= desorcdo térmica indireta é um processo de\\
\\
aio que os custos ¢ beneficios desta técnica de\\
Jogistica locais (Figura 9.2),\\
\\
‘© capturados ¢\\
mentacao de es.\\
tae incineracao,\\
dois estagios,\\
igdes de Operacées e\\
\\
Separacao térmica em\\
Pendem das condi\\
\\
Figura 9.2: Esquema de processo de Separacao térmica de fases (SWACO, 2002).\\
\\
* Incineracao: técnica na qual o cascalho é aquecido a alta temperatura em um inci-\\
nerador a céu aberto. Apesar do uso de filtros Para gases, a incineracao nao é empre-\\
\\
gada atualmente devido ao alto Grau de toxicidade dos gases gerados na queima dos\\
contaminantes existentes no cascalho e As restricdes severas impostas pelos érgios\\
bineis\\
\\
Vantagens e desvantagens do método de desor¢ao térmica indireta\\
\\
(@) Permite a Tecuperacao do hidrocarboneto e seu reuso;\\
©) Possibitita a reciclagem da Agua ou sua reutilizacdo no processo;\\
(©) Baixo custo operacional: US\\$ 24,00/m3; cad\\
(d) Possibilidade de utilizagio dos cascalhos tratados em pavimentagio de estra\\
ou na indastria cimenteira; ;\\
® Pequena area Para implantacio da unidade (cerca de 900 m°).\\
2 Desvantagens:\\
\\
ili " de manga);\\
'9) Utiizacio de equipamento tipo “fim de tubo" (filtro\\
\\
Digitalizada com CamScanner\\

\\section*{Pagina 263}

v INDUSTRIA\\
246 — SEGURANCA. MeIo AMBIENTE E SAUDE NA\\
SOx devido a combustio de gas ou éleo .\\
issi iculados, NOx, € ;\\
(b) Emissio de part\\
rocessoO;\\
) an alterar a composi¢ao dos hidrocarbonetos recuperados;\\
(c ie i\\
(d) Alto custo de implantaca0 do projeto;\\
) Necessidade de transporte dos cascalhos até a unidade de tratamento;\\
(e) Necessi\\
\\
transportar oS cascalhos tratados até o ponto de destinacao final,\\
\\
(f) Necessidade de\\
9.3.1.4 Selecao de métodos de disposigao de rejeitos\\
E necessario haver uma conjuga¢4o de fatores econémicos, técnicos e ambientais na to-\\
mada de decisio sobre a técnica de disposicao final dos residuos da perfuracao terrestre\\
de pocos de petréleo que ser empregada. O aterro com diluigdo é a melhor técnica que\\
atende A legislacao ambiental exigida; além disso, é de facil aplicacao.\\
\\
A técnica mais segura e completa do ponto de vista ambiental é a do COp Supercritico.\\
A impermeabilizacao de diques, apesar de seu baixo custo de implantacao e operacao, é\\
caracterizada por ser um processo no qual os rejeitos sao dispostos sem tratamento, e seu\\
\\
uso deve ser limitado.\\
\\
9.3.2 Impactos ambientais no contexto das politicas de pro-\\
tecao e seus requisitos\\
\\
O grau ou potencial de impactos causados pelas atividades petroliferas depende forte-\\
mente do tipo de operacdo e ambiente. A maioria dos impactos pode ser minimizada\\
ou mitigada através da aplicacao de técnicas de gestao e melhores praticas ambientais de\\
acordo com as normas internacionais, da legislacao local e legislagao globais de protecao.\\
Trés principios fundamentais servem de base para a legislacao, seja ela de ambito nacional\\
\\
ou internacional:\\
* Principio do desenvolvimento sustentavel:\\
* Principio da precaucio;\\
* Principio do poluidor pagador,\\
Principais Leis Regulat6rias\\
de Angola sobre o meio ambiente e as atividades petroliferas:\\
ba . zs Constituigéo da Republica de Angola;\\
* Lei n®S de 19 de\\
Junho de 1998 \\textasciitilde{} Lei de Bases do Ambiente;\\
\\
Digitalizada com CamScanner\\

\\section*{Pagina 264}

ASPECTOS DE SEGURANCA \\& SAUDE SCUPACIONAY\\
\\
247\\
ipecreto Executive n® 12 de 12 de Janeiro de 2005, Gestio de Desc aTBas Operay jonas\\
\\
‘ °\\
pecreto Bxecutivo n°11 de 12 de Janeiro de 2005, Notificagao de Derrames\\
’ on 8;\\
\\
+ Decreto Executive n°08 de 5 de Janeiro de 2005, Procedimentos de Gestio, Remocig\\
¢ Depisito de Desperdicios,\\
\\
9.3.2.1 Gestéo ambiental da Industria Petrolifera\\
\\
Nesta secdo, abordaremos de forma resumida as medidas administrativas ¢ Operacionais\\
de gestio ambiental na inddstria petrolifera de Angola.\\
\\
1. Medidas administrativas de gest4o ambiental:\\
\\
a) Integracdo da gestao de satide, Seguranga e ambiente\\
\\
(>) Prevengdo, na fonte, da Seracdo de residuos pelo uso\\
da poluicao ¢ reutilizagdo dos componentes residuais;\\
\\
(© Consideragdo de todos os com\\
vel estratégico ¢ operacional;\\
\\
(d) Avaliagdo das alternativas com b:\\
ambientais;\\
\\
num unico Programa;\\
\\
de técnicas de Prevencio\\
\\
Ponentes ambientais na tomada de decisdes a ni-\\
ase NO custo/beneficio/risco, incluindo valores\\
\\
(¢) Melhoria continua e reducdo do uso dos recursos naturais.\\
\\
2 Medidas operacionais de Best4o ambiental (Técnicas de Bestdo e disposicio final dos\\
fesiduos):\\
(8) Métodos de reducdo na fonte;\\
(>) Reducdo do volume;\\
‘9 Redugio da toxicidade:\\
(@) Reciclagem e recolha;\\
(©) Reutilizacio;\\
\\{f) Tratamento: métodos térmicos, métodos fisicos,\\
\\
biolbgicos.\\
\\
métodos quimicos, métodos\\
\\
ad Aspectos de seguranca e satide ocupacional\\
\\
“808 Centros de\\
‘e0la(CRMay\\
"ga ¢\\
Ned ‘aide\\
\\
i Je\\
; edo Angol i Formagio Maritima ¢\\
do ceunn ha rie an ramo de segu\\
: orpeore vida carne ds, Raoamans famos a todos os interes\\
Pacional tividades petroliferas. Recomend\\
mad son ie ionados a este campo de conhecimento\\
i Medidas on ¢m seus conhecimentos relaciona\\
\\
Resto t Jem ser resumicas\\
ranca idades petroliteras poc vai\\
\\
: pee von cio de emergeéncia €\\
\\
ny Jeitos a au ‘Ses especiais, Este é um wipico q\\
\\
Jades petroliteras\\
\\
vit\\
? a sua Sensibilidade e importancia para 0 sucesso das ativid\\
\\
Digitalizada com CamScanner\\

\\section*{Pagina 265}

BIENTE E SAUDE NA INDUSTRIA\\
\\
248 SEGURANGA, MEIO AM\\
¢ do Trabalho (SGSST) proporciona um eon,\\
ncia da gestao dos riscos de Seguranga e Saiide\\
tividades da industria do petréleo. Este ‘is\\
\\
do sistema de gestao de qualquer orga.\\
\\
le Gestio da Seguranga ¢ Saud\\
mentas que melhoram 4 eficié\\
ST) relacionados a todas as al ;\\
lerado como parte integrant\\
\\
O Sistema d\\
junto de ferrar\\
do Trabalho (S\\
tema deve ser consid\\
nizacio.\\
OSGSST é b\\
\\
do na politica de SST e deve incluir os seguintes aspectos:\\
aseado I\\
\\
* Definigao da estrutura operacional;\\
* Estabelecimento das atividades de planejamento;\\
\\
* Definicao das responsabilidades;\\
\\
* Alocacao dos recursos necessarios;\\
\\
= Estabelecimento das praticas € procedimentos;\\
\\
* Garantia da identificagao de perigos ¢ avaliacao e controle de riscos.\\
\\
Apés a definicao da politica de SST, é necessdrio projetar um sistema de gestao que\\
abranja desde a estrutura operacional até a disponibilizacao de recursos, passando pelo\\
planejamento, definicao de responsabilidades, praticas, procedimentos e processos. Esses\\
aspectos sio fundamentais para a gestao eficaz da seguranga e satide ocupacional e devem\\
ser integrados horizontalmente em toda a organizacao.\\
\\
E importante destacar que 0 cumprimento da politica de SST da organizacao deve\\
ser assegurado pela alta administracao e deve ser revisado periodicamente e sempre que\\
necessirio. O sistema deve estar focado na gestao de riscos, incluindo a identificacao de\\
perigos e a avaliacdo e controle de riscos.\\
\\
Assim como outros aspectos da industria do petréleo, as questdes de satide ocupacio-\\
nal sio regulamentadas por normas elaboradas por instituicdes crediveis, como 0 APle\\
a Organizacio Internacional de Normalizacio (IS' f jacé -\\
vencdes da Organizacio Mundial de Sau 40 (ISO), além das recomendagées das Con\\
Trabalho (OIT). Essas normas e pels ota a le aren pee “\\
\\
q s deve i e\\
adogao das medidas que sero descritas a ote ‘ill\\
\\
1. Equipament a siti\\
sderaeho pecan Individual (BPIs). A selecao dos EPIs deve levar em con:\\
trabalho, as teitee ds ae 0 trabalhador estaré exposto, incluindo as condicoes de\\
trabalhadores, TPO a serem protegidas e as caracteristicas individuais dos\\
\\
2 Ambiente de trabalho,\\
3. Protecio contra ruido,\\
\\
4. Protecio contra vibracio,\\
\\
5. Protegio contra incén, dios,\\
\\& Radiagbes lonizantes,\\
\\
Digitalizada com CamScanner\\

\\section*{Pagina 266}

PRINCIPAIS ACIDENTES, REGISTRADOS 249\\
\\
uso de Produtos Quimicos:\\
\\
(a) Estado liquido.\\
\\
(b) Estado gasoso.\\
© Estado solido.\\
\\
94.1 Medidas construtivas ou de engenharia\\
\\
Asegurangaéa principal prioridade nas atividades da industria d\\
ie as operacdes Ocorram com o maximo de seguran¢a possive!\\
yersas medidas, que envolvem intervencdes no desempenho e na\\
¢ crucial realizar testes rigorosos em todas as instala\\
\\
Jos em funcionamento ou utiliz4-los em qualquer\\
\\
0 petréleo. Para garantir\\
I, € necessario adotar di-\\
construgao. Além disso,\\
i¢des e equipamentos antes de coloca-\\
Operacao. Dois aspectos sao de suma\\
\\
Aresponsabilidade pela seguranca no desenvolvimento da:\\
wolifera recai sobre todos, desde a base até a alta administracao. A preservacio da vida, da\\
saiide dos trabalhadores e do meio ambiente é uma meta inegociavel, visando a evitar aci-\\
dentes, doencas e danos ao meio ambiente. Isso é Particularmente crucial para empresas\\
que lidam com processos nos quais 0 risco é constante. Portanto, é essencial garantir que\\
forga de trabalho da empresa seja altamente capacitada, abordando questées comporta-\\
mentais, de responsabilidade e conceitos técnicos basicos em Seguranca, Meio Ambiente\\
t Saiide (SMS) e também em conformidade com a legislacao do pais.\\
\\
Todas as empresas devem estabelecer sistemas ou diretrizes contratuais de Seguranca,\\
Meio Ambiente, Satie e Responsabilidade Social, que definem os requisitos, deveres e\\
tesponsabilidades das partes contratantes. Essas diretrizes também estabelecem orienta-\\
ges, requisitos e procedimentos relacionados a SMS e Responsabilidade Social (RS). O\\
objetivo dessas diretrizes é Proteger as pessoas, equipamentos e instalacGes da empresa e\\
\\
do contratante, Promovendo a preservacao do meio ambiente e a satide dos trabalhado-\\
*s, em decorréncia da execucio dos servicos contratados.\\
\\
S atividades da industria pe-\\
\\
95 Principais acidentes registrados\\
\\
Nesta secio, descreveremos alguns dos principais acidentes registrados na industria do\\
Petéleo, Alguns deles podem ter sido causados por erro humano, falta de manutengao\\
adequada, oy até mesmo de forma deliberada. E impossivel listar todos os incidentes,\\
msacteditamos que os los a seguir so os mais significativos. :\\
9 acidente pa arbi ‘Alpha em 1988, no Mar do Norte, foi ° desastre we\\
— Scorrido em plataformas de petréleo. A plataforma estava leona a —\\
we talémetros da costa da Escécia. A causa do acidente foi um Lita: ee ona\\
84s natural que se incendiou, provocando uma explosio rapt Se\\
canta, causada pelo continuo fornecimento de gis, envolveu toda . Les penis\\
ore athe opens: explosto dewrniy's prea pace treet pessoas, a maioria\\
one Tes muito dificil. Esse acidente resultou na mor! mire\\
“evido a inalacao de fumaga, e apenas 62 trabalhadores so\\
\\
Digitalizada com CamScanner\\

\\section*{Pagina 267}

250 —SEGURANGA, MEIO AMBIENTE E SAUDE NA INDUSTRIA\\
historia ocorreu em Kuwait, Golfo Pérsico (ja-\\
Embora nao seja um acidente propriamente\\
rcas iraquianas abriram as valvulas de pocos de petréleo ¢\\
ao se retirarem do Kuwait, causando danos a vida marinha no Golfo Pérsico,\\
madas cerca de 1 milhao e 360 mil toneladas de petréleo,\\
forma Enchova da Petrobras, situada na Bacia de Cam-\\
ou 23 feridos. A explosdo ocorreu durante\\
petréleo, seguida de um grande incéndio. Durante a ten-\\
\\
de aco da baleeira em que alguns trabalhadores estavam\\
fazendo com que a embarcacao despencasse\\
de ter sido causado pela falta de manutencao\\
\\
azamentos de petrdleo da\\
\\
Um dos piores v:\\
Guerra do Golfo.\\
\\
neiro/1991), durante a\\
dito, foi deliberado. As fo:\\
oleodutos\\
Estima-se que tenham sido derrai\\
Em 1984, um acidente na platai\\
pos, resultou na morte de 37 pessoas e deix\\
a perfuragio de um poco de\\
tativa de fuga, um dos cabos\\
ficou preso, enquanto © outro se rompeu,\\
de uma altura de 30 metros. Este acidente po’\\
adequada em alguns equipamentos da plataforma.\\
Baia de Campeche, Golfo do México (ju-\\
\\
A plataforma mexicana Ixtoc 1 rompeu na\\
nho/1979), derramando cerca de 454 mil toneladas de petrdéleo no mar. Essa enorme\\
\\
maré negra afetou as costas de uma 4rea de mais de 1.600 km? por mais de um ano.\\
Estima-se que tenham sido derramados 454 mil toneladas de petréleo.\\
\\
A explosao de um poco no Vale da Fergana é considerada um dos maiores acidentes\\
terrestres jé registrados. Este acidente ocorreu em Uzbequistao (mar¢o/1992) e afetou\\
uma das dreas mais densamente povoadas e agricolas da Asia Central. Estima-se que\\
tenham sido derramadas 285 mil toneladas de petrdleo.\\
\\
Em 1980, a plataforma norueguesa Alexander Kielland, localizada no campo de Eko-\\
fisk, Noruega, sofreu um grave acidente quando um dos principais bracos horizontais de\\
sustenta¢ao de sua estrutura se partiu, causando outras rupturas que inclinaram a plata-\\
forma a um Angulo de 35°. Dos 212 homens a bordo, 123 perderam a vida.\\
\\
Dois superpetroleiros gigantescos colidiram préximo 4 ilha caribenha de Tobago du-\\
rante uma tempestade tropical em julho de 1979, resultando na morte de 26 membros da\\
tripulacao e no derramamento de milhGes de litros de petréleo bruto no mar. Estima-se\\
que tenham sido derramadas 287 mil toneladas de petréleo.\\
\\
Oacidente ocorrido em 2010 com a plataforma Deepwater Horizon, no Golfo do Mé-\\
xico, foi o segundo maior desastre na historia dos acidentes petroliferos, apds a Guerra do\\
pani Lappe eee) Beyond Petroleum, explodiu com 126 pes-\\
e acredita-se que tenham sido derrai pe a iabaladores perderam a vida,\\
\\
mados cerca de 206 milhdes de galdes de petréleo no\\
\\
mar.\\
\\
9.6 Fontes alternativas do petréleo\\
\\
A demand: ‘ 3\\
RSet Pei nae aes te € a preocupagao com as emissdes de didxido de carbono\\
giergia: Baten’ oo hares levado a um aumento na busca por fontes alternativas de\\
d ; ‘ombustiveis tradicionais, principalmente para o setor de transporte.\\
O eeeaaiia a gasolina e o Ultra Low Sulfur Diesel (ULSD)\\
A gasolina aii é ead .\\
hae poteinlrapech aii cieian em relagio as fontes alternativas devido a sua alta den-\\
gética ¢ facilidade de controle em temperatura e pressio ambiente. O ULSD\\
\\
Digitalizada com CamScanner\\

\\section*{Pagina 268}

FONTES ALTORNATIVAS DO PETROL —5y\\
cum combustivel de queima mais limpa usado em novos Motores a diesel, ineluinde\\
‘elculos, 0 que resulta em uma significativa melhoria na qualidade do ar\\
\\
No entanto, a Uniio Buropeia estabeleceu metas Para reduzir as emissbes de CO) em\\
40\\% até 2030, com 0 objetivo de alcangar a neutralidade de earbono em 2050, Nesse con\\
texto, estio sendo feitos esforgos para antecipar a proibigdo da venda de carros a gasolina\\
e diesel, com a proposta de implementagio j4 em 2030, em vez de 2050, como Original\\
mente planejado.\\
\\
Paises como Reino Unido, Noruega, Dinamarca e Espanha ja se comprometeram a\\
climinar os carros a combustio das estradas até 2040, Na Alemanha, por exemplo, a\\
partir de 2030, todos os carros vendidos deverio ser movidos a eletricidade, hidrogénio\\
ou outras fontes de energia limpa. Além disso, esté ocorrendo um aumento significative\\
no investimento em biomassa, que envolve a conversio de residuos de varias plantas em\\
fontes de energia.\\
\\
Atualmente, fabricantes de automéveis estio focados no desenvolvimento de veiculos\\
que utilizem energias alternativas como parte da transi¢io gradual para substituir os car.\\
\\
tos a combustao. Abaixo, apresentamos algumas fontes alternativas Para os combustiveis\\
fosscis:\\
\\
* Reaproveitamento de res\\{duos: no lixo, as bactérias decompéem materiais como\\
\\
alimentos e papel, liberando um gas que é composto por 60\\% de metano. Esse\\
metano pode ser coletado ¢ usado como fonte de energia.\\
\\
* Biocombustiveis: os biocombustiveis sio produzidos a partir da transformagio de\\
acuicares e amido de culturas como milho e cana-de-agiicar em etanol, ou transfor\\
mando soja, canola, linhaga ¢ outras espécies em biodiesel. Também é possivel obter\\
metanol a partir do processamento de madeira ou residuos agricolas. No entanto, o\\
desafio ¢ garantir que o cultivo dessas plantas nao afete o suprimento alimentar.\\
\\
* Hidrogénio ¢ energia solar; carros movidos a hidrogénio podem usar células de\\
combustivel ou motores de combustio in\\
em vez de gasolina. O hidrogénio\\
usando energia solar, Isso resulta\\
fontes de energia renovavel por ex\\
\\
* Oleo vegetal: ¢ possivel adaptar motores de carros para funcionar com 6leo vegetal\\
Esse tipo de dleo pode ser obtido 4 partir de varias plantas ou a partir do residuo de\\
6leo vegetal usado na cozinha,\\
\\
terna adaptados para queimar hidrogénio\\
pode ser produzido a partir da divisio da agua\\
\\
ria em veiculos movidos a agua ¢ energia solar,\\
celéncia,\\
\\
* Células de metanol: cientistas estao desenvolvendo pequenas células de combustivel\\
capazes de gerar eletricidade para carregar baterias usando metanol como fonte. A\\
maior parte do metanol é obtida a partir do gas natural.\\
\\
Muitas dessas fontes alternativas exigem cultivos significativos, 0 que pode representar\\
um desafio para equilibrar a questo da seguranga alimentar global ¢ o combate A pobreza\\
Por exemplo, a Pprodugio de milho para biocombustiveis pode afetar o suprimento de\\
€m certos paises. Além disso, a disponibilidade de dleo vegetal para uso cme\\
combustivel em larga escala, vindo de restaurantes, também \\& uma preocupaglo a set\\
considerada.\\
\\
Digitalizada com CamScanner\\

\\section*{Pagina 269}

gaupe WA INDUSTRIA\\
\\
Conclusao\\
mpactos que podem ser positivos OU Negatives. Oy\\
dade de exploragio, produgio, transporte ¢ refi im\\
icadas medidas mitigadoras, a fim de reduzir acm\\
a\\
\\
dos impactos indissocidveis da fase de produc\\
40,\\
\\
minar os impactos\\
\\
sugerimos que sejam aplicadas medidas compensatorias. Uma das medidas compensa\\
rias que podem ser consideradas é 0 investimento de capital, proveniente do poieman\\
dedor, para agdes sociais ¢ de infraestrutura voltadas 4 populagio. Neste context en\\
aplicagio de um sistema de gestio ambiental eficiente ¢ capaz de evitar a geragio de, ‘\\
guns aspectos ambientais causadores de impactos, € as boas priticas de produgio prt ‘\\
evitar mas atitudes, como 0 despejo de efluentes industriais contendo residuos de 6| \\_\\
sem tratamento, diretamente no mar -™\\
\\
dos pela ativi\\
sejam apl\\
\\
negativos. Para 0 caso\\
\\
Digitalizada com CamScanner\\

\\end{document}
